% UNINTEGRATED ARITHMETIC15 FOUNDATION PART, 25 October2026.
% Same existing arithmetic paper family; source preparation only.
% First batch: accepted general dihedral3/19 theorem and D14 consumer.
% DRAFT pending independent and coordinating review; no execution/compilation.
\section{A strengthened ordinary dihedral comparison}
\label{app:arithdihedral}

This section supplies the ordinary rank-two comparison needed beside
\(D_{14}\).  Its component theorem is general real geometry.
Arithmeticity enters only when the already established arithmetic
\(D_{14}\) prefix bound is used in the final corollary.
Throughout, \(p\) is the ORIGINAL vertex farthest from one strict
interior point.  The theorem is not asserted at an arbitrary
neighbour with a weakened comparison factor.

\begin{proposition}\label{prop:ordinarydihedralthreenineteenths}
Let a compact hyperbolic Coxeter polytope have an actual farthest
vertex \(p\), and write \(z=E(p-x)\), \(\|x\|<1\).
Let one WHOLE component of the incident finite diagram be
\(I_2(m)\), \(m\ge4\), and let \(e\) be the squared norm of the
projection of this SAME \(x\) onto that component.
If \(e\le3/19\), then \(m=7\) and one actual opposite wall
closes both component directions, has a single old contact of
label three there, and has no contact with any other old component.
This is precisely the actual \(O_7\) profile.
Consequently every ordinary \(I_2(m)\), \(m\ge4\), has energy
strictly greater than \(3/19\), with arbitrary finite spectators.
\end{proposition}

\begin{proof}
Use outward unit incident roots, finite Gram \(B\), inverse
\(C=B^{-1}\ge0\), duals \(\omega_i\), and unit inward dual
directions \(-u_i=-\omega_i/\sqrt{C_{ii}}\).
For an actual outgoing edge of length \(L_i\), maximality at
the original \(p\) gives
\[
 x\cdot u_i\ge\tau_i=\tanh(L_i/2)>0,\qquad
 r_i=\tanh L_i={2\tau_i\over1+\tau_i^2}.
\]
The old dihedral inverse is
\[
 C_I=d\begin{pmatrix}1&k\\k&1\end{pmatrix},\qquad
 k=\cos(\pi/m)\ge1/\sqrt2,\quad d=(1-k^2)^{-1}.
\]
Thus each \(\tau_i^2\le e\), and the increasing function
\(4t/(1+t)^2\) on \(0\le t<1\) gives
\[
                         r_i^2\le57/121=:r .
\tag{Dih1}\label{eq:arithdihedralr}
\]

We recall the actual-wall calculation forcing a common opposite.
For an outside actual wall \(c=Ap-v\), let
\(\beta_j=-c\cdot n_j\ge0\) and put
\[
 T_i(c)={(C\beta)_i\over A\sqrt{C_{ii}}}.
\]
For the actual opposite \(c_i\) of edge \(i\),
\(T_i(c_i)=1/r_i\).
The matrix
\(C-C_{\cdot j}C_{j\cdot}/C_{jj}\)
is entrywise nonnegative: its nonzero block is the inverse of
the positive principal Gram obtained by deleting \(j\).
Actual acuteness between two distinct opposites therefore gives
\(T_j(c_i)\le r_j\), while that same nonnegative matrix gives
\(T_j(c_i)\ge k/r_i\) within this old component.
Hence distinct opposites imply \(r_ir_j\ge k\).
But~\eqref{eq:arithdihedralr} gives
\(r_ir_j\le57/121<1/\sqrt2\le k\).
One actual wall \(c\) therefore closes BOTH directions.

Write its full decomposition as
\[
 c=Ap-v_I-v_S,\quad h=A^2>0,\quad
 b=-c\cdot n_1\ge0,\quad g=-c\cdot n_2\ge0,\quad
 Q=\|v_S\|^2\ge0 .
\]
Here \(S\) is the entire orthogonal spectator forest.
Every positive contact is a native magnitude at least \(1/2\);
none of its finite, parallel or dotted possibilities is omitted
before the actual endpoint constraints are imposed.
At least one of \(b,g\) is positive.

Suppose both are positive.
At the actual endpoint releasing direction one, the retained
white Gram has residual \(\delta=1-Q-g^2>0\).
Its actual edge formula is
\[
 \cosh L_1={\sqrt d(b+kg)\over\sqrt{1-Q-g^2}}
 \ge{\sqrt d(1+k)\over\sqrt3}
 ={\,\cot(\pi/(2m))\over\sqrt3}
 \ge{1+\sqrt2\over\sqrt3}>{11\over8}.
\]
The last comparison is \(128\sqrt2>171\), whose positive
squared margin is \(32768-29241\).
It follows that
\[
 e\ge\tanh^2(L_1/2)>
 {11/8-1\over11/8+1}={3\over19},
\]
a contradiction.  Hence exactly one old contact \(\beta>0\)
remains.

We need one whole-spectator fact at small time.
If \(h<1/4\), then \(c\) cannot contact any foreign component.
Indeed a positive foreign contact gives a normalized derivative
\[
 {(C_S\beta)_j\over\sqrt{(C_S)_{jj}}}\ge1/2 .
\]
This uses the full spectator inverse, with nonnegative entries
and diagonal at least one.
Its ratio to \(\sqrt h\) exceeds one.
The actual-wall shadow inequality therefore forces this SAME
\(c\) to be the opposite of that foreign edge as well.

For clarity, the required orthogonal obstruction is specific to
the original farthest vertex.  Suppose a wall closed orthogonal
dual directions \(u_i,u_j\).
Clear it positively off all other incident roots \(W\):
\[
 c_F=c+\sum_{s\in W}\alpha_sn_s,\qquad
 \alpha=B_W^{-1}(-c\cdot n_W)\ge0.
\]
Strict interiority gives \(z\cdot c_F<0\).
In the retained Lorentz three-space its exact form is
\[
 c_F=A p-(A/r_i)u_i-(A/r_j)u_j .
\]
Therefore
\[
 {x\cdot u_i\over r_i}+{x\cdot u_j\over r_j}<1.
\]
Original maximality gives the opposite bound
\(\tau_i/r_i+\tau_j/r_j
 =1+(\tau_i^2+\tau_j^2)/2>1\).
This proves the obstruction, including arbitrary contacts in \(W\).
Thus \(h<1/4\) really forces \(Q=0\).

Put \(q=d\beta^2\).
The smaller normalized derivative square is \(k^2q\).
The common-wall edge formula and~\eqref{eq:arithdihedralr} imply
\[
 q-1\le h=q+Q-1\le rk^2q .
\]
Substituting \(k^2=1-1/d\) and \(r=57/121\) gives the exact
strengthened restriction
\[
                         \beta^2(64d+57)\le121 .
\tag{Dih2}\label{eq:arithdihedralcut}
\]
This is derived at the \(3/19\) threshold; it is not the older
\(1/7\) inequality with an unchanged coefficient.

Since \(\beta^2\ge1/4\), this forces \(d\le427/64\).
For every old \(m\ge8\),
\(d\ge4+2\sqrt2>34/5>427/64\), impossible.
The complete remaining old alphabet is \(4,5,6,7\).
At old \(m=5\),
\(d=2+2/\sqrt5>185/64\):
equivalently \(128>57\sqrt5\), with squared margin139.
Since \(d\) increases with \(m\), if
\(\beta\ge1/\sqrt2\) at any of old \(5,6,7\),
\eqref{eq:arithdihedralcut} would give the opposite
\(d\le185/64\).
Thus all three rows have exactly \(\beta=1/2\).
Every larger finite contact, parallel and dotted contact has
been excluded by the native jump, rather than a sampled label list.

At old four, \(d=2\) and
\(\beta^2\le121/185<\cos^2(\pi/5)\).
The strict comparison is \(413<185\sqrt5\), with squared
margin556.  Only contacts three and four remain, giving
\(q=1/2\) or \(1\).
Consequently \(0<h\le57/242<1/4\).
The small-time argument just proved gives \(Q=0\), whereas
\(h=q-1\le0\), a contradiction.

At old five only \(\beta=1/2\) remains.
Since \(d<3\), \(q<3/4\) and
\(0<h\le r(d-1)/4<57/242<1/4\).
Again \(Q=0\), and \(h=q-1<0\).

At old six, \(\beta=1/2\), \(d=4\), \(q=1\).
If \(Q=0\), then \(h=0\), impossible.
For any \(Q>0\), the actual positive sum
\[
                          \nu=c+2n_1+\sqrt3n_2
\]
is a nonzero future null: all three selected products vanish.
The coefficient order is for OLD6/NEW3, since
\(c\cdot n_1=-1/2\), \(n_1\cdot n_2=-\sqrt3/2\),
\(c\cdot n_2=0\).
Every other actual wall pairs nonpositively termwise with \(\nu\).
It is a feasible ideal ray, contrary to compactness.
No foreign contact was discarded in this affine exclusion.

At old seven only \(\beta=1/2\) remains, and
\[
 0\le Q=1+h-d/4
 \le{427-64d\over484}<{107\over484}<{1\over4},
\]
using \(d_7>5\).
Every positive spectator cost is at least \(1/4\), by inverse
positivity and the native contact lower bound.  Thus \(Q=0\).
We have exactly the ACTUAL isolated-new-three wall in old
\(I_2(7)\), closing both directions: the stated \(O_7\).
Its time square is \(d_7/4-1\in(0,1/3)\);
the bounds \(5<d_7<16/3\) follow from
\(6/5<2\cos(2\pi/7)<5/4\).
Its separate inherited energy floor \(>9/100\) is unchanged.

All weak inequalities included \(e=3/19\), so no ordinary equality
case is omitted.  The proof has retained every finite spectator
component, all its contacts and the same original \(x\).
No arithmetic hypothesis or weakened-neighbour comparison was used.
\end{proof}

\begin{corollary}\label{cor:D14ordinarydihedral}
At the original farthest vertex in arithmetic dimension sixteen,
a whole \(D_{14}\) cannot have an ordinary irreducible
rank-two spectator.
\end{corollary}
\begin{proof}
The fresh true-rank-\(14\) argument in the proof of
Proposition~\ref{prop:f16} gives
\[
                         e_D>.84816={5301\over6250}
\]
for its actual whole rank-two spectator.
Its bootstrap uses only the independently available spectator
floor \(>.09\), so it does not depend on the present strengthening.
For an ordinary \(I_2(m)\), \(m\ge4\), the proposition gives
\[
 \|x\|^2=e_D+e_S>
 {5301\over6250}+{3\over19}
 =1+{719\over118750}>1 .
\]
For ordinary \(A_2\), use its geometric companion bound
\(319/2000\), which exceeds \(3/19\) by \(61/38000\).
Both contradict \(\|x\|<1\).
The actual typed \(N_7/O_7\) alternatives retain their separate
costs and arguments; they have not been charged as ordinary.
\end{proof}

% Next batches may append reviewed arithmetic proof bodies here.
% Preserve this exact first-batch source before any such expansion.

\section{Arithmetic chain windows and simultaneous actual feet}
\label{app:arithmeticchain}

The endpoint arguments below share one arithmetic chain calculation.
We prove it before specializing to the whole \(A_{16}\) and
\(A_{17}\) cases.  The geometric finite-white classification is the
one in the general companion; the arithmetic restrictions and the
simultaneous treatment of isolated feet are given here.
No component dimension is obtained by deleting a node from a
previous ambient theorem.

Let \(N=16,17,18\), let the original incident forest contain the
WHOLE component \(A_{N-1}\), and retain all finite spectators.
Use its zero-sum \(\mathbb R^N\) frame and unit old roots
\(n_i=(e_i-e_{i+1})/\sqrt2\).  Then
\[
 D_i=\omega_i^2={2i(N-i)\over N},\qquad
 C_{ij}=\omega_i\cdot\omega_j={2i(N-j)\over N}\quad(i\le j).
\]
All heights \(h_i=x_A\cdot\omega_i\) concern the SAME original
farthest vector.  Set
\[
             \ell_{16}=.601,\qquad\ell_{17}=\ell_{18}=.6.
\]

\subsection{The complete central inventory at the new losses}
\label{app:arithchaininventory}

For \(5\le i\le N-5\), call an index weak when
\(h_i\le\sqrt{D_i}-\ell_N\).
We claim that every central constraint is either strong, an
isolated-within-\(A\) cap with its actual cleared parameter
\(0<u<1\), or a pure golden end-helper.
The last alternative can occur at at most ONE central index.

First establish the full inventory, without assuming a contrary
component energy.  Deleting \(i\) retains two old paths of ranks
at least four.  Their complete one-node finite extensions are
isolation; one admissible end, \(D\)- or \(E\)-position three;
one terminal four; or two endpoint threes, one in each path.
Two contacts in a single path make a cycle.
A connected two-path completion has rank \(N-1>8\), so finite
classification forces the ordinary path bridge, with no extra
finite foreign arm.

The zero-own repair depends only on side order \(m\ge5\).
At a far-end three the time square is less than its full positive
foreign cost; the next-to-far three has foreign remainder \(<2/m\)
and the same obstruction.
If such a positive foreign contact existed, its actual shadow
would force the same opposite to close an orthogonal component,
contrary to the original-farthest obstruction proved above.
Terminal four has foreign remainder \(<1/m\le1/5\), so no positive
spectator cost, and then negative time.
The other zero-own \(D/E\), internal-four and two-white roles
have their actual positive affine stars or cycles.
Their time is positive through the actual cap and every omitted
wall product is nonpositive.  Own contact is therefore positive.

Write \(q=\sqrt{D_i}\).  Here \(5/2<q\le3\).
The actual edge ratio needed for a loss \(\ell\) is
\[
 R_\ell(q)={1\over\ell}
                -{q-\ell\over q(2q-\ell)}.
\]
It increases with \(q\) and decreases with \(\ell\) on this
rectangle.  Thus
\[
 R_\ell(q)\le R_{.6}(3)=41/27<761/500.
\]
The finite-white residual formulas give the following ordinary
lower ratios, with side order at least five:
near-three/own at least five has
\((101/125+2/5)/\sqrt{3/5}\);
a \(D\)-position has
\((101/125+1/5)/\sqrt{2/5}\);
the \(E\), terminal-four and two-end-three rows are stronger.
All exceed \(761/500\).
The smaller owns have the same actual affine principals.
For the two exceptional roles WITH positive foreign cost,
the complete clearings give ratio
\(>227/(25\sqrt{35})>761/500\), or ratio square \(>121/35\).
These estimates retain the ENTIRE contacted spectator cost.
Hence no additional ordinary nonisolated row is weak.

Unlike the older \(3/2\) test, the present ratio threshold retains
near-own-four also at side order seven.  Up to reflection all
remaining such profiles are
\[
          v=\omega_i/\sqrt2+\omega_{i-1}/2,\qquad i=5,6,7 .
\]
The simple-three helper normalizes the cap to square two over the
actual field, with old product \(-1\); its own product is
\(-\sqrt2\), so \(\sqrt2\in K\).
At the embedding reversing \(\sqrt2\), the time numerators at
\(N=16\) are
\[
        126-88\sqrt2,\quad143-100\sqrt2,\quad154-108\sqrt2 .
\]
They are positive since \(\sqrt2<17/12\).
Their increments with \(N\) are
\((i-1)(3-2\sqrt2)>0\).
The conjugate new-root Schur complement is the negative of this
numerator divided by \(N\), contradicting positive definiteness.
This removes all three side orders and both orientations.

The remaining nonisolated possibility is pure:
\[
 c=Ap-\gamma\omega_i-\omega_e/2,\qquad
 e=1\text{ or }N-1,\quad Q=0,
\]
with own five, finite own at least seven, or dotted.
The exact helper-side identity is
\[
 {(\gamma+1/(2m))^2\over(m+1)/(2m)}-(4\gamma-2)
 ={(\gamma-1-1/(2m))^2\over(m+1)/(2m)}\ge0.
\]
The ratio threshold and native-five gap therefore give
\[
             202/125<t=2\gamma<3139/1458<54/25.
\tag{Astrip}\label{eq:arithAhelperphysical}
\]

For isolated-within-\(A\) contact, retain the whole foreign
projection \(v_S\) and write
\[
 c=Ap-\gamma\omega_i-v_S,\quad Q=\|v_S\|^2,\quad
 u=(1-Q)/\gamma^2,\quad A^2=\gamma^2(D_i-u).
\]
Here \(Q<1\) by finite-white positivity.
Before splitting \(Q\), owns three, four, six and parallel
are excluded for EVERY foreign cost.
Own three has the actual affine \(E_7\) principal with arms
\(1,3,3\); own four has the affine \(4,3,3\) star;
own six has the affine \(6,3\) path.
Each positive affine kernel has positive time through the cap,
and every omitted original wall, including every foreign wall,
has nonpositive product termwise.
Parallel similarly has the feasible positive null pair.
Compactness forbids all these nulls.
Thus \(\gamma\ge\varphi/2\) before any foreign-cost estimate.
When \(Q>0,u>1\), the actual positive clearing has loss at most
\[
 {1\over\sqrt{D_i}+\sqrt{D_i-1}}+\sqrt{u-1}
 <.2+.383<\ell_N .
\]
This uses \(Q\ge1/4\), \(\gamma\ge\varphi/2\), and
\(u-1\le5-3\varphi<.383^2\).
At \(u=1\) the positive cleared vector is a forbidden feasible
future null.
The interval \(0<u<1\) is RETAINED, including \(Q>0\);
we do not charge a strong floor there.

For the remaining PURE isolated estimates, own five has the
actual \(H_4\) projected-point loss
\(<25/76+5/24<.6\).
For own seven, the increasing cubic
\(8a^3-4a^2-4a+1\) gives \(.9<a<.902\).
The actual positive \(A_2\) core gives
\[
 h_i>\sqrt{D_i-a^{-2}}-\sqrt{4/3-a^{-2}},
\]
and hence loss
\[
 <{100/81\over2.6+2.37}+\sqrt{17/150}
 <.25+.337=.587 .
\]
Pure own eight is arithmetically impossible:
\(4\gamma^2=2+\sqrt2\), and its opposite conjugate would require
\(D_i(2-\sqrt2)<4\), whereas
\(D_i\ge55/8>4+2\sqrt2\).
Every embedding of this minimal field extends to the actual
totally real field, so this contradiction applies there.
For all finite owns at least nine, the exact singleton finite
projection gives
\(h_i>\sqrt{D_i-1}-\tan(\pi/m)\).
Since \(\tan(\pi/9)<1386/3727<2/5\), the loss is \(<.6\).
This covers the infinite finite-own tail, not a finite label sample.
All positive finite and negative-norm projections just used are
positive sums of ACTUAL roots; their selected signs, positive
time and all omitted walls are retained as in the companion's
actual-point lemmas.

\subsection{Every pure weak helper is golden over the actual field}
\label{app:arithgoldenfield}

Normalize old \(A\) roots to square two.
If \(r\) is a \(K\)-root on the actual cap line and its simple-three
helper is \(a_e\) of square two, then
\(r^2=2(r,a_e)^2\).  Rescaling by \(-(r,a_e)^{-1}\in K\)
therefore makes the cap square two with helper product \(-1\).
Its own product is \(-t\).
The path cycle through old simple-three bonds makes \(t\) integral.
This compares only connected root lines; no disconnected norm
squareclasses are equated.

Every nonidentity embedding of \(\mathbb Q(t)\) extends to \(K\).
Strict conjugate positivity of the full old-component Schur
complement gives
\[
                    D_it_\sigma^2+2C_{ie}t_\sigma+D_e<4 .
\]
After reversing \(i\le N/2\), the cheaper costs for negative and
positive \(t_\sigma\) are respectively
\[
 Q_-(s)=2is^2+4s+2-2(is+1)^2/N,\qquad
 Q_+(s)=2is^2+2-2(is-1)^2/N .
\]
At the smallest \(N=16,i=5\), their values at \(-4/5,12/25\)
exceed four by \(3/40,59/1000\), respectively.
The costs at those endpoints increase with \(i\) and \(N\);
their values at zero are below four.
Convexity of the quadratics therefore gives the full strict strip
\[
                         -4/5<t_\sigma<12/25 .
\tag{Agoldenstrip}\label{eq:arithAgoldenstrip}
\]
This covers EVERY nonphysical conjugate, not chosen ones.

If the degree is at least five, the nonzero integral norm has
magnitude \(<(54/25)(4/5)^4<1\).
At degree four or three its magnitude is below two.
Any positive small conjugate would make the magnitude below one,
so all small conjugates are negative and their absolute product
is \(1/t\).
At degree four the integral trace lies in \((-1,0)\):
its lower bound is \(t-12/5>-1\), and its upper bound
\(t-3t^{-1/3}<0\) follows from \((54/25)^4<27\).
At degree three the trace lies in \((0,1)\), since
\(t-8/5>0\) and \(t-2/\sqrt t<1\);
the latter follows from
\((54/25)(29/25)^2<4\).
Both degrees are impossible.
At degree two the integral trace is one or two.
Trace two makes the nonzero norm \(t(2-t)\) lie in \((-1,1)\).
Trace one forces \(t<9/5\) and \(t(1-t)\in(-2,0)\);
the norm is \(-1\), and \(t=\varphi\).
Degree one leaves only \(t=2\), the actual parallel pair with
a feasible future null.
Thus every pure weak helper has \(\gamma=\varphi/2\).

The old component and this foreign-free cap have a physically
Lorentzian Gram over \(\mathbb Q(\varphi)\).
If \(K\) were a proper extension, a nonphysical real embedding
fixing that subfield would preserve this indefinite principal
Gram, impossible.
Consequently the helper forces the ACTUAL field
\(K=\mathbb Q(\varphi)\), rather than a convenient overfield.

\subsection{The complete golden packing calculation}
\label{app:arithgoldenpacking}

For a golden cap reversed to \(5\le i\le N/2\), its time
numerator \(M=2NA^2\) is
\[
 M_i^{\rm near}=i(N-i)-(N+1)+(i+2)(N-i)\varphi,\quad
 M_i^{\rm far}=i(N-i)-(N+1)+i(N-i+2)\varphi .
\]
These formulas retain BOTH helper ends.
At \(N=18\), conjugate positivity excludes the midpoint nine
and opposite-end indices six through nine.
The endpoint costs are \((173-84\varphi)/36\) and
\((197-99\varphi)/36\), both greater than one, and the
opposite-end cost is concave in \(i\).
At \(N=17\), only far seven is excluded, by
\((170-84\varphi)/34>1\) using \(\varphi<34/21\).
At \(N=16\) retain every orientation.
The resulting COMPLETE necessary tables are
\[
\begin{array}{c|c|r}
 N&\text{type and }M&\mathrm N(M)\\ \hline
16&\mathrm{near}\ 5:\ 38+77\varphi&-1559\\
 &\mathrm{near}\ 6:\ 43+80\varphi&-1111\\
 &\mathrm{near}\ 7:\ 46+81\varphi&-719\\
 &\mathrm{near}\ 8:\ 47+80\varphi&-431\\
 &\mathrm{far}\ 5:\ 38+65\varphi&-311\\
 &\mathrm{far}\ 6:\ 43+72\varphi&-239\\
 &\mathrm{far}\ 7:\ 46+77\varphi&-271\\ \hline
17&\mathrm{near}\ 5:\ 42+84\varphi&-1764\\
 &\mathrm{near}\ 6:\ 48+88\varphi&-1216\\
 &\mathrm{near}\ 7:\ 52+90\varphi&-716\\
 &\mathrm{near}\ 8:\ 54+90\varphi&-324\\
 &\mathrm{far}\ 5:\ 42+70\varphi&-196\\
 &\mathrm{far}\ 6:\ 48+78\varphi&-36\\
 &\mathrm{far}\ 8:\ 54+88\varphi&-76\\ \hline
18&\mathrm{near}\ 5:\ 46+91\varphi&-1979\\
 &\mathrm{near}\ 6:\ 53+96\varphi&-1319\\
 &\mathrm{near}\ 7:\ 58+99\varphi&-695\\
 &\mathrm{near}\ 8:\ 61+100\varphi&-179\\
 &\mathrm{far}\ 5:\ 46+75\varphi&-59 .
\end{array}
\]
At \(N=16\), near and far eight are reflections at the same
central index, so one row suffices.

Why do these finite tables restrict coexistence?
Subtracting a cap's \(K\)-valued projection onto the FULL old
forest leaves a nonzero vector on the SAME one-dimensional
time line over \(K\).
The caps have no foreign component.
Thus any two time coefficients have ratio in \(K\);
their \(M\)-ratio must be a square IN \(K\), and its rational
norm must be a rational square.
This argument retains all spectators.

At \(N=16\) and \(18\), the positive norm magnitudes in each
table are pairwise coprime nonsquares; their ratios are therefore
not rational squares.  These are direct small integer checks of
the displayed complete tables.
At \(N=17\), different norm squareclasses \(-1,-19,-179\)
cannot mix.
In class \(-1\), write the four entries as
\(42\varphi^3,18\varphi^5,14\varphi^5,6\varphi^7\).
After squares are removed the rational coefficient classes are
\(42,2,14,6\).
Every pair ratio has an odd valuation at three or seven;
none is a rational square or five times a rational square.
Those are exactly the positive rational squares in
\(\mathbb Q(\sqrt5)\), as expansion of
\((a+b\sqrt5)^2\) shows.
For the repeated \(-19\) class,
\[
 {M_{\rm far8}\over M_{\rm near6}}
   ={25+33\varphi\over76}={83+33\sqrt5\over152}
\]
has norm \(1/16\).
A square would require
\(a^2=121/304\) or \(45/304\), both with odd denominator
valuation at nineteen.
Thus different types cannot coexist, including the repeated
norm classes where a norm-only test would be insufficient.

At a fixed index the actual opposite is unique.
Reflected copies of the same noncentral type would have
\[
 -2c_L\cdot c_R=
 { (N-2i)(i+(i+2)\varphi)-(N+2)\over N}
\]
for near helpers, and the same formula with \(i-2\) replacing
\(i+2\) for far helpers.
Their simple-three helper bonds and old path make this
parameter integral by the invariant cyclic criterion, even for
a dotted mutual.
But the constant numerators for \(i=5,6,7\) at \(N=16\) are
\(12,6,-4\); for \(i=5,6,7,8\) at \(17\) they are
\(16,11,2,-11\); and at \(18\) they are \(20,16,8,-4\).
The far rows have the same respective constants.
None is divisible by \(N\), contrary to
\(\mathcal O_{\mathbb Q(\varphi)}=\mathbb Z[\varphi]\).
The central \(N=16\) row cannot occupy two different indices.
Hence at most ONE central index has a golden helper.
This is an all-native-five packing statement, independent of
which of those caps are actually weak at the new loss.

\subsection{Sum every retained isolated foot in one reference}
\label{app:arithsummedfeet}

The intervals \(0<u<1\) retained earlier cannot be declared
strong by the older \(.61\) estimate.
For such a cap clear the entire spectator projection by its
positive actual-root combination.  The actual vector
\[
 \eta=c+v_S+\gamma u n_i
      =\gamma\{\sqrt{D_i-u}\,p-(\omega_i-un_i)\}
\]
has square \(-\gamma^2u(1-u)\), positive time, zero contacted
spectator products, negative or zero \(c,n_i\) products, and
all other actual products nonpositive.
Its normalization is a genuine point of the whole chamber.
Original maximality gives
\[
 x_A\cdot(\omega_i-un_i)
 \ge\sqrt{D_i-u}-\sqrt{u(1-u)}.
\]
If \(\lambda_i=\sqrt{D_i}-\sqrt{D_i-u}+\sqrt{u(1-u)}\)
denotes its loss, then
\[
 \lambda_i\le a_i u+\sqrt{u(1-u)},\qquad
 a_i={1\over\sqrt{D_i}+\sqrt{D_i-1}}.
\]
The rational constants needed below are
\[
\begin{array}{c|ccc}
N&a_i<a&\text{positive reference coefficient }\beta_i\ge b&
 \ell\\ \hline
16&.1982&.022&.601\\
17&.1955&.0205&.6\\
18&.194&.018&.6 .
\end{array}
\]
The \(a\) bounds are positive-square comparisons at
\(D_{\min}=55/8,120/17,65/9\).
For every \(0<u<1\),
\[
 a u+\sqrt{u(1-u)}-\frac b2(2u-u^2)<\ell.
\tag{Aabsorb}\label{eq:arithAabsorb}
\]
For \(u\le1/2\) use \(a/2+1/2\), giving
\(.5991,.59775,.597\).
For \(u\ge1/2\), use \(2u-u^2\ge3/4\) and
\(\max(au+\sqrt{u(1-u)})=(a+\sqrt{1+a^2})/2\).
Radical upper bounds \(1.0195,1.019,1.019\) give respectively
\(.6006,.5995625,.59975\), all strictly below their \(\ell\).

Let \(H\) be a concave reference with available nonnegative
coefficients \(\beta_i\), positive ones at least \(b\):
\[
 w_j=(H_j-H_{j-1})/\sqrt2,\quad
 \beta_i=w\cdot n_i,\quad
 w=\sum_i\beta_i\omega_i,\quad
 E_H=\|w\|^2=\sum_i\beta_iH_i .
\]
Require \(H_i+\ell<\sqrt{D_i}\) at its positive central
coefficients.  Omit the possible golden index.
If \(S\) is the set of retained isolated indices with
\(\beta_i>0\), define ONE vector
\[
                       Z=w-\sum_{i\in S}\beta_i u_i n_i .
\]
Sum all actual point pressures BEFORE estimating its norm:
\[
 x_A\cdot Z\ge E_H-\sum_{i\in S}\beta_i(\lambda_i-\ell).
\]
Because \(w\cdot n_i=\beta_i\) and
\(n_i\cdot n_j\le0\) for \(i\ne j\),
\[
 \|Z\|^2\le E_H-\sum_{i\in S}\beta_i^2(2u_i-u_i^2).
\]
It follows that
\[
 e_A\ge2x_A\cdot Z-\|Z\|^2
 \ge E_H+\sum_{i\in S}\beta_i
       \{\beta_i(2u_i-u_i^2)-2(\lambda_i-\ell)\}>E_H.
\tag{Asummed}\label{eq:arithAsummed}
\]
Strictness follows from~\eqref{eq:arithAabsorb} if \(S\ne\varnothing\),
and from the strict available constraints otherwise.
Adjacent isolated indices and positive foreign costs are allowed.
No nonadjacency is assumed for this new isolated set, and no
independently optimized energies are added.

The references used by the endpoint arguments are
\[
\begin{array}{c|c|c|c}
N&H_5,H_6,\ldots\text{ to midpoint}&E_H&
 \text{maximum internal golden omission loss}\\ \hline
16&2.021,2.137,2.205,2.227&.8354522&.000576\\
17&2.056,2.186,2.269,2.310&.8708972&.00055225\\
18&2.087,2.228,2.324,2.381,2.399&.9037838&.00050625 .
\end{array}
\]
Reflect \(H_{N-i}=H_i\), use linear end tails, and at \(N=17\)
take \(H_9=H_8\).
The positive curvatures have exactly the lower bounds \(b\) above.
All floor inequalities are positive-square comparisons:
the new heights are the old \(.61\) heights shifted by
\(.61-\ell\), so \(H_i+\ell\) is unchanged.
A single internal omission sets its coefficient zero by neighbour
averaging and increases the adjacent coefficients.
An end omission replaces the first tail by the chord \(0\)--\(6\),
or its reflection.  These are changes in ONE reference.
In particular the full and omitted references preserve the
summed-foot argument whenever their positive central coefficients
have the displayed lower bounds; subsequent endpoint grafts must
check that condition again rather than assume it.

\subsection{A golden end hole pays for its own endpoint}
\label{app:arithgoldenend}

For later proper-component consumers we record how the exceptional
end omission is closed.  Suppose its central index is five,
reversing if necessary, and put \(\gamma=\varphi/2<.81\).
For either helper end \(e\), with its adjacent old node \(j\),
the actual finite \(H_4\) correction is
\[
 b_c=\gamma n_5+\tfrac23n_e+\tfrac13n_j,\qquad
 \eta=c+b_c,\quad
 \eta^2=\tfrac23-\gamma^2<.111^2 .
\]
The actual positive finite projected point has pure spatial
displacement \(v-b_c\).
Its normalization and original maximality give
\(x\cdot(v-b_c)\ge A-.111\).
Since \(x\cdot b_c>0\), the helper-end signed coordinate satisfies
\[
                  x_{\rm end}>
             \sqrt2(A-.111-\gamma h_5)>.49 .
\]
Both orientations are covered: the smaller opposite-end time
squares, at \(N=16,17,18\), are
\[
 (38+65\varphi)/32,\quad(42+70\varphi)/34,\quad
 (46+75\varphi)/36.
\]
They give \(A>2.115,2.136,2.155\).
Weakness gives \(h_5<2.022,2.057,2.088\).
Using \(\sqrt2>1.4\), the corresponding signed-coordinate lower
bounds are \(.512652,.502362,.493808\).
The nearer-end time is larger.
Here the signed coordinate is \(x_1\) at helper one and
\(-x_N\) at helper \(N-1\).
No strong floor at the missing index five has been inferred.

Use the chord \(0\)--\(6\) at the omitted end.
At the helper end the resulting reference is flat on \(L=6\)
coordinates in the near case and \(L=5\) in the far case,
with signed mean \(.25<b_0<.3\).
Replace that block by
\[
             .49,\quad
             {Lb_0-.49\over L-1},\ldots,
             {Lb_0-.49\over L-1}.
\]
Its sum is unchanged.  The latter coordinates lie in
\((.19,.262)\); the next old coordinate is below \(.101\).
Thus concavity persists, the new central boundary coefficient
is \(>.062>b\), and all other positive central coefficients
keep their previous lower bound.
The new endpoint coefficient is positive and is paid by the
actual endpoint pressure just proved.
The one-reference energy increases by
\[
                {L\over L-1}(.49-b_0)^2>.04332 .
\]
The pre-graft chord energy is exactly
\[
 E_H-\frac12\{H_5^2/5+(H_6-H_5)^2-H_6^2/6\}.
\]
At \(N=17,18\) it is respectively
\(12569249/15000000\) and \(13079276/15000000\);
at \(N=16\) it exceeds \(.8008\).
These direct rational formulas give, after the graft and
simultaneous absorption~\eqref{eq:arithAsummed}, the conservative
component bounds
\[
                 e_{A_{15}}>.8348,\qquad
                 e_{A_{16}}>.8703,\qquad
                 e_{A_{17}}>.9032 .
\]
If the single golden hole is internal instead, subtract the
internal losses in the reference table:
the energies are \(>.8348762,.87034495,.90327755\).
If there is no hole they are stronger.
The three conclusions therefore include every omission location,
all retained isolated feet, and arbitrary finite spectators.
They are the shared component estimates; no historical
ambient-dimension consumer has been imported.

\section{The actual whole-\(A_{17}\) endpoint reduction}
\label{app:wholearithA17}

We now prove the geometric and arithmetic inclusion needed before
using the recorded finite \(A_{17}\) propositions.
The old incident diagram is WHOLE \(A_{17}\), so \(N=18\) and
there is no spectator.
Keep the original \(p,z,x\) throughout.
Use the \(N=18\) reference and simultaneous-foot argument of
Appendix~\ref{app:arithmeticchain}, with energy
\(E_0=.9037838\), central coefficients at least \(.018\),
and first/last five-coordinate signed means
\(.2951<b<.2952\).
The next signed coordinate is below \(.1\).
A possible single golden central hole is omitted before its
pressure is estimated.

\subsection{Actual endpoint lengths and the full finite-white list}
\label{app:arithA17endpointlist}

Even with a golden end hole, the chord reference has energy
\[
 E_{\min}={13079276\over15000000}>.8719 .
\]
Each signed end block has length \(L=5\) or \(6\) and mean
\(.25<b<.3\).
If its actual endpoint coordinate were at least \(.63\), replace
the block by
\[
 .63,\quad (Lb-.63)/(L-1),\ldots,(Lb-.63)/(L-1).
\]
Its sum and every available central height are unchanged.
The remaining coordinates exceed \(.155\), and the following
coordinate is \(<.1\); hence the changed positive central
coefficient is \(>.055/\sqrt2>.018\).
All other positive central coefficients still meet \(.018\).
Thus the SAME summed-foot proof applies to this ONE new
reference, whose energy increases by
\[
 {L\over L-1}(.63-b)^2>{6\over5}(.33)^2=.13068.
\]
This exceeds one after adding \(E_{\min}\).
Therefore both signed actual endpoint coordinates are \(<.63\).

For either actual endpoint edge,
\[
 x_{\rm end}\ge\sqrt{17/18}
                        \sqrt{(C-1)/(C+1)} .
\]
The bound \(C\ge5/2\) would give endpoint square
\(\ge17/42>.63^2\), impossible.
If the ACTUAL field is nongolden, no central golden hole exists.
Use the full reference and graft instead to \(.575=23/40\).
The remaining block coordinates exceed \(.225125\);
all changed central coefficients still exceed \(.018\).
The energy gain is
\((5/4)(.575-.2952)^2=.09786005\), giving total
\(>1.00164385\).
Since \(C\ge21/10\) would give endpoint square
\(\ge187/558>(23/40)^2\), we obtain
\[
 \text{always }C<5/2,\qquad
 \text{outside }\mathbb Q(\varphi): C<21/10 .
\tag{A17C}\label{eq:arithA17C}
\]

Release old node one.
The retained white path is \(A_{16}\) on nodes two through
seventeen.  Finite one-node classification gives isolation;
helper three at \(j=2,3,16,17\); or helper four at \(j=2,17\).
A second contact would make a cycle; no exceptional finite
rank-seventeen diagram contains that path.
There is no foreign contact in this whole-diagram case.
Parallel own always has its actual future-null pair.

For helper three put \(t=2\gamma\).
Connected simple-three normalization gives a square-two cap,
helper product \(-1\), own product \(-t\), and integral \(t\).
Its full Lorentz star is defined over \(\mathbb Q(t)\);
all-place definiteness gives ACTUAL \(K=\mathbb Q(t)\).
The unit time square and edge coshes are exactly
\[
 A^2={17t^2+2(18-j)t+j(18-j)-36\over36},
\qquad
\begin{array}{c|c}
j&C\\ \hline
2&(17\gamma+8)/9\\
3&(17\gamma+15/2)/(3\sqrt2)\\
16&(17\gamma+1)/(3\sqrt2)\\
17&(17t+1)/18 .
\end{array}
\]
These are obtained by clearing ALL retained \(A_{16}\) roots.
For helper four the two formulas are
\[
 j=2:\ C=(17\gamma+8\sqrt2)/3,\qquad
 j=17:\ C=(17\gamma+1/\sqrt2)/3 .
\]
The raw white residuals are \(18/17\) at a helper-three
endpoint, \(4/17\) at a next-to-end three, and \(1/17\)
at a helper-four endpoint.

At helper \(j=3\), every positive own violates \(C<5/2\).
Own zero has the ACTUAL feasible affine-\(E_8\) null
\[
 2n_1+4n_2+6n_3+5n_4+4n_5+3n_6+2n_7+n_8+3c .
\]
Its future time is \(3A>0\), and all omitted products are
nonpositive.
At \(j=16\), own zero has \(A^2=-1/9\), own three has
the actual all-three cycle, and own at least four violates
\(C<5/2\).
The helper-four rows fail already at own zero for \(j=2\)
and own at least three for \(j=17\); the latter own-zero
time is \(-1/18\).
At \(j=2\), own zero has negative time and
\(C<5/2\) gives \(\gamma<29/34\), leaving owns three,
four or five.
Own three is an actual affine triangle.
Own four generates \(\mathbb Q(\sqrt2)\), but its edge exceeds
\(21/10\), contrary to the nongolden bound.
Only own five in the ACTUAL golden field can remain at this end.

\subsection{The far-helper strip and the golden exceptions}
\label{app:arithA17goldenends}

For the far helper \(j=17\), positive time gives \(t>1\).
Outside the golden field, \eqref{eq:arithA17C} gives
\[
 1<t<184/85,\qquad -19/17<t_\sigma<1\quad(\sigma\ne\mathrm{id}).
\]
Set \(f(s)=s(s^2-1)\).
On this nonphysical interval
\(\lvert f(s)\rvert\le2/(3\sqrt3)\):
on \([-1,1]\) use
\(4/27-q(1-q)^2=(1-3q)^2(4-3q)/27\);
on \([-19/17,-1]\) the bound \(1368/4913\) is smaller.
Physically \(0<f(t)<8\).
Degree at least four would give nonzero integer norm of
magnitude \(<8[2/(3\sqrt3)]^3<1\).
At degree three the magnitude is \(<32/27<2\), so
\(t,t-1,t+1\) would all be units.
The signs at one and minus one leave exactly
\[
 X^3-X-1,\quad X^3-2X^2-X+1,\quad
 X^3+X^2-2X-1,\quad X^3-X^2-2X+1 .
\]
The first has discriminant \(-23\), so is not totally real.
The second is negative at \(11/5\), with increasing positive
branch there, so its physical root exceeds \(11/5>184/85\).
The latter two have positive values \(5347/4913,2899/4913\)
at \(-19/17\), contrary to all three roots lying above that
endpoint.

In degree two the integral trace is \(0,1,2,3\), and the nonzero
norm is \(\pm1,\pm2\).
For trace zero the negative conjugate strip forces
\(1<t<19/17\), so \(-t^2\) lies strictly between \(-2,-1\).
For trace one the norm \(t(1-t)\) leaves only \(-1\), giving
\(t=\varphi\), or \(-2\), giving the rational \(t=2\).
For trace two the norm \(t(2-t)\) is strictly between \(-1,1\),
with zero excluded in degree two.
For trace three the conjugate strip requires \(t>2\), and
\(t(3-t)\) lies strictly between one and two.
Thus no nongolden quadratic survives.
In degree one only \(t=2\) is possible, the forbidden parallel.
Hence any far helper forces the actual golden field.

In \(K=\mathbb Q(\varphi)\), the unconditional bound gives
\(1<t<44/17\) and the same nonphysical strip.
Writing \(t=a+b\varphi\) in the FULL ring yields \(0<b<2\);
thus \(b=1\), \(a=0\) or one, and \(a=1\) violates the
physical bound.  Again \(t=\varphi\).
Every possible nonisolated endpoint is therefore a near- or
far-simple-three helper with own five in \(K=\mathbb Q(\varphi)\).

We now exclude that actual field in full, retaining its isolated
alternatives.  The unit near/far helper times are \(M_N/36,M_F/36\),
where
\[
 M_N=13+49\varphi,\quad \mathrm N(M_N)=-1595,\qquad
 M_F=-2+19\varphi,\quad \mathrm N(M_F)=-395 .
\]
Two far helpers have positive actual product \((1+4\varphi)/9\).
Near/far time-norm ratio \(319/79\) is not a rational square.
Two near helpers have invariant two-cycle
\((500+660\varphi)/81\notin\mathbb Z[\varphi]\).
Thus both endpoints cannot be nonisolated.

For an isolated endpoint put \(U=4\gamma^2\).
Rescale its actual \(K\)-root line to raw square \(2U\) and
sole old product \(-U\).
The actual two-cycle makes \(U\) integral and the full Lorentz
star gives \(K=\mathbb Q(U)\).
Its raw time square is
\[
 H^2=U(17U-36)/18,\quad
 36/17<U<225/17,\quad0<U'<36/17 .
\]
The FULL ring models \(U=a+b\varphi\) are exactly
\[
\begin{array}{c|r|r|r}
U&\mathrm N(U)&-\mathrm N[U(17U-36)]&
 \text{time norm squareclass}\\ \hline
1+\varphi&1&251&251\\
2+\varphi&5&1595&5\cdot11\cdot29\\
2+2\varphi&4&4880&5\cdot61\\
3+2\varphi&11&4631&11\cdot421\\
2+3\varphi&1&2699&2699\\
3+3\varphi&9&14499&179\\
3+4\varphi&5&16895&5\cdot31\cdot109\\
4+4\varphi&16&22784&89\\
4+5\varphi&11&38291&11\\
5+5\varphi&25&16475&659 .
\end{array}
\]
Indeed \(b\sqrt5<225/17\) gives \(b=1,\ldots,5\);
the strict conjugate strip gives respectively
\(a=1,2;\ 2,3;\ 2,3;\ 3,4;\ 4,5\).
The entries follow from
\[
 \mathrm N[U(17U-36)]
 =n(289n-612T+1296),\quad n=\mathrm N(U),\quad T=\mathrm{Tr}(U).
\]
The displayed squareclasses are distinct; their small prime
factorizations can be checked by trial divisibility through the
respective square roots.

Actual roots on the SAME \(K\)-time line require
\(U(17U-36)/M_N\) or \(/M_F\) to be a square IN \(K\).
Their rational norms exclude every model except \(U=2+\varphi\)
with a near helper; none has the far-helper class \(5\cdot79\).
For that exception \(U(17U-36)=M_N\) exactly.
Its equal positive raw time coefficients and spatial product
\(U^2/18\) give
\[
 g=-(4+22\varphi)/9,\qquad
 g^2/U=(168+164\varphi)/81,
\]
whose nonintegral two-cycle excludes it.

Finally two isolated golden endpoints cannot use different rows,
since their rational time-norm squareclasses differ.
With the same row their equal positive raw times give
\[
 g=2U(9-4U)/9,\qquad g^2/U^2=4(9-4U)^2/81\in\mathcal O_K .
\]
Integral closedness makes \(2(9-4U)/9\) integral, hence
\(8U/9\) integral; Bezout \(8\cdot8-7\cdot9=1\)
then makes \(U/9\) integral.
But \(0<\mathrm N(U)<8100/289<29\), whereas
\(\mathrm N(U/9)=\mathrm N(U)/81\) would be a positive integer.
Thus the actual golden field is impossible.

Both endpoints are now isolated in a nongolden actual field,
with independent integral generators satisfying
\[
 K=\mathbb Q(U)=\mathbb Q(V),\quad
 36/17<U,V<3969/425,\quad
 0<U^\sigma,V^\sigma<36/17\quad(\sigma\ne\mathrm{id}).
\tag{A17isolated}\label{eq:arithA17isolated}
\]
No equality of their parameters, finite-own assumption or common
disconnected raw scale was used.

\subsection{One two-end graft gives actual degree at most five}
\label{app:arithA17degree}

There is no golden hole after the preceding argument.
If BOTH signed endpoint coordinates were at least \(1/2\),
modify the first and last five-coordinate blocks simultaneously.
Replace each signed block mean \(b\) by
\[
             a,r,r,r,r,\quad a=1/2,\quad r=(5b-a)/4 .
\]
Both sums and every central height stay unchanged.
The supports are disjoint and each residual has sum zero,
so the two residuals are orthogonal to one another and to the
old flat blocks.
Moreover \(r>.243875>.1\), making every changed positive
central coefficient \(>.018\).
The new endpoint coefficients are paid by the two ACTUAL
coordinates assumed at least \(a\).
The summed-foot proof therefore applies to this ONE new reference,
with energy
\[
 >.9037838+2(5/4)(.5-.2952)^2
 =1.0086414 .
\]
Thus at least one actual endpoint coordinate is \(<1/2\).
For its generator \(U\), the edge formula
\(C^2=17U/36\) and the farthest endpoint comparison give
\[
                  36/17<U<66564/10625<157/25.
\tag{A17shortU}\label{eq:arithA17shortU}
\]
The other generator need not satisfy this sharper upper bound;
the chosen one still generates the ENTIRE actual field.

Let \(F(U)=U(U-1)(U-2)\) and \(d=[K:\mathbb Q]\).
For \(d\ge2\) all factors are nonzero algebraic integers.
Physically
\[
                    0<F(U)<142 .
\]
At any other embedding put \(t=U^\sigma-1\in(-1,19/17)\).
On \([-1,1]\), \(|t^3-t|\le2/(3\sqrt3)<2/5\);
on \((1,19/17)\) it is at most \(1368/4913<1/3\).
If \(d\ge7\), the nonzero integral norm would have magnitude
\(<142(2/5)^6=9088/15625<1\).
Hence \(d\le6\).
For \(d=6\), the bound \(142(2/5)^5<2\) makes its absolute
norm one, and all three separate factor norms have absolute
value one.

We give the complete degree-six seal.
Set \(s=U-1\), and let
\(r_i=|U^{\sigma_i}-1|\), \(i=1,\ldots,5\).
Irreducibility excludes \(r_i=0,1\).
The unit norms and~\eqref{eq:arithA17shortU} give
\[
 0<r_i<19/17<9/8,\quad
 \prod r_i=1/s>3/16,\quad s<132/25,\quad
 s^2-1<16799/625<27.
\]
The sharper \(s<132/25\), rather than merely \(s<16/3\),
is needed in the last inequality.

If \(m\) positive numbers \(q_i<1\) have
\(p=(\prod q_i)^{2/m}\), ordinary AM--GM applied separately
to \(q_i^2\) and \(1-q_i^2\) gives
\[
       p+\{\prod(1-q_i^2)\}^{1/m}\le1 .
\]
Let \(k\) be the number of \(r_i>1\), so \(m=5-k\).
Each large \(r_i\) contributes
\(|1-r_i^2|<72/289<1/4\).
The following six cases exhaust all possibilities:
\[
\begin{array}{c|c|c|c}
k&\prod_{\rm small}r_i\text{ lower}&
 p\text{ lower}&\prod_{i=1}^5|1-r_i^2|\text{ upper}\\ \hline
0&3/16&1/2&1/32\\
1&1/6&2/5&81/2500\\
2&4/27&1/4&27/1024\\
3&32/243&1/8&49/4096\\
4&\text{not needed}&\text{not needed}&1/256\\
5&\text{empty product}&\text{not needed}&1/1024 .
\end{array}
\]
For the product bounds use
\(\prod_{\rm small}r_i>(3/16)(8/9)^k\).
The four root comparisons for \(p\) follow respectively from
\(9/256>1/32\), \(1/6>4/25\), \(16/729>1/64\),
\(32/243>1/8\); each exponent has been cleared positively.
Multiplying the last column by \(27\) gives
\[
 27/32,\quad2187/2500,\quad729/1024,\quad
 1323/4096,\quad27/256,\quad27/1024,
\]
all below one.
But the unit norm of \(U(U-2)\) would be
\[
             1=(s^2-1)\prod_{i=1}^5|1-r_i^2|<1 .
\]
This excludes degree six and proves the ACTUAL degree bound
\(d\le5\), without a field list or an equality assumption on the
two endpoints.

\subsection{The whole two-dotted interval and the derived finite fields}
\label{app:arithA17anchors}

Suppose both isolated owns are dotted.
For the left endpoint put \(u=\gamma^{-2}\in(0,1)\).
The positive actual core
\[
 c+\gamma^{-1}n_1
   =\gamma\{\sqrt{17/9-u}\,p-(\omega_1-un_1)\}
\]
has negative square \(u-1\), zero \(c\) product and negative
\(n_1\) product.
Every other original wall, INCLUDING the other endpoint opposite,
has nonpositive product with both summands.
After normalization it is a genuine chamber point and gives
\[
 x\cdot Z_L\ge h(u)=\sqrt{17/9-u}-\sqrt{u(1-u)},\quad
 Z_L=\omega_1-un_1 .
\]
The right endpoint has its independent parameter \(v\in(0,1)\)
and reflected vector \(Z_R\).

Let
\[
 Y_L=Z_L-\omega_5/5,\qquad
 R(u)=\|Y_L\|^2=8/5-2u+u^2 .
\]
This is a zero-sum vector supported in the first five coordinates.
Its central-root products are zero except \(-1/5\) at five.
Its reflected partner is supported in the LAST five coordinates.
Consequently \(w\cdot Y_L=w\cdot Y_R=Y_L\cdot Y_R=0\), while
\(w\cdot Z_L=w\cdot Z_R=.4174\).
Pay both actual core pressures at \(\mu=.3\):
\[
 w'=w+\mu Y_L+\mu Y_R,\quad
 \beta'_5=\beta'_{13}=.1382-.06=.0782,\quad
 \beta'_i=\beta_i\ (6\le i\le12).
\]
Every positive central coefficient is still at least \(.018\),
and \(w'\cdot n_i=\beta'_i\) EXACTLY.

Retain every central isolated foot, with its OWN parameter
\(v_i\in(0,1)\), using ONE vector
\[
                       V=w'-\sum_{i\in S}\beta'_iv_in_i .
\]
Sum the endpoint and central actual pressures before estimating
the norm.  The identities just proved and
\(n_i\cdot n_j\le0\) give
\[
 \begin{split}
 x\cdot V&\ge E_0+
 \mu\{h(u)-.4174+h(v)-.4174\}
       -\sum_S\beta'_i(\lambda_i-.6),\\
 \|V\|^2&\le E_0+\mu^2\{R(u)+R(v)\}
       -\sum_S(\beta'_i)^2(2v_i-v_i^2).
 \end{split}
\]
Every central summand is nonnegative by the already proved
absorption inequality, since \(\beta'_i\ge.018\).
Thus
\[
 \|x\|^2\ge E_0+G(u)+G(v),\qquad
 G(u)=.6\{h(u)-.4174\}-.09(8/5-2u+u^2).
\]

This whole continuous interval has a strict rational seal.
Concavity and the endpoint squares give
\(\sqrt{17/9-u}>1.3743-.4315u\) on \([0,1]\).
Therefore
\[
 G(u)>.43014-.0789u-.09u^2-.6\sqrt{u(1-u)}.
\]
For \(u\le1/2\), use \(u^2\le u/2\) and the circle inequality
\[
 au+b\sqrt{u(1-u)}\le(a+\sqrt{a^2+b^2})/2 .
\]
With \(a=.1239,b=.6\), \(\sqrt{a^2+b^2}<.6127\)
gives \(G(u)>.06184\).
For \(u\ge1/2\), use \(u^2\le1.5u-.5\);
\(a=.2139,b=.6\), with radical \(<.637\), gives
\(G(u)>.04969\).
The positive square margins are \(.00005008,.00001579\),
respectively.
Hence throughout \((0,1)\),
\[
          \|x\|^2>.9037838+2(.04969)=1.0031638,
\]
impossible.  Parallel owns already have their actual future null.
At least one endpoint is a finite native own.

At that endpoint \(U=4\cos^2(\pi/m)=2+2\cos(2\pi/m)\)
generates the ACTUAL field.
The degree bound makes \(\phi(m)\le10\).
The COMPLETE elementary totient list is
\[
 3,4,5,6,7,8,9,10,11,12,14,15,16,18,20,22,24,30 .
\]
Indeed only prime divisors \(2,3,5,7,11\) are possible;
their permissible prime powers and products give that list.
Labels three and four have \(U\le2<36/17\).
Six gives the rational field, excluded by Meyer and compact
anisotropy at form rank eighteen.
Five and ten give the already excluded golden field.
For the following nine labels choose the indicated coprime
cyclotomic embedding:
\[
\begin{array}{c|c|c|c}
m&k&\cos(2k\pi/m)=\sin\theta&2\theta/\pi\\ \hline
9&2&\theta=\pi/18&1/9\\
14&3&\theta=\pi/14&1/7\\
15&2&\theta=7\pi/30&7/15\\
16&3&\theta=\pi/8&1/4\\
20&3&\theta=\pi/5&2/5\\
24&5&\theta=\pi/12&1/6\\
30&7&\theta=\pi/30&1/15\\
11&2&\theta=3\pi/22&3/11\\
22&3&\theta=5\pi/22&5/11 .
\end{array}
\]
Every last entry exceeds \(1/17\).
Strict concavity on \(0<\theta<\pi/2\) gives
\(\sin\theta>2\theta/\pi\), so that conjugate would have
\(U^\sigma>36/17\), a contradiction.
Label thirty was included before its conjugate test.
Only labels \(7,8,12,18\) remain.
Their actual fields are, respectively,
\(K_7,\mathbb Q(\sqrt2),\mathbb Q(\sqrt3),K_9\);
the cubic parameters have polynomials
\[
 X^3-5X^2+6X-1,\qquad X^3-6X^2+9X-3,
\]
with discriminants \(49,81\) and irreducibility by the
rational-root test.  This proves the actual field reduction,
rather than assuming a list of possible fields.

\subsection{Actual inclusion in the recorded finite domains}
\label{app:arithA17datainclusion}

The isolated endpoint normalization in
\eqref{eq:arithA17isolated} has raw square \(2U\), sole old
product \(-U\) and time square \(F(U)/18\), where now
\(F(X)=X(17X-36)\).
For the two actual endpoint roots, subtracting their \(K\)-valued
old projections leaves one common \(K\)-time line.
Thus
\[
 R=F(U)F(V)\in K^2,\quad
 g={UV-\sqrt R\over18},\quad
 L={g^2\over UV}.
\]
The square root is PHYSICALLY POSITIVE.
The long old-three cycle gives \(g\in\mathcal O_K\);
the actual two-cycle gives \(L\in\mathcal O_K\).
Acuteness gives \(g\le0\).
If \(g\ne0\), every nonphysical place satisfies \(0<L^\sigma<4\);
the branch \(g=L=0\) is retained.
The physical owns and mutual products retain their native/dotted
alphabet, with a parallel pair excluded by its actual future null.
These are precisely the stated necessary tests in the quadratic
domain; the endpoint parameters need not be equal.

If \(K\) is quadratic, this actual pair therefore enters the FULL
quadratic domain of Proposition~\ref{prop:pair-output}(i),
including its maximal-order classes, square signs and zero mutual.
That recorded finite-output proposition says the domain is empty.
If \(K=K_7\) or \(K_9\), the finite endpoint proved above supplies
the actual anchor seven or eighteen, respectively.
The same original time-line square, cyclic and native conditions
put the other endpoint in the full cubic finite-anchor domain
of Proposition~\ref{prop:pair-output}(ii).
Its complete ring/box/decision arguments are given in
Section~\ref{sec:algorithms}, and the recorded output is empty.

Thus, at those explicitly stated EXISTING finite-output premises,
the whole \(A_{17}\) at the original arithmetic rank-seventeen
vertex is impossible.
The geometric domain inclusion and complete hand arguments are
proved in this appendix; no new enumeration, independent replay,
runtime PASS, higher-dimensional arithmetic tail or geometric
Lean formalization is inferred by transcribing them.

\section{The separate whole-\(A_{16}\) endpoint proof}
\label{app:wholearithA16}

The whole \(A_{16}\) case has \(N=17\) and its own constants.
It is not a consequence of reranking the preceding \(A_{17}\)
theorem.  We retain one original farthest vector throughout.
The common chain proof gives central reference energy
\[
 E_0=.8708972,\quad
 (H_5,H_6,H_7,H_8)=(2.056,2.186,2.269,2.310),\quad H_9=H_8,
\]
with positive coefficients at least \(.0205\).
Its worst single end-hole chord has energy
\(12569249/15000000\).
All grafts below preserve the common simultaneous-foot proof
by checking those coefficients afresh.

\subsection{Actual lengths and all nonisolated endpoint roles}
\label{app:arithA16endpointlist}

On either signed end block, of length five or six and mean
\(.25<b<.3\), grafting to \(.67\) leaves each remaining
coordinate \(>.145\); the following coordinate is below \(.1\).
The central coefficient stays \(>.045/\sqrt2>.0205\).
The gain \(>(6/5)(.37)^2=.16428\), even after the worst
chord omission, exceeds the available norm budget.
Both signed endpoint coordinates are therefore \(<.67\).
The actual edge comparison is now
\[
 x_{\rm end}\ge\sqrt{16/17}\sqrt{(C-1)/(C+1)}.
\]
At \(C=17/6\) its square \(176/391\) exceeds \(.67^2\).
Hence both endpoints have \(C<17/6\).
If both signed coordinates were at least \(.56\), disjoint
grafts on the SAME reference would gain
\(>2(6/5)(.26)^2=.16224\), again exceeding one.
Since \(16/51>.56^2\), at least one endpoint has \(C<2\).

Outside the actual golden field there is no hole and the
full block mean lies in \((.2907,.2908)\).
A single graft to \(.613\) gains \(>.12976605\), making total
\(>1.00066325\).
The retained coordinates are \(>.210125\), so the coefficient
threshold remains valid.
As \(32/85>.613^2\), both endpoints then have \(C<7/3\).
Two disjoint grafts to \(.52\) gain \(>.1313316\), giving total
\(>1.0022288\).
The exact comparison
\(1296/4777>(13/25)^2\) therefore proves
\[
 \text{always both }C<17/6,\ \min C<2;\qquad
 K\ne\mathbb Q(\varphi):
        \text{ both }C<7/3,\ \min C<181/100 .
\tag{A16lengths}\label{eq:arithA16lengths}
\]

Release node one and retain the whole white \(A_{15}\).
Its complete finite one-node extensions are isolation;
helper three at \(j=2,3,15,16\); or helper four at \(j=2,16\).
There is no foreign component or exceptional rank-sixteen
completion.  A second contact makes a cycle.
For helper three, the actual connected normalization gives
integral \(t=2\gamma\), ACTUAL \(K=\mathbb Q(t)\), raw
time square \(M_j/17\), and
\[
 M_j=16t^2+2(17-j)t+j(17-j)-34,\qquad
\begin{array}{c|c}
j&C\\ \hline
2&(16t+15)/17\\
3&(16\gamma+7)/\sqrt{17}\\
15&(16\gamma+1)/\sqrt{17}\\
16&(16t+1)/17 .
\end{array}
\]
For helper four the lengths are
\[
 (16\gamma+15/\sqrt2)/\sqrt{17/2},\qquad
 (16\gamma+1/\sqrt2)/\sqrt{17/2}.
\]
These follow by eliminating ALL retained roots; the unit white
residuals are \(17/32,1/8,1/16\) in the endpoint-three,
next-three and endpoint-four rows.

Helper \(j=3\) with positive own gives
\(C\ge15/\sqrt{17}>17/6\).
Its own-zero case has the same actual affine-\(E_8\) vector
printed in the preceding endpoint proof.
At \(j=15\), own zero has unit time square \(-2/17\), own three
has the actual all-three cycle, and owns at least four give
\(C>17/6\).
Helper four at \(j=2\) fails even at own zero.
At \(j=16\), its own-zero time is \(-1/17\), and every positive
own gives \(C\ge(8\sqrt2+1)/\sqrt{17}>17/6\).
At near \(j=2\), own zero has negative time and own three
the actual affine triangle.
At far \(j=16\), positive time requires \(t>1\), since
\(M_{16}=(t-1)(16t+18)\).
Parallel has the positive actual future-null pair.
All these nulls use UNIT actual normals and are feasible
against every omitted wall.

In a nongolden field the near helper has \(t<37/24\),
so its only remaining native own is four, \(t=\sqrt2\).
The far helper satisfies
\[
                    1<t<29/12,\qquad -9/8<t_\sigma<1.
\]
As before \(f(t)=t(t^2-1)\) has other-place magnitude at most
\(2/(3\sqrt3)\); on the additional interval
\([-9/8,-1]\), the maximum is \(153/512\), smaller.
Physically \(0<f(t)<12\).
The integral norm excludes degrees at least four.
At degree three \(t,t-1,t+1\) would all be units.
The same complete four cubics as before occur.
The negative-discriminant cubic and the two with a conjugate
below \(-9/8\) fail; the only survivor is
\[
                P(X)=X^3-2X^2-X+1,
\]
whose physical root \(a>2\) generates \(K_7\).
For degree two, trace \(0,1,2,3\) and norm
\(-2,-1,1,2\) exhaust the strip.
Trace zero would require \(\sqrt2\), with conjugate below
\(-9/8\); trace one leaves only \(\varphi\) or rational two;
trace two leaves \(1+\sqrt2\), since \(1+\sqrt3\) is too large;
trace three's irrational positive root exceeds \(29/12\).
The rational two is parallel.
Thus the nongolden helper roles are exactly near silver,
far silver, or far \(K_7\).
Their actual lengths are respectively
\[
 (16\sqrt2+15)/17,\quad (17+16\sqrt2)/17,\quad
 (16a+1)/17,
\]
all \(>181/100\).
The short endpoint selected in~\eqref{eq:arithA16lengths}
is consequently ISOLATED.

\subsection{The full golden endpoint table}
\label{app:arithA16golden}

For any isolated endpoint, its OWN \(K\)-rescaling gives
raw norm \(2U\), sole old product \(-U\), integral \(U=4\gamma^2\)
generating ACTUAL \(K\), and
\[
 F(U)=U(16U-34),\quad H_U^2=F(U)/17,\quad C^2=8U/17,
 \quad U>17/8,\quad0<U^\sigma<17/8 .
\tag{A16U}\label{eq:arithA16U}
\]
No raw time itself is assumed \(K\)-valued; the ratios along
the SAME one-dimensional \(K\)-time line are \(K\)-valued.

Suppose \(K=\mathbb Q(\varphi)\).
The near helper has \(t<199/96\) and conjugate
\(-2<t'<1/8\).  The finite native list leaves \(\varphi\);
a dotted full-ring \(a+b\varphi\) has \(b=1\) and no integer
\(a\) in \(2<t<199/96\).
The far helper satisfies \(1<t<283/96\), \(-9/8<t'<1\),
giving \(t=\varphi\) or \(\varphi^2\).
All three types have raw times \(M/17\):
\[
\begin{array}{c|c|r}
\text{helper}&M&\mathrm N(M)\\ \hline
\text{near }\varphi&12+46\varphi&-1420\\
\text{far }\varphi&-2+18\varphi&-356\\
\text{far }\varphi^2&16+50\varphi&-1444 .
\end{array}
\]
The norm squareclasses \(355,89,1\) forbid different types.
Reflected near copies have nonintegral actual long-cycle
\(g=-(7+41\varphi)/17\).
Reflected far-\(\varphi\) copies have positive product
\(g=(4+15\varphi)/17\).
Reflected far-\(\varphi^2\) copies have
\(g=(19-15\varphi)/17\in(-1,0)\), forbidden for two
distinct square-two actual roots by the native gap.
Thus both endpoints cannot be nonisolated.

For the isolated golden end, the universal length cap gives
\[
                       17/8<U<4913/288,\qquad0<U'<17/8 .
\]
Writing \(U=a+b\varphi\) gives \(b=1,\ldots,7\).
The complete conjugate intervals give
\(a=1,2;\ 2,3;\ 2,3;\ 3,4;\ 4,5;\ 4,5;\ 5,6\),
and the physical upper bound removes only \((6,7)\).
The complete thirteen rows are
\[
\begin{array}{c|r|r|r}
U&\mathrm N(U)&-\mathrm N(F(U))&\text{squareclass}\\ \hline
1+\varphi&1&220&55\\
2+\varphi&5&1420&355\\
2+2\varphi&4&4336&271\\
3+2\varphi&11&4180&1045\\
2+3\varphi&1&2396&599\\
3+3\varphi&9&12924&359\\
3+4\varphi&5&15020&3755\\
4+4\varphi&16&20416&319\\
4+5\varphi&11&34100&341\\
5+5\varphi&25&15100&151\\
4+6\varphi&4&21744&151\\
5+6\varphi&19&50996&12749\\
5+7\varphi&11&58036&14509 .
\end{array}
\]
For \(n=\mathrm N(U),T=\mathrm{Tr}(U)\), every entry follows from
\(\mathrm N(F(U))=n(256n-544T+1156)\).
The listed classes are squarefree:
\(1045=5\cdot11\cdot19\), \(3755=5\cdot751\),
\(319=11\cdot29\), \(341=11\cdot31\),
\(12749=11\cdot19\cdot61\), \(14509=11\cdot1319\),
\(355=5\cdot71\), \(55=5\cdot11\);
the remaining small primality checks are direct trial division.
The only repeated class is151.
Precisely the FIRST SIX rows have \(U<17/2\), equivalently \(C<2\),
and their classes occur nowhere else in this complete table.

If one endpoint is near-\(\varphi\) or far-\(\varphi^2\),
its \(C>2\), so the opposite isolated end is in the first six rows.
Class one matches none.
The near class355 matches only \(U=2+\varphi\);
there \(F(U)=12+46\varphi\), so the actual positive raw
times are equal.
Their spatial product is \(U(\varphi+2)/17=U^2/17\),
giving the same nonintegral long cycle
\(g=-(7+41\varphi)/17\).
For a far-\(\varphi\) helper the opposite end may be ANY of
the thirteen rows, but none has class89.
All helper/isolated possibilities are excluded.

Both golden endpoints must now be isolated.
One belongs to the first six rows by the selected \(C<2\).
The SAME time-line square test forces equality of the two
\(U\)-models; the repeated151 class cannot occur at the selected end.
Equal positive raw times give
\[
 g=U(34-15U)/17,\qquad
 L=g^2/U^2=((34-15U)/17)^2\in\mathcal O_K .
\]
Integral closedness makes \((34-15U)/17\), and then \(U/17\),
integral, using Bezout for15 and17.
But the selected norms \(1,5,4,11,1,9\) are smaller than289,
so \(\mathrm N(U/17)=\mathrm N(U)/289\) cannot be a
nonzero integer.
The ACTUAL golden field is impossible.

\subsection{Both actual endpoints are isolated}
\label{app:arithA16helpers}

The actual field is now nongolden, and the selected endpoint
of \(C<181/100\) is isolated.
At \(K=\mathbb Q(\sqrt2)\), its full possible isolated strip
from \(C<7/3\) is
\[
 17/8<U<833/72,\quad0<U'<17/8 .
\]
Writing \(U=A+B\sqrt2\), the complete models are
\[
 2+s,\ 3+s,\ 3+2s,\ 4+2s,\ 5+3s,\ 6+3s,\qquad s=\sqrt2 .
\]
Indeed \(B=1,\ldots,4\);
the conjugate strip permits respectively
\(A=2,3;\ 3,4;\ 5,6;\ 6,7\), and the physical upper
bound removes both \(B=4\) possibilities.
Their signed \(F\)-norms are
\[
                   -1016,-2212,-1852,-9184,-17444,-13752 .
\]
The near-silver helper has \(M=28+30s\), norm
\(-1016=-8\cdot127\).
Only \(U=2+s\) has a compatible norm squareclass:
none of the other five norms is divisible by127.
For that exception \(F(U)=M\) and the actual long cycle
\(g=-(22+26s)/17\) is not integral.
The far-silver helper has \(M=32+34s\), norm
\(-1288=-8\cdot7\cdot23\).
None of the six time norms is divisible by23.
Both silver helper roles are excluded.

For the remaining far helper in \(K_7\), let \(a>2\) be
the root of \(P(X)=X^3-2X^2-X+1\).
Its three roots have disjoint rational brackets
\[
 2.246<a<2.247,\quad -.803<b<-.801,\quad .554<c<.556 .
\]
Signs and derivatives of \(P\) isolate them.
The full ring is \(\mathbb Z[a]\), as justified below.
The selected isolated endpoint has
\(17/8<U<7\), \(0<U_b,U_c<17/8\).
Write \(U=A+Ba+Ca^2\).
Lagrange interpolation expresses \(C\) as the sum of the
three values divided by the derivatives at the three roots.
The two positive denominators exceed five and four;
the magnitude of the negative denominator exceeds \(17/8\).
Consequently \(-1<C<7/5+(17/8)/4<2\), so \(C=0,1\).
The small conjugate difference gives
\[
                 |B+C(2-a)|<(17/8)/(c-b)<8/5 .
\]
If \(C=0\), positive physical difference forces \(B=1\),
and the conjugate interval forces \(A=1\).
If \(C=1\), the physical/conjugate difference \(<7\)
removes \(B\ge1\), leaving \(-1,0\).
For \(B=-1\), the \(c\)-value requires \(A\ge1\), while
the \(b\)-value requires \(A\le0\).
For \(B=0\), \(A=0,1\).
Thus the COMPLETE selected models are
\[
                               U=1+a,\quad a^2,\quad1+a^2 .
\]
Their positive time norms are \(5048,32248,82264\).
The helper time numerator is
\((a-1)(16a+18)\), of norm \(7496=8\cdot937\).
The prime937 divides none of these three norms;
therefore no actual time ratio can be a \(K\)-square.
This excludes the last nongolden helper.
BOTH actual endpoints are isolated.

\subsection{Degree, the whole two-dotted domain and finite anchors}
\label{app:arithA16anchors}

The selected isolated endpoint generates the actual field, with
\[
 17/8<U<556937/80000<7,\qquad0<U^\sigma<17/8 .
\]
For \(F_3(U)=U(U-1)(U-2)\), physical value \(0<F_3(U)<210\).
At a nonphysical embedding \(y=U^\sigma-1\in(-1,9/8)\).
The same cubic bound gives
\(|y^3-y|<2/5\), including the upper extra interval where
the maximum is \(153/512\).
A degree at least seven would give a nonzero integer norm
\(<210(2/5)^6=.86016<1\).
Hence ACTUAL degree is at most SIX.
No degree-five refinement from the \(A_{17}\) case is used.

Suppose both isolated owns were dotted.
Their ACTUAL endpoint cores give, for independent \(u,v\in(0,1)\),
\[
 x\cdot Z_L\ge h(u)=\sqrt{32/17-u}-\sqrt{u(1-u)},\qquad
 Z_L=\omega_1-un_1 ,
\]
and its reflection.
The core is future timelike, zero on its cap, negative on the
released old root, and nonpositive against EVERY other actual
wall, including the opposite endpoint.
The residual \(Y_L=Z_L-\omega_5/5\) has zero sum,
support in the first five coordinates, norm
\(R(u)=8/5-2u+u^2\), and sole central product \(-1/5\)
at five.
The reflected residual subtracts \(\omega_{12}/5\), and is
orthogonal to both the first residual and the old flat block.
Here \(w\cdot Z_L=w\cdot Z_R=.4112\).

Pay each core at weight \(\mu=.32\).
The individual outer coefficients change from \(.1406\) to
\(.1406-.064=.0766>.0205\);
every other central coefficient is unchanged.
In particular the changed reference has its EXACT products
equal to these new coefficients.
Subtract all retained central isolated-foot terms with their
own parameters, and apply the common simultaneous absorption
proof to this ONE vector.
All endpoint/central and adjacent-central crosses are retained,
giving
\[
 \|x\|^2\ge E_0+G(u)+G(v),\quad
 G(u)=.64\{h(u)-.4112\}-.1024(1.6-2u+u^2).
\]
The following COMPLETE three-chord certificate covers \([0,1]\):
\[
\begin{array}{c|c|c|c|c}
u\text{ interval}&L(u)<\sqrt{32/17-u}&K&a&
 \sqrt{a^2+.64^2}\text{ upper}\\ \hline
[0,1/2]&1.3719-.3924u&.451008&.097536&.6475\\
[1/2,3/4]&1.3989-.4464u&.506688&.208896&.6733\\
[3/4,1]&1.4385-.4992u&.570432&.293888&.7043 .
\end{array}
\]
The strict endpoint values \(1.3719,1.1757,1.0641,.9393\)
have squares below \(32/17,47/34,77/68,15/17\).
Concavity proves each chord lower bound.
On \([l,r]\) use \(u^2\le(l+r)u-lr\);
then \(G(u)>K-au-.64\sqrt{u(1-u)}\).
The circle inequality used earlier gives the three lower values
\(.07849,.06559,.071338\).
The squared radical margins are
\(.000142978704,.000095351184,.000068333456\).
Thus \(G(u)>.06559\) on the whole dotted interval, and
\[
                  \|x\|^2>.8708972+2(.06559)=1.0020772 .
\]
Both owns cannot be dotted.
Parallel has its actual future null, so at least one endpoint
is a finite native own.

Its integral \(U=2+2\cos(2\pi/m)\) generates the actual field,
and the degree bound gives \(\phi(m)\le12\).
The full totient list is
\[
 3,4,5,6,7,8,9,10,11,12,13,14,15,16,
 18,20,21,22,24,26,28,30,36,42 .
\]
Only prime factors \(2,3,5,7,11,13\) and their allowed prime
powers occur; checking products with totient at most twelve
gives exactly this list.
Labels three and four have \(U\le2<17/8\).
For own six the ACTUAL unit-root null is
\(\sqrt3c+2n_1+n_2\), a nonzero future ray feasible against
all other walls.
Five and ten give the excluded golden field.
The remaining forbidden labels have explicit coprime embeddings:
\[
\begin{array}{c|c|c}
m&k&\cos(2k\pi/m)\text{ lower}\\ \hline
9,11,13,15,21&2&\sin(\pi/18)>1/9\\
14,16,20,22,26,28&3&\sin(\pi/14)>1/7\\
24,36,42&5&\sin(\pi/12)>1/6\\
30&7&\sin(\pi/30)>1/15 .
\end{array}
\]
All exceed \(1/16\), by strict sine concavity.
Their conjugate \(U\) would exceed \(17/8\).
Only finite anchors seven, eight, twelve and eighteen remain,
with ACTUAL fields \(K_7,\mathbb Q(\sqrt2),\mathbb Q(\sqrt3),K_9\).

We use FULL integer rings.
For \(K_7\), \(P(X+3)=X^3+7X^2+14X+7\) is Eisenstein
at seven, and the polynomial discriminant is49.
For \(K_9\), put \(Q(X)=X^3-3X-1\);
\(Q(X+1)=X^3+3X^2-3\) is Eisenstein at three and the
discriminant is81.
The Eisenstein uniformizer criterion proves maximality at the
only possible index primes, while the discriminants exclude
every other index prime.
Equivalently these are the full-ring proofs already given in
Section~\ref{sec:algorithms}, under translation.
Thus the bases are the full \(\mathbb Z[a]\), \(\mathbb Z[k]\);
the two quadratic rings are \(\mathbb Z[\sqrt2]\) and
\(\mathbb Z[\sqrt3]\).

\subsection{All twenty-one actual finite-anchor time norms}
\label{app:arithA16times}

The opposite isolated endpoint \(V\), whether finite or dotted,
has
\[
                    17/8<V<833/72,\qquad0<V^\sigma<17/8 .
\]
For \(\mathbb Q(\sqrt2)\), its complete six models were printed
above.
For \(\mathbb Q(\sqrt3)\), the full ring gives
\(B=1,2,3\), and conjugate intervals
\(A=2,3;\ 4,5;\ 6,7\).
The physical upper bound removes only \((A,B)=(7,3)\).
The EXACT five models are
\[
                    2+r,\ 3+r,\ 4+2r,\ 5+2r,\ 6+3r,
                    \qquad r=\sqrt3.
\]

For \(K_7\), the leading-coefficient interpolation bound now
gives \(-1<C<3\), so \(C=0,1,2\) in \(V=A+Ba+Ca^2\).
The same small conjugate difference gives
\(|B+C(2-a)|<8/5\).
The complete integer possibilities are
\[
\begin{array}{c|c|c}
C&B&A\\ \hline
0&1&1\\
1&-1&\text{none}\\
1&0&0,1\\
1&1&1\\
2&-1,0&0 .
\end{array}
\]
For \(C=2,B\ge1\), the physical/conjugate difference exceeds
\(3.047(1+2(1.444))=11.846736>833/72\).
This excludes every remaining larger \(B\);
the conjugate-difference bound excludes the smaller ones.
The resulting EXACT six models are
\[
                 1+a,\ a^2,\ 1+a^2,\ 1+a+a^2,\ 2a^2-a,\ 2a^2 .
\]

For \(K_9\), the three roots have brackets
\[
 1.879<k<1.880,\quad -.348<b<-.347,\quad -1.533<c<-1.532 .
\]
The two positive Lagrange denominators exceed \(7.5,4\),
and the magnitude of the negative one exceeds \(17/8\).
Thus again \(C=0,1,2\).
The full conjugate difference gives \(|B-Ck|<9/5\).
The constant-term interval cases are
\[
\begin{array}{c|c|c}
C&B&A\\ \hline
0&1&2\\
1&1&1\\
1&2&1,2\\
1&3&3\text{ fails the physical upper bound}\\
2&2&\text{none}.
\end{array}
\]
For \(B=3,C=1\), the physical value is
\(>3+3(1.879)+1.879^2=12.167641>833/72\).
For \(C=2,B\ge3\), the physical difference exceeds
\(3.411(3+.694)>833/72\).
At \(C=2,B=2\), the \(b\)-strip forces \(A\ge1\) and
the \(c\)-strip \(A\le0\).
Hence the EXACT four models are
\[
                     2+k,\ 1+k+k^2,\ 1+2k+k^2,\ 2+2k+k^2 .
\]
Every listed accepted model satisfies all strict strips by the
root brackets; every omitted integer coefficient has been excluded.
These are complete handwritten domains, not unexecuted worker results.

The original raw times of two isolated roots are \(F(U)/17\),
\(F(V)/17\).
The same actual \(K\)-time line makes their ratio a \(K\)-square.
A necessary condition is a square rational norm; BOTH possible
signed norms of any putative square root satisfy that condition.
Here are ALL twenty-one signed time norms, in the stated model order:
\[
\begin{array}{c|l}
\text{field and finite anchor}&\mathrm N(F(V))\\ \hline
\sqrt2,\ 2+\sqrt2&
 -1016,-2212,-1852,-9184,-17444,-13752\\
\sqrt3,\ 2+\sqrt3&
 -764,-3432,-8688,-12428,-27612\\
K_7,\ 1+a&
 5048,32248,82264,110552,1784,330688\\
K_9,\ 2+k&
 16872,66168,79032,293624 .
\end{array}
\]
For a quadratic model the reproducible identity is
\[
 \mathrm N(F(V))=n(256n-544T+1156).
\]
For a cubic model with coefficient data \((T,S,n)\), it is
\[
 \mathrm N(F(V))
      =n(4096n-8704S+18496T-39304).
\]
The six \(K_7\) triples are
\[
 (5,6,1),(6,5,1),(9,20,13),(11,24,13),(10,17,1),(12,20,8),
\]
and the four \(K_9\) triples are
\[
                     (6,9,3),(9,18,9),(9,6,1),(12,27,17).
\]
They follow by multiplication in the displayed full defining bases.

The anchors have odd valuations respectively at the primes
\(127,191,631,19\):
\[
 1016=8(127),\quad764=4(191),\quad5048=8(631),\quad
 16872=8\cdot3\cdot19\cdot37.
\]
NONE of the other time norms in its field row is divisible by
that prime.  The exact nonzero residues of the ABSOLUTE time norms
are respectively
\[
 (53,74,40,45,36),\quad(185,93,13,108),\quad
 (67,234,127,522,44),\quad(10,11,17).
\]
Thus every UNEQUAL anchor/model time ratio has nonsquare rational
norm and cannot be a square in \(K\).
No square norm is used as a sufficient test, and no native,
orthogonal, parallel or dotted mutual case is omitted by this
necessary invariant.

Only \(V=U\) at the actual finite anchor remains.
The endpoint walls are distinct, and their canonical raw norms
are \(2U\).
Equal time squares give equal POSITIVE physical times;
their spatial product is \(U^2/17\).
Therefore the actual two-cycle is
\[
 g={U(34-15U)\over17},\qquad
 L={g^2\over U^2}=\left({34-15U\over17}\right)^2
                     \in\mathcal O_K .
\]
The zero branch would force the irrational anchor \(U=34/15\),
impossible, but integral closedness applies in either case.
It makes \(15U/17\), hence \(U/17\), integral:
\(15\cdot8-17\cdot7=1\).
The four anchor norms are \(2,1,1,3\), in degrees \(2,2,3,3\);
their quotients by the corresponding powers of seventeen
are not integers.
This final contradiction excludes every diagonal pair too.

The whole-\(A_{16}\) exclusion is complete: actual endpoint
lengths and helper/golden cases, BOTH isolated, degree at most six,
full two-dotted interval, finite anchors and actual fields, complete
full-ring domains, every unequal time ratio, and the actual
diagonal two-cycle.
It is manual at the printed finite inequalities and norm tables,
with no new enumeration-output premise, altered dimension bound
or assumed equality of disconnected root scales.

\section{One actual rank-one opposite}
\label{app:arithrankone}

For the full forests \(A_{15}+A_1,B_{15}+A_1,D_{15}+A_1\),
the correct first split concerns the SAME actual opposite \(d\)
of the old rank-one direction \(e\).
Either it contacts the old large component, or all those products
vanish.  In the isolated case physical unit normalization gives
\[
                 d=\sqrt h\,p-\sqrt{1+h}\,e,\qquad h>0 .
\]
No unrelated wall is chosen in the second case, and this identity
does not identify disconnected field-root norm squareclasses.
We first close the whole \(A_{15}+A_1\) case and record its
separate \(B_{15}\) nonisolated corollary.

\subsection{Nonisolated rank-one opposites beside \(A_{15}\)}
\label{app:arithA15nonisolated}

Let the original full forest be \(A_{15}+A_1\), with old
unit \(A\)-roots in zero-sum \(\mathbb R^{16}\), and write
\[
 d=Ap-v_A-\gamma e,\quad A>0,\quad\gamma>0,\quad
 Q=\|v_A\|^2,\quad A^2=\gamma^2+Q-1.
\]
At the actual endpoint of the rank-one edge, old \(A_{15}+d\)
is a finite rank-sixteen white diagram.
Its COMPLETE nonisolated extensions are
\[
\begin{array}{c|c|c}
\text{role}&v_A&Q\\ \hline
\text{endpoint three}&\omega_j/2,\ j=1,15&15/32\\
\text{endpoint four}&\omega_j/\sqrt2,\ j=1,15&15/16\\
\text{next-to-end three}&\omega_j/2,\ j=2,14&7/8 .
\end{array}
\]
These are \(A_{16},B_{16},D_{16}\), including both orientations.
Two old-path contacts make a nonpositive cycle; no exceptional
rank-sixteen completion contains this long path.

We must have \(\gamma>1\).
If \(\gamma<1\), the full positive actual replacement
\[
                 r=p+{A\over1-\gamma^2}(d+\gamma e)
\]
has negative norm, positive time, zero products with \(e,d\)
and all retained old \(A\)-roots except the contacted node \(j\).
Every unknown wall is nonpositive termwise, and its omitted
\(n_j\) product is strictly negative.
It is therefore the ACTUAL endpoint of old edge \(j\).
The SAME \(d\) would close both that direction and the
orthogonal old rank-one direction, contradicting the
original-farthest orthogonal obstruction.
For \(\gamma=1\), \(d+e\) is an actual future null.
Thus \(Q<A^2\) and the entire surviving own domain is dotted.

The actual rank-one edge length is
\(C=\gamma/\sqrt{1-Q}\), with energy at least \((C-1)/(C+1)\).
At \(Q\ge7/8\), \(C>\sqrt8>7/3\) gives energy \(>2/5\).
Together with the independently proved same-vector
\(e_{A_{15}}>.8348\), this excludes endpoint-four and
next-three throughout their unbounded own domains.
For endpoint-three, if \(\gamma\ge41/40\),
\[
                  C^2\ge1681/850>49/25
\]
gives rank-one energy \(>1/6\), and
\(.8348+1/6=1+11/7500\).
Only \(1<\gamma<41/40\) remains.

Here the ACTUAL positive core
\[
                  d+\gamma e=Ap-\omega_j/2
\]
is future timelike of square \(1-\gamma^2\), zero on \(e\),
negative on \(d\), and nonpositive on every other original wall.
Original maximality gives
\(h_j/2\ge A-\sqrt{\gamma^2-1}\).
The right side decreases with \(\gamma\).
At \(41/40\), \(A^2=831/1600>.72^2\) and
\(\sqrt{\gamma^2-1}=9/40\).
Thus the appropriate ACTUAL signed \(A_{15}\) endpoint satisfies
\[
                         x_{\rm end}>\sqrt2(.72-.225)>.69 .
\]

Rebuild its reference before using the new pressure.
Take the common N16 raw heights
\(2.021,2.137,2.205,2.227\), reflected,
with their possible single golden hole.
Internal omission loses at most \(.000576\).
End omission replaces the first tail by \(0\)--\(6\), giving
\[
                         E_{\rm end}=48050651/60000000>.80084 .
\]
Thus every RAW reference has energy \(>.79\), before any
earlier golden-endpoint graft.
At the selected endpoint its flat block has length five or six,
mean \(.25<b<.3\), and next coordinate \(<.101\), including
an adjacent internal omission.
Graft this block afresh to \(.69\).
The other coordinates are \(>(L(.25)-.69)/(L-1)\ge.14\);
the changed central coefficient is
\(>(.14-.101)/\sqrt2>.027>.022\).
Every other positive central coefficient retains \(.022\);
the golden hole remains coefficient zero.
The exact sum-preserving zero-sum change has zero old block cross,
so the new ONE reference has energy
\[
                    >.79+(6/5)(.69-.3)^2=.97252 .
\]
Absorb EVERY retained actual isolated foot again into THIS
reference, using the common N16 loss inequality and its NEW
coefficients.
This gives \(e_A>.97252\), including adjacent and foreign feet.
No earlier optimized endpoint gain or foot energy was added.

Finally the original rank-one edge has
\(C>\sqrt{32/17}>11/9\), so its distinct orthogonal energy is
\(>1/10\).
Adding that energy ONCE gives
\(\|x\|^2>.97252+.1>1\).
This excludes the entire nonisolated actual rank-one opposite
beside \(A_{15}\), over every actual field.

For later use the corresponding \(B_{15}\) argument is separate.
A finite rank-sixteen extension of old \(B_{15}\) can only be
\(B_{16}\), attaching by three at its LONG endpoint.
Thus \(v_B=\Omega_1/2\), \(Q=1/2\), and
\(A^2=\gamma^2-1/2\).
The same reciprocal actual projection proves \(\gamma>1\).
The original rank-one energy is then
\(>(\sqrt2-1)/(\sqrt2+1)=3-2\sqrt2>.17\).
The already written fresh \(B_{15}\) bulk in the proof of
Proposition~\ref{prop:f16}, with arbitrary spectators and both
initial-five cases, gives \(e_B>E_{13}>.88989\).
Their sum exceeds \(1.05989\).
Thus a nonisolated actual rank-one opposite is impossible
beside \(B_{15}\) as well.  No \(D_{15}\) conclusion is
transported from this corollary.

\subsection{The isolated \(A_{15}\) case and its invariant cycles}
\label{app:arithA15isolatedcycles}

Now suppose the SAME actual opposite is isolated:
\(d=\sqrt h\,p-\sqrt{1+h}\,e\).
If \(h\ge1\), its actual edge has energy at least
\(3-2\sqrt2>.17\); together with \(e_A>.8348\) this
already contradicts the norm budget.
Assume \(0<h<1\) and put \(V=4h\).
The actual two-cycle makes \(V\) integral and all nonphysical
places satisfy \(-4<V^\sigma<0\).
Physically \(0<V<4\).
Consequently \(V/4\) CANNOT be integral, since its nonzero norm
would have magnitude below one.
The actual field has degree at least two by the stated
Meyer/compact-anisotropy input for the ENTIRE rank-seventeen form.

We need pressures only at
\[
 J=\{5,6,7,9,10,11\},\qquad
 \ell_i=.49\ (i\text{ odd}),\quad
 \ell_6=\ell_{10}=.496 .
\]
The midpoint eight will never be assigned a pressure or coefficient.

Let an actual central opposite be
\[
 c=Ap-v_A-\beta e,\quad Q=\beta^2<1,\quad
 \lambda=-c\cdot d=A\sqrt h-\beta\sqrt{1+h}\ge0.
\]
The INVARIANT actual two/three cycles are
\[
 L=4\lambda^2,\quad W=4\beta^2,\quad
T=8\beta\sqrt{1+h}\lambda,\qquad L,W,T\in\mathcal O_K .
\]
Zero \(\beta\) or \(\lambda\) is included.
Direct expansion cancels the mixed radical:
\[
                  L+T=VA^2-W(1+V/4).
\tag{Rankonecycle}\label{eq:arithrankonecycle}
\]
No individual unit coefficient or disconnected raw time is
assumed \(K\)-valued.

For an isolated-within-\(A\) cap,
\(v_A=\gamma\omega_i\), \(B=4\gamma^2\),
this gives the exact full-foreign identity
\[
              L+T+V+W={i(16-i)BV\over32}\in\mathcal O_K.
\tag{A15isocycle}\label{eq:arithA15isocycle}
\]
At EVERY nonphysical embedding, positive definiteness on the
whole old \(A_{15}\) and this cap gives
\(0<B^\sigma<4/D_i\); its additional foreign square is
nonnegative and only strengthens the bound.
Together with \(-4<V^\sigma<0\), this gives
\(\lvert B^\sigma V^\sigma\rvert<16/D_i\).
For one white-three contact,
\(v_A=\gamma\omega_i+\omega_j/2\),
\(t=2\gamma\), and
\[
 M=i(16-i)t^2+2\min(i,j)(16-\max(i,j))t+j(16-j)-32,
\]
the corresponding identity is
\[
                         L+T+W=MV/32\in\mathcal O_K.
\tag{A15whitecycle}\label{eq:arithA15whitecycle}
\]
Thus the foreign radical is retained through its triangle term.
Dropping that term would not prove these congruences.

\subsection{Finite native parameters and all dyadic places}
\label{app:arithA15dyadic}

For \(a=2\cos(\pi/m)\), \(m\ge4\), in any containing field,
every dyadic valuation normalized by \(v(2)=1\) satisfies
\[
                         0\le v(a),v(a+1)\le1/2 .
\]
Here is the full root-of-unity argument.
Write \(a=\zeta+\zeta^{-1}\), \(\zeta\) primitive \(2m\)-th.
Prime-to-two roots keep their order on reduction: a nontrivial
odd-order root reducing to one would make its zero geometric
sum reduce to a nonzero odd integer.
Thus \(v(a)>0\) occurs only when \(2m=2^s\).
For \(s\ge3\),
\(v(a)=v(\zeta^2+1)=1/2^{s-2}\le1/2\).
The formula follows because conjugate primitive \(2^r\)-root
differences have equal valuation and product two.
Likewise
\[
                a+1=\zeta^{-1}(\zeta^3-1)/(\zeta-1)
\]
has positive valuation only when \(2m=3\cdot2^s\);
for \(m\ge4\), \(s\ge2\), and
\(v(a+1)=1/2^{s-1}\le1/2\).
For \(m=3\), \(a=1\) and \(v(a+1)=1\); that exception is
removed by an actual affine cycle BEFORE this lemma is used.
In particular \(B=4\cos^2(\pi/m)=a^2\) satisfies \(v(B)\le1\)
for every finite \(m\ge3\).
If \(a\notin K\), use the cyclotomic compositum and descend
the normalized valuation of \(B\in K\).
These bounds persist under every actual field extension.

For a finite one-white-three own, the actual three-cycle
along the old path excludes own three.
For every other finite own, apply the bounds to \(t\).
Put \(a_i=i(16-i)\) and \(C=\min(i,j)(16-\max(i,j))\).
The following four parity cases prove \(v(M)\le3\)
at EVERY dyadic place.

If \(i,j\) are odd, rewrite in \(s=t+1\):
\[
 M=a_is^2+2(C-a_i)s+[r(16-r)-32],\qquad r=|i-j|.
\]
The last two coefficients are divisible by four and \(a_i\)
is odd.  The leading term has valuation at most one, while
the others have valuation at least two.
If \(i\) is odd and \(j\) even, use \(t\) instead:
\(a_i\) is odd, \(C\) even and \(a_j\) divisible by four.
The same unique-minimum argument gives \(v(M)\le1\).
If \(i\) is even and \(j\) odd, the constant is odd and the
other coefficients are even, giving \(v(M)=0\).
If both are even, \(i=6,10\) makes \(k=i/2\) odd.
With \(l=j/2\),
\[
 M/4=k(8-k)t^2+2\min(k,l)(8-\max(k,l))t+l(8-l)-8 .
\]
For even \(l\) use the \(t\)-argument; for odd \(l\) use
\(s=t+1\), whose constant
\(r(8-r)-8\), \(r=|k-l|\), is divisible by four.
Thus \(v(M/4)\le1\) and \(v(M)\le3\).
No midpoint-eight assertion has been made.

Equation~\eqref{eq:arithA15whitecycle} now forces \(v(V)\ge2\)
at every dyadic place.
At all other finite places \(V\) is already integral.
Hence \(V/4\) would be integral, impossible.
This excludes EVERY finite one-white-three cap at the six
required sites, including foreign contacts and both orientations.
The same proof applies to a dotted \(t\) whenever both
\(v(t),v(t+1)\le1/2\) have been established.

\subsection{The fresh six-site inventory at losses \(.49,.496\)}
\label{app:arithA15freshinventory}

At each \(i\in J\), both retained old paths have ranks at
least four.
Their COMPLETE finite-white list is no old-\(A\) contact;
one three at an admissible end or \(D/E\) position;
one four at a path end; or one endpoint three in each path.
The two-path bridge has rank fifteen, so it is \(A_{15}\):
a third arm through the new root would have two old arms at
least four, incompatible with finite \(D\), and no exceptional
finite type has this rank.
The old \(e\)-contact \(\beta\) is retained throughout.
It only restricts the white list further.

Own zero is not silently discarded.
Without old-\(A\) contact its time \(Q-1\) is negative.
For a white three at an internal old node, the full old diagram
and cap have actual affine arms \((1,2,5)\) or \((1,3,3)\),
unless the contact is at \(1,2,14,15\).
At global ends one or fifteen, the full finite white path
has rank at least five; adding \(e\) allows only
\(\beta=0,1/2,1/\sqrt2\).
Then \(A^2=15/32+Q-1<0\).
At two or fourteen, \(\beta=0\) gives time \(-1/8\).
A positive foreign contact in that finite branched white
component forces \(\beta=1/2\), and the full actual roots
then have the affine arms \((1,2,5)\).
An internal white four has the \(4/3/3\) star.
At a global end white four, any positive \(e\)-contact would
put a three BEFORE the four in a path of rank at least six,
or make a nonfinite \(4,4\) principal; hence \(Q=0\) and
time \(15/16-1<0\).
Two white threes at own zero give the actual all-three cycle.
Every null is a positive sum of actual roots, future through
the cap time, and nonpositive against ALL omitted walls.

For every positive own, the parallel value is excluded BEFORE
the remaining case dispatch: \(c+n_i\) is its actual feasible
future null, including positive foreign contacts.
For white three on a side of order \(m=i\) or \(16-i\),
number its node \(j\) from that side's far end.
The actual residual and numerator are
\[
 \delta=1-j(m-j)/(2m),\qquad \sigma=\gamma+j/(2m).
\]
For white four at an end they are
\(\delta=1/m\), \(\sigma=\gamma+j/(\sqrt2m)\).
For two white threes, \(\delta=1/D_i\),
\(\sigma\ge\gamma+1/D_i\).
The ENTIRE old-\(e\) cost changes the denominator to
\(\delta-Q\), and
\[
                   C_i={\sigma\sqrt{D_i}\over\sqrt{\delta-Q}} .
\]
If the desired \(\ell_i\)-pressure fails, this implies
\[
 {\sigma\over\sqrt{\delta-Q}}
       \le R_{\ell_i}(\sqrt{D_i})<1.881 .
\tag{A15ratio}\label{eq:arithA15ratio}
\]
For completeness the ratio target increases with \(q\) and
decreases with \(\ell\) on \(q\le\sqrt8<2.829\),
\(\ell\ge.49\).
At \(q=2.829,\ell=.49\), positive denominator clearing gives
the strict margin119649968.
Thus this is a fresh changed-loss inventory, rather than
reuse of the old \(.601\) weak list.

Finite one-white-three owns have already been excluded by
the invariant parity argument, including own three's affine cycle.
Thus only DOTTED \(\gamma>1\) remains there.
At a \(D\)-position the ratio is
\(>(1+1/m)/\sqrt{2/m}\ge6/\sqrt{10}>1.881\).
At an \(E\)-position \(\delta\le1/4,\sigma>1\), stronger still.
Foreign cost decreases the denominator and cannot weaken this.
For near white four, every positive own gives ratio
\[
            \ge(1/2+4/(5\sqrt2))\sqrt5>2.38 .
\]
For far white four, owns three, four and six have actual
affine principals using old roots at most three steps on each
side of \(i\), avoiding the helper at one or fifteen.
The remaining own is at least \(\varphi/2>.809\), giving
\(.809\sqrt5+1/\sqrt{10}>2.12\).
For two white threes, own three has an actual cycle.
Every other positive own is at least \(1/\sqrt2>.7\),
and its ratio exceeds \(.7(2.6)+.35>1.881\).
Every finite-white role except dotted near/far three has now
been included and removed.

\subsection{Every remaining dotted helper fails over every field}
\label{app:arithA15dottedhelpers}

For near white three the ratio is
\[
 {\,\gamma+(m-1)/(2m)\over\sqrt{(m+1)/(2m)-Q}}.
\]
At \(m\ge7,\gamma>1\), it exceeds \(10/\sqrt{28}>1.881\).
For \(m=5,6\), weakness gives
\(2<t=2\gamma<2.11555<53/25\).
These are \(i=5,j=4\), \(i=6,j=5\), and reflections.
Lift EVERY nonidentity embedding of \(\mathbb Q(t)\) to \(K\).
Strict full-old-\(A\) positivity gives respectively
\[
 (55/8)t_\sigma^2+11t_\sigma+2<0,\qquad
 (15/2)t_\sigma^2+(25/2)t_\sigma+23/8<0 .
\]
The foreign square is nonnegative, strengthening these conditions.
Both negative quadratic intervals lie in \((-1.4,-.2)\);
check positive endpoint values at \(-7/5,-1/5\).

If the minimal degree is at least three and some small conjugate
is at least \(-1\), the nonzero integral norm of \(t(t+1)\)
has magnitude
\[
              <6.6144(1/4)(.56)^{d-2}\le.926016<1 .
\]
If all small conjugates are below \(-1\), the norm of \(t+1\)
has magnitude \(<3.12(.4)^{d-1}<1\).
In degree two the integral trace is one, and the norm
\(t(1-t)\) lies strictly between \(-3,-2\).
In degree one no integer lies in \((2,2.12)\).
Thus every weak near dotted helper is impossible, including
foreign contacts and proper actual fields.

For far white three, if \(\gamma\ge34/25\) its squared ratio
is at least its \(m=5\) value \(5329/1500>1.881^2\).
To check the minimum, put \(r=1/m\).
The derivative of \(2(\gamma+r/2)^2/(1+r)\) has sign
\(1-\gamma+r/2<0\), and the ratio increases with \(\gamma\).
Hence a weak far helper has \(2<t<68/25\).
Conjugate positivity gives the entire strip
\[
                            -4/5<t_\sigma<12/25 .
\]
The N16 endpoint cost checks at \(-.8,.48\) have margins
\(3/40,59/1000\); they increase over the reflected half-window.
Again positive foreign cost only strengthens them.

For \(f(t)=t(t^2-1)\), small conjugate magnitude is at most
\(2/(3\sqrt3)\).
Degree at least four would give norm magnitude
\[
 {68\over25}\left\{\left({68\over25}\right)^2-1\right\}
                         {8\over81\sqrt3}<1 .
\]
For an explicit positive seal, eight times the physical factor
is139.229184, below \(81(1.732)=140.292\).
Thus the minimal degree is at most three.
At degree three, \(|\mathrm N(t)|<2.72(.8)^2<2\), so it is
a unit.
Two positive small conjugates would give norm below one.
For opposite signs their absolute product \(1/t\) forces
\[
 t>125/48,\quad
 t_{\rm positive}>125/272,\quad
 |t_{\rm negative}|>625/816,
\]
putting the integral trace strictly between two and three.
For two negative conjugates, the trace is one and \(t<2.6\).
The integral second coefficient
\(1/t-t(t-1)\) lies in \((-4,-1.5)\), hence is \(-2\) or \(-3\).
The two cubics are
\(X^3-X^2-2X-1\), of discriminant \(-31\), and
\(X^3-X^2-3X-1=(X+1)(X^2-2X-1)\).
Neither is the minimal polynomial of the totally real cubic.
Degree three is excluded.

Degree two has trace two or three.
Trace two has norm in \((-1.9584,0)\), hence \(-1\);
trace three has positive norm below two, hence one.
The exact possibilities are
\[
                         t=1+\sqrt2,\qquad t=\varphi^2 .
\]
Degree one is the forbidden parallel two.
For \(1+\sqrt2\), \(t\) is a dyadic unit and
\(v(t+1)=1/2\).
For \(\varphi^2\), both \(t\) and \(t+1=\varphi+2\) have
odd norm and are dyadic units.
These facts persist in every actual field extension.
Thus the already proved four-parity argument applies to \(M\)
for BOTH remaining dotted possibilities.
Equation~\eqref{eq:arithA15whitecycle} forces \(V/4\) integral,
a contradiction.
All weak dotted helpers are excluded.

\subsection{The isolated central caps through the full \(h<1\) interval}
\label{app:arithA15centralisolated}

For isolated-within-\(A\) contact, use
\eqref{eq:arithA15isocycle}.
At odd sites \(i(16-i)\) is coprime to32, so integer Bezout
gives \(BV/32\in\mathcal O_K\).
At six and ten, the gcd is four, giving \(BV/8\in\mathcal O_K\).
For ANY finite own the dyadic bound \(v(B)\le1\) then forces
\(v(V)\ge4\) or \(\ge2\), respectively.
Thus every finite isolated own would make \(V/4\) integral,
impossible; own zero has negative time.
Parallel was already excluded globally by its actual null.

For dotted own, \(u=(1-Q)/\gamma^2\in(0,1)\).
Its positive ACTUAL core
\(c+\beta e+\gamma u n_i\)
has square \(-\gamma^2u(1-u)\), positive time, zero \(c,e\)
products, negative \(n_i\) product and all other wall products
nonpositive.
It gives
\[
 h_i>\sqrt{D_i-u}-\sqrt{u(1-u)},\qquad
 \text{loss}\le au+\sqrt{u(1-u)} ,
\]
with \(a<.2\) at the odd sites and \(a<.19\) at six and ten.
If \(u\le1/4\) at an odd site, the loss is
\(<1/20+\sqrt3/4<39/80<.49\).
Thus an odd weak cap forces \(u>1/4\) and \(B<16\).
The integral \(BV/32\) has physical magnitude \(<2\) and
EVERY other magnitude \(<1/(2D_i)\le4/55\).
Since the actual degree is at least two, its nonzero norm has
magnitude \(<8/55<1\), impossible.

At six or ten, for \(u\le.27\) the loss is
\(<.0513+.444=.4953<.496\).
An even weak cap therefore has \(u>.27\), \(B<400/27\).
The integral \(BV/8\) has physical magnitude \(<200/27\),
and all other magnitudes \(<4/15\).
At degree at least three its norm is
\(<(200/27)(4/15)^2=128/243<1\).
At degree two, instead use the rational integer norms:
\[
       0<\mathrm N(B)<640/81<8,\qquad
       0<|\mathrm N(V)|<16 .
\]
The first is positive and the second negative before taking
absolute values.
Integrality of \(BV/8\) requires64 to divide their product.
But an integer from one through seven has dyadic valuation
at most two, and one through fifteen at most three.
Their sum is less than six, a contradiction.
This argument permits arbitrary ramification and full maximal order.

Combining the complete own-zero, finite-white, finite, parallel,
dotted-helper and isolated cases proves, for EVERY \(i\in J\),
\[
                            h_i>\sqrt{D_i}-\ell_i .
\]
There is still NO midpoint-eight assertion.

\subsection{One final reference for every isolated rank-one height}
\label{app:arithA15isolatedcomposition}

Choose the rational heights
\[
 H_5=H_{11}=2.132,\quad H_6=H_{10}=2.2426,\quad
 H_7=H_8=H_9=2.3162
\]
and linear outer tails.
Their six positive constraints lie strictly below the just
proved pressures by direct squares at \(55/8,15/2,63/8\).
The midpoint coefficient is zero, the slopes decrease, and
\[
                 E=2.132^2/5+.1106^2+.0736^2=.92673412 .
\]
For \(1/2\le h<1\), the separate rank-one edge energy is
at least \(5-2\sqrt6>.1\); hence the original norm exceeds one.
Together with the earlier \(h\ge1\) argument this leaves
only \(0<h<1/2\).

We retain BOTH actual \(A_{15}\) endpoint pressures in that case.
Put
\[
 T=\sqrt{1+4h}-2\sqrt h,\qquad
 u={\sqrt h\over\sqrt{1+h}+1},\qquad
 \alpha=\sqrt{15/16}.
\]
The actual rank-one edge gives energy \(\ge u^2\), \(0<u<1/3\),
and direct substitution gives
\((1-u^2)(1-T^2)=8uT\).
The rank-free actual-\(I\) wall comparison
\eqref{eq:Ibarrier} applies with every foreign contact retained.
Explicitly, for another actual opposite with time \(A\)
and old-\(e\) contact \(\beta=0\) or \(1/2\le\beta<1\),
its chosen spatial derivative square is at most
\(1+A^2-\beta^2\).
For \(\beta=0\), the native gap gives \(A^2\ge1/(4h)\);
for \(\beta>0\), acuteness gives
\(A^2\ge(1+h)\beta^2/h\).
Both imply full tangent square \(\ge1/(1+4h)\).
In the positive-contact case the comparison, after positive
denominators are cleared, has numerator
\(h((5+4h)\beta^2-1)>0\).
Thus its normalized half-edge height is at least \(T\).
At BOTH signed \(A_{15}\) endpoints this gives
\(x_1,-x_{16}\ge\alpha T\).
This is the original wall comparison, not a new field assumption.

The first and last five-coordinate reference means are
\(b=2.132/(5\sqrt2)<.30152<.302\).
Since \(T>\sqrt3-\sqrt2>.317\) and \(\alpha>.968\),
\(\eta=\alpha T>b\).
Graft BOTH disjoint blocks to \(\eta\).
Their other coordinates
\((5b-\eta)/4>1/8\) exceed the next old coordinate
\(.1106/\sqrt2<.079\).
Thus all dual coefficients remain nonnegative, supported only
on the six proved pressures and the two ACTUAL endpoint bounds;
the midpoint coefficient remains zero.
The single rebuilt reference gives
\[
                         e_A>E+(5/2)(\eta-b)^2 .
\]

For completeness the scalar barrier \(T+3u/2>.76\) follows
directly from the displayed relation.
At \(T_0=.76-1.5u>0\), its left-minus-right is bounded below by
\[
                  264/625-(19/5)u+(21419/2500)u^2>0 .
\]
The quadratic discriminant is \(-55964/1562500<0\);
the cubic term was bounded using \(u<1/3\).
The relation decreases strictly with positive \(T\), proving
the barrier.
With \(D=\eta-b>0\),
\[
 D+(3/2)\alpha u>.76\alpha-b>.43368.
\]
Weighted Cauchy gives
\[
 (5/2)D^2+u^2>
 {.43368^2\over2/5+(9/4)\alpha^2}
 ={.43368^2\over803/320}>.074 .
\]
The positive comparison is
\(.1880783424>.18569375\).
Adding the original rank-one energy ONCE to the rebuilt
\(A_{15}\) reference yields
\[
                     \|x\|^2>.92673412+.074=1.00073412>1 .
\]
All \(h>0\), including boundaries \(1/2,1\), are excluded
for the isolated opposite.

Together with the complete nonisolated argument, this excludes
the ENTIRE original \(A_{15}+A_1\) family, with either ordinary
nonshort rank one or actual short \(I\).
The dichotomy concerns one actual wall throughout, and the two
case contradictions were never added to one another.
No historical frontier-count subtraction is needed.
The separate \(B_{15}\) nonisolated corollary only supplies its
isolation reduction; its isolated ordinary and actual-\(I\)
branches remain to be proved.
\subsection{A fresh two-tail comparison for \(D_{15}\)}
\label{app:arithD15body}

We first supply the missing body estimate for \(D_{15}+A_1\).
The conclusion is an energy bound, rather than unconditional
strength of the last two knots.
Use actual unit roots
\[
 n_j=(e_j-e_{j+1})/\sqrt2\quad(1\le j\le14),\qquad
 n_{15}=(e_{14}+e_{15})/\sqrt2 .
\]
For \(P_i=e_1+\cdots+e_i\), the long dual is
\(\Omega_i=\sqrt2P_i\), \(i\le13\).
The two fork duals are
\[
 \Omega_{14}=(P_{14}-e_{15})/\sqrt2,\qquad
 \Omega_{15}=(P_{14}+e_{15})/\sqrt2,\qquad
 \|\Omega_{14}\|^2=\|\Omega_{15}\|^2=15/2 .
\]
The old rank-one unit root \(e\) is orthogonal to this space.
Retain the SAME original decomposition
\[
 z=E(p-x-y e),\quad x_1>\cdots>x_{14}>|x_{15}|,\quad
 e_D=\|x\|^2,\quad e_1=y^2,\quad e_D+e_1<1 .
\]
Set \(f_i=\sqrt i-.61/\sqrt2\) and
\[
 E_K=f_5^2/5+\sum_{j=5}^{K-1}(\sqrt{j+1}-\sqrt j)^2.
\]
The shared geometric bulk inventory and
Lemma~\ref{lem:literal-helper} give strong \(6,\ldots,10\)
and \(e_D>E_{10}>.82436\), with either initial-five reference.
The generic parameter is \(M=15-6=9\), so its last
established index is ten.  Eleven and twelve need a fresh check.

For their actual opposite write
\[
 c=Ap-v_D-\beta e,\quad A>0,\quad
 \gamma=-c\cdot n_i\ge0,\quad Q=\beta^2<1.
\]
The full retained white Gram is positive definite.
Every nonzero foreign cost satisfies \(Q\ge1/4\).
If the inverse calculation gives residual \(\delta>0\) and
\(\cosh L_i=\sqrt{2i}\,\sigma/\sqrt\delta\),
the original edge comparison makes the index strong if
\[
 {\sigma\over\sqrt\delta}>T(\sqrt{2i}),\qquad
 T(q)={100\over61}-{q-.61\over q(2q-.61)}.
\]
On this interval \(T\) increases and \(\sqrt{2i}<5\), so
\[
 T(\sqrt{2i})<T(5)={442721\over286395}<{31\over20};
\]
the final cleared numerator is \(23825>0\).
This actual threshold, rather than \(100/61\), retains
the full bridge list at this rank.

Deleting eleven leaves \(A_{10}+D_4\), with long leaf twelve,
centre thirteen and forks fourteen and fifteen.
Deleting twelve leaves \(A_{11}+D_3\), with middle thirteen
and endpoints fourteen and fifteen.
Finite classification gives the COMPLETE list:
on the left, endpoint-three, endpoint-four or next-to-end-three;
on \(D_4\), simple three at any of its three leaves;
on \(D_3\), middle-three or endpoint-three, four or five.
A two-piece bridge is simple three at a left endpoint and
any listed simple-three right contact.
Both left orientations and both fork orientations are retained.
There is no exceptional finite rank-fifteen diagram.
A bridge cannot have an extra old-\(e\) contact:
that adds two forks to its \(D_{15}\) completion or the
nonfinite arms \(11,3,1\) to its \(A_{15}\) completion.
Thus ONLY for bridges \(Q=0\).

For a left simple three at \(j\), counted from its far end,
\[
 \sigma=\gamma+{j\over2i},\qquad
 \delta={i+1\over2i}\text{ at endpoints},\qquad
 \delta={2\over i}\text{ at next-to-end sites}.
\]
For endpoint four,
\(\sigma=\gamma+j/(\sqrt2i)\), \(\delta=1/i\).
Foreign \(Q\) subtracts from every displayed residual.
The remaining exact inverse-Gram formulas are
\[
\begin{array}{c|c|c}
 \text{right contact}&\delta\text{ before }Q&\sigma\\ \hline
 D_4\text{ long three or }D_3\text{ middle three}
       &1/2&\gamma+1/2\\
 D_4\text{ fork three}&1/2&\gamma+1/4\\
 D_3\text{ endpoint three}&5/8&\gamma+1/4\\
 D_3\text{ endpoint four}&1/4&\gamma+1/(2\sqrt2)\\
 D_3\text{ endpoint five}&(5-3\varphi)/8&\gamma+\varphi/4
\end{array}
\]
and
\[
\begin{array}{c|c|c}
 i,\text{ right contact}&\delta&
       \sigma-\gamma\text{ for near/far left endpoint}\\ \hline
11,\text{ long three}&1/22&21/22,\ 6/11\\
11,\text{ fork three}&1/22&31/44,\ 13/44\\
12,\text{ middle three}&1/24&23/24,\ 13/24\\
12,\text{ endpoint three}&1/6&17/24,\ 7/24 .
\end{array}
\]
The middle-three \(D_{15}\) bridge is retained separately;
treating \(D_3\) only as an end-contact path would miss it.

Every affine exclusion below applies with arbitrary foreign
contacts.  Its positive combination of actual roots has
positive future time and nonpositive product with every
unselected original wall termwise.
It would be a feasible null direction, impossible in a
compact polytope.
For a single three at global long node \(j\ge2\), own zero
has the affine \(D\) principal on the cap, \(n_{j-1}\),
trunk \(n_j,\ldots,n_{13}\), and BOTH forks;
the trunk coefficients are two and the four terminal
coefficients one.
At a fork, own zero has affine \(E_8\) arms \(1,2,5\):
the other fork is the length-one arm and the contacted fork
plus cap is the length-two arm.
Far-three at one with own zero has
\(A^2=Q-1/2\) and \(Q<(i+1)/(2i)\le6/11\);
the only native costs are \(0,1/4,1/2\), so \(A^2\le0\).
Far-four with own zero forces \(Q<1/i<1/4\),
hence \(Q=0\) and zero time.
Near-four with own zero has the affine \(4,3,3\) star.
Right endpoint-four with own zero has affine \(B_4\)
kernel \((\sqrt2,2,2,1,1)\) on the cap, contacted fork,
centre, old twelve and other fork.
Right endpoint-five remains in the direct estimates.

Own three with a single-three helper has the actual
all-three affine cycle.
Own six has the affine \(6,3\) path using an old neighbour
other than the helper.
Own four with a right fork helper has affine \(F_4\)
path \(n_{i-2},n_{i-1},n_i,c,\text{fork}\),
whose bonds are \(3,3,4,3\).
With a nonadjacent left helper it has the \(4,3,3\) star.
The ADJACENT near-three/own-four profile is retained;
it is not declared an affine star.
An isolated-within-\(D\) cap has the affine \(D\),
star and \(6,3\) principals at owns three, four and six.
Parallel own always has the actual feasible null \(c+n_i\).

The remaining direct estimates are exhaustive.
Right long-three and middle-three have ratio at least
\(1+1/\sqrt2>31/20\).
Right fork-three and endpoint-three have
\(\gamma\ge\varphi/2>.8\); if \(Q>0\), their largest
residual \(3/8\) gives
\(1.05/\sqrt{3/8}>31/20\).
Positive-own right endpoint-four has ratio at least
\(1+1/\sqrt2\).
Endpoint-five has residual \(<.02\) and numerator \(>.4\),
hence ratio \(>2.8\), even for own zero.
Left near-three with \(\gamma\ge\varphi/2\) has ratio
\[
 {\,\sigma\over\sqrt\delta}>
 {\, .8+5/11\over\sqrt{6/11}}>31/20 .
\]
Its extra own-four case has \(\gamma>.7\),
\(\sqrt\delta<.74\), ratio \(>635/407>31/20\).
Next-to-end-three has \(\delta\le2/11\) and \(\gamma>.8\).
Positive-own endpoint-four has ratio at least
\(\sqrt{11}/2>31/20\).
Every positive-own bridge has ratio \(>1.9\):
even its weakest numerator is \(19/24\) over \(\sqrt{1/6}\).
Own-zero bridges have the actual affine cycle.

The only unresolved left far-three is the literal packet
\(v_D=\gamma\Omega_i+\Omega_1/2\).
Weakness gives \(\gamma<(31/20)\sqrt{6/11}<1.147\);
the own removals give
\(202/125<\varphi\le2\gamma<12/5\).
Lemma~\ref{lem:literal-helper} excludes it at eleven
and twelve, including foreign contacts.
That lemma is not applied to the right-fork packet.

We also need all isolated-within-\(D\) caps, without
assuming \(Q=0\).
Their remaining own satisfies \(\gamma\ge\varphi/2\)
by the ALL-foreign-cost affine removals already given.
Put \(D_i=2i\) and \(u=(1-Q)/\gamma^2\).
For \(0<u<1\), the actual positive core
\[
 \xi=c+\beta e+{1-Q\over\gamma}n_i
\]
has selected products zero except
\(\xi\cdot n_i=\gamma(u-1)<0\), square
\(-\gamma^2u(1-u)\), and positive time.
The original maximum gives, using \(x\cdot n_i>0\),
\[
 x\cdot\Omega_i>
       \sqrt{D_i-u}-\sqrt{u(1-u)} .
\]
Since \(q_i-\sqrt{D_i-u}<.108u\), its loss is at most
\[
                {.108+\sqrt{1+.108^2}\over2}<.557<.61 .
\]
At \(u=1\) the same core is a forbidden feasible null.
For \(Q>0,u>1\), positively project onto the ACTUAL
finite set \(c,n_i,e\), with residual
\(\gamma^2(u-1)>0\).
All correction coefficients are positive; clearing the
old-\(e\) contact leaves a pure-\(D\) displacement.
Here
\[
 u\le{3\over\varphi^2}=6-3\varphi<7/6,\qquad
 \sqrt{u-1}\le2-\varphi<.4,\qquad
 q_i-\sqrt{D_i-u}<.13 ,
\]
so the full loss is \(<.53\).

For \(Q=0\), native five uses the LEFT actual \(H_4\)
correction
\[
 w_5=\gamma(3n_i/2+n_{i-1}+n_{i-2}/2),\qquad
 \delta_5=1-3\gamma^2/2>0 .
\]
Its actual finite projection gives
\(x\cdot(\gamma\Omega_i-w_5)>A-\sqrt{\delta_5}\).
The divided displacement is
\(3\Omega_{i+1}/4+\Omega_{i-3}/4\);
long-arm concavity bounds its pressure above by
\(x\cdot\Omega_i\).
With \(u=\gamma^{-2}<25/16\), the loss is
\(<25/144+1/4<.61\).
Native seven uses the positive LEFT \(A_2\) negative
core with correction
\(\gamma(4n_i/3+2n_{i-1}/3)\).
Its divided displacement is
\(2\Omega_{i+1}/3+\Omega_{i-2}/3\).
Using \(1.21<u<100/81\), its full loss is
\(<.14+.36=.5\).
All finite natives at least eight use the actual singleton
projection, with loss
\(<.14+\tan(\pi/8)<.555\).
Dotted own uses the preceding core and loss \(<.557\).
All corrections use original roots, with \(i+1\le13\);
neither fork is treated as another long root.

Thus a new weak tail can ONLY be the pure packet
\[
 v_D=\gamma\Omega_i+\Omega_f/2,\qquad
 i=11\text{ or }12,\quad f=14\text{ or }15,\quad Q=0 .
\]
Weakness gives
\(\gamma+1/4<(31/20)\sqrt{5/8}<1.24\).
For \(t=2\gamma\),
\[
                    202/125<t<1.98<12/5 .
\]
Normalize the old \(D\) roots to raw square two over
the ACTUAL field \(K\).
The actual simple-three fork contact permits the same
normalization for the cap; its actual path cycle makes
\(t\) integral.
The raw inverse entries \(i,i/2,15/4\) give
\[
             4i\,t_\sigma(t_\sigma+1)+15<8,
             \qquad -1<t_\sigma<0
\]
at EVERY nonphysical embedding.
Every embedding of \(k=\mathbb Q(t)\) extends to \(K\).
The nonzero integer \(t(t+1)\) has physical value
\(<204/25\) and other conjugates of magnitude at most \(1/4\).
Degree at least three would give a nonzero integral norm
of magnitude \(<51/100\).
In degree two the trace is one; trace two would require
\(t>2\), contrary to the sharper weak strip.
Then \(-2<t(1-t)<-.995456<0\) forces norm \(-1\),
and \(t=\varphi\).
There is no physical rational integer in the strip.
The new weak tail is therefore native five.
Arithmetic positivity by itself does not remove this packet.

For that packet put
\[
 \alpha=\sqrt2(\gamma+1/4),\quad
 \alpha^2=\varphi+5/8>56/25,\quad
 \delta=(5-3\varphi)/8,\quad
 {\sqrt\delta\over\alpha}<.091 .
\]
The last comparison follows, for example, from
\(\delta/\alpha^2<73/8972<.091^2\).
The full time square is
\[
 A^2=i\alpha^2-\kappa,\qquad
 \kappa=1/2\ (i=11),\quad \kappa=5/8\ (i=12).
\]
The LEFT \(H_4\) correction gives the actual finite point
\(p+(A/\delta)(c+w_5)\).
It has positive future time, negative square and zero
selected products.
EVERY omitted original wall, including the OTHER fork
and old \(e\), pairs nonpositively termwise.
For \(Y=v_D-w_5\), \(V=Y/\alpha\), \(K=i-3\),
the SAME original maximum yields
\[
 \begin{gathered}
 x\cdot V>G=\sqrt{i-\kappa/\alpha^2}-.091,\qquad
 V_j=1\quad(1\le j\le K),\\
 V^2=(A^2+\delta)/\alpha^2<i .
 \end{gathered}
\]
All fifteen coordinates are retained in this norm.

Choose the shared positive prefix reference \(w\)
on its first \(K=8\) or \(9\) coordinates.
If five is strong it uses knots \(5,\ldots,K\).
Otherwise use the actual weak-five endpoint,
strong six and terminal \(K\).
Both have \(w^2\ge E_K\), \(w\cdot P_K\le f_K\),
and terminal INDIVIDUAL coefficient respectively
\[
 \sqrt K-\sqrt{K-1}\quad\hbox{or}\quad
 {1\over\sqrt K+\sqrt6},
\]
both greater than \(1/6\).
Pay that coefficient BEFORE using the terminal floor:
\(W=w+(V-P_K)/6\).
The FULL cross with \(w\) is zero since \(V-P_K\)
vanishes on its first \(K\) coordinates.
Its residual norm is \(<i-K=3\).
One square gives
\[
                 e_D>E_K+{G-f_K\over3}-{1\over12}.
\]
At eleven,
\(3.282^2<11-25/112\) gives \(G>3.191\);
\(f_8<2.398\), \(E_8>.76864\) yield \(e_D>.94964\).
At twelve,
\(3.423^2<12-125/448\) gives \(G>3.332\);
\(f_9<2.569\), \(E_9>.79805\) yield \(e_D>.96905\).
If either tail is weak its body bound is stronger than
\(E_{12}>.86989\).
If neither is weak, the ordinary or weak-five reference
through twelve gives this last bound directly.
We have proved the UNCONDITIONAL conclusion
\[
                         e_D>.86989 .
\tag{R1D}\label{eq:arithD15body}
\]
We have not proved unconditional strength of eleven or twelve.

Let \(d\) now be the actual opposite of the old rank-one edge.
If it contacts \(D_{15}\), classification of the FULL
white diagram \(D_{15}+d\) permits only the \(D_{16}\)
completion: simple three at the long endpoint.
Thus
\[
 d=Ap-\Omega_1/2-\gamma e,\qquad
 Q_D=1/2,\qquad A^2=\gamma^2-1/2 .
\]
For \(\gamma<1\), the positive actual projection
\[
                  p+{A\over1-\gamma^2}(d+\gamma e)
\]
is future and feasible against ALL original walls.
All retained products vanish, and its omitted \(n_1\)
product is negative.
It is a finite endpoint on old long edge one, so the
SAME \(d\) closes both that edge and the orthogonal
old rank-one direction.
This contradicts the original-farthest orthogonal
common-opposite obstruction proved above.
For \(\gamma=1\), \(d+e\) is the forbidden future
feasible null.
Hence \(\gamma>1\) and
\(e_1>3-2\sqrt2>.17\).
Together with~\eqref{eq:arithD15body} this exceeds one.
The ACTUAL rank-one opposite is isolated from the
ENTIRE \(D_{15}\).

\subsection{One ordinary rank-one argument for \(B_{15}\) and \(D_{15}\)}
\label{app:arithBD15ordinary}

For \(B_{15}\) use actual roots
\(n_j=(e_j-e_{j+1})/\sqrt2\), \(j<15\), \(n_{15}=e_{15}\);
for \(D_{15}\) retain the preceding roots.
The long duals through thirteen agree.
Their actual old rank-one opposites have been isolated
by separate proofs and therefore have the form
\[
                d=\sqrt h\,p-\sqrt{1+h}\,e,\qquad h>0.
\]
Ordinary rank one is the range \(h\ge1/2\), INCLUDING
its boundary, so its original edge gives
\[
 e_1\ge q(h)={\sqrt{1+h}-1\over\sqrt{1+h}+1}>.1,
 \qquad q(1/2)=5-2\sqrt6>.1 .
\]
This disjoint original energy is counted once.

The \(B_{15}\) short-layer inventory gives strong
\(6,\ldots,13\), with \(e_B>E_{13}>.88989\)
in either initial-five case.
If \(h\ge.57\), then \(\sqrt{1+h}>5/4\)
and \(e_1>1/9>.111\), a contradiction.
Thus \(h<.57\).
For \(D_{15}\),~\eqref{eq:arithD15body} instead gives
\(h<.7\): otherwise \(e_1>3/23\), whereas
\[
                         .86989+3/23-1={747\over2300000}>0 .
\]
In this ORDINARY branch only, either weak tail is now
impossible: its pure-body bound \(>.94964\) or \(>.96905\)
already exceeds the remaining \(1-.1\) budget.
Thus \(D_{15}\) has strong \(6,\ldots,12\).

Before selecting the head, eliminate weak five in either
component.
Its actual pure-component point gives \(x_1>.5965\);
strong six gives \(S_6>f_6>2.01815\).
Interpolate through these pressures and EVERY later
strong knot up to \(J=13\) in \(B\), \(J=12\) in \(D\).
The first-six energy is
\[
 .5965^2+{(2.01815-.5965)^2\over5}
        =.7600299945>.7600 .
\]
Its middle slope is \(>.2843\), above the next ordinary
increment \(<.197\), so all individual duals are positive.
As \(E_6<.6971\), this SINGLE rebuilt reference gives
\(e_T>E_J+.0629\), \(T=B,D\).
Even for \(D\) this is \(>.93279\); adding \(e_1>.1\)
is impossible.
Five is therefore strong.

Fix the positive ordinary reference \(w\) through all
knots \(5,\ldots,J\):
\[
 \begin{array}{c|c|c|c}
 T&J&E_J\text{ lower bound}&h\text{ upper bound}\\ \hline
 B&13&.88989&.57\\
 D&12&.86989&.7
 \end{array}
\]
and
\[
 w_1=\cdots=w_5=b=f_5/5<.361,\qquad
 \text{individual }P_5\text{ coefficient}=
 b-(\sqrt6-\sqrt5)>.14 .
\]
No previous tail or weak-five gain is added to this \(w\).
For \(.49\le X\le.643\), replace its first five
coordinates by \(X,(f_5-X)/4,\ldots,(f_5-X)/4\).
The other four stay \(>.2904\), above the next
coordinate \(<.214\).
Every individual dual remains nonnegative and later
knots are unchanged.
Whenever ACTUAL \(x_1>X\), the single reference square gives
\[
                   e_T>E_J+{5\over4}(X-b)^2
                       >E_J+{5\over4}(X-.361)^2 .
\tag{R1G}\label{eq:arithBD15graft}
\]

Let \(c\) be the ACTUAL opposite of old long edge one.
It is distinct from \(d\), which has no \(B/D\) component.
Deleting one retains \(B_{14}\) or \(D_{14}\) and old \(e\).
The COMPLETE finite extension in that connected white tree
is no contact or simple three at its LONG endpoint two.
For \(D_{14}\), fork attachment gives nonfinite arms
\(11,2,1\), and an internal attachment adds a second fork.
For \(B_{14}\), another weighted bond or fork cannot
give a finite classical extension.
Two contacts to either tree form a nonpositive cycle;
exceptional finite ranks are too small.
Foreign contacts to \(e\) are still retained.

In the nonisolated branch, with \(g=-c\cdot n_1\) and
\(Q=\beta^2=(-c\cdot e)^2\), the unit inverse Gram gives
\[
 c=Ap-g\Omega_1-\Omega_2/2-\beta e,\quad
 A^2=2g^2+2g+Q,\quad
 C_1={2g+1\over\sqrt{1-2Q}},\quad Q<1/2 .
\]
Own zero and \(Q=0\) has zero time.
Own zero and \(Q>0\) has the ACTUAL positive null
\[
                         c+\Omega_2/2=\sqrt Q\,p-\beta e .
\]
Its clearing expansions must be kept separate:
\[
 {\Omega_2\over2}=
 \begin{cases}
 n_1+2(n_2+\cdots+n_{14})+\sqrt2 n_{15},&B,\\
 n_1+2(n_2+\cdots+n_{13})+n_{14}+n_{15},&D .
 \end{cases}
\]
Thus BOTH \(D\) forks and the ACTUAL \(B\) short root
have their correct coefficients.
Selected products vanish; the old-\(e\) product and
every unknown original-wall product are nonpositive.
The null is future feasible, contrary to compactness.
Own three has the affine triangle \(c,n_1,n_2\).
Every other positive own has \(g\ge1/\sqrt2\), so
\(C_1\ge1+\sqrt2\) and
\(x_1>\sqrt{\sqrt2-1}>.643\).
The legal graft~\eqref{eq:arithBD15graft} adds \(>.099405\);
even the \(D\) total is
\(.86989+.099405+.1=1.069295>1\).
The head is therefore isolated WITHIN its entire
classical component, but it may still contact old \(e\).

Its remaining form is
\[
 c=Ap-g\Omega_1-\beta e,\quad
 A^2=2g^2+Q-1,\quad
 C_1^2={2g^2\over1-Q}=1+{A^2\over1-Q}.
\]
Because its actual finite endpoint retains \(e\),
the COMPLETE foreign alphabet splits into
\[
                   Q=0,\qquad Q=1/4,\qquad 1/2\le Q<1 .
\]
The last interval includes EVERY finite foreign label
at least four; no typed-seven alphabet is imported.
Actual acuteness with \(d\) gives
\[
 \lambda=-c\cdot d=A\sqrt h-\beta\sqrt{1+h}\ge0,
 \qquad h(2g^2-1)\ge Q .
\tag{R1N}\label{eq:arithBD15native}
\]
If \(\lambda>0\), the ACTUAL native gap gives
\(\lambda\ge1/2\), including parallel and dotted mutuals.
The squared inequality excludes \(g<1/\sqrt2\);
equality requires \(Q=A=0\).
Owns zero, three and four are therefore impossible.
Own six has the affine \(6,3\) kernel \((\sqrt3,2,1)\)
on \(c,n_1,n_2\), and parallel own has the future
feasible null \(c+n_1\).
Both apply with every foreign contact present.
For \(Q=0\), native five has
\(A^2=(\varphi-1)/2<.31\); even the weaker \(D\)
cut \(h<.7\) gives
\[
                         0<\lambda^2=A^2h<.217<1/4 ,
\]
forbidden by the native gap.
The remaining PURE owns are finite labels at least
seven and the ENTIRE dotted interval.

With \(u=g^{-2}\), their actual pressure table is
\[
\begin{array}{c|c|c|c}
 \text{own}&Z&H&R=Z^2\\ \hline
 \text{finite }\ge7&
 (1-u/2)(e_1+e_2)+(u-1)e_3&
 [\sqrt{2-u}-\sqrt{7u-4-3u^2}]/\sqrt2&
 3-4u+3u^2/2\\
 \text{dotted}&
 (1-u/2)e_1+(u/2)e_2&
 [\sqrt{2-u}-\sqrt{u(1-u)}]/\sqrt2&
 1-u+u^2/2 .
\end{array}
\]
Both \(Z\) are nonnegative, sum to one, and use genuine
first-three long-root coordinates.
For the finite row use the positive combination
\[
 \xi=c+g^{-1}n_1+2(g^{-1}-g)n_2,\qquad
 \xi^2=-(7-4g^2-3g^{-2})<0 .
\]
Its selected products are \(0,0,(3-4g^2)/(2g)<0\).
For the dotted row use
\(\xi=c+g^{-1}n_1\), of square \(-(1-u)<0\).
Both have positive time \(A\).
All unselected original walls, including \(e,d\) and
BOTH \(D\) forks, pair nonpositively termwise.
Their future normalizations and the SAME original
maximum give \(x\cdot Z\ge H\).

We prove the scalar estimate on the FULL domain:
\[
                         .8H-.16R>.2893 .
\tag{R1S}\label{eq:arithBD15purescalar}
\]
For native seven the elementary cubic gives
\(.9009<g<.901\).
Positive rational squaring gives \(A>.7894\),
\(\sqrt{7-4g^2-3/g^2}<.2395\), \(g\sqrt2<1.2743\),
hence \(H>.4315\).
Also \(1.2318<u<1.233\), where \(R\) decreases,
so \(R<.349\), and the expression is \(>.28936\).

For every finite native at least eight,
\(1<u\le4-2\sqrt2<1.172\).
Put \(D=(u-1)(4-3u)\) and \(a=2\sqrt2/5\).
For \(F=.8H-.16R\),
\[
 F'(u)=a\left(-{1\over2\sqrt{2-u}}
             -{7-6u\over2\sqrt D}\right)+.64-.48u .
\]
On \((1,7/6]\) the first contribution is \(<-.28\),
the second is nonpositive and the last is at most \(.16\).
On \([7/6,1.172)\), \(D>.08066\), so \(\sqrt D>.284\);
the now positive second contribution is \(<.032\),
while the others total \(<-.28+.08\).
Thus \(F\) decreases.
At the native-eight endpoint,
\(\sqrt{2-u}>.910\), \(\sqrt D<.289\), \(R<.373\),
giving \(H>.4388\) and \(F>.29136\).

For dotted \(0<u<1\), put \(a=.8/\sqrt2\),
so \(.5656<a<.5657\).
For \(u\le1/2\), concavity gives
\(\sqrt{2-u}\ge1.4142-.379u\), \(R\le1-.75u\).
Consequently
\[
 F>.63987152-.0944u-.5657\sqrt{u(1-u)}
   >.30587152 .
\]
Indeed \(\sqrt{.0944^2+.5657^2}<.5736\), and
\[
 bu+c\sqrt{u(1-u)}
 \le{b+\sqrt{b^2+c^2}\over2}.
\]
For \(u\ge1/2\), use the chord
\(\sqrt{2-u}\ge1.4494-.4494u\), \(R\le.75-.25u\);
then
\[
 F>.69978064-.2142u-.5657\sqrt{u(1-u)}
   >.29023064,
\]
using \(\sqrt{.2142^2+.5657^2}<.6049\).
The two chord endpoint inequalities are positive
rational squarings.
This covers the unbounded dotted-own interval and
proves~\eqref{eq:arithBD15purescalar}.

Use ONE vector
\[
                         W=w+{2\over5}(Z-P_5/5).
\]
Its debit \(2/25=.08\) is below the INDIVIDUAL
\(P_5\) coefficient \(>.14\), paid BEFORE the floor.
The residual has zero sum on the first five coordinates,
hence FULL cross zero with \(w\), and norm \(R-1/5\).
The actual pressure, remaining individual floors and
one square give
\[
 \begin{split}
 e_T&\ge2x\cdot W-W^2\\
    &>E_J+.8H-.16R-.8b+.032\\
    &>E_J+.2893-.8(.361)+.032 .
 \end{split}
\]
Even for \(D\) this is \(>.90239\).
Adding original \(e_1>.1\) is impossible.
This PURE displacement is never used for a mixed head.

For EVERY \(Q\ge1/2\),
\eqref{eq:arithBD15native} gives
\[
 C_1^2>{1+Q/h_*\over1-Q},\qquad
 h_*=\begin{cases}.57,&B,\\ .7,&D.\end{cases}
\]
The expression increases on \([1/2,1)\).
For \(B\) it is at least \(214/57>3\),
so \(C_1>41/25\), \(x_1^2>8/33>.49^2\).
For \(D\) it is at least \(24/7>(37/20)^2\),
so \(x_1^2>17/57>.54^2\).
Their legal graft gains are
\[
 {5\over4}(.49-.361)^2=.02080125\quad(B),\qquad
 {5\over4}(.54-.361)^2=.04005125\quad(D).
\]
Their original totals exceed \(1.01069125\)
and \(1.00994125\), respectively.
This covers ALL finite foreign labels at least four,
finite/dotted head owns and zero or positive mutuals.

For \(Q=1/4,\lambda>0\), the native gap gives
\[
                 A\sqrt h\ge1/2+(1/2)\sqrt{1+h}>1.1 .
\]
For \(B\), \(1.45^2(.57)<1.21\) gives \(A>1.45\),
hence \(C_1>1.95\), \(x_1^2>19/59>.56^2\).
Its legal gain \(.04950125\) is contradictory.
For \(D\), \(1.3^2(.7)<1.21\) gives \(A>1.3\),
\(C_1^2>244/75>(9/5)^2\), \(x_1^2>2/7>.53^2\).
Its legal gain \(.03570125\) makes the total
\(>1.00559125\).

Finally \(Q=1/4,\lambda=0\) gives
\[
                         g^2={1+1/(4h)\over2}.
\]
For \(B\), \(41/57<g^2\le3/4\);
for \(D\), \(19/28<g^2\le3/4\).
Thus
\(\cos(\pi/5)<.81<g\le\cos(\pi/6)\);
more precisely \(g>.84\) in \(B\), \(g>.82\) in \(D\).
The full native alphabet leaves only own six at the
INCLUDED \(h=1/2\) boundary.
Its actual affine null was already excluded.
Parallel and dotted \(g\ge1\) are also impossible.

This covers every original wall, head own, foreign
contact and mutual parameter for \(h\ge1/2\).
Neither \(B_{15}\) nor \(D_{15}\) can occur with an
ORDINARY nonshort rank-one component.
The actual short-\(I\) cases \(0<h<1/2\) still need
their separate arguments; this ordinary bootstrap
does not discard them.

\subsection{The short-\(I\) cycle identity and the \(B_{15}\) long window}
\label{app:arithB15Iwindow}

We now enter \(0<h<1/2\).
Keep the ACTUAL isolated opposite
\(d=\sqrt h\,p-\sqrt{1+h}\,e\), and put \(V=4h\).
Its actual two-cycle and strict nonphysical positivity give
\[
 V\in\mathcal O_K,\qquad 0<V<2,\qquad -4<V_\sigma<0
             \quad\hbox{at EVERY nonphysical place}.
\]
In particular \(V/4\) cannot be integral: every conjugate
of that nonzero number has magnitude \(<1\).

For ANY actual classical-edge opposite
\(c=Ap-v-\beta e\), put \(Q=\beta^2<1\) and
\(A_0^2=A^2-Q=\|v\|^2-1\).
Actual acuteness gives
\(\lambda=-c\cdot d=A\sqrt h-\beta\sqrt{1+h}\ge0\).
The actual cycles
\[
 L_c=4\lambda^2,\qquad W_I=4\beta^2,\qquad
 T_c=8\beta\sqrt{1+h}\lambda
\]
are integral.
The full radical cancellation is
\[
                  L_c+T_c+W_I=V(A^2-\beta^2)=VA_0^2
                        \in\mathcal O_K .
\tag{IC}\label{eq:arithBD15Icycle}
\]
This retains \(\beta=0\), \(\lambda=0\) and every
positive foreign contact.
It does not assume integrality of an individual
Cartan entry or identify disconnected root scales.

For \(B_{15}\), strong \(6,\ldots,13\) are already
available.  If five is weak, its COMPLETE geometric
profile is pure:
\[
 v=\gamma\Omega_5+\Omega_1/2,\qquad
 t=2\gamma\in\mathcal O_K,\quad
 202/125<t<12/5,\quad
 A^2=(5/2)t^2+t-1/2.
\]
Every other conjugate of \(t\) obeys
\(5t_\sigma^2+2t_\sigma-1<0\), hence
\(-.69<t_\sigma<.29\).
Degree at least four contradicts
\(2.4(.69)^3<1\).
In degree three the norm is a unit; a positive small
conjugate would make its magnitude \(<2.4(.29)(.69)<1\).
Writing the other conjugates \(-a,-b\),
the trace is one: it is \(>1.616-1.38>0\) and
at most \(t-2/\sqrt t<2\).
Then \(a+b=t-1<1.38\), \(t<2.38\), and the
integer second coefficient \(1/t-t(t-1)\) is \(-1\)
or \(-2\).
The corresponding polynomials
\[
                  X^3-X^2-X-1,\qquad X^3-X^2-2X-1
\]
have discriminants \(-44,-31\), so are not totally real.
In degree two the norm is \(\pm1\) and the trace one
or two.
Trace two gives \(t(2-t)\in(-.96,.621)\), impossible;
trace one gives \(t=\varphi\).
The only rational value, two, is an actual parallel null.
Thus weak five is golden.

Its connected raw old-\(B\)/cap Gram is defined over
\(\mathbb Q(\varphi)\), with old long and cap squares
two and old short square one.
It is physically Lorentzian.
A proper actual-field extension would give a
nonphysical embedding fixing that indefinite principal.
Hence ACTUAL \(K=\mathbb Q(\varphi)\).
Here \(A^2=(4+7\varphi)/2\), and
\(\mathrm N(4+7\varphi)=-5\) is odd.
Equation~\eqref{eq:arithBD15Icycle} makes \(V/2\)
integral at every dyadic place, and therefore everywhere.
The strict strips force
\(\mathrm N(V)=-4\).
For \(T=V/2\), \(\mathrm N(T)=-1\), \(0<T<1\),
\(-2<T'<0\).
Thus \(T'=-1/T\), \(T>1/2\), and the integral trace
\(T-1/T\in(-1.5,0)\) is \(-1\).
It follows that
\[
          V=2\varphi-2,\quad h=(\varphi-1)/2>.3,\quad
          q(h)>.065 .
\]
But the same weak-five reference gives
\(e_B>E_{13}+.7600-.6971>.95279\).
Adding \(e_I>.065\) is impossible.
Thus five is strong, without a long-head hypothesis.

We next classify the ACTUAL long-fourteen opposite.
Its retained white pieces are \(A_{13}\) on one
through thirteen, the genuine short root \(n_{15}\),
and old \(e\).
The COMPLETE finite possibilities within \(B\) are:
isolation; left-only endpoint-three/four or next-to-end-three;
short-only any finite native contact;
or a left endpoint-three bridge to short-three/four.
Two left contacts form a cycle.
A next-to-end branch joined to short has nonfinite
arms \(1,2,11\); a left-four bridge puts its double
bond internally in rank fifteen; short-five or larger
cannot complete that long finite path.
Every old-\(e\) contact remains, including an old-\(e\)
three arm on a short-three bridge.

Write \(\gamma=-c\cdot n_{14}\), \(Q=\beta^2\).
Weakness at fourteen implies
\[
 {\sigma\over\sqrt{\delta-Q}}\le T(\sqrt{28})<{25\over16}.
\]
Indeed \(\sqrt{28}<5.292\), \(T\) increases,
and the rational comparison at \(5.292\) has positive
cleared numerator \(76368925\).

For isolation, own zero has \(A^2=Q-1<0\).
Own three has the affine star on \(c,n_{14},n_{13},n_{15}\)
with kernel \((1,2,1,\sqrt2)\);
own four has the \(4,4\) null, own six the \(6,3\)
null, and parallel has \(c+n_{14}\).
Every kernel is future feasible against ALL original
walls, regardless of foreign contacts.
For the rest put \(u=(1-Q)/\gamma^2\).
For \(0<u<1\), the actual core
\(c+\beta e+\gamma u n_{14}\) gives
\[
 x\cdot\Omega_{14}>
          \sqrt{28-u}-\sqrt{u(1-u)}.
\]
As \((\sqrt{28}+\sqrt{27})^{-1}<.1\), its loss is
\(<(.1+1.005)/2<.61\).
The \(u=1\) core is null.
For \(Q>0,u>1\), \(Q\ge1/4\),
\(\gamma\ge\varphi/2>.8\), so \(u<75/64\).
The actual finite projection onto \(c,e,n_{14}\)
gives pressure \(>\sqrt{27}-\sqrt{u-1}\),
with loss \(<.1+5/12<.61\).
For pure finite owns at least seven the singleton
projection has loss \(<.1+\tan(\pi/7)<.6\);
one may use \(\tan(\pi/7)<1078/2159<1/2\).
Native five uses the actual LEFT \(H_4\) correction on
\(n_{14},n_{13},n_{12}\), with residual
\((5-3\varphi)/8<.136^2\).
Since \(\gamma>.8\), \(\gamma^{-2}<25/16\),
its loss is \(<5/32+.17<.61\).
Thus all isolated fourteen profiles are strong.

For left-three at \(j=1,13,2,12\),
\[
              \delta=1-{j(14-j)\over28},\qquad
              \sigma=\gamma+j/28.
\]
For left-four at \(j=1,13\),
\(\delta=1/14\), \(\sigma=\gamma+j/(14\sqrt2)\).
Subtract \(Q\) in each denominator.
At left-three \(j\ge2\), own zero has the actual
affine \(B\) kernel
\[
       c+n_{j-1}+2(n_j+\cdots+n_{14})+\sqrt2n_{15}.
\]
At \(j=1\) finite classification permits only
\(\beta=0,1/2,1/\sqrt2\), hence
\(A^2=Q-1/2\le0\).
Own-zero left-four uses the positive affine \(C\)
kernel
\(c+\sqrt2(n_1+\cdots+n_{14})+n_{15}\) at one,
and the local squared-cosine star
\(1/2,1/4,1/4\) at thirteen.
For a left-three, own three has its actual affine cycle
and own four the \(4,4\) null.
At far one, own six has the \(6,3\) null; parallel
always has its null pair.
The remaining far weakness has precisely
\(\gamma\Omega_{14}+\Omega_1/2\) and
\(202/125<2\gamma<12/5\);
Lemma~\ref{lem:literal-helper} excludes it, including \(Q>0\).
At near thirteen every remaining own is at least
\(1/\sqrt2\), with ratio
\(>(.707+13/28)/.732>25/16\).
At next-to-end two/twelve the ratio is
\(>(.7+1/14)/\sqrt{1/7}>25/16\).
All positive-own left-four rows have ratio at least
\((.5+1/(14\sqrt2))\sqrt{14}>25/16\).
Thus every left-only role is strong.

For short-only contact \(\eta=-c\cdot n_{15}>0\),
introduce the actual integral cycles
\[
 B=4\gamma^2,\quad W=4\eta^2,\quad
 T=4\sqrt2\gamma\eta,\qquad
 A_0^2=7B+7T+15W/4-1.
\]
Equation~\eqref{eq:arithBD15Icycle} implies \(WV/4\)
integral, because fifteen is odd.
For short labels three, five, six and seven,
\(W\) is respectively \(1,\varphi^2,3\), or the unit
root of \(X^3-5X^2+6X-1\).
Each is a unit at every dyadic place.
Thus \(V/4\) is integral, impossible.
This argument includes EVERY own, including zero,
and every old-\(e\) contact.
For short four, only \(V/2\) follows.
Own zero has the actual \(4,4\) null.
For own three, normalize the cap to raw square two;
both products with old fourteen and fifteen are \(-1\).
Its rational old-\(B\)/cap principal is physically
Lorentzian, with Schur residual \(-55\).
Hence actual \(K=\mathbb Q\), making \(0<V/2<1\)
an integer, impossible.
Every other positive own has
\[
 {\sigma\over\sqrt{\delta-Q}}\ge
 {\,\gamma+1/2\over\sqrt{1/2}}\ge1+1/\sqrt2>25/16.
\]
For every short label at least eight, even own zero
has ratio at least
\(\cot(\pi/8)/\sqrt2=1+1/\sqrt2>25/16\).
All higher finite and unbounded dotted owns are included.

A short-four bridge has residual \(1/28\) and
\(\sigma=\gamma+j/28+1/2\), \(j=1,13\);
even own zero has ratio \(>2.8\).
For a short-three bridge,
\[
 v=\gamma\Omega_{14}+\Omega_j/2+\Omega_{15}/2 .
\]
Its original cycles make \(a=2\gamma\) integral
and \(\sqrt2\) integral and \(K\)-valued.
Directly
\[
 4A_0^2=M=28a^2+(4j+28\sqrt2)a+2j+11+2j\sqrt2
             \in1+2\mathcal O_K .
\]
Equation~\eqref{eq:arithBD15Icycle} gives \(MV/4\)
integral.
Writing \(M=1+2Z\) and multiplying by \(1-2Z\)
makes \(V/4\) integral.
This is valid with arbitrary dyadic ramification;
no quotient ring is presumed a field.
It includes \(a=0\), both bridge orientations,
and the retained old-\(e\) three arm.

We have proved strong EVERY \(5,\ldots,14\) and
\[
 e_B>E_{14}>.9084,\quad
 b=f_5/5<.361,\quad f_{14}<3.311 .
\tag{BW}\label{eq:arithB15Iwindow}
\]
The genuine reference has individual \(P_5\) price
\(>.14\), terminal \(P_{14}\) price
\(\sqrt{14}-\sqrt{13}>.13\), and \(w_{15}=0\).
Its first-coordinate graft is legal for
\(.49\le X\le.67\): the other first-five coordinates
are \(>.2836>.214\).
Therefore
\[
             x_1>X\quad\Longrightarrow\quad
             e_B>.9084+{5\over4}(X-.361)^2 .
\tag{BG}\label{eq:arithB15Igraft}
\]

\subsection{The actual \(B_{15}\) head in the short-\(I\) branch}
\label{app:arithB15Ihead}

The finite \(B_{14}\) head list and the own-zero/triangle
clearings in Section~\ref{app:arithBD15ordinary} apply
without any ordinary-\(h\) premise.
A nonisolated head would have \(x_1>.643\).
Now~\eqref{eq:arithB15Igraft} alone gives
\(e_B>.9084+.099405>1\), without charging \(I\).
Thus the ACTUAL head is isolated within \(B_{15}\):
\[
 c_H=A_Hp-g_H\Omega_1-\beta_H e,\quad
 A_H^2=2g_H^2+Q_H-1,\quad Q_H=\beta_H^2 .
\]
As before,
\[
 \lambda_H=A_H\sqrt h-\beta_H\sqrt{1+h}\ge0,\qquad
               (2g_H^2-1)h\ge Q_H .
\]
This removes owns zero, three and four.
The actual \(6,3\) and parallel nulls remove six and one.
Pure native five would give
\(0<\lambda_H^2<.31(.5)<1/4\), impossible.

For pure finite own at least seven, use the actual
finite head point and the SINGLE paid reference
of Section~\ref{app:arithBD15ordinary}:
\[
                e_B>E_{14}+F(u)-.8b+.032,\qquad
                F(u)=.8H-.16R>.2893 .
\]
This is not a separately optimized head energy.
Since \(g_H<1\), its positive mutual satisfies
\(\lambda_H<1/\sqrt2\); the native alphabet forces
\(\lambda_H=1/2\).
Thus \(h=[4(2g_H^2-1)]^{-1}>1/4\).
At native seven or eight,
\(h\ge\sqrt2/4\), so \(q(h)>.0755\):
use \(\sqrt2>1.4142\), \(\sqrt{1+h}>1.1634\)
and \(.1634/2.1634>.0755\).
The original total is then \(>1.0164\).
For all finite natives at least nine,
the largest-root cubic gives \(g_H>.939\), hence \(u<1.135\).
The already proved monotonicity of the finite scalar gives
\[
 F(u)>F(1.135)>.302416;
\]
at that rational endpoint,
\(\sqrt{2-u}>.93\), \(\sqrt{(u-1)(4-3u)}<.284\),
\(H>.4565\), \(R<.3924\).
Also \(q(h)>.0557\), since \(\sqrt{1.25}>1.118\).
The single original total is \(>1.009716\).
Thus the only pure head left is DOTTED.

For \(Q_H\ge1/2\), acuteness gives
\(C_H^2\ge2+1/h\).
If \(h<.2\), then \(C_H^2>7\), \(x_1>.67\),
and~\eqref{eq:arithB15Igraft} gives body \(>1.02775125\).
If \(.2\le h<.5\), then \(C_H>2\), \(x_1>.57\),
and \(q(h)>.045\), using \(\sqrt{1.2}>1.095\).
The original total is \(>1.00800125\).
For \(Q_H=1/4,\lambda_H>0\), native gap gives
\(A_H\sqrt h>1\).
For \(h<.2\) this again gives \(C_H^2>7\).
For \(h\ge.2\), use \(\sqrt{1+h}>1.08\) to get
\(A_H\sqrt h>1.04\), so
\[
 C_H^2=1+(4/3)A_H^2>
             1+(4/3)(1.04^2/.5)>1.95^2.
\]
Then \(x_1>.56\); its graft plus \(q(h)>.045\)
makes the total \(>1.00290125\).
These intervals cover all positive foreign labels
and every positive finite/parallel/dotted mutual.

The only mixed possibility left is
\(\beta_H=1/2,\lambda_H=0\).
For \(U=4g_H^2\) its exact relation is
\[
                           V(U-2)=2,\qquad U>3 .
\tag{BQ}\label{eq:arithB15Iquarter}
\]
Initially both finite and dotted own are retained here.

We shall need a point valid for BOTH the pure dotted
and quarter-foreign heads.
Put \(\delta_H=1-Q_H\), \(u=\delta_H/g_H^2\).
For the pure dotted head \(0<u<1\).
For the quarter head,~\eqref{eq:arithB15Iquarter}
and \(0<h<1/2\) give \(g_H^2>3/4\), again \(0<u<1\).
The positive ACTUAL combination
\[
             \xi=c_H+\beta_H e+{\delta_H\over g_H}n_1
\]
has square \(-\delta_H(1-u)<0\), products zero
with \(c_H,e\), and product
\(-g_H+\delta_H/g_H<0\) with \(n_1\).
Its time is \(A_H>0\).
All omitted original walls, including \(d\), have
nonpositive product termwise.
Its old-\(e\) displacement cancels EXACTLY.
Its normalized pressure is the SAME dotted-row
\(Z,H,R\) used in~\eqref{eq:arithBD15purescalar},
now genuinely derived for the mixed head as well.

If this quarter head had finite own, then
\(u>3/4\), \(h>1/4\).
For the dotted-row scalar on \([3/4,1)\),
the derivative comparison reduces to
\(-4u^3+13u^2-10u+2>0\).
This polynomial is \(1/8\) at \(3/4\) and increases;
the remaining derivative term \(.16(1-u)\) is positive.
Hence
\[
 F(u)>F(3/4)={2\sqrt{10}-\sqrt6\over10}-.085>.3024.
\]
Together with \(q(h)>.0557\) and the same paid
head reference, the total exceeds \(1.0097\).
Thus BOTH remaining head families are dotted.

For any actual short-fifteen pressure \(S_{15}>T\),
use the SINGLE vector
\[
                  W=w+.4(Z-P_5/5)+.13e_{15}.
\]
The individual debits are \(.08\) at \(P_5\) and
\(.13\) at \(P_{14}\), paid before using either floor.
The head residual is zero-sum on the first five
coordinates and hence orthogonal to \(w\).
The terminal \(e_{15}\) is orthogonal to both.
The FULL square gives
\[
 e_B>E_{14}+F(u)-.8b+.032
                    +.26(T-f_{14})-.0169 .
\tag{BT}\label{eq:arithB15Ijoint}
\]
If the actual short cosh were at least fifteen,
its original edge would give
\(S_{15}\ge\sqrt{105/8}>3.62\),
and the body alone would exceed
\(.9084+.0325+.06344=1.00434\).
Thus \(C_S<15\).
Also the legal \(X=.632\) head graft would give
\(.9084+(5/4)(.632-.361)^2=1.00020125\);
hence \(x_1<.632\), and the actual head comparison
gives \(C_H<7/3\).
These are necessary ACTUAL edge bounds, not cutoffs
imposed on formal roots.

\subsection{Every actual short-fifteen white profile}
\label{app:arithB15Ishort}

The finite endpoint of the actual short-fifteen edge
retains \(A_{14}\) on one through fourteen and old \(e\).
Finite classification gives exactly SEVEN profiles,
with both endpoint orientations kept:
\[
\begin{array}{c|c|c}
 \text{contact in }A_{14}&v_B&\text{allowed }\beta\\ \hline
 \text{none}&\gamma P_{15}&\text{any finite native or zero}\\
 \text{endpoint three }j=1,14&
       \gamma P_{15}+P_j/\sqrt2&0,1/2,1/\sqrt2\\
 \text{next-to-end three }j=2,13&
       \gamma P_{15}+P_j/\sqrt2&0\\
 \text{endpoint four }j=1,14&
       \gamma P_{15}+P_j&0 .
\end{array}
\]
An extra old-\(e\) arm at a next-to-end branch gives
nonfinite long \(E\) arms.
At a white four it puts the double bond internally
in rank sixteen.
An endpoint-three old-\(e\) arm permits only three
or four, giving \(A_{16}\) or \(B_{16}\).
Thus the displayed foreign list is complete.
The own \(\gamma=-c\cdot n_{15}\) is initially zero,
finite native, parallel or dotted.

Own zero is impossible in every row.
Isolation has \(A^2=Q-1<0\).
Endpoint-three at one has \(A^2=Q-1/2\le0\).
For every other white-three \(j\ge2\), use the actual
affine \(B\) null
\[
              c+n_{j-1}+2(n_j+\cdots+n_{14})+\sqrt2n_{15}.
\]
For white-four at one use
\(c+\sqrt2(n_1+\cdots+n_{14})+n_{15}\);
at fourteen use \(c+\sqrt2n_{14}+n_{15}\).
Every selected product vanishes; its positive cap
coefficient gives future time, and all unselected
original products are nonpositive termwise.
Parallel own has the universal null \(c+n_{15}\).

For positive own the time squares, before adding \(Q\), are
\[
\begin{array}{c|c}
 \text{profile}&A_0^2\\ \hline
 \text{isolated}&15\gamma^2-1\\
 \text{white three},\ a=2\sqrt2\gamma&
               (15a^2+4ja+4j-8)/8\\
 \text{white four},\ t=2\gamma&
               (15t^2+4jt+4j-4)/4 .
\end{array}
\]
These are positive for every \(\gamma\ge1/2\).
Normalize the old \(B\) long roots to square two and
short root to square one.
A white-three cap has square two, white product \(-1\)
and short product \(-a/2\);
a white-four cap has square one, white product \(-1\)
and short product \(-t/2\).
These scalings follow from the actual native bonds over
\(K\); no disconnected scale is guessed.
The old-\(B\)/cap principal is physically Lorentzian
and defined over \(\mathbb Q(a)\) or \(\mathbb Q(t)\).
A proper actual-field extension would fix that
indefinite principal at a nonphysical embedding.
Hence actual \(K\) equals that generating field.
For isolation normalize \(c/\gamma\) to get
actual \(K=\mathbb Q(B)\), \(B=4\gamma^2\), in the same way.

For an isolated finite own,
\(0<B_\sigma<4/15\) at every nonphysical place.
For \(B=2+2\cos(2\pi/m)\), degree \(d=\phi(m)/2\)
and trace \(2d+\mu(m)\) give
\(26d<56-15\mu(m)\), hence \(d\le2\).
The complete finite list is \(m=3,4,5,6,8,10,12\).
Rational cases three/four/six force actual \(K=\mathbb Q\),
excluded by the compact WHOLE-form Meyer obstruction.
For five/eight/ten the other \(B\) values
\[
 (3-\sqrt5)/2,\quad2-\sqrt2,\quad(5-\sqrt5)/2
\]
all exceed \(4/15\).
For twelve the other value \(2-\sqrt3>4/15\),
using \(\sqrt3<26/15\).
Thus no isolated finite own survives, regardless of \(\beta\).

For white three, the actual path makes \(a\) integral;
the own two-cycle makes \(a^2/2\) integral.
Write \(M_j=15a^2+4ja+4j-8\).
The cycle identity gives \(M_jV/8\) integral.
Let \(t=2\cos(\pi/m)\).
The root-of-unity dyadic lemma proved in
Section~\ref{app:arithA15dyadic} says \(t\) is a unit
at every dyadic place unless \(m\) is a power of two.
For a non-power-two label, \(a=\sqrt2t\) has
valuation \(1/2\), normalized by \(v(2)=1\).
The leading term of \(M_j\) has valuation one;
the other terms have valuation at least two.
Thus \(v(M_j)=1\), forcing \(V/4\) integral, impossible.
For \(m=4\), \(a=2\) gives a rational Lorentzian
principal, so actual \(K=\mathbb Q\).
For \(m=2^k\ge8\), put \(B=t^2\).
Its even degree \(d\) and norm two imply
\(\mathrm N(2B)=2^{d+1}\) is nonsquare.
Hence \(\mathbb Q(a)/\mathbb Q(B)\) has the sign
automorphism \(a\mapsto-a\).
At this nonphysical embedding,
\[
                  M_j(-a)\ge15a^2-56a+48>0
\]
because \(a>12/5\), the right polynomial vanishes
at \(12/5\) and increases thereafter.
This contradicts positive definiteness.

For white four the actual native-four path makes
\(2t\) \(K\)-valued, and the own two-cycle makes
\(t^2\) integral, so \(t\in\mathcal O_K\).
For \(M_j=15t^2+4jt+4j-4\), the cycle identity
gives \(M_jV/4\) integral.
A non-power-two own makes \(t\), and hence \(M_j\),
a dyadic unit, again impossible.
For the remaining power-two owns, the near \(j=14\)
short edge already has \(C\ge15/2+14>15\).
At far one the sign automorphism \(t\mapsto-t\)
exists for every power-two \(m\ge4\), and its time
numerator \(15t^2-4t>0\) contradicts admissibility.
Thus EVERY finite own has been removed.

For completeness, the actual cosh with a white contact
of unit magnitude \(\eta\) is
\[
 C_S={\sqrt{15}\gamma+\sqrt2\eta j/\sqrt{15}
                   \over\sqrt{\delta-Q}},\qquad
 \delta=1-{2\eta^2j(15-j)\over15}.
\]
White-four has pure \(C_S=15\gamma+j\);
its dotted range violates \(C_S<15\) at both ends.
White-three at near next-to-end thirteen has pure
\(C_S=15a/4+j/2>17\).
A foreign cost only increases these coshes.

For the remaining far next-to-end three at \(j=2\),
\[
 2\sqrt2<a<56/15,\qquad
 -8/15<a_\sigma<0,\qquad K=\mathbb Q(a),\ a\in\mathcal O_K .
\]
Degree at least four contradicts
\((56/15)(8/15)^3<1\).
In degree three the norm has magnitude \(<2\), hence
is one.
For other roots \(-b,-c\), \(bc=1/a\) and
\(b,c<8/15\) force \(a>225/64\).
The trace is then three, so
\(b+c=a-3<11/15\), whereas
\((b+c)^2\ge4/a>15/14\), impossible.
Degree one is excluded by the whole-form rational obstruction.
Degree two has norm \(-1\), trace three and
\(a=(3+\sqrt{13})/2\).
Its time numerator \(15a^2+8a=53a+15\)
has odd norm \(-199\).
It is a unit at EVERY dyadic place;
\(M_jV/8\) would make \(V/8\) integral.

Only end-three dotted roles remain.
Far one has
\[
 2\sqrt2<a<118/15,\qquad -2/3<a_\sigma<2/5 .
\]
For near fourteen put \(S=a+2\), giving time numerator
\(15S^2-4S-4\) and
\[
 2+2\sqrt2<S<122/15,\qquad -2/5<S_\sigma<2/3 .
\]
Both \(a,S\) are integral and have square divided
by two integral.
If their actual degree \(d\ge3\), their integer norms
are divisible by four, while their magnitudes are
\[
                   <(122/15)(2/3)^{d-1}\le488/135<4,
\]
impossible.
In degree two both trace and norm are EVEN:
\(\mathrm N(z)^2/4\) integral forces even norm, and
\(\operatorname{Tr}(z^2/2)
 =\operatorname{Tr}(z)^2/2-\mathrm N(z)\)
forces even trace.

The far strips allow trace \(4,6,8\), with
\(-6<\mathrm N(a)<4\), nonzero.
The near strips allow trace \(6,8\), with
\(-4<\mathrm N(S)<6\).
The COMPLETE finite hand tables are
\[
\begin{array}{c|ccc}
 \text{far trace}&N=-4&N=-2&N=2\\ \hline
4&2+2\sqrt2\ (\text{wrong conjugate})&2+\sqrt6&
        2+\sqrt2\ (\text{wrong conjugate})\\
6&3+\sqrt{13}&3+\sqrt{11}&3+\sqrt7\\
8&4+2\sqrt5\ (\text{too large})&
        4+3\sqrt2\ (\text{too large})&4+\sqrt{14}
\end{array}
\]
and
\[
\begin{array}{c|ccc}
 \text{near }S\text{ trace}&N=-2&N=2&N=4\\ \hline
6&3+\sqrt{11}&3+\sqrt7&3+\sqrt5\ (\text{wrong conjugate})\\
8&4+3\sqrt2\ (\text{too large})&4+\sqrt{14}&4+2\sqrt3 .
\end{array}
\]
For their nine retained rows, the integral time-numerator
factor \(M_j/2^k\) and its odd norm are
\[
\begin{array}{c|c|c|r}
 \text{role and }a&k&M_j/2^k&\mathrm N(M_j/2^k)\\ \hline
 \text{far }2+\sqrt6&1&77+32\sqrt6&-215\\
 \text{far }3+\sqrt{13}&2&(169+47\sqrt{13})/2&-39\\
 \text{far }3+\sqrt{11}&1&154+47\sqrt{11}&-583\\
 \text{far }3+\sqrt7&1&124+47\sqrt7&-87\\
 \text{far }4+\sqrt{14}&1&231+62\sqrt{14}&-455\\
 \text{near }1+\sqrt{11}&1&142+43\sqrt{11}&-175\\
 \text{near }1+\sqrt7&1&112+43\sqrt7&-399\\
 \text{near }2+\sqrt{14}&1&215+58\sqrt{14}&-871\\
 \text{near }2+2\sqrt3&3&50+29\sqrt3&-23 .
\end{array}
\]
Every factor is a unit at ALL dyadic places,
including ramified places in the actual field.
The seven \(k=1\) rows make \(V/4\) integral.
The \(k=2\) row makes \(V/2\) integral.
In a quadratic field such a \(W=V/2\) has
\(0<W<1\), \(-2<W'<0\), norm \(-1\)
and trace \(-1\) or zero.
Trace zero is reducible or a boundary;
trace \(-1\) gives \(\mathbb Q(\sqrt5)\),
not the actual \(\mathbb Q(\sqrt{13})\).

Only the final near row remains.
Here actual \(K=\mathbb Q(\sqrt3)\) and
\(\mathcal O_K=\mathbb Z[\sqrt3]\).
The full strict \(V\) strips give
\(V\in\{\sqrt3-1,\sqrt3\}\):
writing \(V=a+b\sqrt3\) gives \(0<b<\sqrt3<2\),
so \(b=1\), then \(a=-1,0\).
Hence \(h>.18\), \(q(h)>.04\).
Its pure short cosh is
\(15(2+2\sqrt3)/8+7/2>13\);
foreign \(Q\) only increases it.
Thus \(S_{15}>\sqrt{90/7}>3.58\).
The SAME joint vector~\eqref{eq:arithB15Ijoint},
with \(I\) counted once, gives
\[
                   e_B+e_I>.9084+.0325+.05304+.04
                                =1.03394>1.
\]
Every nonisolated short profile is excluded.
The actual short-fifteen opposite is ISOLATED
within \(B_{15}\) and DOTTED.

\subsection{The quarter-foreign head cannot meet an isolated short cap}
\label{app:arithB15Iquarter}

Suppose \(\beta_H=1/2,\lambda_H=0\).
Let \(U=4g_H^2>4\), so \(V(U-2)=2\), and write the
isolated dotted short cap
\[
 c_S=A_Sp-g_SP_{15}-\beta_Se,\qquad
 B=4g_S^2>4,\quad Q=\beta_S^2.
\]
The actual short edge and the nonphysical principal give
\[
 B<60(1-Q),\qquad
 0<B_\sigma<{4(1-Q_\sigma)\over15}.
\]
All finite native \(\beta_S\), including zero, remain.

The native-three head/old-\(e\) link fixes the actual
\(K\)-root squares:
\[
 r={c_H\over\sqrt2g_H},\quad f={e\over\sqrt2g_H},
 \quad r^2=f^2={2\over U},\quad r\cdot f=-1/U,\quad
 s={c_S\over g_S},\quad s^2=4/B .
\]
The old short product of \(s\) is \(-1\).
Its old-\(e\) coefficient \(q_0\in K\) is nonnegative,
with \(s\cdot f=-2q_0/U\), and
\[
 q_0^2={2UQ\over B}\text{ is a FULL }K\text{-square},\qquad
 r\cdot s=1+q_0/U-J=-D,\quad J>0,\ D\ge0 .
\]
These are actual-root normalizations, not assigned
relative lengths for disconnected old components.

The ENTIRE long old-\(B\) path cycle, the two-cycle,
the head/old-\(e\)/short triangle and the long
old-\(e\)/head/old-\(B\)/short cycle give, respectively,
\[
              G=UBD,\quad L_c=UBD^2/2,\quad
              T_c=q_0BD,\quad P_c=2q_0B
                         \quad\hbox{integral}.
\]
All formulas retain mutual zero.
The two time squares are
\((2U-3)/(2U)\) and \(15+4(Q-1)/B\).
Expanding \(J=D+1+q_0/U\) gives
\[
 4(Bq_0+G+T_c+L_c)
       =28UB-45B+8UQ-16Q-8U+12 .
\tag{BQ2}\label{eq:arithB15Iquarterparity}
\]
Since \(4Bq_0=2P_c\), \(4Q\) is integral and
forty-five is odd, this equality forces \(B/2\)
integral at every dyadic place, with arbitrary
ramification.
For \(Q=0\) its stronger divisibility forces \(B/4\)
integral.
Degree at least three is impossible because
\[
                 0<\mathrm N(B/2)
                       <30(2/15)^{d-1}\le8/15<1 .
\]
Degree one is excluded by compact rational anisotropy.
In degree two, \(Q=0\) likewise gives
\(0<\mathrm N(B/4)<1\).
Only positive \(Q\) remains.

The actual head principal is Lorentzian over \(\mathbb Q(U)\),
so actual \(K=\mathbb Q(U)\).
The earlier actual head bound \(C_H<7/3\) gives
\[
                   4<U<49/6,\qquad 0<U'<3/2 .
\]
From \(V=2/(U-2)\) integral,
\(\mathrm N(U-2)\) is a negative divisor of four,
of magnitude greater than one.
It is \(-2\) or \(-4\), and the trace of \(U\) is
\(5,\ldots,9\).
The COMPLETE surviving irreducible rows are
\[
\begin{array}{c|c|c}
 \mathrm N(U-2)&U&\mathrm N(U)\\ \hline
 -2&3+\sqrt3&6\\
 -2&(7+\sqrt{17})/2&8\\
 -4&(5+\sqrt{17})/2&2\\
 -4&3+\sqrt5&4\\
 -4&4+2\sqrt2&8\\
 -4&(9+\sqrt{41})/2&10 .
\end{array}
\]
For verification, writing \(T=\operatorname{Tr}(U)\)
gives \(\mathrm N(U)=2T-6\) or \(2T-8\).
The other \(T=5,\ldots,9\) rows are rational,
boundaries, or violate \(U'<3/2\).
Thus this is a derived list of actual quadratic
fields, not a prior field restriction.

Put \(Z=B/2\), \(M_0=\mathrm N(Z)\in\mathbb Z_{>0}\).
The full short bounds give \(M_0<4\mathrm N(1-Q)\).
A finite \(\beta_S\) in a quadratic field has
label \(3,4,5,6,8,10,12\).
For labels four, six, eight, ten and twelve, the
right side is respectively at most
\(1,1/4,1/2,1/4,1/4\), and the inequality is strict.
Only
\[
 Q=1/4,\ M_0=1\text{ or }2;\qquad
 Q=\varphi^2/4,\ M_0=1,\ K=\mathbb Q(\sqrt5)
\]
remain; native five has \(4\mathrm N(1-Q)=5/4\).
In both branches \(\mathrm N(Q)=1/16\).
The FULL square \(q_0^2=UQ/Z\) requires
\(\mathrm N(U)/M_0\) to be a rational square.
In the entire six-row table this leaves only
\(\mathbb Q(\sqrt2)\) or \(\mathbb Q(\sqrt{17})\)
with \(M_0=2\), or
\(\mathbb Q(\sqrt5)\), \(U=3+\sqrt5\), \(M_0=1\).

For \(Q=1/4,M_0=2\), the full bounds give
\(20<Z<45/2\), \(0<Z'<1/10\).
Its integer trace is twenty-one or twenty-two,
with discriminant \(433\) or \(476\).
Neither has squareclass two or seventeen.
For \(\mathbb Q(\sqrt5),M_0=1,Q=1/4\),
\[
                   10<Z<45/2,\qquad 0<Z'<1/10.
\]
The trace is \(11,\ldots,22\).
For \(T^2-4\) to be five times a square, only
\(T=12,13,17,18,22\) need inspection;
the quotients are \(28,33,57,64,96\).
Only eighteen survives, giving
\[
 Z=9+4\sqrt5=\varphi^6,\quad
 U=3+\sqrt5=2\varphi^2,\quad
 q_0^2={1\over2\varphi^4}.
\]
This is not a square in \(\mathbb Q(\sqrt5)\),
because two is not a square there.
Finally \(Q=\varphi^2/4,M_0=1\) gives
\(8.29<Z<10.365\), \(0<Z'<1/8\):
use \(Z<30(3-\varphi)/4\),
\(Z'<(2+\varphi)/30\), and
\(1.618<\varphi<1.6181\).
Its only integer traces are nine and ten;
discriminants \(77,96\) are not five times squares.
Every mixed quarter-head branch is excluded.

\subsection{The remaining pure dotted pair: a complete norm argument}
\label{app:arithB15Ipurepair}

The actual remaining head is PURE dotted and the short
cap is isolated-within-\(B\) dotted:
\[
 c_H=A_Hp-\sqrt2g_He_1,\qquad
 c_S=A_Sp-g_SP_{15}-\beta e,\qquad g_H,g_S>1 .
\]
Put \(U=4g_H^2\), \(B=4g_S^2\), \(Q=\beta^2\),
\(\delta=1-Q\).
The original Lorentzian connected principals give
\(K=\mathbb Q(U)=\mathbb Q(B)\).
At all places the complete necessary strips are
\[
 \begin{gathered}
 4<U<98/9,\quad 4<B<60\delta,\quad 0<V<2,\\
 0<U_\sigma<2,\quad
 0<B_\sigma<4\delta_\sigma/15,\quad
 -4<V_\sigma<0,\quad 0<\delta_\sigma\le1 .
 \end{gathered}
\]
All \(U,B,V\) are integral.
Equation~\eqref{eq:arithBD15Icycle} gives
\[
                       UV/2,\ BV/4\in\mathcal O_K.
\]
For example subtracting the integral \(V\) term
from \(V(U/2-1)\) gives the first statement;
for \(V(15B/4-1)\), the odd coefficient fifteen
is invertible at dyadic places and all other
places already allow division by four.
For the nonzero integer \(Y=BV/4\),
\[
             0<Y<30,\qquad |Y_\sigma|<4/15 .
\]
Degree at least four contradicts
\(30(4/15)^3<1\).
The whole-form rational obstruction removes degree one.
Thus actual degree is two or three.
Also \(K=\mathbb Q(V)\): any proper extension would have
a nonphysical embedding fixing positive physical \(V\),
contrary to the strict \(V_\sigma<0\) strip.

The one paid head reference already proves
\[
                         e_B>.9409 .
\]
The same head-plus-short vector gives
\[
                 e_B>.9084+.0325+.26(T-3.311)-.0169
              \quad\hbox{if }S_{15}>T .
\tag{BP}\label{eq:arithB15Ipurebudget}
\]
In particular \(S_{15}>3.61\) forces \(e_B>1.00174\).
There are two uniform dispatches.
If \(V>6/5\), then \(h>.3\), \(q(h)>.065\),
and \(e_B+e_I>1.0059\).
If \(B>225/4\), the ACTUAL isolated short edge has
\[
 C_S^2={15B\over4\delta}\ge15B/4>3375/16>(29/2)^2 .
\]
Its original comparison gives
\(S_{15}>\sqrt{405/31}>3.61\), contradicting
\eqref{eq:arithB15Ipurebudget} without charging \(I\).
These dispatches use one reference and its two distinct
individual anchor prices.

In degree three every conjugate of \(B\) is positive, so
\[
        \mathrm N(B)\in\{1,2,3,4\},\qquad
        \mathrm N(V)\in\{1,\ldots,31\}.
\]
Indeed \(\mathrm N(B)<60(4/15)^2=64/15\),
and \(V\) has two negative conjugates with norm \(<32\).
Integrality of \(BV/4\) requires
\(64\mid\mathrm N(B)\mathrm N(V)\).
The entire finite list forces
\(\mathrm N(B)=4,\mathrm N(V)=16\).
As both nonphysical \(B\) values are \(<4/15\),
the physical \(B>4/(4/15)^2=225/4\).
The high-short-edge dispatch excludes the entire
cubic branch, with every finite \(\beta\) retained.

In degree two write \(K=\mathbb Q(\sqrt D)\),
\(D>1\) squarefree.
In the full quadratic order, the gap between
\(V,V'\) is \(<6\).
For \(D\equiv1\pmod4\), write
\(V=(a+b\sqrt D)/2\), same-parity integers,
\(b>0\); then \(b\sqrt D<6\).
The complete possible \(D\) are \(5,13,17,21,29,33\).
Otherwise \(V=a+b\sqrt D\), \(b>0\), and
\(b\sqrt D<3\), so \(D=2,3,6,7\).
Only \(D=5\), respectively \(D=2\), admits \(b=2\);
all other possible \(b\) equal one.
Imposing BOTH strict strips gives exactly
\[
\begin{array}{c|c|r}
 D&V&-\mathrm N(V)\\ \hline
2&\sqrt2-1,\ \sqrt2,\ 2\sqrt2-1&1,\ 2,\ 7\\
3&\sqrt3-1,\ \sqrt3&2,\ 3\\
5&(\sqrt5-1)/2,\ (\sqrt5+1)/2,\ \sqrt5-1&1,\ 1,\ 4\\
6&\sqrt6-1&5\\
7&\sqrt7-1&6\\
13&(\sqrt{13}-3)/2,\ (\sqrt{13}-1)/2&1,\ 3\\
17&(\sqrt{17}-3)/2,\ (\sqrt{17}-1)/2&2,\ 4\\
21&(\sqrt{21}-3)/2,\ (\sqrt{21}-1)/2&3,\ 5 .
\end{array}
\]
For full coverage: at \(D=5,b=1\), \(a=-1,1\);
at \(b=2\), \(a=-2\).
At \(D=13,17,21\), \(a=-3,-1\).
At \(D=29,33\) the strips would require both
odd \(a\le-3\) and odd \(a\ge-1\), impossible.
At \(D=2,b=1\), \(a=-1,0\), and at \(b=2\),
\(a=-1\).
At \(D=3\), \(a=-1,0\), and at \(D=6,7\),
\(a=-1\).
Thus no coefficient, actual field or strict boundary
has been silently omitted.

Now \(0<\mathrm N(B)<16\).
Writing \(n=-\mathrm N(V)\), integrality of \(BV/4\)
requires \(16\mid\mathrm N(B)n\).
Odd \(n\) is impossible.
For \(n=2,6\), \(\mathrm N(B)=8\);
for \(n=4\), it is \(4,8,12\).
Only SIX \(V\) rows remain:
\[
 \sqrt2,\quad\sqrt3-1,\quad\sqrt5-1,\quad\sqrt7-1,\quad
 (\sqrt{17}-3)/2,\quad(\sqrt{17}-1)/2 .
\]
Four exceed \(6/5\), and are excluded by the
head-plus-\(I\) dispatch.
The entire quadratic branch is reduced to
\[
 K=\mathbb Q(\sqrt3),\ V=\sqrt3-1,\ \mathrm N(B)=8;
 \qquad
 K=\mathbb Q(\sqrt{17}),\
 V=(\sqrt{17}-3)/2,\ \mathrm N(B)=8 .
\]
In \(\mathbb Z[\sqrt3]\), norm eight would require
\(a^2-3b^2=8\), impossible modulo three.
In the full \(\sqrt{17}\) order put
\(B=(T+s\sqrt{17})/2\), same-parity integers, \(s>0\).
Its norm and strict strip give
\[
               30<B<60,\quad31\le T\le60,\quad
               T^2-17s^2=32 .
\]
Modulo seventeen \(T\equiv7,10\).
In that COMPLETE integer interval the only candidates
are \(T=41,44,58\), giving \(s^2=97,112,196\).
Only the last is a square.
Thus \(B=29+7\sqrt{17}>57>225/4\),
and the high-short-edge dispatch is contradictory.
This excludes BOTH quadratic remnants and the entire
pure dotted pair.

Together with the ordinary branch, the full actual
\(B_{15}+A_1\) family is now excluded by the written
arguments, whether the isolated rank-one wall has
ordinary \(h\ge1/2\) or actual short \(0<h<1/2\).
No numerical-output premise, field census or realization
is used in this closure.

\subsection{The odd-\(D\) fork cycle and the actual \(D_{15}\) centre}
\label{app:arithD15Icentre}

Retain the original \(D_{15}\) roots, both signed forks
and the actual short \(I\) opposite.
We use~\eqref{eq:arithBD15Icycle}, with every foreign
\(\beta\) and its radical-cancelling triangle retained.

First, in WHOLE \(D_r+A_1\) with \(r\ge5\) ODD,
an actual long-edge opposite cannot have exactly one
extra \(D\) contact, at a fork with ANY finite native
magnitude \(\eta\).
Indeed for
\(v_D=\gamma\Omega_i+\eta\Omega_f\) put
\[
 U=4\gamma^2,\qquad B=4\eta^2,\qquad Z=4\gamma\eta .
\]
Actual two-cycles and the entire old-three path cycle
make these integral, also at \(\gamma=0\).
The exact inverse products give
\[
                   8A_0^2=4iU+4iZ+rB-8=:N .
\]
At every dyadic place, normalized by \(v(2)=1\),
a finite native \(B\) has \(0\le v(B)\le1\).
To see this in its minimal field, write
\(B=(1+\zeta)(1+\zeta^{-1})\),
\(\mathrm N(B)=\Phi_m(-1)\).
This norm is odd unless \(m\) is a power of two,
when it equals two; the elementary factorization
\(n=\prod_{d\mid n,d>1}\Phi_d(1)\) proves the assertion.
The valuation bound persists in EVERY actual extension.
As \(r\) is odd, \(rB\) is the unique lowest term in \(N\).
Thus \(v(N)=v(B)\le1\).
Integrality of \(VA_0^2\) forces \(v(V)\ge2\)
at every dyadic place, hence \(V/4\) integral,
impossible.
This covers own zero, every finite/parallel/dotted
own and every old-\(I\) contact.
Native four is included: its own two-cycle is essential
and must not be dropped.

At \(D_{15}\), this removes the ONLY weak right-fork
packets at eleven and twelve already classified in
Section~\ref{app:arithD15body}.
Together with the shared bulk, strong \(6,\ldots,12\)
are available without ordinary rank-one energy.
It remains to prove thirteen directly.

Deleting thirteen leaves \(A_{12}\), BOTH forks,
and old \(e\).
Finite classification gives the COMPLETE table:
\[
\begin{array}{c|l}
 \text{left contact}&\text{fork contacts}\\ \hline
 \text{none}&
 \text{none; one arbitrary finite label; }
                 (3,3),(3,4),(3,5)\\
 \text{endpoint three at }1,12&
 \text{none; one fork three/four; both forks three}\\
 \text{next-to-end three at }2,11&\text{none}\\
 \text{endpoint four at }1,12&\text{none}.
\end{array}
\]
Two left contacts form a cycle.
Next-to-end with a new fork has nonfinite unequal
long arms.
With no left contact, two finite fork contacts must
form \(A_3,B_3,H_3\); with an endpoint contact,
finite long \(A/B/D\) completions give precisely the
displayed three/four choices.
Adding old \(e\) can restrict but cannot enlarge this list.
For left coefficient \(\eta_L\) at \(j\), fork
coefficients \(a,b\), and foreign \(Q=\beta^2\), put
\[
 q_L={2\eta_L^2j(13-j)\over13},\quad
 \Delta=1-q_L-a^2-b^2-Q>0,\quad
 \sigma=\gamma+{\eta_Lj\over13}+{a+b\over2}.
\]
The actual inverse formula is
\(A^2=26\sigma^2-\Delta\),
\(C_{13}=\sqrt{26}\sigma/\sqrt\Delta\).
The exact threshold satisfies
\[
 T(\sqrt{26})<{100\over61}-{4490\over48909}<{31\over20},
\]
with final positive cleared numerator \(146719\).
Foreign \(Q\) only increases every direct ratio.

For isolation, own zero has negative time.
Own three has the affine four-leaf-three star on
\(c,n_{13},n_{12},n_{14},n_{15}\);
own four has the \(4,3,3\) star; own six the \(6,3\)
path; parallel its actual null pair.
All apply with arbitrary foreign cost.
Every remaining finite own is excluded arithmetically.
For \(B_0=4\gamma^2\), actual reflection normalization
gives old raw squares two, cap square \(2B_0\) and
marked product \(-B_0\).
The original connected principal has Schur residual
\(B_0(2-13B_0)\).
Thus every nonphysical \(B_0\) obeys
\(0<(B_0)_\sigma<2/13\).
Cyclotomic trace \(2d+\mu(m)\) and physical \(B_0<4\)
force \(d\le2\).
The complete labels are \(3,4,5,6,8,10,12\).
The nonrational other values at five/eight/ten/twelve,
\[
 (3-\sqrt5)/2,\quad2-\sqrt2,\quad
 (5-\sqrt5)/2,\quad2-\sqrt3,
\]
all exceed \(1/4>2/13\).
The rational cases were already null.
For every dotted isolated own, the actual core
\(c+\beta e+\gamma u n_{13}\),
\(u=(1-Q)/\gamma^2\in(0,1)\), gives
\[
 x\cdot\Omega_{13}>
           \sqrt{26-u}-\sqrt{u(1-u)}.
\]
Its full loss is \(<(.1+\sqrt{1+.1^2})/2<.553\).
Every isolated row is therefore strong or impossible.

With a left contact and no fork, far-three own zero
has \(A^2=Q-1/2\), \(Q<7/13\), so its native
costs \(0,1/4,1/2\) give nonpositive time.
Far-four own zero forces \(Q=0\).
Near/next-three own zero has affine \(D\) trunk
coefficients two and cap/previous-node/both-fork
coefficients one.
Near-four own zero has the affine \(4,3,3\) star.
Positive own three has its all-three cycle,
own four the fork star, own six the fork \(6,3\)
path, and parallel its null pair.
The remaining near/next-three owns have
\(\gamma\ge\varphi/2>.8\), with ratios at least
\[
 {\, .8+6/13\over\sqrt{7/13}}>1.7,\qquad
 {\, .8\over\sqrt{2/13}}>2.
\]
Positive endpoint-four has ratio at least
\(\sqrt{13}/2>1.8\).
The only left-only possible weak role is the literal
far helper at thirteen, with
\[
       202/125<2\gamma<2(31/20)\sqrt{7/13}<12/5.
\]
Lemma~\ref{lem:literal-helper} excludes it including
foreign contacts.

One fork and no left contact is excluded by the
odd-\(D\) cycle just proved.
For fork pair \((3,3)\), own zero is the actual
four-cycle and own three the local triangle;
all other positive owns have ratio
\(\ge1+1/\sqrt2>31/20\).
For pair \((3,5)\), even own zero has ratio
\(\varphi^3/2>2.1\).
For \((3,4)\), every positive own has ratio \(>2.2\).
At own zero finite whiteness forces \(Q=0\).
The original four-cycle puts \(\sqrt2\) in \(K\),
and
\[
                      A_0^2=(37+26\sqrt2)/8 .
\]
The numerator has odd norm seventeen, so is a unit
at EVERY dyadic place, also in proper extensions.
The cycle identity forces \(V/4\) integral.

For a left endpoint-three and one fork-three,
\(\Delta\le15/52\), \(\sigma\ge\gamma+15/52\).
Own zero has its actual cycle, own three its triangle;
all other owns have ratio \(>1.8\).
For one fork-four,
\(\Delta\le1/26\), and even own zero has ratio
\((1/26+1/(2\sqrt2))\sqrt{26}>1.9\).
For both forks-three,
\(\Delta\le1/26\), \(\sigma\ge7/13\), ratio \(>2.7\).
Both left and fork orientations are included.
Each affine kernel has positive cap coefficient and
time, with all omitted ORIGINAL walls nonpositive
termwise.  No foreign wall was discarded.
This exhausts the table and proves strong thirteen.

Hence \(6,\ldots,13\) are strong.
Both initial-five references give \(e_D>E_{13}>.88989\).
If five is weak, its curved reference gives the stronger
\(e_D>E_{13}+.7600-.6971>.95279\).
Therefore weak five implies \(h<.21\), since
\(q(.21)=1/21\) already exceeds the remaining budget.
Its literal far-helper inverse entries are \(5,1,1\),
the SAME entries used in Section~\ref{app:arithB15Iwindow}.
The complete minimal-degree and actual-field golden
argument there therefore applies over old \(D_{15}\),
without using a \(B\) rank conclusion.
It forces \(V=2\varphi-2\), \(h>.3\), contrary to
\(h<.21\).
Five is strong.
Fix the genuine reference through \(5,\ldots,13\):
\[
 w^2=E_{13}>.88989,\quad
 b<.361,\quad f_{13}<3.175,\quad
 w_{14}=w_{15}=0.
\]
Its INDIVIDUAL \(P_5\) and \(P_{13}\) prices exceed
\(.14\).  The first-five graft is legal on
\(.49\le X\le.67\), as in the \(B\) calculation.

\subsection{The actual \(D_{15}\) head with short \(I\)}
\label{app:arithD15Ihead}

The COMPLETE actual \(D_{14}\) white head list and
both-fork positive clearing were proved in
Section~\ref{app:arithBD15ordinary}.
For a nonisolated head,
\[
 C_1^2={(2g+1)^2\over1-2Q}
       =1+{2A^2\over1-2Q}.
\]
If \(Q=0\), native gap with \(d\) gives \(A^2h\ge1/4\).
If \(Q>0\), acuteness gives
\(A^2h\ge Q(1+h)\ge1/4\).
For \(h<.05\), \(C_1^2>11>7\), \(x_1>.67\);
the legal graft gives \(e_D>1.00924125\).
For \(h\ge.05\), the complete own-zero and own-three
nulls leave \(g\ge1/\sqrt2\), hence \(x_1>.643\).
Now \(q(h)>.012\), using \(\sqrt{1.05}>1.0245\).
The original total is
\(.88989+.099405+.012=1.001295>1\).
Thus the head is isolated WITHIN \(D_{15}\).

Its form and full foreign split are
\[
 c_H=A_Hp-g_H\Omega_1-\beta_He,\quad
 A_H^2=2g_H^2+Q_H-1,\quad
 Q_H=0,\ 1/4,\text{ or }[1/2,1),
\]
with \((2g_H^2-1)h\ge Q_H\).
The original nulls and native gap remove owns zero,
three, four, six, parallel, and pure five.
For a pure finite head use the SAME proved actual
finite point and the ONE paid head reference:
\[
              e_D>.88989+F(u)-.2888+.032 .
\]
For finite own, native gap forces
\(h=[4(2g_H^2-1)]^{-1}>1/4\).
The complete four native cases are
\[
\begin{array}{c|c|c|c}
 \text{native}&F\text{ lower bound}&q(h)\text{ lower bound}
                         &\text{original total lower bound}\\ \hline
7&.2893&.083&1.00539\\
8&.291984&.0755&1.000574\\
9&.302416&.069&1.004506\\
\ge10&.31690336&.0557&1.00569336 .
\end{array}
\]
Here are the tight fresh scalar details.
For seven, \(g_H<.901\) and
\(4(2(.901)^2-1)(.4)<1\), so \(h>.4\).
For eight, put \(s=\sqrt2\),
\[
 H_8={\sqrt{2s-2}-\sqrt{34s-48}\over s},\quad
 R_8=23-16s,\quad h_8=s/4.
\]
The brackets \(1.4142135<s<1.4142136\) give
\(\sqrt{2s-2}>.91017\),
\(\sqrt{34s-48}<.28856\), \(H_8>.4395\),
\(R_8<.3726\).
Thus \(F_8>.291984\), and
\(\sqrt{1+h_8}>1.1634\) gives \(q(h_8)>.0755\).
For nine, the cubic gives \(.939<g_H<.94\),
\(u<1.135\), \(h>.32\).
The earlier finite scalar at \(1.135\) gives
\(F>.302416\); \(\sqrt{1.32}>1.1488\) gives
\(q(h)>.069\).
For EVERY native at least ten,
\(u\le2-(2/5)\sqrt5<1.106\).
At that rational bound,
\(\sqrt{.894}>.9455\), \(\sqrt{.072292}<.269\),
\(H>.4783\), \(R=.410854\),
giving \(F>.31690336\).
The finite scalar's already proved monotonicity supplies
the entire infinite native tail.
Each table row applies to ONE reference and charges \(I\)
once.

For \(Q_H\ge1/2\), \(C_1^2\ge2+1/h\).
The complete \(h\) split is
\[
\begin{array}{c|c|c|c}
 h\text{ range}&x_1\text{ bound}&q(h)\text{ bound}
                            &\text{total lower bound}\\ \hline
(0,.2)&.67&\text{unused}&1.00924125\\
[.2,.3)&.62&.045&1.01874125\\
[.3,.5)&.57&.065&1.00949125 .
\end{array}
\]
The edge squares are respectively \(>7\),
\(>16/3>(9/4)^2\), and \(>4\).
For \(Q_H=1/4,\lambda_H>0\), native gap gives
\(A_H\sqrt h\ge.5+.5\sqrt{1+h}\).
The same three ranges give, respectively,
\(x_1>.67,.63,.57\), and totals
\(>1.00924125,1.02534125,1.00949125\).
For the middle row use \(A_H\sqrt h>1.04\),
\(C_1^2>49/9\); for the last use
\(A_H\sqrt h>1.07\), \(C_1^2>4\).
The \(I\) bounds follow from
\(\sqrt{1.2}>1.095,\sqrt{1.3}>1.14\).
All grafts lie in the PROVED \(.49,.67\) interval;
these are alternative contradictions, never stacked.

Thus the only remaining heads are pure dotted or
\(\beta_H=1/2,\lambda_H=0,V(U_H-2)=2\).
The second still includes finite own.
For BOTH families \(u=(1-Q_H)/g_H^2\in(0,1)\),
and the actual cleared point
\[
                  c_H+\beta_He+{1-Q_H\over g_H}n_1
\]
is future negative, selected-wall nonpositive and
feasible against EVERY original wall, including both
forks and \(d\).
The detailed computation in
Section~\ref{app:arithB15Ihead} uses only these
first long roots and applies here verbatim.
It gives the pure first-two-coordinate pressure
\(x\cdot Z\ge H\), \(F=.8H-.16R>.2893\).
It does not declare the finite quarter-head dotted.

\subsection{Both signed forks and the complete isolated-cap argument}
\label{app:arithD15Iforkisolated}

The genuine fork duals are
\(\Omega_\pm=(P_{14}\pm e_{15})/\sqrt2\);
their squared norm is \(15/2\).
The ORIGINAL edge comparison gives, independently
of the opposite cap's contact pattern,
\[
 T_\pm=x\cdot(P_{14}\pm e_{15})
       \ge\sqrt{15}\sqrt{(C_\pm-1)/(C_\pm+1)} .
\]
For either ACTUAL signed pressure \(T_\pm>T\), use
\[
          W_\pm=w+.4(Z-P_5/5)
                     +.14(P_{14}\pm e_{15}-P_{13}).
\]
Its individual debits are \(.08\) from \(P_5\)
and \(.14\) from \(P_{13}\), both affordable BEFORE
using their floors.
The head residual is zero-sum on the first five
coordinates.
The signed residual \(e_{14}\pm e_{15}\) has norm
two and is disjoint from the head residual.
Both are orthogonal to \(w\).
The FULL one-vector square therefore gives
\[
 e_D>.88989+.0325+.28(T-3.175)-.0392 .
\tag{DF}\label{eq:arithD15Iforkbudget}
\]
If either \(C_\pm\ge14\), then
\(T_\pm\ge\sqrt{13}>3.605\), making \(e_D>1.00359\).
Thus BOTH actual fork coshes are \(<14\),
without assuming a fork-cap whitelist.

Fix either genuine fork \(f\), and suppose its actual
opposite is isolated within \(D_{15}\):
\[
                  c=Ap-g\Omega_f-\beta e,\quad Q=\beta^2<1.
\]
Own zero has negative time; parallel own has
the actual future feasible null \(c+n_f\).
For every remaining positive own put \(U=4g^2\);
it is integral and
\(A_0^2=15U/8-1>0\).
Here is an explicit actual \(K\)-normalization,
valid for finite as well as dotted own.
Choose old raw roots \(\alpha_i=\sqrt2n_i\).
The actual reflection difference
\[
                       r=s_c(\alpha_f)-\alpha_f=2\sqrt2g\,c
\]
is a \(K\)-vector on the actual cap line.
Its Gram has
\(r^2=2U\), \(r\cdot\alpha_f=-U\), other old products zero.
The connected original principal is defined over
\(\mathbb Q(U)\), with physical Schur residual
\[
                     2U-(15/4)U^2=-2UA_0^2<0.
\]
Any proper actual extension would have a nonphysical
embedding fixing that indefinite principal.
Hence ACTUAL \(K=\mathbb Q(U)\), and every nonphysical
place has \(0<U_\sigma<8/15\).
The actual finite fork endpoint gives
\[
              C_f^2={15U\over8(1-Q)}<196,\qquad
              0<U<1568/15 .
\]
The cycle identity makes \(15UV/8\) integral;
as \(15=16-1\), \(Y=UV/8\) is integral too.
Its full strips are
\[
                 0<Y<392/15,\qquad |Y_\sigma|<4/15.
\]
Degree at least four contradicts
\((392/15)(4/15)^3<1\).
Degree one is excluded by the compact WHOLE-form
rational-anisotropy obstruction.

In degree three,
\(0<\mathrm N(U)<100352/3375<32\) and
\(0<\mathrm N(V)<32\).
Integrality of \(UV/8\) requires
\(512\mid\mathrm N(U)\mathrm N(V)\).
Each positive integer below thirty-two has ordinary
two-adic valuation at most four, so their product
cannot have valuation nine.
This excludes EVERY cubic field and ramification pattern.

In degree two \(0<\mathrm N(U)<12544/225<56\).
The actual \(V\) generates \(K\), and the COMPLETE
quadratic strict-\(I\) list is precisely the sixteen
values already derived in
Section~\ref{app:arithB15Ipurepair};
that derivation used only \(0<V<2,-4<V'<0\),
not a \(B\) geometric hypothesis.
Write \(n=-\mathrm N(V)\).
Now \(64\mid\mathrm N(U)n\).
Odd \(n\) is impossible.
For \(n=2,6\), \(\mathrm N(U)=32\);
for \(n=4\), it is \(16,32,48\).
The same SIX surviving \(V\) values are
\[
 \sqrt2,\quad\sqrt3-1,\quad\sqrt5-1,\quad\sqrt7-1,\quad
 (\sqrt{17}-3)/2,\quad(\sqrt{17}-1)/2.
\]
If \(\mathrm N(U)\ge32\), the strict conjugate strip
gives physical \(U>60\).
Then \(C_f^2>112.5>10.6^2\), and
\(T_f>\sqrt{360/29}>3.52\).
Equation~\eqref{eq:arithD15Iforkbudget} gives
\(e_D>.97979\).
Every displayed \(V\) is \(>14/25\), hence \(h>.14\),
\(\sqrt{1+h}>16/15\), \(q(h)>1/31>.032\).
The original total is \(>1.01179\).
If \(\mathrm N(U)=16\), necessarily \(n=4\),
so \(V=\sqrt5-1\) or \((\sqrt{17}-1)/2\),
both \(>6/5\).
Thus \(q(h)>.065\).
Also \(U>30\), \(C_f>15/2\),
\(T_f>\sqrt{195/17}>3.38\).
The SAME signed reference gives \(e_D>.94059\)
and total \(>1.00559\).
No head or fork energy optimized separately has
been added to this reference.

Thus neither ACTUAL fork opposite can be isolated
within \(D_{15}\), with its entire finite old-\(I\)
contact and every own range retained.
The remaining nonisolated cases are now classified
at their own actual finite endpoints.

\subsection{Every nonisolated fork cap}
\label{app:arithD15Iforknonisolated}

Fix fork \(f\) and denote the other fork by \(k\).
Deleting \(f\) retains the ACTUAL \(A_{14}\) path
\(n_1,\ldots,n_{13},n_k\), plus old \(e\).
The COMPLETE finite contact and foreign table is
\[
\begin{array}{c|c}
 \text{retained }D\text{ contact}&\text{allowed }\beta\\ \hline
 \text{simple three at }n_1\text{ or }n_k&0,1/2,1/\sqrt2\\
 \text{simple three at }n_2\text{ or }n_{13}&0\\
 \text{native four at }n_1\text{ or }n_k&0.
\end{array}
\]
Two path contacts form a cycle.
Other branches would need finite exceptional rank
at most eight, while this completion has rank fifteen.
A next-to-end old-\(I\) arm creates nonfinite long
\(E\) arms.
An old-\(I\) arm at white four makes its double
bond internal.
At endpoint three only old-\(I\) labels three
or four give finite \(A_{16}/B_{16}\).
This retains BOTH path ends and next-to-end sites.

Write \(v_D=g\Omega_f+\eta\Omega_j\).
Own zero is completely dispatched as follows.
Far-three at \(n_1\) has \(A^2=Q-1/2\le0\);
next-three at \(n_2\) has \(A^2=0\).
Near-next-three at \(n_{13}\) has the actual affine
\(D_4\) null
\[
                       c+n_{12}+n_f+n_k+2n_{13}.
\]
Near-end-three at \(n_k\) has affine \(E_8\) kernel
\[
 2c+4n_k+6n_{13}+3n_f+5n_{12}
                 +4n_{11}+3n_{10}+2n_9+n_8 .
\]
Its arms have lengths \(1,2,5\).
Far-four at one and near-four at \(k\) have respectively
\[
 \sqrt2c+2(n_1+\cdots+n_{13})+n_f+n_k,\qquad
 \sqrt2c+2n_k+2n_{13}+n_f+n_{12}.
\]
Every selected product vanishes; its cap coefficient
gives positive future time.
Every unselected ORIGINAL wall, including the entire
old-\(I\) arm, has nonpositive product termwise.
Thus these are forbidden feasible nulls.
Parallel own has \(c+n_f\).

For white three put \(t=2g\).
The actual closed three-path cycle makes \(t\)
integral and \(K\)-valued.
Actual reflection differences normalize cap and old
\(D\) roots to square two, with own product \(-t\)
and white product \(-1\).
The physically Lorentzian old-\(D\)/cap principal
then proves ACTUAL \(K=\mathbb Q(t)\), with time
numerators
\[
 \begin{array}{c|c}
 j=1,2,13&8A_0^2=15t^2+4jt+4j-8\\
 j=k&8A_0^2=15t^2+26t+7 .
 \end{array}
\]
For white four put \(a=2\sqrt2g\).
The original reflection difference at its white-four
neighbour equals \(2c\), fixing the cap's \(K\)-line.
Use raw cap square one, white product \(-1\)
and own product \(-a/2\).
The actual cycles give \(a\in\mathcal O_K\),
\(a^2/2\in\mathcal O_K\), and actual \(K=\mathbb Q(a)\).
Now
\[
 \begin{array}{c|c}
 j=1&16A_0^2=15a^2+8a\\
 j=k&16A_0^2=15a^2+52a+44 .
 \end{array}
\]
All these \(A_0^2\) are positive for positive native
or dotted own, so the actual-field descent uses
an original Lorentzian principal in its precise domain.

For white three with own three, the original old-three
path closes an all-three affine cycle; its all-ones
kernel remains future feasible with foreign contacts.
For every other FINITE own, the dyadic lemma of
Section~\ref{app:arithA15dyadic} gives
\(v(t),v(t+1)\le1/2\) at EVERY dyadic place.
In the long-\(j\) numerator the term \(15t^2\)
has unique smallest valuation \(\le1\).
At \(j=k\), putting \(s=t+1\) gives
\(15s^2-4s-4\), with the same unique-lowest-term
argument.
Integrality of \(VA_0^2\) forces \(v(V)\ge2\)
at all dyadic places, contrary to \(V/4\) nonintegral.

For white four and a non-power-two finite own,
\(t=2\cos(\pi/m)\) is a dyadic unit,
so \(v(a)=1/2\).
The leading term in either numerator has unique
valuation one; the others have valuation at least two.
The cycle identity makes \(V/8\) integral, impossible.
For \(m=4\), \(a=2\) gives a rational Lorentzian
principal, so actual \(K=\mathbb Q\), impossible.
For higher powers of two, the SAME norm-two
half-square argument used in
Section~\ref{app:arithB15Ishort} gives an
automorphism \(a\mapsto-a\).
The far conjugate numerator \(15a^2-8a\) is positive.
The near numerator \(15a^2-52a+44\) has roots
\(22/15,2\) and is positive since
\(a^2\ge4+2\sqrt2>6>(12/5)^2\).
Both contradict strict nonphysical positivity.
Every finite own is now excluded.

Only \(g>1\) remains.
The complete pure fork coshes, which only increase
on adding foreign \(Q\), are
\[
\begin{array}{c|c}
 \text{white contact}&C_f\text{ at }Q=0\\ \hline
 n_1\text{ three}&(15t+2)/8\\
 n_k\text{ three}&(15t+13)/8\\
 n_2\text{ three}&(15t+4)/4\\
 n_{13}\text{ three}&(15t+26)/4\\
 n_1\text{ four}&15a/4+1\\
 n_k\text{ four}&15a/4+13/2 .
\end{array}
\]
They follow from
\(\Omega_f^2=15/2\),
\(\Omega_f\cdot\Omega_k=13/2\),
\(\Omega_f\cdot\Omega_j=j\),
and retained-path inverse costs
\(7/15,13/15,14/15\).
Since \(C_f<14\), \(t>2\), \(a>2\sqrt2\),
near-next-three and near-four are immediately impossible.

Far-four has
\[
            2\sqrt2<a<52/15,\qquad -8/15<a_\sigma<0.
\]
Degree at least three contradicts
\((52/15)(8/15)^2=3328/3375<1\).
In degree two the norm magnitude is \(<2\)
but must be EVEN since \(a^2/2\) is integral.
In degree one an even integer in the physical
interval does not exist.
Far-next-three has \(t=2g\) in the FULL dotted strip
\(2<t<52/15\), with \(-8/15<t_\sigma<0\)
at every nonphysical place.
No half-square condition or \(2\sqrt2\) lower bound
is imposed on \(t\).
Degree at least three is again impossible;
degree one is rational, and degree two has norm
\(-1\), trace two or three.
Thus \(t=1+\sqrt2\) or \((3+\sqrt{13})/2\).
Their time numerators \(15t^2+8t\) have odd norms
\(-79,-199\), so are units at all dyadic places.
The cycle identity makes \(V/8\) integral.
Both next-three and both white-four dotted roles
are now exhausted.

The only roles left are endpoint-three dotted.
For the far one put \(z=t\), and for the near one
put \(z=t+1\).
The COMPLETE actual strips and time numerators are
\[
\begin{array}{c|c|c|c}
 \text{role}&\text{physical }z&\text{every other }z_\sigma&M=8A_0^2\\ \hline
 \text{far}&2<z<22/3&-2/3<z_\sigma<2/5&15z^2+4z-4\\
 \text{near}&3<z<38/5&-2/5<z_\sigma<2/3&15z^2-4z-4 .
\end{array}
\]
Here actual \(K=\mathbb Q(z)\), \(z\) is integral,
and \(MV/8\) is integral.
If \(\mathrm N(z)\) were odd, \(z\) and \(M\)
would be units at every dyadic place, forcing
\(V/8\) integral.
Thus its nonzero integer norm is EVEN.

We need a small ideal argument to exclude absolute
norm two WITHOUT assuming a half-square or unramified
dyadic primes.
If \(|\mathrm N(z)|=2\), then \((z)=\mathfrak p\)
for one prime of norm two, with exponent one.
Let its ramification index be \(e\).
For \(e\ge2\), \(v_{\mathfrak p}(z)=1/e\);
the leading \(15z^2\) in \(M\) has unique smallest
valuation \(2/e\le1\).
Thus \(v_{\mathfrak p}(V)\ge3-2/e\ge2\).
At all other dyadic primes \(z\), hence \(M\),
is a unit, giving \(v(V)\ge3\).
So \(V/4\) is integral.
For \(e=1\), the other dyadic primes have
total ramification-residue degree \(d-1\).
At each, \(M\) is a unit and
\(\operatorname{ord}_{\mathfrak q}(V)\ge3e_{\mathfrak q}\).
Hence
\[
             2^{3(d-1)}\mid |\mathrm N(V)|,
 \qquad 0<|\mathrm N(V)|<2\,4^{d-1}=2^{2d-1}.
\]
For \(d\ge2\), \(3(d-1)\ge2d-1\), a contradiction.
This proves the norm-two exclusion in EVERY
ramification pattern using the Dedekind ideal norm
and \(\sum e_{\mathfrak p}f_{\mathfrak p}=d\).

For \(d\ge5\), the full strips give
\(|\mathrm N(z)|<(38/5)(2/3)^4<2\),
contrary to nonzero even norm.
For \(d=3,4\) its magnitude is \(<152/45<4\);
only two could remain, already excluded.
Degree one is rational and impossible.
Thus actual degree is two, with
\(|\mathrm N(z)|=4\).
Its trace is EVEN as well.
If odd, \(z\) cannot be nonunit at every dyadic
prime: conjugation permutes those primes, making
both \(z,z'\) nonunit there and their integer
trace even.
Norm four also ensures a nonunit prime exists.
A quadratic field would therefore need TWO
dyadic primes, both with \(e=f=1\), one unit
and one nonunit for \(z\).
At the unit prime \(M\) is a unit, so eight
divides \(|\mathrm N(V)|\), contrary to its strict
quadratic bound \(<8\).

For far \(z\), positive norm four is impossible
since \(zz'<(22/3)(2/5)<4\).
Thus norm is \(-4\); \(z'>-2/3\) gives \(z>6\).
Its even trace is six, so
\[
 z=3+\sqrt{13},\qquad
 M=4(169+47\sqrt{13})/2 .
\]
The displayed factor in the FULL quadratic order
has odd norm \(-39\).
Thus \(MV/8\) integral makes \(V/2\) integral.
As proved in the short-\(B\) analysis, an integral
\(W=V/2\) with \(0<W<1,-2<W'<0\) must
generate \(\mathbb Q(\sqrt5)\), not the actual
\(\mathbb Q(\sqrt{13})\).
No far endpoint-three dotted cap exists.

For near \(z\), negative norm \(-4\) is impossible
since \(|zz'|<(38/5)(2/5)<4\).
Norm is four, \(z'>0\), \(z>6\);
its even trace is eight.
Thus
\[
 z=4+2\sqrt3,\quad t=3+2\sqrt3,\quad
 M=8(50+29\sqrt3),\quad K=\mathbb Q(\sqrt3).
\]
The full \(V\) strips in \(\mathbb Z[\sqrt3]\)
give \(V=\sqrt3-1\) or \(\sqrt3\), so
\(h>.18\), \(q(h)>.04\).
Its actual pure fork cosh is
\[
                  (15t+13)/8=(29+15\sqrt3)/4>13 .
\]
Any foreign \(Q\) increases it.
The original signed edge gives
\(T_f>\sqrt{90/7}>3.58\).
The SINGLE reference~\eqref{eq:arithD15Iforkbudget}
gives \(e_D>.99659\) and
\(e_D+e_I>1.03659\).
This excludes the LAST endpoint-three dotted role.

Every actual fork opposite is either isolated or
in the complete nonisolated table.
Both classes have now been contradicted for either
signed fork, both remaining head families and
every finite foreign old-\(I\) contact.
This proves the short-\(I\) \(D_{15}\) branch.
With its separately written ordinary argument and
the unconditional actual rank-one isolation, the
whole \(D_{15}+A_1\) family is covered for every \(h>0\).
There is no field-census output premise and no
new global bound asserted by this unintegrated draft.

\section{Actual typed-seven spectators beside \(B_{14}\) and \(D_{14}\)}
\label{app:arithBD14typed}

Here the ENTIRE old forest is \(B_{14}+Z\) or
\(D_{14}+Z\), where \(Z\) is an ACTUAL \(N_7\)
or \(O_7\) profile.
We do not apply these conclusions to a bare old
seven bond.
The actual defect field descent in
Section~\ref{subsec:shared-seals} gives
\(K=K_7=\mathbb Q(\theta)\), where
\[
 \theta=2\cos(2\pi/7),\quad
 \theta^3+\theta^2-2\theta-1=0,\quad
 4g_7^2=2+\theta,\quad g_7=\cos(\pi/7).
\]
Indeed the ACTUAL old-plus-defect-wall \([7,3]\)
principal is Lorentzian over this cubic field;
any nonphysical embedding fixing it would preserve
its indefinite Gram.
The normalized defect wall \(c_T\) is orthogonal
to the WHOLE classical component, with time squares
\[
                     h_N=(\theta-1)/3,\qquad
                     h_O={\theta-1\over2-\theta}.
\]
Keep the one original decomposition \(z=E(p-x-y)\),
where \(x\) is classical and \(y\) is the entire
old rank-two projection.
Then
\[
                   e_T+e_Z=\|x\|^2+\|y\|^2<1,\qquad e_Z>.09 .
\]
The already proved actual-defect geometry and
propagation are reused at their precise whole-component
domain; no duplicate propagation theorem is introduced.

The three cubic roots lie in
\[
              (31/25,5/4),\quad(-1/2,0),\quad(-2,-9/5).
\]
The polynomial signs at those six endpoints are
\(-559/15625,1/64,1/8,-1,-1,1/125\).
Thus \(0<h_N<1/12\), \(0<h_O<1/3\).
The order discriminant is forty-nine, so
\([\mathcal O_K:\mathbb Z[\theta]]\) is odd and
\(K_7\) is unramified at three.

For an opposite retaining the whole old spectator,
let \(Q\) be its FULL inverse unit-Gram cost.
The COMPLETE possible list is
\[
                       O_7:\ Q=0;\qquad
                       N_7:\ Q=0,1/3,2/3 .
\tag{TQ}\label{eq:arithBD14foreign}
\]
Even one old-\(I_2(7)\) contact costs
\(1/(4\sin^2(\pi/7))>1\).
For old \(A_2\), two contacts contain the affine
all-three triangle or a stronger nonpositive principal.
A single contact permits only labels three/four/five,
with costs \(1/3,2/3,\varphi^2/3\).
Five is incompatible with the cubic actual field.
These are whole-spectator costs, not isolated
bond estimates.

\subsection{A common actual head argument}
\label{app:arithBD14typedhead}

The fresh rank-fourteen prefix gives strong
\(6,\ldots,12\) for \(B\), and \(6,\ldots,11\)
for \(D\).
The \(B\) conclusion uses its actual B3/F4 and
B2 sign-conjugate rows, not a larger-rank tail.
For \(D\), the actual D4/D3 sites are ten and eleven;
their only extra weak fork packets have the actual
left-\(H_4\) ONE-reference body bounds
\(>.92531\) and \(>.94697\).
Adding the separate spectator \(>.09\) excludes them.
The exact local inventory and reference proof are
in Proposition~\ref{prop:f16}; their terminal
indices remain ten and eleven here.
No floor at the unclassified \(D\) centre twelve
is assumed.

If five is weak, its actual pure point gives
\(x_1>.5965\).
Interpolate through that pressure, strong six and
every later genuine strong knot.
The first-six energy is \(>.7600\), versus \(E_6<.6971\),
with middle slope \(>.2843\), above the next \(<.197\).
Even the weaker \(D\) total is
\(.84816+.0629+.09=1.00106>1\).
Five is therefore strong BEFORE selecting the head.
The positive reference \(w\), zero after \(J\), has
\[
 \begin{array}{c|c|c}
 T&J&w^2=E_J\\ \hline
 B&12&>.86989\\
 D&11&>.84816
 \end{array},
 \qquad b=f_5/5<.361,\qquad
 \text{individual }P_5\text{ coefficient}>.14 .
\]
For \(X=.64,.643\), the first-five graft is legal,
preserves every later knot and gives gain
\(>(5/4)(X-.361)^2\).

Deleting the actual long-head direction retains
\(B_{13}\) or \(D_{13}\) and the ENTIRE spectator.
The COMPLETE classical white list is no contact,
or simple three at its long endpoint two.
In the latter branch
\[
 c=Ap-g\omega_1-\omega_2/2-v_Z,\quad
 A^2=2g^2+2g+Q,\quad
 C_1={2g+1\over\sqrt{1-2Q}},\quad Q<1/2 .
\]
Own zero at \(Q=0\) has zero time;
at \(Q=1/3\) it has the actual future feasible null
\(c+\omega_2/2=\sqrt Q\,p-v_Z\).
Its original clearing is
\[
 {\omega_2\over2}=
 \begin{cases}
 n_1+2(n_2+\cdots+n_{13})+\sqrt2n_{14},&B,\\
 n_1+2(n_2+\cdots+n_{12})+n_{13}+n_{14},&D .
 \end{cases}
\]
The foreign products are nonpositive because its
spatial foreign part is the negative of a positive
inverse-Gram combination.
All omitted ORIGINAL wall products are nonpositive
termwise.
Own three has its actual affine triangle.
Every other positive own gives
\(C_1\ge1+\sqrt2\), \(x_1>.643\).
Its legal graft adds \(>.099405\), already exceeding
one with \(e_Z>.09\), even at \(E_{11}\).
The head is isolated WITHIN its entire classical
component, while foreign contacts are retained.

Write the remaining head
\[
 c=Ap-g\omega_1-v_Z,\quad
 A^2=2g^2+Q-1,\quad
 C_1^2={2g^2\over1-Q},\quad U=4g^2\in\mathcal O_K.
\]
Own zero has negative time.
Own three has negative time at \(Q=0,1/3\);
at \(Q=2/3\), the ACTUAL chain
\(n_2,n_1,c,z_1,z_2\), labels \(3,3,4,3\),
has affine kernel \((1,2,3,2\sqrt2,\sqrt2)\).
Own four at \(Q=0\) has zero time;
at \(Q=1/3\), the \(3,4,3,3\) chain has kernel
\((\sqrt2,2\sqrt2,3,2,1)\);
at \(Q=2/3\), \(n_1,c,z_1\) is the \(4,4\) null.
Own five is absent in \(K_7\); own six has the
\(6,3\) kernel \((\sqrt3,2,1)\);
parallel has \(c+n_1\).
Every affine combination has positive cap coefficient,
positive future time and nonpositive product with
ALL omitted original walls.
Thus only finite owns at least seven or dotted remain.

The COMPLETE finite native invariant alphabet in \(K_7\) is
\[
           4\cos^2(\pi/m)\in
             \{1,2,3,2+\theta,1+\theta+\theta^2\},
 \qquad m=3,4,6,7,14 .
\tag{TA}\label{eq:arithK7native}
\]
Indeed \(\phi(m)/2\) divides three, giving the
complete totient list \(3,4,6,7,9,14,18\).
The nine/eighteen field is ramified at three:
its polynomial \(X^3-3X+1\), translated by minus one,
becomes Eisenstein \(X^3-3X^2+3\).
It cannot equal \(K_7\), whose discriminant divides
forty-nine.
The last invariant follows from
\(2\cos(\pi/7)=\theta^2+\theta-1\).
Native four is retained; an unsquared \(\sqrt2\)
need not belong to \(K\) when root-line squareclasses differ.

For \(Q>0\), a full nonphysical Schur complement
requires \(U_\sigma/2+Q<1\).
Normalize the head to raw square \(2U\), with own
product \(-U\) against an old square-two long root;
the full Schur complement is \(2U(1-U/2-Q)\).
This compares no disconnected root scales.
The remaining finite own seven has a conjugate
\(U_\sigma>3/2\), contradicting \(Q\ge1/3\).
Fourteen has one \(U_\sigma>61/25>2\).
Every finite mixed head is excluded.

For a dotted mixed head with \(Q=2/3\),
\(C_1^2>6\), so \(C_1>12/5\), \(x_1>.64\).
For \(Q=1/3\) write \(v_Z=\omega_b/2\).
The actual \(N_7\) wall gives
\[
 c\cdot c_T=q-A\sqrt{h_N},\qquad
 q=2g_7/3\text{ in the same old }A_2\text{ direction},
 \quad q=g_7/3\text{ in the other}.
\]
In the first orientation acuteness yields
\(A^2h_N\ge q^2>16/45\), so \(A^2>64/15\).
The CORRECT identity is
\[
                         C_1^2={3\over2}(A^2+2/3)>37/5 .
\]
In the other orientation \(g>1\) gives
\[
 A^2h_N-q^2>{4\over3}h_N-{2+\theta\over36}
                  ={5\theta-6\over12}>0 .
\]
The product is strictly negative.
Its actual native gap forces
\(A\sqrt{h_N}\ge q+1/2>.8\), hence
\(A^2>192/25\), \(C_1^2>313/25\).
Both orientations give \(x_1>.64\).
The one legal graft and \(e_Z>.09\) contradict
even the \(D\) total:
\[
                  .84816+.09+{5\over4}(.64-.361)^2>1.03546 .
\]
Every remaining head is PURE.

For a pure head its product with the actual defect
wall is \(-A\sqrt{h_T}<0\).
Their actual two-cycle gives
\[
              M=4A^2h_T=2(U-2)h_T\in\mathcal O_K,\qquad M\ge1 .
\]
For \(N_7\), \(h_N<1/12\) forces \(U>8\).
For \(O_7\), \(M\ge4\) likewise forces \(U>8\).
For its finite mutual alphabet, use
\eqref{eq:arithK7native}.
Put \(k=\theta^2+2\theta-1\).
The cubic gives \(1/h_O=k\), so \(U=2+Mk/2\).
The ODD order index means integrality of \(Mk/2\)
requires its three polynomial coefficients to be even.
The COMPLETE table is
\[
\begin{array}{c|rrr}
 M&[1]Mk&[\theta]Mk&[\theta^2]Mk\\ \hline
 1&-1&2&1\\
 2&-2&4&2\\
 3&-3&6&3\\
 2+\theta&-1&5&3\\
 1+\theta+\theta^2&1&6&2 .
\end{array}
\]
Only \(M=2\) remains.
Then \(U=(1+\theta)^2\) and
\[
             .79<64/81<u=4/U<625/784<.8 .
\]
All pure heads, including every original finite own,
have thus reduced to DOTTED owns with complete domains
\[
                            0<u<1/2
                   \quad\hbox{or}\quad .79<u<.8 .
\]

The actual positive point \(c+g^{-1}n_1\) has square
\(-(1-u)\), positive time, product zero with \(c\),
negative product with \(n_1\), and nonpositive product
with EVERY other original wall, including \(c_T\).
It gives the same dotted-row \(Z,H,R\) pressure as before.
On these TWO full intervals its scalar improves to
\[
                             F=.8H-.16R>.3055 .
\]
For \(u\le1/2\), use \(a=2\sqrt2/5\),
\(.56568<a<.5657\), the chord
\(\sqrt{2-u}>1.4142-.379u\) and \(R\le1-.75u\).
The circle inequality gives
\[
 F>.639984656-
       {\, .0944+\sqrt{.0944^2+.5657^2}\over2}
       >.305984656,
\]
since \(.0944^2+.5657^2<.5736^2\).
For \(.79<u<.8\), direct positive squares give
\(\sqrt{2-u}>1.0954\),
\(\sqrt{u(1-u)}<.4074\),
\(H>.4864\), \(R<.52205\).
Thus \(F>.305592\).
Both continuous domains, including unbounded dotted
own, are covered.

Pay ONLY \(.08\) from the individual \(P_5\) price:
\(Y_H=Z-P_5/5\), \(W=w+.4Y_H\).
The residual is zero-sum on the first five coordinates,
has FULL cross zero with \(w\), and norm \(R-1/5\).
The effective gain is
\[
                        F-.8b+.032>.0487 .
\tag{TG}\label{eq:arithBD14typedgain}
\]
For \(B_{14}\), adding spectator ONCE already gives
\(.86989+.0487+.09=1.00859>1\).
This completes BOTH typed \(B_{14}\) cases.
For \(D_{14}\), the weaker \(E_{11}\) reference
still needs its own centre-twelve proof.

\subsection{The full \(D_{14}\) centre at its actual rank}
\label{app:arithD14typedcentre}

Use actual roots
\(n_i=(e_i-e_{i+1})/\sqrt2\), \(i\le13\),
\(n_{14}=(e_{13}+e_{14})/\sqrt2\).
The centre is twelve, forks thirteen and fourteen.
The signed coordinate \(x_{14}\) is retained.
At this point ONLY \(5,\ldots,11\) are strong;
\(E_{12}\) is a scalar value, not yet an available floor.
The reference through eleven has individual terminal
price \(\sqrt{11}-\sqrt{10}>1/7\),
\(f_{11}<2.886\), and the head residual \(Y_H\)
is supported on the first five coordinates.

Deleting twelve retains \(A_{11}\), both genuine
fork \(A_1\) roots and the entire old spectator.
The COMPLETE finite-white list is:
without forks, isolation, simple three at
\(j=1,2,10,11\), or endpoint four at \(j=1,11\);
with left contact and forks, only endpoint-three
at one/eleven with one fork-three/four or both
forks-three;
without left contact, one arbitrary finite fork
label, or fork pairs \((3,3),(3,4),(3,5)\).
Five is absent in \(K_7\).
Two left contacts or two contacts within old \(A_2\)
form a nonpositive cycle.
Next-to-end plus a fork has nonfinite arms \(1,2,9\);
a four internal to a long path is nonfinite.
Every left/fork bridge is foreign-sealed:
an extra old-\(A_2\) arm makes arms \(11,1,2\)
or adds another branch to a full \(D_{14}\) tree.
All other foreign contacts retain
\eqref{eq:arithBD14foreign}.

For the actual opposite write
\[
 c=Ap-g\omega_{12}-\eta\omega_j
                   -\delta\omega_{13}-\epsilon\omega_{14}-v_Z .
\]
A missing contact has coefficient zero.
The exact actual Schur and outgoing data are
\[
 \begin{gathered}
 q_L={2\eta^2j(12-j)\over12},\quad
 \Delta=1-q_L-\delta^2-\epsilon^2-Q>0,\\
 \sigma=g+\eta j/12+(\delta+\epsilon)/2,\qquad
 C_{12}=\sqrt{24}\sigma/\sqrt\Delta .
 \end{gathered}
\]
The \(.61\) sufficient ratio is \(<25/16\).
For \(a=\sqrt{24}\), \(4.8<a<5\), its exact
expression is
\[
 {\,2-2(.61)/a+(.61/a)^2\over1.22-.3721/a}
                 <{1.773\over1.142}<{25\over16}.
\]
Thus any ratio above \(25/16\) proves a true
strong-twelve floor.

For no fork and a left contact, own-zero far-three
would require \(1/2<Q<13/24\), next-to-end-three
would require \(0<Q<1/6\), and far-four would
require \(0<Q<1/12\).
None occurs in the COMPLETE foreign list; the latter
two have zero time at \(Q=0\).
Near/next-three own zero has the affine double-fork
principal: preceding old node, path to twelve,
both forks and cap.
Near-four has the actual kernel
\((\sqrt2,2,2,1,1)\) on
\(c,n_{11},n_{12},n_{13},n_{14}\).
Own three with white three has the actual cycle;
own four the star at twelve with leaves cap and
both forks; own six the \(6,3\) path to a fork;
parallel the null pair.
Own five is absent.
Positive endpoint-four ratios are at least
\(\sqrt{12}/2+1/\sqrt{24}>25/16\);
remaining next-to-end-three ratios exceed \(2.4\),
and near-three ratios exceed \(1.84\).
A possible weak far-three is precisely
\(g\omega_{12}+\omega_1/2\), with
\(.9<g<6/5\).
It is excluded by Lemma~\ref{lem:literal-helper}
at twelve, which includes arbitrary foreign contacts.

For no fork and isolation, own zero has negative time.
Own three has the actual four-leaf-three star;
four/six/parallel have the preceding star/path/null,
and five is absent.
For \(Q>0\) after these removals, or for dotted
\(Q=0\), put \(u=(1-Q)/g^2<1\).
The original positive clearing
\(c+v_Z+gu n_{12}\) is future negative and gives
\[
 x\cdot(\omega_{12}-un_{12})
                \ge\sqrt{24-u}-\sqrt{u(1-u)} .
\]
Here \(v_Z\) is a positive inverse-Gram combination
of ACTUAL old spectator normals.
At \(u=1\) this is a forbidden feasible null.
For \(0<u<1\), the loss below \(\sqrt{24}\) is
\(<.104u+\sqrt{u(1-u)}<.555<.61\);
use \(\sqrt{24}+\sqrt{23}>9.69\) and
\(\sqrt{1+.104^2}<1.006\).
Because \(x\cdot n_{12}>0\), this proves strength.
For every pure finite own at least seven,
\(1<u=g^{-2}<4/3\), use instead
\[
              \xi=c+g^{-1}n_{12}+2(g^{-1}-g)n_{11}.
\]
It has square \(-g^2(7u-4-3u^2)<0\), products zero
with cap and twelve, negative with eleven and both
forks, positive future time, and all omitted
original-wall products nonpositive termwise.
Its displacement is
\(\omega_{12}-un_{12}-2(u-1)n_{11}\);
the pressure is
\(\sqrt{24-u}-\sqrt{7u-4-3u^2}\).
The first radical loses \(<.139\), while
\[
              7u-4-3u^2=1/12-3(u-7/6)^2<.289^2 .
\]
The total loss is \(<.428<.61\).
This covers the ENTIRE finite native tail.

For left endpoint-three plus forks, \(Q=0\).
With both forks-three,
\(\Delta=1/24\), \(\sigma\ge1/2+1/24\);
with one fork-four,
\(\Delta=1/24\), \(\sigma\ge1/(2\sqrt2)+1/24\).
Both are directly strong even at own zero.
With one fork-three,
\(\Delta=7/24\); own zero has its actual cycle,
own three its contacted-fork triangle,
and every own at least four is directly strong.
The far orientation is used in each lower bound,
so near orientations are also covered.

With no left contact and positive own, a single
fork-three/own-three has its affine triangle.
Single fork at least six with own three is directly
strong.
Single fork-four/own-three with \(Q>0\) has
\(\Delta\le1/6\) and is strong; its \(Q=0\)
case is retained below.
Own four has the star at twelve with leaves cap,
eleven and the UNCONTACTED fork.
Own six has the \(6,3\) path through eleven;
parallel has its null pair.
For every remaining own at least seven or dotted,
fork at least four is directly strong.
Fork-three with \(Q>0\) has \(\Delta\le5/12\)
and is strong.
At \(Q=0\), a possible weak fork-three has
\(.9<g<6/5\).
Its ACTUAL triangle makes \(t=2g\) integral.
Normalize the connected simple-three cap to raw
square two; all nonphysical \(K_7\) embeddings obey
\[
                  48t_\sigma(t_\sigma+1)+14<8,
                  \qquad -1<t_\sigma<0 .
\]
Its nonzero integer \(t(t+1)\) has norm
\(<(204/25)(1/4)^2=.51<1\), impossible.
For two forks-three, own three has its triangle
and owns at least four are directly strong.
Pair three/four is directly strong for every
positive own.

With no left contact and own zero, single
fork-three has the actual affine \(E_8\) kernel:
centre twelve coefficient six, uncontacted fork
three, contacted fork four, cap two, and old
eleven through seven coefficients five through one.
Single fork-four has kernel
\((\sqrt2,2,2,1,1)\) on
cap, contacted fork, twelve, eleven and other fork.
Single fork-six has the \(6,3\) principal.
Every null has positive cap time and correct signs
against ALL omitted ORIGINAL walls.
Single fork-five is absent.
For any finite fork label at least seven,
\(\delta>.9\); finite whiteness
\(\delta^2+Q<1\) forces \(Q=0\).
Two forks-three have the actual four-cycle;
pair three/four has \(Q=0\) and remains below.
Thus the ONLY nonstrong possibilities are:
own-three/single-fork-four \(Q=0\);
own-zero/pair-three-four;
own-zero/single-finite-fork-at-least-seven.
All other own, foreign and unbounded dotted cases
have been covered.

\subsection{Reciprocal endpoints and the remaining continuous fork family}
\label{app:arithD14typedexceptions}

For own-three/single-fork-four, the SAME actual cap
is also the contacted fork's opposite.
Deleting that fork retains \(A_{13}\); adding the
cap at its next-to-end node by three gives finite \(D_{14}\).
For own-zero/pair-three-four, release the fork-three
direction instead; the retained \(A_{13}\) and cap
give finite \(B_{14}\).
In both cases keep the ENTIRE orthogonal spectator.
Positive inverse finite-Gram projection from \(p\)
gives an actual future vertex on that original edge.
Every unknown wall remains strict: its product with
\(p\) is strict and all correction coefficients
are nonnegative.
The omitted old fork is strict because the cap
coefficient and that fork contact are positive.
This proves an ACTUAL endpoint, not a formal
completion or a new farthest vertex.

Both reciprocal fork coshes are
\[
                         C_f=6+7\sqrt2/2>10 .
\]
In the first case
\(A^2=(17+12\sqrt2)/2\),
\(v\cdot\omega_f=6+7/\sqrt2\),
\((v\cdot\omega_f)^2-7A^2=1\).
In the second,
\(A^2=(17+12\sqrt2)/4\),
\(v\cdot\omega_f=7/2+3\sqrt2\),
and the difference is \(1/2\).
These give the displayed cosh by the original
edge formula.
For either sign, the genuine spatial direction
\(V=P_{13}\pm e_{14}\), \(V^2=14\), gives
\[
                 x\cdot V>\sqrt{126/11}>3.38 .
\]
Pay \(\mu=1/7\) at the individual \(P_{11}\)
coefficient \(>1/7\).
The residual \(V-P_{11}\) vanishes on all first
eleven coordinates, has square three and full
cross zero with \(w\).
One original square WITHOUT any head credit gives
\[
          e_D+e_Z>.84816+.09+
                      {2(3.38-2.886)\over7}-{3\over49}>1.018.
\]
Both exceptional reciprocal rows are excluded.

Finally let own be zero and the single finite fork
contact have square \(b_f=\delta^2\in(.81,1)\).
Put \(u=1/b_f\in(1,100/81)\); \(Q=0\) was proved,
rather than assumed.
The actual positive combination
\[
               \xi=c+\delta^{-1}n_f
                           +2(\delta^{-1}-\delta)n_{12}
\]
has zero cap and contacted-fork products, negative
centre and other-fork products, and square
\[
                   -k,\qquad k={(1-b_f)(4b_f-3)\over b_f}>0 .
\]
Its time is positive and every unknown original-wall
product is nonpositive termwise.
It is therefore a feasible future point of the
ENTIRE polytope.
For \(\alpha=\delta/\sqrt2\), its divided spatial
displacement and pressure are
\[
 \begin{gathered}
 V=(1^{11},3-2u,u-1,\ \pm(1-u)),\\
 x\cdot V\ge G=\sqrt{14-2u}-\sqrt{-8+14u-6u^2},\\
 R=V^2=22-16u+6u^2 .
 \end{gathered}
\]
This retains the ACTUAL fork sign; it never assumes
\(x_{14}>0\).
The residual \(V-P_{11}\) vanishes on the entire
reference support and is orthogonal to \(Y_H\).

On the ENTIRE interval,
\[
                             14G-R>30.58 .
\]
Here is a continuous rational seal.
The chord
\(\sqrt{14-2u}\ge3.756-.292u\)
has positive endpoint-square margins
\(11/15625\), \(574919/410062500\).
Put \(z=u-7/6\).
Then
\[
 -8+14u-6u^2=1/6-6z^2,\qquad
 R=23/2-2z+6z^2,
\]
and
\(\sqrt{1/6-6z^2}\le1/\sqrt6-3\sqrt6z^2\).
Consequently
\[
 \begin{split}
 14G-R
 &\ge13618/375-14/\sqrt6-2.088z+(42\sqrt6-6)z^2\\
 &>13618/375-5.716-2.088z+96.8748z^2\\
 &>13618/375-5.716-.012=2294/75>30.58 .
 \end{split}
\]
Use \(\sqrt6>2.4494\), \(5.716(2.4494)>14\)
and \(2.088^2<.048(96.8748)\) for the final
completed square.

Now combine the TWO primitive pressures BEFORE
taking the norm:
\[
                         W=w+.4Y_H+(V-P_{11})/7 .
\]
Its INDIVIDUAL prices are \(.08\) at \(P_5\)
and \(1/7\) at \(P_{11}\).
Both remaining coefficients are positive.
Both residuals are orthogonal to \(w\) AND to
each other.
Using the head gain~\eqref{eq:arithBD14typedgain}
and the same spectator ONCE gives
\[
 \begin{split}
 e_D+e_Z
 &>E_{11}+.09+.0487+
                  {14G-R-14f_{11}+11\over49}\\
 &>.84816+.09+.0487+
                  {30.58-40.404+11\over49}
   =1.01086>1 .
 \end{split}
\]
This is one original reference, not the addition
of independent head and centre optima.

Every ACTUAL centre-twelve cap is now strong
or impossible.
The positive curve through twelve is available
with \(E_{12}>.86989\).
Its first-five mean and individual \(P_5\) price
are unchanged.
Apply the SAME actual head point and its primitive
\(.08\) payment once to this NEW reference:
\[
                       e_D+e_Z>.86989+.0487+.09=1.00859>1.
\]
This excludes BOTH \(D_{14}\) actual typed-seven
branches.
The order was fresh \(6,\ldots,11\), then five,
then the actual head, THEN centre twelve and its
reciprocal/continuous exceptions.
Neither fork was assigned an unsupported prefix
floor, and no actual short-\(I\) barrier or
arithmeticity of a face was used.

\section{A whole \(D_{13}\) component with any rank-three remainder}
\label{app:arithD13rankthree}

This argument covers EVERY finite rank-three spectator
forest beside a whole \(D_{13}\) at the original
farthest vertex in dimension sixteen.
It does not depend on subtracting rows of a
previous frontier.
Retain actual unit \(D_{13}\) roots, long spine
one through eleven, forks twelve and thirteen,
and signed coordinates
\(x_1>\cdots>x_{12}>|x_{13}|\).
Write \(e_D+e_S<1\) for the SAME original vector
and its entire orthogonal rank-three spectator.

Every such spectator has
\[
                                e_S>.21 .
\tag{S3}\label{eq:arithD13spectator}
\]
Here is the COMPLETE partition and exception argument.
With ordinary components only, partitions three,
two-plus-one and one-plus-one-plus-one cost
\(>1/4\), \(>1/7+.1\) and \(>.3\).
For actual \(I\), the remaining rank two is one
ordinary irreducible \(T_2\), or two ordinary \(A_1\)'s.
The existing rank-free joint costs are
\(I+T_2>.21\) and \(I+A_1>.20204\);
the latter plus one further \(A_1\) is \(>.30204\).
The actual \(I\) is counted once.
A \(J\) pair plus the last ordinary \(A_1\)
costs \(>.34+.1\).
An actual typed rank-two defect costs \(>.09\).
For \(N_7\), existing propagation gives the last
\(A_1>.2\); for \(O_7\), it gives \(A_1>2/15\).
The compatibility theorem permits no additional
\(I/J\) beside a typed defect and at most the
single stated short-rank exception.
These are every rank-three partition and exception.
The propagation/joint theorems are the SAME whole
component statements already used in the shared seals;
no new field refinement is needed for this cost.

For \(D_{13}\), choose generic parameter \(M=7\).
Its exact rank hypothesis \(13=M+6\) holds.
Deleting \(i=5,\ldots,8\) retains
\(A_{i-1}+D_{13-i}\), with both necessary side
ranks at least four/five.
The full finite connected replacement has rank
thirteen, so exceptional high-rank alternatives
do not arise; the existing small \(E_6,E_7,E_8\)
fork attachments are retained in the local list.
The complete generic weak row is the pure far
helper, excluded at \(i\ge6\) by
Lemma~\ref{lem:literal-helper}.
Thus six, seven and eight are strong before
either short layer or strength at five.

We need only two earlier anchors, \(K=6,7\).
If five is strong use the positive ordinary
reference with norm \(E_K\).
If five is weak, its actual pure point gives
\(x_1>.5965\), \(S_6>2.01815\).
At \(K=6\) use coordinates
\((.5965,.28433,\ldots,.28433)\).
Their norm is \(.7600299945>E_6\), their sum
\(2.01815<f_6\), and terminal coefficient
\(.28433>1/6\).
At seven append \(\sqrt7-\sqrt6\);
the endpoint is \(<f_7\), energy \(>E_7\)
and terminal coefficient
\(1/(\sqrt7+\sqrt6)>1/6\).
All individual duals are nonnegative.
The strong-five terminal coefficients at six/seven
are also \(>1/6\).
In EITHER initial branch we therefore have an
actual supported reference \(w\) on the first \(K\)
coordinates satisfying
\[
 w^2\ge E_K,\quad w\cdot P_K\le f_K,\quad
        \text{individual }P_K\text{ coefficient}>1/6 .
\]
The exact scalar seals are
\[
 E_6>.69685,\ E_7>.73531,\ E_{10}>.82436,\qquad
 f_6<2.019,\ f_7<2.215.
\]

The fresh short layers are \(i=9\), with
white \(A_8+D_4\), and \(i=10\), with
white \(A_9+D_3\).
Their COMPLETE left-only finite extensions are
endpoint-three/four and next-to-end-three.
Ranks nine/ten preclude a new exceptional type.
Two contacts to either retained component form
a forbidden cycle.
The shared local inverse formulas and original
spectator clearing apply at these ACTUAL side ranks.
We give the right and bridge cases explicitly:
\[
\begin{array}{c|c|c}
 \text{right role}&\delta\text{ before }Q&\sigma\\ \hline
 D_4\text{ long three}&1/2&\gamma+1/2\\
 D_4\text{ fork three}&1/2&\gamma+1/4\\
 D_3\text{ middle three}&1/2&\gamma+1/2\\
 D_3\text{ end three}&5/8&\gamma+1/4\\
 D_3\text{ end four}&1/4&\gamma+1/(2\sqrt2)\\
 D_3\text{ end five}&(5-3\varphi)/8&\gamma+\varphi/4 .
\end{array}
\]
Every foreign \(Q\) is the ENTIRE spectator cost,
not an individual bond, and subtracts from the residual.

The long/middle-three own-zero rows have the actual
two-fork/four-leaf affine kernel, own three their
triangle; every remaining own at least four has
ratio \(\ge1+1/\sqrt2>100/61\).
End-four own zero has the actual star of squared
costs \(1/4,1/4,1/2\); every positive own has
the same direct lower ratio.
End-five is the sealed \(H_4\), strong even at
own zero, since its residual is \(<1/40\)
and numerator \(>2/5\).
A left endpoint-three/right-simple-three bridge
completes \(D_{13}\), or, for \(D_3\) end-three,
\(A_{13}\).
The former has residual \(1/(2i)\) and positive-own
numerator at least \(\gamma+1/4+1/(2i)\);
the latter at ten has residual \(7/40\) and
positive-own numerator \(>.75\).
Both ratios exceed \(100/61\).
Their own-zero cycles are affine.
The bridge cap is internal in a rank-thirteen
classical tree, so has no finite spectator arm;
the \(A_{13}\) bridge's old arms have lengths
nine and three.
An end-four/five cannot join the long left path:
its non-three bond would be internal in a
nonfinite rank-thirteen tree.

For the left-only small owns, internal-three and
endpoint-four own zero retain their actual affine
fork principals.
Own three with a white-three contact has its cycle.
Own four with no right contact has the affine
principal obtained from the retained right
\(D_5/D_4\) INCLUDING \(n_i\):
its unit long-leaf inverse diagonal is two,
so a four attachment has cost one.
Own six has its local \(6,3\) path.
Parallel always has its future feasible null pair.
Far-left-three at own zero with foreign contacts
has \(A^2=Q-1/2<1/(2i)<1/4\).
The original small-time shadow argument proved
in the dihedral section would make the SAME cap
close an orthogonal spectator direction, contrary
to the original-farthest common-opposite obstruction.
For \(Q=0\) its time square is negative.
Thus that own-zero row also is impossible.
Every null here is a positive combination of ACTUAL
walls with positive cap time and nonpositive product
at EVERY omitted original wall.
The nonzero left-only ratios and foreign clearings
are the legal shared prefix formulas at \(i=9,10\).
The only potentially weak left row is the literal
far helper and is removed by
Lemma~\ref{lem:literal-helper}.

For isolation, the original singleton and LEFT
\(H_4/A_2\) projected points and negative cores
remain available through \(i=r-3\), with \(D_i=2i\ge18\).
They cover arbitrary foreign cost and all finite/
parallel/dotted own, without a fictitious long root.
For clarity their scalar margins still suffice here:
the \(0<u<1\) core has full loss
\[
             <{.12+\sqrt{1+.12^2}\over2}<.57<.61,
\]
using \((\sqrt{18}+\sqrt{17})^{-1}<.12\).
For foreign \(u>1\), \(u<7/6\), the radical
loss is \(<.14+.4<.61\).
Pure native five's LEFT \(H_4\) correction has
loss \(<.2+.25<.61\);
native seven's LEFT \(A_2\) correction has
loss \(<.15+.36<.61\);
all finite natives at least eight use the
singleton loss \(<.15+.415<.61\).
The underlying divided displacements and concavity
are exactly those printed in
Section~\ref{app:arithD15body};
only these freshly stated \(D_i\ge18\) scalar
bounds are used, not its \(i=11,12\) estimates.

For the retained single-fork-three/end-three roles,
own zero, three, four and six have respectively
the actual \(E_8\), cycle, \(F_4\) and \(6,3\)
principals; all required old arms exist at nine/ten.
Parallel has its actual null pair.
Thus \(\gamma\ge\varphi/2\).
If \(Q>0\), then \(\delta\le3/8\),
\(\sigma>21/20\), giving ratio \(>42/25>100/61\).
For \(Q=0,\gamma\ge6/5\), the ratio again exceeds
\(100/61\).
ALL finite and dotted own values in between
are retained.
The ONLY possible weak short caps are therefore
\[
 v=\gamma\omega_i+\omega_f/2,\quad Q=0,\quad
 i=9,10,\ f=12,13,\quad 202/125<2\gamma<12/5.
\]

Their actual simple-three contact normalizes the
cap to raw square two over ACTUAL \(K\).
The original path cycle makes \(t=2\gamma\) integral.
At EVERY nonphysical embedding,
\[
                     4it_\sigma(t_\sigma+1)+13<8 .
\]
Hence \(-1<t_\sigma<0\).
All embeddings of the minimal field lift to \(K\).
The integer \(t(t+1)\) has physical value \(<204/25\)
and other absolute conjugates at most \(1/4\).
Degree at least three contradicts its nonzero
integral norm bound \(<.51\).
In degree two the trace is one or two.
Trace one gives norm \(-1\) and \(t=\varphi\);
trace two gives norm in \((-24/25,0)\), impossible.
The rational value two is an actual parallel null.
Every weak fork helper is therefore golden,
without a field census or a scale comparison
with spectators.

For each golden packet use the ACTUAL LEFT finite
\(H_4\) correction
\[
 w_{\rm old}=\gamma(3n_i/2+n_{i-1}+n_{i-2}/2),\quad
 \alpha^2=\varphi+5/8>56/25,\quad
 \delta=(5-3\varphi)/8,\quad \sqrt\delta/\alpha<.091 .
\]
The actual finite point is the normalization of
\(p+(A/\delta)(c+w_{\rm old})\).
It has square \(-1-A^2/\delta\), positive future
time and zero selected products.
Every omitted original-wall product is nonpositive
by actual acuteness and the positive coefficients.
The correction alone is not a hyperbolic point.
For \(Y=v-w_{\rm old}\), \(V=Y/\alpha\), \(K=i-3\),
the fresh inverse calculations give
\[
 \begin{gathered}
 A^2=i\alpha^2-\kappa,\quad
 \kappa=1/2\ (i=9),\quad \kappa=5/8\ (i=10),\\
 V_j=1\ (j\le K),\quad V^2<i,\quad
 x_D\cdot V>G=\sqrt{i-\kappa/\alpha^2}-.091 .
 \end{gathered}
\]
Use the already established \(K=6/7\) reference
in EITHER initial-five branch.
Pay its INDIVIDUAL terminal coefficient \(1/6\)
BEFORE using its floor.
For \(W=w+(V-P_K)/6\), the residual vanishes
on the first \(K\) coordinates, so its FULL
cross with \(w\) is zero, and its norm is \(<3\).
One square gives
\[
                         e_D>E_K+{G-f_K\over3}-{1\over12}.
\]
The COMPLETE fresh table is
\[
\begin{array}{c|c|c|c|c|c}
 \text{layer}&i&K&G>&f_K<&e_D>\\ \hline
 D_4&9&6&2.871&2.019&.8975\\
 D_3&10&7&3.026&2.215&.92231 .
\end{array}
\]
The two radical seals are
\(2.962^2<9-25/112\),
\(3.117^2<10-125/448\).
Adding ONLY the actual disjoint spectator
\eqref{eq:arithD13spectator} contradicts the
original norm in BOTH cases.
Thus all \(6,\ldots,10\) are now genuinely strong.

The final ordinary or legitimate weak-five
reference through ten gives \(e_D>E_{10}>.82436\).
Consequently
\[
                         e_D+e_S>.82436+.21=1.03436>1.
\]
This excludes WHOLE \(D_{13}\) with EVERY rank-three
spectator forest at the actual farthest vertex in
dimension sixteen.
Every pressure belonged to one ORIGINAL vector;
the spectator was counted once.
No actual-I barrier, typed field lock, field census,
arithmeticity of a face or unexecuted computation
was used to enlarge this theorem's scope.

\section{Whole \(B_{16}\): geometry before the finite pair domain}
\label{app:arithwholeB16}

Now the ENTIRE incident forest is \(B_{16}\);
there is no spectator to delete or to charge.
Use actual unit roots
\(n_i=(e_i-e_{i+1})/\sqrt2\), \(i\le15\),
\(n_{16}=e_{16}\), and the original vector
\(z=E(p-x)\), \(x_1>\cdots>x_{16}>0\), \(\|x\|<1\).
The legal fresh bulk and B3/B2 terminal checks
give strong \(6,\ldots,14\) independently of five
and either endpoint.
We first prove all physical reductions needed
to enter the exact quadratic algorithm.

\subsection{The actual head and a tight short-edge length bound}
\label{app:arithB16edges}

For \(a=.643\) or \(.648\), the reference through
\((1,a),(6,f_6),\ldots,(14,f_{14})\) has positive
decreasing coordinates
\[
 a,\quad(f_6-a)/5\text{ five times},\quad
 \sqrt j-\sqrt{j-1}\ (7\le j\le14),\quad0,0 .
\]
Its first middle slope exceeds \(.274\), hence
the next square-root slope \(<.2\).
If actual \(x_1>a\), ONE primitive square gives
\[
 \|x\|^2>
 a^2+(f_6-a)^2/5+
             \sum_{i=6}^{13}(\sqrt{i+1}-\sqrt i)^2.
\]
The last sum is \(>.2117\).
A direct rational proof avoids a numerical receipt:
for \(A=4i+2\), put
\(L=(A^2-1)/(A(A^2-2))\).
The small root \(d=(\sqrt{i+1}-\sqrt i)^2\)
satisfies \(d(A-d)=1\), whereas
\[
                 L(A-L)-1=-{1\over A^2(A^2-2)^2}<0 .
\]
Thus \(d>L\).
For \(i=6,\ldots,13\), strict lower numerators
of \(L\) at denominator \(10^7\) are
\[
 385186,333704,294372,263340,238230,217494,200080,185248 .
\]
Their sum is \(2117654\).
With \(f_6>2.01815\), the two reference norms are
\(>1.0033565045\), \(>1.0070662045\).

The ACTUAL head's retained white diagram is \(B_{15}\).
Finite classification permits isolation or simple
three at its LONG endpoint two, with no symmetric
short-end alternative.
For the latter,
\(A_H^2=2g^2+2g\), \(C_H=2g+1\).
Own zero has zero time, own three has its affine
triangle, and parallel has the actual null pair.
Every remaining own has \(g\ge1/\sqrt2\), so
\(x_1>\sqrt{\sqrt2-1}>.643\), impossible by the
first reference.
The head is isolated:
\[
 c_H=A_Hp-\sqrt2g_He_1,\quad
 U=4g_H^2\in\mathcal O_K,\quad A_H^2=U/2-1,\quad U>2 .
\]
If \(U\ge12\), its ACTUAL edge has
\(C_H=\sqrt{U/2}\ge\sqrt6\), giving \(x_1>.648\),
impossible by the second reference.
The strict rational seal uses
\(\sqrt6>2.449\),
\(.648^2(3.449)<1.449\).
Thus \(2<U<12\).

Let \(C_S\) be the ACTUAL short-sixteen edge cosh.
Its actual finite endpoint and original comparison give
\(S_{16}\ge4\sqrt{(C_S-1)/(C_S+1)}\).
If \(C_S\ge16\), this is \(>3.7573\).
For strong five, interpolate through true knots
five through eight and this terminal pressure.
The last slope is \(<.172<\sqrt8-\sqrt7\).
Using
\[
 E_8>.69685+.0385186+.0333704=.768739,\qquad f_8<2.3971,
\]
its norm is
\[
                      >.768739+1.3602^2/8=1.000007005 .
\]
The upper seal for \(f_8\) follows from
\(\sqrt8<2.828428\), \(\sqrt2<1.414214\),
\(.431328(1.414214)<.61\); also \(f_8>2.39\)
gives the claimed last-slope bound.
For weak five use the actual pressures
\(x_1>.59\), \(S_6>2\) and \(S_{16}>3.7573\).
The decreasing slopes \(.59,.282,.17573\) give
one reference of norm
\[
                     .59^2+1.41^2/5+1.7573^2/10
                            =1.054530329>1 .
\]
Both alternatives contradict the SAME original
norm without adding a head gain.
Therefore \(C_S<16\).

\subsection{All seven actual short-wall profiles}
\label{app:arithB16shortwhite}

The retained ACTUAL \(A_{15}\) finite-white list is
isolation; helper three at \(j=1,2,14,15\);
or helper four at \(j=1,15\).
Two contacts give a cycle.
Other simple branches or non-three bonds cannot
give a finite rank-sixteen tree.
There are no foreign components.

Normalize old long roots to raw square two and
the genuine short root to square one:
\(\alpha_i=e_i-e_{i+1}\), \(s_{16}=e_{16}\).
All \(e_i\) are then \(K\)-valued.
A helper-three cap normalizes to square two,
helper product \(-1\), and short product \(-T/2\),
where \(T=2\sqrt2\gamma\).
The ACTUAL cycles give
\(T,T^2/2\in\mathcal O_K\), not \(T/2\) integral.
Its exact formulas are
\[
 \begin{gathered}
 c=Ap-\gamma P_{16}-P_j/\sqrt2,\quad
 2A^2=F_j(T)=4T^2+jT+j-2,\\
 C_S={8T+j\over\sqrt{32-j(16-j)}} .
 \end{gathered}
\]
The original old-plus-cap Gram is Lorentzian over
\(\mathbb Q(T)\), so actual \(K=\mathbb Q(T)\),
and every other \(F_j(T_\sigma)<0\).
A helper-four cap has raw square one and helper
product \(-1\), with
\[
       A^2=16\gamma^2+2j\gamma+j-1,\quad
       C_S=16\gamma+j,\quad K=\mathbb Q(\gamma).
\]
These field equalities follow from nonphysical
embeddings fixing an original indefinite principal,
not from arithmeticity of a face.

Own zero is retained before dispatch.
Helper-three at one/two has \(A^2=-1/2,0\).
At fifteen its actual null is
\(c+n_{14}+2n_{15}+\sqrt2n_{16}\);
at fourteen it is
\(c+n_{13}+2n_{14}+2n_{15}+\sqrt2n_{16}\).
Helper-four at one has zero time;
at fifteen use \(c+\sqrt2n_{15}+n_{16}\).
Isolation has \(A^2=-1\).
Every null has positive cap time and selected
products zero; ALL omitted original-wall signs
are nonpositive termwise.
Parallel own always has \(c+n_{16}\).

Near helper-four positive own has \(C_S\ge23\).
Far helper-four has \(\gamma<15/16\), so its
complete positive native alphabet is
three through eight.
The native-nine cubic at \(15/16\) is
\(-17/512\), so all larger finite labels, parallel
and dotted are beyond this cutoff.
Own three is rational and excluded by WHOLE-form
compact anisotropy/Meyer.
For owns four through eight use respectively
the nonphysical conjugates
\[
 -\sqrt2/2,\quad(1-\varphi)/2,\quad-\sqrt3/2,\quad
 \cos(3\pi/7),\quad-\cos(\pi/8).
\]
The resulting time squares are
\(8-\sqrt2\), \(9-5\varphi\), \(12-\sqrt3\),
\(16\gamma_\sigma^2+2\gamma_\sigma\), and
\(16\gamma_\sigma^2-2\gamma_\sigma\), all positive.
They contradict admissibility in the ACTUAL
generating field.

Far next-end-three has
\(T<15/4\), every other \(-1/2<T_\sigma<0\).
In degree \(d\ge2\), its nonzero integral norm
has magnitude \(<15/8<2\), hence one.
But \(T^2/2\) integral would make
\(\mathrm N(T)^2/2^d\) an integer, impossible.
Degree one is rational.
Near next-end-three has \(T<9/4\), leaving only
own three \(T=\sqrt2\) or own four \(T=2\).
Own five already has \(T^2=2\varphi^2>81/16\).
The rational case is impossible; the other
conjugate time numerator is \(20-14\sqrt2>0\).
Both next-end orientations and both helper-four
roles are exhausted.

For the remaining far-end-three put \(S=T\);
for near-end-three put \(S=T+2\).
Their time polynomials are \(4S^2+S-1\) and
\(4S^2-S-1\), with respective nonphysical intervals
\[
 { -1-\sqrt{17}\over8}<S_\sigma<{ -1+\sqrt{17}\over8},
 \qquad
 {1-\sqrt{17}\over8}<S_\sigma<{1+\sqrt{17}\over8}.
\]
Both have \(0<S<42/5\), \(|S_\sigma|<.641\)
and \(S^2/2\) integral.
For degree \(d\ge3\), the nonzero integer norm
is divisible by four but has magnitude
\(<(42/5)(.641)^2=3.4514004<4\).
Degree one is rational; actual degree is two.
Trace and norm are EVEN by the same half-square
argument already proved above.

The hand table uses the EXACT physical short-edge
cutoff \(C_S<16\), rather than only the coarser
\(S<42/5\) used to bound the degree:
\[
\begin{array}{c|c}
 \text{far}&\sqrt2\le S<(16\sqrt{17}-1)/8<65/8\\
 \text{near}&2+\sqrt2\le S<(16\sqrt{17}+1)/8<67/8 .
\end{array}
\]
Indeed the actual edge formulas are
\(C_S=(8S+1)/\sqrt{17}\) and
\(C_S=(8S-1)/\sqrt{17}\), respectively;
the lower bounds retain the native own gap.
For example the far integer \(4+3\sqrt2\)
passes the coarse bound but violates this
exact physical cutoff and is not a table row.

The far strips give traces \(2,4,6,8\),
\(-6<\mathrm N(S)<4\).
The near strips give traces \(4,6,8\),
\(-4<\mathrm N(S)<6\).
Solving ALL even trace/norm pairs and imposing
the two strict strips gives the COMPLETE eleven rows:
\[
\begin{array}{c|c|r}
 \text{role}&T&\mathrm N(F_j(T))\\ \hline
 \text{far}&2+\sqrt6&-53\\
 &3+\sqrt{13}&-25\\
 &3+\sqrt{11}&-151\\
 &3+\sqrt7&-19\\
 &4+\sqrt{14}&-117\\ \hline
 \text{near}&\sqrt2&-9\\
 &1+\sqrt{11}&-43\\
 &1+\sqrt7&-103\\
 &2+3\sqrt2&-137\\
 &2+\sqrt{14}&-229\\
 &2+2\sqrt3&-83 .
\end{array}
\]
For explicit coverage, the far norm columns are
\(-4,-2,2\).
Trace two gives only inadmissible conjugates or
nonreal roots; trace four retains \(2+\sqrt6\);
trace six retains all three \(3+\sqrt{13}\),
\(3+\sqrt{11}\), \(3+\sqrt7\);
trace eight retains only \(4+\sqrt{14}\).
The near \(S\) norm columns are \(-2,2,4\).
Trace four retains only \(2+\sqrt2\);
trace six retains \(3+\sqrt{11},3+\sqrt7\);
trace eight retains \(4+3\sqrt2,4+\sqrt{14},4+2\sqrt3\).
Subtract two to recover the printed near \(T\).
The omitted rows violate a strict conjugate/physical
bound or are repeated rational roots.
All displayed norms are negative and their physical
times positive, so the other time signs are exactly
the required negative ones.
These are necessary actual models, never examples.

\subsection{The original connected-basis squares remove every helper}
\label{app:arithB16helpersquares}

Normalize the actual head to the \(K\)-root
\(h=c_H/(\sqrt2g_H)\), with
\(h\cdot\alpha_1=-1\), \(h^2=2/U\).
Every nonphysical head Schur complement gives
\(0<U_\sigma<2\).
For any of the eleven helper-three rows, the
sixteen old raw roots and \(r=\sqrt2c\) are an
ORIGINAL \(K\)-basis.
The coefficient of \(r\) in \(h\), computed from
its positive time coefficients, therefore gives
\[
                      R_0={U-2\over U F_j(T)}\in K^2 .
\]
This derives a common time-line square in one
actual connected basis.  It does not put separate
unit roots or square roots in \(K\).

For quadratic \(K\), the full head strip gives
\[
                     0<\mathrm N(U)<24,\qquad
                     0<-\mathrm N(U-2)<20.
\]
For all five NEAR DOTTED rows and the far
\(\sqrt6,\sqrt{11}\) rows, the absolute time
norm is a prime at least forty-three.
Its valuation in \(\mathrm N(R_0)\) is odd
because neither of the two bounded integer norms
can contain it.
For far \(\sqrt7\), the FULL ring
\(\mathbb Z[\sqrt7]\) gives the complete head list
\[
 U=3+\sqrt7,\ 4+\sqrt7,\ 6+2\sqrt7,
\quad \mathrm N(U)=2,9,8,\quad
 \mathrm N(U-2)=-6,-3,-12 .
\]
Writing \(U=a+b\sqrt7\) shows \(b=1,2\);
the strict head strips give precisely these \(a\)'s.
None contains prime nineteen.
For far \(\sqrt{14}\), its full strip similarly
gives
\[
 U=4+\sqrt{14},\ 5+\sqrt{14},\quad
 \mathrm N(U)=2,11,\quad\mathrm N(U-2)=-10,-5 .
\]
The odd prime thirteen in time norm \(-117\)
therefore cannot cancel.
These are nine eliminated rows.

The near \(T=\sqrt2\) row needs its FULL square
test rather than a norm-only rejection.
Here actual \(K=\mathbb Q(\sqrt2)\), because its
connected old-plus-cap Gram is physically Lorentzian
over that field.
Put \(s=\sqrt2\), \(\epsilon=1+s\).
The ACTUAL cap has
\[
 A_F^2=(21+15s)/2=3\epsilon^3/2,\qquad
 c_F=A_Fp-P_{16}/2-P_{15}/s .
\]
The complete head strip in \(\mathbb Z[s]\) is
\[
         U=2+s,\ 3+s,\ 3+2s,\ 4+2s,\ 5+3s,\ 6+3s,\ 6+4s .
\]
Indeed \(U=a+bs\) has \(1\le b\le4\);
the allowed pairs are
\(b=1:a=2,3\);
\(b=2:a=3,4\);
\(b=3:a=5,6\);
\(b=4:a=6,7\), with the last \(a=7\) beyond twelve.
For \(R_0=(U-2)/(3\epsilon^3U)\), the norms are
\[
                      1/9,\ 1/63,\ 7/9,\ 1/18,\ 1/7,\ 1/81,\ 4/9 .
\]
The four nonsquares are impossible.
At \(U=2+s\), \(R_0=1/(3\epsilon^4)\);
at \(U=6+4s\), \(R_0=2/(3\epsilon^4)\).
Neither \(1/3\) nor \(2/3\) is a square in
\(\mathbb Q(s)\): a rational square in that field
requires the rational number itself or half of
it to be a rational square.
The last row \(U=6+3s\) REALLY gives a square,
with its positive physical root \(k=1/(3\epsilon)\).
For the ACTUAL mutual magnitude \(M=-c_H\cdot c_F\),
\[
 A_HA_F/g_H=\epsilon^2,\qquad
 M/g_H=(4+3s)/2,\qquad
                    4M^2=(174+123s)/2 .
\]
Its rational norm is \(9/2\), so it is not integral.
This contradicts the actual two-cycle.
The positive time sign is fixed by actual outward
roots, not an optional square-root sign.
Thus the near own-three model also is excluded.

The far \(T=3+\sqrt{13}\) model is last.
Put \(r=\sqrt{13}\), \(\eta=(3+r)/2\).
Its time polynomial is \(F_1(T)=5\eta^3\).
Use the FULL ring \(\mathbb Z[(1+r)/2]\).
The complete head strip gives the SIX rows
\[
\begin{array}{c|r|r|c}
 U&\mathrm N(U)&\mathrm N(U-2)&\mathrm N(R_0)\\ \hline
 (5+r)/2&3&-3&1/25\\
 (7+r)/2&9&-1&1/225\\
 4+r&3&-9&3/25\\
 5+r&12&-4&1/75\\
 (11+3r)/2&1&-17&17/25\\
 (13+3r)/2&13&-9&9/325 .
\end{array}
\]
Writing \(U=(a+br)/2\), same-parity integers,
the strip gives \(b=1,2,3\) with respectively
\(a=5,7;\ 8,10;\ 11,13\).
The last four norm entries are nonsquares.
For the first,
\(R_0=(101-28r)/15\).
A full field square would have rational coefficient
square \(52/15\) or \(49/15\), neither rational-square.
For the second, multiplying by the square \(\eta^2\)
gives
\[
                   \eta^2R_0=(7-r)/90 .
\]
Its rational coefficient square would be \(13/180\)
or \(1/180\), again impossible.
No norm-square is treated as sufficient.
All seven actual white profiles are now exhausted
except isolation.

\subsection{The actual isolated pair and the entire cubic branch}
\label{app:arithB16cubic}

The actual short cap is
\(c_S=A_Sp-g_SP_{16}\).
Own zero and parallel are impossible as above.
Put \(V=4g_S^2\).
Rescale the ACTUAL \(K\)-root to raw norm \(V\),
old-short product \(-V/2\):
physically \(F_S=\sqrt V\,c_S\), spatial
\((V/2)P_{16}\), time square \(V(4V-1)\).
Its original Lorentzian Gram over \(\mathbb Q(V)\)
gives ACTUAL \(K=\mathbb Q(V)\), and EVERY
nonphysical place has \(0<V_\sigma<1/4\).
The short edge bound gives \(V<64\).
Rational \(K\) is excluded by WHOLE-form Meyer.
Finite own would give \(V<4\) and
\(0<\mathrm N(V)<4(1/4)^{d-1}\le1\), impossible.
Thus \(V>4\), the short own is dotted.
Degree at least four gives
\(\mathrm N(V)<64(1/4)^3=1\).
Hence actual degree is two or three.

Normalize \(h=c_H/(\sqrt2g_H)\),
\(r_F=c_S/g_S\), so
\(h^2=2/U\), \(r_F^2=4/V\), spatial parts
\(e_1,P_{16}\).
The original old sixteen roots plus \(r_F\)
form a \(K\)-basis.
The positive coefficient of \(r_F\) in \(h\) gives
\[
                   k^2={(U-2)V\over4U(4V-1)}\in K^2,
                   \qquad k>0 .
\tag{B16K}\label{eq:arithB16square}
\]
Let \(H=Uh\), \(F=(V/2)r_F\).
The ACTUAL full long-path cycle and two-cycle give
\[
 \begin{gathered}
 G=-2H\cdot F=4Uk(4V-1)-UV\ge0,\qquad
 L={G^2\over2UV},\qquad G,L\in\mathcal O_K,\\
 (G+UV)^2=4UV(U-2)(4V-1),\\
 L+G=((15U-32)/2)V-2U+4 .
 \end{gathered}
\tag{B16C}\label{eq:arithB16cycles}
\]
The old-path factors are \(-1\), followed by
\(-V,-G/V\), whose complete product is \(G\).
Orthogonal \(G=L=0\) is retained.
The last identity, subtracting seven integral
copies of \(UV\), implies \(UV/2\) integral.
This is rank sixteen's own half-product, not
the odd-rank-seventeen census invariant.
Thus every actual pair has the COMPLETE domain
\[
 \begin{gathered}
 K=\mathbb Q(U)=\mathbb Q(V),\quad d=2\text{ or }3,\quad
 U,V,UV/2\in\mathcal O_K,\\
 2<U<12,\quad4<V<64,\qquad
 0<U_\sigma<2,\quad0<V_\sigma<1/4 .
 \end{gathered}
\]

Suppose \(d=3\).
The positive norm \(\mathrm N(V)\) is one, two or three.
For odd norm, \(V\) is a unit at every dyadic prime,
so \(UV/2\) integral makes \(U/2\) integral.
For norm two, \((V)\) is a SINGLE prime
\(\mathfrak p\) of norm two, exponent one.
Use integer prime valuations and let
\(e=\operatorname{ord}_{\mathfrak p}(2)\).
If \(U/2\) were nonintegral, its only deficient
prime could be \(\mathfrak p\), with
\(\operatorname{ord}_{\mathfrak p}(U)=e-1\).
Then
\(\operatorname{ord}_{\mathfrak p}(U-2)=e-1\),
\(\operatorname{ord}_{\mathfrak p}(4V-1)=0\).
The nonzero field square~\eqref{eq:arithB16square}
would have valuation
\[
                         (e-1)+1-2e-(e-1)=1-2e ,
\]
ODD, impossible for a square.
Thus \(W=U/2\) is integral in every case, with
physical \(1<W<6\) and both other conjugates
strictly between zero and one.

We prove the small cubic-integer fact needed here.
Any totally real cubic integer with physical root
\(w>1\) and other roots \(a,b\in(0,1)\) has \(w>5\).
Assume \(w\le5\).
Put integer trace \(S=w+a+b\le6\),
positive integer norms \(A=wab\), \(B=(w-1)(1-a)(1-b)\).
Its monic polynomial is
\[
              P(X)=X^3-SX^2+(A+S-B-1)X-A .
\]
It is positive between the two distinct small roots.
For \(0<X<1\), it increases with \(S\) and
decreases with \(A,B\).
If \(A\ge2\), then
\[
 P(X)\le X^3-6X^2+6X-2
                 <-5(X-3/5)^2-1/5<0 .
\]
If \(B\ge2\), similarly
\(P(X)<-5(X-2/5)^2-1/5<0\).
Thus \(A=B=1\).
For \(S\le5\),
\(P(X)\le X^3-5X^2+4X-1<-(2X-1)^2\le0\).
Only \(S=6\) remains, giving
\(P(X)=X^3-6X^2+5X-1\).
But \(P(5)=-1=(5-w)(5-a)(5-b)\) forces \(w>5\),
a contradiction.
Therefore \(U>10\).

The ACTUAL head edge now gives
\(x_1>(\sqrt5-1)/2>.618\).
The strict conjugate short bounds give
\(V>16\,\mathrm N(V)\ge16\), so its ACTUAL
short edge gives \(S_{16}>4\sqrt7/3>3.5276\).
Together with unchanged \(S_6>2.01815\), use
ONE positive reference
\[
 w_0=.33797P_1+.129085P_6+.150945P_{16}.
\]
Its coordinates are \(.618\), then five
\(.28003\)'s, then ten \(.150945\)'s.
The primitive square on the SAME original \(x\)
gives
\[
 \|x\|^2>
 .618^2+1.40015^2/5+1.50945^2/10
             =1.00185193475>1 .
\]
This excludes EVERY actual cubic field, with
all dyadic ramification allowed and no polynomial
field census.
Actual \(K\) is consequently quadratic.

\subsection{Exact inclusion in the distinct whole-\(B_{16}\) census}
\label{app:arithB16domain}

Every hypothetical whole \(B_{16}\) now satisfies
EXACTLY the actual-root strips and cycles of
Proposition~\ref{prop:pair-output}(iv), with the
original scales \(U=4g_H^2,V=4g_S^2\).
The square in~\eqref{eq:arithB16square} is a FULL
\(K\)-square, with its PHYSICAL positive root.
The exact \(G,L\) are
\eqref{eq:arithB16cycles}.
Acuteness gives \(G\ge0\).
For \(G=0\), \(L=L'=0\); otherwise strict
nonphysical two-root positivity gives \(0<L'<4\).
The own \(U\) and mutual \(L\) lie in the actual
quadratic finite native alphabet or are \(>4\).
Parallel four has its actual feasible null and
is excluded; zero mutual is retained.
These are necessary conditions of ACTUAL walls,
not conditions assigned to a formal projection.

For completeness, this algorithm has a finite
full-order domain justified by the physical strip.
For the maximal basis
\(\omega=(1+\sqrt D)/2\) or \(\sqrt D\),
\(U-U'=b\sqrt\Delta<12\), hence
fundamental discriminant \(\Delta<144\).
Trial \(2\le D\le143\), exact squarefreeness and
this discriminant test exhaust all actual fields.
The FULL safe integer boxes are
\[
 U=a+b\omega:\quad1\le a\le7,\ 1\le b\le5;\qquad
 V=a+b\omega:\quad1\le a\le32,\ 1\le b\le28 .
\]
For \(U\), the gap gives \(b\sqrt D<12\)
in the half basis and \(b\sqrt D<6\) otherwise;
\(0<U'<2\) bounds \(a\).
For \(V\), replace twelve/six by sixty-four/
thirty-two and use \(0<V'<1/4\).
These give the printed boxes, followed by EXACT
strict two-place strips.
At most ten head and twenty-eight short models
occur per field.
Visit EVERY ordered head/short pair.
Exact rational-pair arithmetic, both signed-norm
square-root branches, verified full square identities,
integrality, conjugate and native predicates are
the algorithm of Section~\ref{sec:algorithms}.
All loops are finite; there is no unbounded
coefficient search or numerical root cutoff.
The three disjoint discriminant intervals
\([0,48),[48,96),[96,144)\) cover the domain.
A separate complete combine recomputes field and
pair records; it is a check with the SAME
implementation, not an independent algorithm.

At this precisely proved domain we use the recorded
finite-output proposition~\ref{prop:pair-output}(iv)
as an EXPLICIT computed premise.
Its 43 fields,100 head models,85 short models
and273 pairs leave one necessary pair,
\[
                  U=5+2\sqrt6,\qquad V=10+4\sqrt6 .
\]
It is not an empty census or a realized polytope.
For this pair the ACTUAL head/short edges have
\[
 C_H=(2+\sqrt6)/2>11/5,\qquad
 C_S=2(2+\sqrt6)>44/5.
\]
They give \(x_1>.61\), \(S_{16}>3.56\).
With unchanged \(S_6>2.01815\), ONE reference
\[
                  w=.32837P_1+.127445P_6+.154185P_{16}
\]
has positive coefficients and norm
\[
                  .61^2+1.40815^2/5+1.54185^2/10
                              =1.00640742675>1 .
\]
The original completed square is contradictory.
Every actual short profile, all ground fields and
the complete nonempty quadratic pair domain have
now been covered for whole \(B_{16}\).
The finite output remains a declared numerical
premise; this source transcription is not a
new runtime replay or geometric Lean proof.

\section{Whole classical rank seventeen: actual fields and census inputs}
\label{app:arithBD17actual}

The two rank-seventeen classical forests require
DISTINCT geometric reductions.
Their final quadratic parameter domain happens
to agree, but this will be proved from each
actual original-root configuration.
We first record the common isolated-five lattice
interface used in their head reductions.

\subsection{The actual golden long-end star and all invariant lattices}
\label{app:arithgoldenlongstar}

We use the following explicitly EXTERNAL cutoff:
for the literal golden lattice
\[
 L=\mathbb Z[\varphi]^{n+1},\qquad
 f=-\varphi t^2+x_1^2+\cdots+x_n^2,\qquad n\ge8,
\]
the full integral reflection arrangement has no
compact chamber.
Where infinite covolume of its reflection group
is used, that stronger published literal cutoff
is an external input as identified in the source ledger.
The all-lattice step below is INTERNAL; similarity
or commensurability alone would not justify it.

Put \(u=\varphi-1\) and use the unit-equivalent
form \(f=-ut^2+\sum x_i^2\).
A whole \(D_n\) long-end-five star has old raw
roots \(-e_i+e_{i+1}\), \(-e_{n-1}-e_n\),
and cap \(r=e_0+\varphi e_1\), all square two.
Its full Lorentzian Gram forces actual
\(K=\mathbb Q(\varphi)\): any larger field has
a nonphysical embedding fixing this indefinite
principal.
The old group gives all even spatial sign
changes and permutations.
Conjugating the cap reflection gives maps \(R_i\)
whose only nonidentity block, on zero and \(i\), is
\[
                         \begin{pmatrix}\varphi&-\varphi\\
                                          1&-\varphi\end{pmatrix}.
\]
A whole \(B_n\) head-five star has the SAME
configuration and contains those even signed
permutations, so the argument also applies there.
Arithmeticity supplies a full invariant
\(\mathcal O_K\)-lattice \(\Lambda\), for example
by summing finitely many translates of a lattice
stable under an arithmetic finite-index subgroup.

Let \(\mathcal A\) be the algebra generated by
these lattice-preserving maps.
For distinct spatial \(i,j,k\), even sign
changes \(d_{ij}\) give
\((1-d_{ij})(1-d_{ik})=4E_{ii}\).
Permutations give all spatial \(4E_{ij}\).
The displayed \(R_i\) and the unit \(\varphi\)
then give \(4E_{i0},4E_{0i},4E_{00}\);
explicitly
\[
 \begin{gathered}
 4E_{ii}R_i=4E_{i0}-4\varphi E_{ii},\\
 R_i4E_{ii}=-4\varphi E_{0i}-4\varphi E_{ii},\\
 R_i4E_{i0}=-4\varphi E_{00}-4\varphi E_{i0}.
 \end{gathered}
\]
The golden ring is a PID, for example by its
Minkowski bound \(\sqrt5/2<2\).
Scale the coordinate-content ideal of \(\Lambda\)
to the unit ideal.
The matrix units imply
\(4L\subset\Lambda\subset L\).
The FOURfold sandwich is essential for odd \(n\).

Modulo two, \(\mathcal O_K/2=\mathbb F_4\).
Put
\[
 \ell(t,x)=\varphi t+\sum x_i,\qquad
 H=\ker\ell,\qquad v=(\varphi^2,1,\ldots,1).
\]
The only invariant subspaces are
\(0,\langle v\rangle,H,L/2L\).
Indeed any unequal spatial coordinates give
\(e_i+e_j\) by a transposition difference;
their orbit spans the spatial augmentation
hyperplane, and \(R_i\) supplies a time direction,
giving \(H\).
If every spatial coordinate is equal,
\(R_i\)-invariance forces the line \(\langle v\rangle\).
Here \(v\in H\) exactly for odd \(n\).

For the two residue flags
\[
 W_1=(\Lambda+2L)/2L,\qquad
 W_2=((\Lambda\cap2L)/2)/2L,\qquad W_1\subset W_2,
\]
the content normalization makes \(W_1\ne0\).
If \(W_1\) is full, \(\Lambda=L\).
If \(W_1=H\), differences under even sign
changes of a lift of \(e_i+e_k\), divided by two,
put a coordinate vector in \(W_2\), forcing it full.
If \(W_1=\langle v\rangle\), the same differences
force \(H\subset W_2\); for even \(n\), \(W_2\)
is full, while for odd \(n\) it is \(H\) or full.
Thus, writing \(E\) for the inverse image of \(H\),
the COMPLETE lattice list is
\[
           L,\quad E,\quad\mathcal O_Kv+2L,\quad
           \mathcal O_K(v+2z)+2E\quad(n\text{ odd},\ z\in L).
\]
This includes every lift, not merely every residue
subspace.

Intrinsically,
\[
 E=\{x\in L:f(x)\in2\mathcal O_K\},\qquad
 E^*=L+\mathcal O_K(v/2),\qquad
 f(v)=n-1-2\varphi.
\]
The latter norm has dyadic valuation zero for even
\(n\), exactly one for odd \(n\).
Writing \(x=y+av/2\) shows
\[
                     \{x\in E^*:f(x)\in\mathcal O_K\}=L:
\]
if \(a\) is not divisible by two, the final
norm term has negative valuation and the other
terms are integral.
Hence every isometry of \(E\), and by duality of
\(\mathcal O_Kv+2L=2E^*\), preserves \(L\).
For a deep odd lattice put \(\lambda=v+2z\);
its norm still has valuation one.
Writing \(x=a\lambda/2+e\), \(e\in E\), gives
\[
                     \{x\in\Lambda/2:f(x)\in\mathcal O_K\}=E .
\]
Its isometries therefore preserve \(E\) and then \(L\).
Every invariant lattice in the COMPLETE list has
\(O(f,\Lambda)\subset O(f,L)\).
Thus the actual reflection group lies inside the
literal integral group.
Its arrangement chamber refines the assumed
compact chamber, contradicting the external
literal no-compact-chamber cutoff
(or the stronger infinite-covolume statement).
This excludes the actual whole-\(B/D\) isolated
long-end-five configuration, with every compatible
invariant lattice retained.

\subsection{Whole \(B_{17}\): the fresh window and actual short wall}
\label{app:arithB17physical}

Use actual \(B_{17}\) roots and
\(x_1>\cdots>x_{17}>0\) from one original
farthest vector.
The generic parameter \(M=12\) proves the bulk
through thirteen; literal far helpers at indices
at least six are excluded arithmetically.
At fourteen the right \(B_3\) extra short-three
extension is the sealed \(F_4\): own zero has
its affine principal, every positive own has
ratio \(\ge1+1/\sqrt2\).
At fifteen the right \(B_2\) extra own-three/
short-three profile has the opposite-\(\sqrt2\)
Schur cost
\(15(3-2\sqrt2)/4+1/2>1\).
Its zero/four/six and parallel owns have actual
nulls; larger owns and the full \(B_{17}\)
bridge are directly strong.
The bridge residual is \(1/(2i)\), and even
own zero has ratio
\(1/\sqrt{2i}+\sqrt i/2>100/61\).
All required left corrections use genuine
long node sixteen.
Thus six through fifteen are strong.

If five is weak with finite own, its ACTUAL
finite helper point gives \(x_1>.65\).
The reference through one, six and every later
strong knot gives
\[
 e_B>E_{15}+.65^2+(2.01815-.65)^2/5-E_6
                       >.92564+.7968-.7>1 .
\]
For dotted weak five put \(t=2g\in(2,12/5)\).
Its actual path cycle makes \(t\) integral,
with all other conjugates in \((-.69,.29)\).
Degree at least four contradicts
\(2.4(.69)^3<1\); degree one has no integer;
degree two has trace two and norm
\(t(2-t)\in(-24/25,0)\).
In degree three a positive small conjugate gives
norm magnitude \(<1\).
Two negative small conjugates force norm one;
trace two violates \((t-2)^2\ge4/t\).
Trace one forces \(t>9/4\), where the integer
second coefficient \(1/t+t-t^2\) lies strictly
between \(-3,-2\).
Thus five is strong too.

The ACTUAL retained \(B_{16}\) head list is
isolation or long-endpoint-three.
The nonisolated case gives \(x_1>.643\),
whose legal first-five graft gains \(>.099405\)
over \(E_{15}>.92564\), impossible.
For an isolated head, small own/nulls remove
zero, three, four, six and parallel;
the golden long-star lemma removes five.
For finite own at least seven or dotted, the
already proved actual first-three-coordinate
point gives \(x\cdot Z\ge H\) and
\(F=.8H-.16R>.2893\).
Paying \(.08\) at the INDIVIDUAL \(P_5\)
coefficient gives gain \(>.0325\).

Independently, a legal first-five graft at
\(x_1\ge.605\) would give
\(E_{15}+(5/4)(.605-.361)^2>1.00006\).
Hence \(x_1<.605\).
For the actual head parameter \(U=4g_H^2\),
\[
 K=\mathbb Q(U),\quad U\in\mathcal O_K,\quad
               2<U<28/3,\quad0<U_\sigma<2 .
\]
The raw head has square \(2U\), spatial \(Ue_1\),
time square \(U(U-2)\).
This original full-rank Lorentzian principal
proves ACTUAL field equality.
FINITE HEADS ARE RETAINED until excluded below.

We first derive a preliminary short bound without
using that head.
If actual \(C_S\ge15\), then \(S_{17}>3.856\).
For strong five, the positive reference at
\((5,1.804),(17,3.856)\) has norm \(1.0017752\).
For weak five, use \((1,.59),(6,2),(17,3.856)\);
its norm is \(.74572+3.444736/11>1.05\).
Therefore \(C_S<15\).
The actual retained \(A_{16}\) short-white list
is exactly isolation, helper three at
\(j=1,2,15,16\), or helper four at \(j=1,16\).

For helper three put \(a=2\sqrt2\gamma\),
raw cap square two and helper product \(-1\).
Actual cycles make \(a,a^2/2\) integral.
Its complete formulas are
\[
 A_j^2=(17a^2+4ja+4j-8)/4,\qquad
 C_S={17a+2j\over2\sqrt{34-j(17-j)}} .
\]
Here \(A_j\) is the RAW cap time coefficient.
For helper four use raw cap square one,
time square \(17\gamma^2+2j\gamma+j-1\),
\(C_S=17\gamma+j\).
The original connected Lorentzian principal gives
ACTUAL \(K=\mathbb Q(a)\) or \(\mathbb Q(\gamma)\).

Own zero at helper-three one/two has negative/
zero time; near fifteen/sixteen has the actual
positive null
\[
 c+n_{14}+2n_{15}+2n_{16}+\sqrt2n_{17},\qquad
 c+n_{15}+2n_{16}+\sqrt2n_{17},
\]
respectively.
Helper-four own zero at one has zero time;
near sixteen already violates the length bound.
Parallel always has the actual future feasible
null pair.
Every null is feasible against ALL original
walls termwise.

At helper-three two,
\(a<56/17\), every other \(-8/17<a_\sigma<0\).
Its norm is a unit, incompatible with
\(a^2/2\) integral.
Near next-end fifteen leaves only own three
\(a=\sqrt2\), with inadmissible conjugate
\((86-60\sqrt2)/4>0\).
Far helper-four has \(\gamma<14/17\), leaving
only owns three/four/five.
The rational case is impossible;
the other time conjugates are
\(17/2-\sqrt2>0\), \((38-21\varphi)/4>0\).
Thus all four non-end roles are excluded.

For far-end-three put \(v=a\); for near-end-three
put \(v=a+2\).
Their nonphysical strips have absolute bound
\(.617\), with smaller one-sided bound \(.382\);
physical \(v<7.7\), and \(v^2/2\) is integral.
Degree at least three would give nonzero integer
norm of magnitude \(<2.932<3\), but its
half-square norm requires divisibility incompatible
with norms one or two.
Hence actual degree is two.
Even trace/norm and the exact negative-time
intervals give the COMPLETE tables
\[
\begin{array}{c|c|c}
 \text{role}&(\operatorname{Tr}v,\mathrm N(v))&a\\ \hline
 \text{far}&(4,-2),(6,2),(6,-2),(6,-4)&
       2+\sqrt6,\ 3+\sqrt7,\ 3+\sqrt{11},\ 3+\sqrt{13}\\
 \text{near}&(4,2),(6,2),(6,-2),(8,4)&
       \sqrt2,\ 1+\sqrt7,\ 1+\sqrt{11},\ 2+2\sqrt3 .
\end{array}
\]
For coverage: far traces \(2,4,6\) and norms
\(-4,-2,2\) exhaust the quadratic formula;
near traces \(4,6,8\) and norms \(-2,2,4\)
do likewise.
All other rows have wrong physical/conjugate
signs or are rational/nonreal.
The exact time signs use
\[
174-72\sqrt6,\ 382-106\sqrt{13},\
348-106\sqrt{11},\ 280-106\sqrt7<0
\]
for far roots, and
\[
90-64\sqrt2,\ 324-98\sqrt{11},\
256-98\sqrt7,\ 456-264\sqrt3<0
\]
for near \(v\)'s.
Thus the full exact intervals, not just safe
decimal overstrips, are imposed.

An isolated short cap has invariant
\(V=4\gamma_S^2\), raw square \(V\), spatial
\((V/2)P_{17}\), time square \(V(17V-4)/4\).
Its original Lorentzian principal gives
\(K=\mathbb Q(V)\), \(0<V_\sigma<4/17\).
Finite own would have nonzero integral norm
\(<16/17\); parallel is its null pair.
Thus it is dotted.
The preliminary \(C_S<15\) gives \(V<900/17\);
degree at least four gives norm
\(<57600/83521<1\).
So the ACTUAL field, whether the short cap is
isolated or one of the helper models, has
degree two or three.
The rational field was separately excluded by
the WHOLE rank-eighteen Meyer obstruction.

\subsection{Whole \(B_{17}\): exact head/short compatibility}
\label{app:arithB17compatibility}

Before excluding any finite head, its actual
degree bound and \(0<U_\sigma<2\) give the
COMPLETE remaining finite list
\[
\begin{array}{c|c}
 \text{native}&U\\ \hline
7&U^3-5U^2+6U-1=0\\
8&2+\sqrt2\\
10&(5+\sqrt5)/2\\
12&2+\sqrt3\\
18&U^3-6U^2+9U-3=0 .
\end{array}
\]
Indeed \(\phi(m)\le6\) gives
\(3,4,5,6,7,8,9,10,12,14,18\);
the earlier null/golden exclusions remove
three/four/five/six.
Nine and fourteen have a nonphysical \(U>2\).
In particular finite native eight in
\(\mathbb Q(\sqrt2)\) MUST still be retained.

The individual \(P_5,P_{15}\) prices exceed
\(.14,.13\).
For any actual short pressure \(S_{17}>T\),
ONE vector
\[
 w_{15}+.4(Z-P_5/5)+.13(P_{17}-P_{15})
\]
spends \(.08,.13\) before using either floor.
The head residual is zero-sum on the first five;
the terminal residual is disjoint with square two.
All FULL crosses with the reference and each
other vanish.
Thus
\[
                 \|x\|^2>
                 E_{15}+.0325+.26(T-f_{15})-.0169(2).
\]
If \(C_S\ge11\), the actual short edge gives
\(S_{17}\ge\sqrt{85/6}>3.763\).
With \(E_{15}>.92564,f_{15}<3.442\),
the right side is \(>1.0078\).
Therefore ACTUAL \(C_S<11\).

Of the eight short-helper models, the actual
coshes
\((17a+2)/(6\sqrt2)\) and
\((17a+32)/(6\sqrt2)\)
leave only far \(a=2+\sqrt6\) or near
\(a=\sqrt2\).
For both remaining rows the raw cap time norm
has prime twenty-three to odd valuation:
\[
\begin{array}{c|c|c}
 \text{role}&H_c^2&-16\mathrm N(H_c^2)\\ \hline
 \text{far}&(87+36\sqrt6)/2&828\\
 \text{near}&(45+32\sqrt2)/2&92 .
\end{array}
\]
The ACTUAL field equals that of \(a\).
The COMPLETE compatible dotted head strips
\(4<U<11,0<U'<2\), with their
\(h_0=-\mathrm N(U(U-2))\), are
\[
\begin{array}{c|l|l}
 K&U&h_0\\ \hline
 \mathbb Q(\sqrt6)&
 3+\sqrt6,\ 4+\sqrt6,\ 5+2\sqrt6,\ 6+2\sqrt6&
 15,20,15,96\\
 \mathbb Q(\sqrt2)&
 3+\sqrt2,\ 3+2\sqrt2,\ 4+2\sqrt2,\ 5+3\sqrt2,\ 6+3\sqrt2&
 7,7,32,63,36\\
 \mathbb Q(\sqrt2)\text{ finite}&2+\sqrt2&4 .
\end{array}
\]
The full rings are the displayed square-root
orders.
Writing \(U=A+B\sqrt d\) gives
\(B<11/(2\sqrt d)\), and
\(A=\lfloor B\sqrt d\rfloor+1,+2\);
the physical strip gives exactly the rows.
There is no finite head in \(\mathbb Q(\sqrt6)\);
the finite \(\sqrt2\) head is explicitly included.
Both raw roots occupy the SAME actual \(K\)-time
line, so the rational time-norm ratio
\(p/(16h_0)\) must be a square.
Its twenty-three valuation is odd in EVERY row.
No nonisolated short cap survives.
The actual short cap is now isolated dotted:
\[
                 K=\mathbb Q(V),\quad
                 4<V<484/17,\quad0<V_\sigma<4/17 .
\]

Derive the B17 cycles in THEIR OWN actual scales:
head raw square \(2U\), time square \(U(U-2)\);
short raw square \(V\), time square \(V(17V-4)/4\).
The actual long path makes
\(a=2(r_H,r_S)\) integral and its actual
two-cycle makes \(L=a^2/(2UV)\) integral.
The original common time line gives
\[
 \begin{gathered}
 R_0=U(U-2)V(17V-4)\in K^2,\qquad
 a=UV-\sqrt{R_0},\quad\sqrt{R_0}>0\text{ physically},\\
 L-a=8UV-17V-2U+4=:d_0,\qquad
 a^2=2UV(a+d_0).
 \end{gathered}
\tag{B17C}\label{eq:arithB17cycles}
\]
No D17 fork normalization was used.

First exclude ALL finite heads.
In degree three,
\(0<\mathrm N(V)<(484/17)(4/17)^2<2\),
so \(V\) is a unit.
At every dyadic prime, \(d_0\equiv-V\) is a unit;
the cycle equation forces \(a\) into the prime
and then
\[
                     2\operatorname{ord}(a)
                      =\operatorname{ord}(2)+\operatorname{ord}(U).
\]
For finite cubic heads seven/eighteen,
the actual field has odd discriminant
forty-nine/eighty-one, and \(\mathrm N(U)=1,3\).
Thus two is unramified and \(U\) a dyadic unit;
the equality reads an even number equals one.
Both finite cubic heads are excluded without
assuming \(U>4\).

For finite quadratic heads eight/ten/twelve,
the FULL short strips have exactly these seven
models and necessary rational time-norm ratios:
\[
\begin{array}{c|c|c}
 \text{head}&V&
 \mathrm N(H_S^2/H_H^2)\\ \hline
8&3+2\sqrt2&103/64\\
8&10+7\sqrt2&383/16\\
12&7+4\sqrt3&647/48\\
12&14+8\sqrt3&61\\
10&2+3\varphi&171/80\\
10&5+8\varphi&919/80\\
10&7+11\varphi&239/16 .
\end{array}
\]
Coverage follows from the full-order coefficient
ranges \(B=1,\ldots,10\) in \(\sqrt2\),
\(1,\ldots,8\) in \(\sqrt3\),
\(1,\ldots,12\) in \(\varphi\).
The corresponding complete ceil arrays for \(A\) are
\[
 \begin{gathered}
 2,3,5,6,8,9,10,12,13,15;\\
 2,4,6,7,9,11,13,14;\\
 1,2,2,3,4,4,5,5,6,7,7,8 .
 \end{gathered}
\]
Exact \(0<V'<4/17\) retains only
\(B=2,7;\ 4,8;\ 3,8,11\), respectively.
Rational brackets for the square roots and
\(\varphi\) decide every ceil and endpoint.
If \(n=\mathrm N(V)\), \(T=\operatorname{Tr}(V)\),
and \(h_0=-\mathrm N(U(U-2))\), each ratio is
\[
                     {n(68T-289n-16)\over16h_0}.
\]
None in the COMPLETE table is a rational square:
the reduced denominator has odd three/five
valuation where indicated, and the remaining
integer numerators \(103,383,61,239\) are
nonsquares.
Thus every finite quadratic head is excluded.
Only NOW conclude \(4<U<28/3\).

In degree three the unit \(V\) and cycle valuation
remain valid.
If two is unramified they force \(U/2\) integral.
Its physical value is in \((2,14/3)\), both other
conjugates in \((0,1)\), contradicting the
complete cubic-integer lemma of
Section~\ref{app:arithB16cubic}, which requires
physical value \(>5\).
If two ramifies, \(V\)'s monic unit polynomial
\(X^3-AX^2+BX-1\) has a repeated factor modulo
two, forcing \(A,B\) ODD.
Its two small positive roots in \((0,4/17)\)
and physical \(V<484/17\) give
\[
 A=19,21,23,25,27,\qquad
 B^2>4(A-8/17),\qquad
 B<{272A+4849\over1156}.
\]
The least and largest possible odd \(B\) are,
respectively,
\[
                (9,7),\ (11,9),\ (11,9),\ (11,9),\ (11,9).
\]
Every row is empty.
This excludes the ENTIRE cubic branch in the
actual B17 domain without a field census.
The actual field is quadratic.

\subsection{Whole \(B_{17}\)'s actual inclusion in the shared quadratic domain}
\label{app:arithB17domain}

Rename ONLY the algebraic variable \(W=V\),
\(g=a\) in~\eqref{eq:arithB17cycles}.
We have proved from the ACTUAL B17 geometry
\[
 \begin{gathered}
 K=\mathbb Q(U)=\mathbb Q(W),\quad U,W\in\mathcal O_K,\\
 4<U<11,\quad0<U'<2,\qquad
 2<W<676/17,\quad0<W'<4/17,\\
 g=UW-\sqrt{U(U-2)W(17W-4)},\quad
 g,L=g^2/(2UW)\in\mathcal O_K .
 \end{gathered}
\]
The physical \(g\) is strictly negative:
\(d_0>(8U-17)2-2U+4=14U-30>0\),
so \(R_0-(UW)^2=2UWd_0>0\).
Thus \(L>0\).
At the other embedding actual positive definiteness
gives \(0<L'<4\).
The physical \(L\) is a true finite native
invariant or \(>4\); parallel four has its
actual null pair.
Every predicate of
Proposition~\ref{prop:pair-output}(iii) is now
satisfied.
This domain inclusion uses B17's short raw
square \(V\), not D17's fork square \(4W\).
The common algebraic formula is a proved
coincidence of necessary invariants.

The complete full-order algorithm and its
explicit output premise are already stated in
Section~\ref{sec:algorithms}.
Its discriminant bound \(\Delta<121\), two
conjugate-strip candidates for each \(U\)
coefficient, one for each \(W\), and
\(b_U\le4,b_W\le17\) cover every field and
ordered pair and make every loop finite.
Both signed-norm square branches are tested;
a rational-square norm alone is not sufficient.
The old 36-field/120-pair output leaves TWO
necessary pairs, in \(\sqrt{13}\) and \(\sqrt{17}\).
The separate old combine authenticates coverage/
counters and survivors; it does NOT independently
recompute all118 rejected square tests.
Those remain part of the explicit source-bound
numerical-output premise.

For the two actual surviving heads,
\(u=4/U\) lies in \((.92,.94)\) and
\((.719,.72)\).
Their actual dotted head points give, respectively,
\[
 H>.535,\ R<.51;\qquad H>.482,\ R<.54 .
\]
Positive rational squaring proves the radical
bounds \(\sqrt{2-u}>1.0295,1.1313\) and
\(\sqrt{u(1-u)}<.272,.4495\), respectively.
Paying \(.36\) at the head costs ONLY
\(.072\) from \(P_5\), and gives gains
\(>.085\), \(>.043\).
Both actual \(W\)'s exceed \(104/5\);
the B17 ACTUAL short edge has
\(C_S^2=17W/4>442/5\), hence
\(S_{17}>\sqrt{357/26}>3.7\).

Use ONE reference through fourteen:
\[
                   w_{14}+.36(Z-P_5/5)+.13(P_{17}-P_{14}).
\]
The individual prices \(.072,.13\) are legal.
The two residuals are disjoint, orthogonal to
the reference and each other; terminal norm is three.
Thus the SAME original vector would satisfy
\[
 \|x\|^2>.9084+.043+.26(3.7-3.311)-.0169(3)
                              =1.00184>1 .
\]
Both complete census survivors are excluded.
This closes whole \(B_{17}\), conditional only
on the explicitly identified external facts and
the recorded complete quadratic output premise.

\subsection{Whole \(D_{17}\): fresh geometry before any field census}
\label{app:arithD17geometry}

Use ACTUAL unit \(D_{17}\) roots, long spine one
through fifteen, forks sixteen/seventeen, and
\[
 x_1>\cdots>x_{16}>|x_{17}|,\qquad\|x\|<1 .
\]
The signed coordinate is retained.
The legal generic \(M=11\) window gives strong
six through twelve.
At thirteen/fourteen the COMPLETE right D4/D3
lists leave only the pure own-five/fork-three
or own-five/end-three packets.
Here completeness uses the SAME actual finite
role list printed in Section~\ref{app:arithD13rankthree},
with the FRESH indices thirteen/fourteen.
The D4 long-three and D3 middle-three roles have
residual \(1/2\), shift \(1/2\);
fork/end-three has residual \(1/2\) or \(5/8\),
shift \(1/4\).
D3 end-four has residual \(1/4\), shift
\(1/(2\sqrt2)\); sealed end-five has residual
\((5-3\varphi)/8\), shift \(\varphi/4\).
The \(D_{17}\) bridge residual is \(1/(2i)\);
the \(A_{17}\) end-three bridge at fourteen
has residual \(1/8+1/28=9/56\).
Their positive-own numerators make them strong,
and their own-zero cycles are affine.
For single-fork/end-three, own zero, three,
four and six retain the actual affine \(E_8\),
cycle, \(F_4\) and \(6,3\) principals.
The own-three cycle has FIVE actual nodes at
D4 and FOUR at D3 end-three; it is not
replaced by a triangle.
The far-helper arithmetic lemma removes the
other nonisolated possibility.
The original LEFT isolated-core/H4/A2/singleton
proof applies with \(2i\ge26\), so its previously
proved \(2i\ge18\) loss margins suffice.
The remaining fork-helper raw Schur constant
is ACTUAL seventeen; the complete norm/trace
argument leaves only its golden own.
There is no spectator or extra long root here.
Their actual LEFT \(H_4\) points have divided
displacement \(V\), flat one through \(K=i-3\),
\(V^2<i\), and pressure
\[
 x\cdot V>G>\sqrt{i-\kappa/(\varphi+5/8)}-.091,\quad
 \kappa=1/2,\ 5/8 .
\]
Every point uses positive actual walls, positive
future time and all original-wall signs.
The relevant rows are
\[
\begin{array}{c|c|c|c|c}
 i&K&\mu&G>&f_K<\\ \hline
13&10&.16&3.48&2.732\\
14&11&.15&3.61&2.886 .
\end{array}
\]
The radical seals are
\(13-25/112>3.571^2\),
\(14-125/448>3.701^2\).
The INDIVIDUAL terminal prices exceed
\(.16,.15\), in both initial-five references.
The residual \(V-P_K\) vanishes on the whole
reference support and has norm \(<3\).

If five is weak, its actual endpoint reference
has energies \(L_{10}>.8870,L_{11}>.9100\).
The two paid squares are
\[
 .8870+.32(3.48-2.732)-.0256(3)>1.0495,\qquad
 .9100+.30(3.61-2.886)-.0225(3)>1.0597 .
\]
If five is strong, first classify the ACTUAL
head: the retained \(D_{16}\) admits only isolation
or simple three at its long endpoint two.
The nonisolated row gives \(x_1>.643\), hence
legal first-five gain \(>.0994\).
An isolated head has small owns/nulls excluded
as before, including five by
Section~\ref{app:arithgoldenlongstar}.
For finite own at least seven, its actual finite
head point gives gain \(>.0135\) at cost \(.1\)
from \(P_5\).
Explicitly for \(1<u<1.232\),
\[
 \eta=H-.361-(R-.2)/4>.0135 .
\]
On \(u\in[1,1.1],[1.1,1.2],[1.2,1.232]\),
the respective rational seals are
\[
 H>.6708-.1871,\quad .63245-.204125,\quad .6196-.2,
 \qquad R\le.5,.415,.36 .
\]
The middle second-radical square has maximum
\(1/24\); all bounds are positive squarings.
They give gains \(>.0477,.013575,.0186\).

For a dotted head use its ACTUAL positive core,
\(Z=(1-u/2)e_1+(u/2)e_2\), \(0<u<1\).
Paying \(.36\) at \(P_5\), cost \(.072\),
gives gain
\[
                .72(H-b)-.1296(R-1/5)>.04 .
\tag{D17H}\label{eq:arithD17dotgain}
\]
Here is the full continuous seal rather than a
borrowed actual-\(I\) theorem.
On \([0,1/2],[1/2,3/4],[3/4,1]\), use lower
chords for \(\sqrt{2-u}\)
\[
1.4142-.379u,\quad1.4381-.4268u,\quad1.4721-.4721u
\]
and upper chords for \(R\)
\[
1-.75u,\quad.8125-.375u,\quad.625-.125u.
\]
With \(.5091<9\sqrt2/25<.5092\), the three
gain bounds reduce to
\[
\begin{array}{l}
 .35636922-.0958u-.5092\sqrt{u(1-u)},\\
 .39283671-.16873u-.5092\sqrt{u(1-u)},\\
 .43444611-.2242u-.5092\sqrt{u(1-u)} .
\end{array}
\]
The circle inequality and squared-root upper
bounds \(.5182,.5365,.5564\) give minima
\(>.04936922,.04022171,.04414611\).
No \(I\) wall or energy enters.

Every head residual is zero-sum on the first
five and orthogonal to the H4 residual beyond
\(K=10/11\).
The bare strong-five H4 bounds are
\(.98692,.99786\).
Adding even the smallest LEGALLY funded head
gain \(>.0135\) excludes both.
All six through fourteen are consequently strong
BEFORE deciding five or the final head.

Weak five with finite own gives actual
\(x_1>.65\); its rebuilt six-through-fourteen
reference exceeds one.
Its dotted case is precisely the complete
minimal-degree norm/trace argument already
proved for the SAME inverse entries in
Section~\ref{app:arithB17physical}.
Thus five is strong.
A nonisolated head now gives
\(E_{14}+.0994>1\), so it is isolated.
The legal graft at \(x_1\ge.632\) gives
\(E_{14}+(5/4)(.632-.361)^2>1.00020125\).
Therefore \(x_1<.632\).
The original head raw normalization gives
\[
 r_H^2=2U,\quad\text{spatial }r_H=Ue_1,\quad
 H_0^2=U(U-2),\quad K=\mathbb Q(U),\quad
               2<U<11,\quad0<U_\sigma<2 .
\]
Finite heads remain possible at this stage.

For completeness the conjugate interval itself
bounds their COMPLETE native alphabet.
For \(U=2+2\cos(2\pi/m)\), a coprime exponent
\(1<a<m/4\) gives a forbidden conjugate \(>2\).
For odd \(m\ge9\) use two.
For even \(m\ge14\) not divisible by three use three.
For \(m=6k\), \(k=4,5\), use five/seven;
for larger even \(k\) use \(k-1\), or \(k-3\)
if \(k\equiv1\pmod3\);
for larger odd \(k\) use \(k-2\), or \(k-4\)
if \(k\equiv2\pmod3\).
Each is coprime and below \(3k/2\).
The remaining labels are
\(3,4,5,6,7,8,10,12,18\).
The prior null/golden exclusions leave exactly
seven/eight/ten/twelve/eighteen, with the
parameters already printed for \(B_{17}\).

For either actual signed fork direction
\(V_f=P_{16}\pm e_{17}\), if \(C_f\ge13\),
the ORIGINAL edge gives
\(x\cdot V_f\ge\sqrt{102/7}>3.817\).
Use the SAME head gain and a separately funded
\(.13(V_f-P_{14})\), whose residual norm is three
and cross with the head/reference is zero.
Then
\[
 \|x\|^2>.9084+.0135+.26(3.817-3.311)-.0169(3)
                                  =1.00276>1 .
\]
Thus BOTH actual fork coshes are \(<13\),
without equating their parameters.

\subsection{The complete actual \(D_{17}\) fork table}
\label{app:arithD17forktable}

Fix either genuine fork \(f\).
Deleting it retains ACTUAL \(A_{16}\).
Finite classification gives isolation, helper three
at \(j=1,2,15,16\), or helper four at \(j=1,16\).
For three put \(t=2\gamma\), raw cap square two,
white product \(-1\), own product \(-t\).
For four put \(a=2\sqrt2\gamma\), raw square one,
white product \(-1\), own product \(-a/2\).
Actual cycles make the parameter integral,
and for four also \(a^2/2\) integral.
The original Lorentzian principal gives ACTUAL
field equality with that parameter.
The COMPLETE raw time and actual edge table is
\[
\begin{array}{c|c|c}
 \text{helper}&H_c^2&C_f\\ \hline
3@1&(17t^2+4t-4)/4&(17t+2)/(6\sqrt2)\\
3@2&(17t^2+8t)/4&(17t+4)/4\\
3@15&(17t^2+60t+52)/4&(17t+30)/4\\
3@16&(17t^2+30t+9)/4&(17t+15)/(6\sqrt2)\\
4@1&(17a^2+8a)/16&(17a+4)/4\\
4@16&(17a^2+60a+52)/16&(17a+30)/4 .
\end{array}
\]
The near endpoint THREE formula differs from
the B17 short formula and is not substituted
for it.
Every other conjugate has negative displayed
time square.
Every null below uses UNIT ACTUAL cap normals;
its positive time and all original-wall signs
follow termwise.

The \(3@15\) own-zero kernel is
\(n_{14}+2n_{15}+n_{16}+n_{17}+c\);
own three has its triangle.
All other owns have \(t\ge\sqrt2\), \(C_f>13\).
The \(4@16\) own-zero kernel is
\(n_{14}+2n_{15}+2n_{16}+n_{17}+\sqrt2c\);
every positive \(a\ge\sqrt2\) gives \(C_f>13\).
At \(3@2\), own zero has zero time; positive
owns satisfy
\(1\le t<48/17\), \(-8/17<t_\sigma<0\).
Degree at least three has nonzero integer norm
magnitude \(<1\).
Degree two leaves norm \(-1\), trace one or two;
the golden root violates the conjugate strip,
leaving \(1+\sqrt2\).
Its raw time norm is \(-47/16\).
For a finite head the only compatible head is
eight, of time norm \(-4\), and the ratio
\(47/64\) is nonsquare.
For a DOTTED head, the complete actual
\(\sqrt2\) head strip gives time norms
\(-7,-7,-32,-63,-36\).
Prime forty-seven again has odd ratio valuation.
Thus this row is excluded under EVERY head.
At \(4@1\) the same norm strips leave only
\(a=1+\sqrt2\), but \(a^2/2=3/2+\sqrt2\)
is not integral.
Its own zero has zero time.
Every four non-end role is exhausted, with
finite, parallel and dotted own retained;
parallel everywhere has \(c+n_f\).
Far-end own zero has negative time, while
near-end own zero gives the forbidden rational
actual field.

For the two remaining end-three roles put
\(v=t\) far, \(v=t+1\) near.
Their raw times are
\[
 H_c^2=(17v^2\pm4v-4)/4,
\]
with upper sign far.
The COMPLETE strips are
\[
\begin{array}{c|c|c}
 \text{role}&\text{physical }v&\text{every other }v_\sigma\\ \hline
 \text{far}&1\le v<6.372&-.617<v_\sigma<.382\\
 \text{near}&2\le v<6.607&-.382<v_\sigma<.617 .
\end{array}
\]
The exact nonphysical endpoints are
\((\mp2\pm6\sqrt2)/17\), and the physical
upper endpoints are
\((78\sqrt2\mp2)/17\).
These exact signs are retained whenever a
safe decimal overstrip admits an extra row.

Let \(g_0\) be the ACTUAL raw head/helper
product.
Its actual triangle or complete old-chain
cycle makes it integral.
The common original time-line identities are
\[
 \begin{split}
 4g_0^2-4g_0U(t+2)+U^2(t+2)^2
                  &=U(U-2)(17t^2+4t-4),\\
 4g_0^2-4g_0U(t+1)+U^2(t+1)^2
                  &=U(U-2)(17t^2+30t+9).
 \end{split}
\]
Reducing modulo \(4\mathcal O_K\) gives the
necessary parity statement
\[
                              Uv^2\in2\mathcal O_K .
\tag{D17P}\label{eq:arithD17helperparity}
\]
It does NOT justify \(Uv/2\) integral at
ramified dyadic primes.
The physical and conjugate bounds belong to
actual \(K=\mathbb Q(v)\), not merely an
unexamined smaller parameter field.

\subsection{Whole \(D_{17}\): every finite head is excluded}
\label{app:arithD17finiteheads}

Assume one of the FIVE finite heads retained
in Section~\ref{app:arithD17geometry}.
Its actual field is the corresponding quadratic
or cubic field, so the complete fork table applies
with that field equality.

For heads seven/ten/eighteen, \(U\) is a dyadic
unit and two is unramified, by discriminants
forty-nine/five/eighty-one.
Parity~\eqref{eq:arithD17helperparity} makes \(v/2\)
integral.
In the far role its nonzero norm has magnitude
\(<3.186(.3085)^{d-1}<1\) for \(d=2,3\).
For the near role the cubic norm is
\(<3.304(.309)^2<1\).
In the remaining quadratic golden near case,
\(v/2=a+b\varphi\) has physical value \(<3.304\),
other value in \((-.191,.309)\).
The gap gives \(b\le1\), \(b>0\).
For \(b=1\) the interval requires
\(.427<a<.927\), no integer.
Both end-helper roles are impossible.

For head eight use \(K=\mathbb Q(s)\), \(s=\sqrt2\),
\(U=2+s\).
At its dyadic prime \(\operatorname{ord}(U)=1\),
\(\operatorname{ord}(2)=2\).
Parity requires \(s\mid v\).
The COMPLETE far strip has
\(t=1+s,3+2s\), neither divisible by \(s\).
The complete near strip has
\(v=2+s,3+2s\); only \(2+s\) survives,
giving \(t=1+s\).
Its raw time square has norm \(-23/4\),
with nonsquare ratio \(23/16\) to the actual
head time norm \(-4\).
For head twelve, \(U=2+\sqrt3\) is a dyadic unit
although two ramifies.
Parity requires \(v\) in its dyadic prime,
where \(\sqrt3=1\) in the residue field.
The COMPLETE far/near strips leave only
\(t=2+\sqrt3\), respectively \(v=2+\sqrt3\);
both have nonzero residue one.
Thus BOTH actual forks are isolated under
EVERY finite head.
The assertion is presently conditional on a
finite head, not carried into the dotted case.

For an isolated positive fork put \(V=4\gamma_f^2\).
Its ACTUAL raw root has square \(2V\),
sole old product \(-V\), time square
\[
              H_f^2=V(17V-8)/4,\quad
              K=\mathbb Q(V),\quad0<V_\sigma<8/17 .
\]
Own zero has no spatial wall part.
Own three has the actual affine \(E_8\) arms
one/two/five; own four has positive affine
vector
\(2n_{15}+2n_f+n_{14}+n_k+\sqrt2c\),
where \(k\) is the other fork.
Own six has its \(6,3\) path; parallel has
the actual null pair.
All are future feasible against EVERY original
wall.
For finite labels at least seven,
degree at least three contradicts
\(\mathrm N(V)<4(1/2)^2\le1\).
In degree two the strict strip forces norm one,
trace four, namely native twelve.
Native five is the other finite candidate.
The cubic/silver head fields cannot equal those
of these candidates.
At golden head ten the fork-five time norm
is \(-55/16\), giving nonsquare ratio \(11/16\)
to head norm \(-5\).
At head twelve the fork-twelve time norm
is \(-191/16\), ratio \(191/48\) to head norm
\(-3\), also nonsquare.
Both forks are DOTTED under every finite head.

The actual head/fork LONG cycle gives
\(g=(r_H,r_f)\) integral; the actual two-cycle
gives \(L=g^2/(UV)\) integral.
Their spatial product is \(UV/2\).
The SAME original time line gives
\[
 (g-UV/2)^2=U(U-2)V(17V-8)/4,\qquad
 L-g=4UV-17V/2-2U+4 .
\]
Hence \(W=V/2\) is integral, by subtracting
eight copies of \(V\) from the integral \(17V/2\).
This uses BOTH cycles without a dyadic-unit
assumption.
The actual fork bound gives
\[
                  2<W<676/17,\quad0<W_\sigma<4/17,\quad
                  H_f^2=W(17W-4).
\tag{D17W}\label{eq:arithD17forkstrip}
\]

For the three quadratic finite head fields,
the COMPLETE full-order strips and time-norm
ratios are
\[
\begin{array}{c|l|l}
 \text{head}&W&\mathrm N(H_f^2/H_0^2)\\ \hline
8&
 3+2\sqrt2,\ 10+7\sqrt2,\ 17+12\sqrt2&
 103/4,\ 383,\ 2007/4\\
10&
 2+3\varphi,\ 5+8\varphi,\ 7+11\varphi,\ 10+16\varphi&
 171/5,\ 919/5,\ 239,\ 5104/5\\
12&
 7+4\sqrt3,\ 14+8\sqrt3&
 647/3,\ 976 .
\end{array}
\]
For coverage, the positive coefficient bounds
are fourteen/seventeen/eleven, and for each
coefficient the small-conjugate width \(4/17<1\)
permits at most one integer constant term.
The physical bound removes the otherwise
small-conjugate \(20+14\sqrt2\).
For \(n=\mathrm N(W),T=\operatorname{Tr}(W)\),
\(\mathrm N(17W-4)=289n-68T+16\)
gives every printed ratio.
Each is a nonsquare by an odd denominator
three/five valuation or the nonsquare integer
numerator; \(976=16(61)\).
The original common \(K\)-time line therefore
excludes every row.

For finite heads seven and eighteen the FULL
rings are \(\mathbb Z[U]\).
Their polynomials have discriminants \(49,81\);
the former translated by four is seven-Eisenstein,
the latter is three-Eisenstein.
The distinct-valuation uniformizer argument
already proved for the cubic bases in
Section~\ref{app:arithA16anchors} excludes the
only possible index prime.
Both are irreducible modulo two, so two is inert.
The positive norm of \(W\) is \(<10816/4913<3\).
An even norm would be divisible by eight;
hence \(\mathrm N(W)=1\), \(W>289/16\).

Let \(\alpha,\beta,\chi\) be physical and two
small conjugates of the actual finite-head generator.
For head seven the exact root and derivative
bounds are
\[
 \begin{gathered}
3.246<\alpha<3.248,\quad1.554<\beta<1.556,\quad
 .198<\chi<.199,\\
5.149<f'(\alpha)<5.167,\quad|f'(\beta)|>2.289,\quad
 f'(\chi)>4.128,\\
1.752<\beta+\chi<1.754,\qquad\beta-\chi>1.355 .
 \end{gathered}
\]
They follow by polynomial signs and products
of the disjoint root gaps.
Writing \(W=A+BU+CU^2\) in the FULL ring,
Lagrange interpolation gives
\[
 C=\sum {W_i\over f'(U_i)},\qquad
 B+C(\beta+\chi)={W_\beta-W_\chi\over\beta-\chi}.
\tag{D17L}\label{eq:arithD17interpolation}
\]
The complete strip implies \(3<C<8\) and
\(|B+C(\beta+\chi)|<.174\).
At \(C=4,5,6,7\), only \(C=4,B=-7\) survives.
Then \(7\chi-4\chi^2\in(1.229,1.235)\),
so the small strip requires \(1.229<A<1.471\),
no integer.
For head eighteen the corresponding bounds are
\[
 \begin{gathered}
3.879<\alpha<3.880,\quad1.652<\beta<1.653,\quad
 .467<\chi<.469,\\
7.590<f'(\alpha)<7.605,\quad|f'(\beta)|>2.63,\quad
 f'(\chi)>4,\quad
2.120<\beta+\chi<2.121,\quad\beta-\chi>1.183 .
 \end{gathered}
\]
Interpolation gives \(2<C<6\),
\(|B+C(\beta+\chi)|<.199\).
For \(C=3,4,5\), the products lie in
\((6.360,6.363),(8.480,8.484),(10.600,10.605)\);
none is within \(.199\) of an integer.
No coefficient \(B\) exists.
Both cubic finite heads are excluded.

All FIVE finite heads have now been eliminated
at their ACTUAL fields.
Only NOW conclude that the head is DOTTED,
with \(4<U<11\), \(0<U_\sigma<2\).
Fork isolation and the dotted fork assumptions
proved conditionally in this subsection are NOT
exported to that remaining head.
We next prove them anew in the dotted case.

\subsection{Whole \(D_{17}\): degree reduction under a dotted head}
\label{app:arithD17degrees}

The complete fork table remains valid.
For an end helper, the integral \(v\) has norm
magnitude \(<6.607(.617)^{d-1}\).
Degree at least five gives norm \(<1\).
At degree four it is a unit;
parity~\eqref{eq:arithD17helperparity} then makes
\(Y=U/2\) integral.
With physical \(2<Y<11/2\) and three other
conjugates in \((0,1)\),
\[
                 0<|\mathrm N(Y(1-Y))|<99/256<1,
\]
impossible.
For an isolated fork the actual LONG+TWO
cycles above remain valid under a dotted head
BEFORE assigning its own type.
They again give integral \(W=V/2\), with
\(0<W<676/17\) and \(0<W_\sigma<4/17\)
at every other embedding.
The lower bound \(2<W\) is not used before
the finite-own exclusion below.
Degree at least four contradicts
\((676/17)(4/17)^3<1\).
Thus actual degree is two or three.
Rational degree was excluded by the named
WHOLE-form Meyer input.

We first remove cubic nonisolated forks.
The end-helper norm is \(\pm1,\pm2\).
Opposite signs of the two small roots give
magnitude \(<6.607(.617)(.382)<2\).
Two positive small roots in the far case,
or two negative in the near case, give magnitude
\(<1\).
Norm magnitude two can therefore occur only
as norm \(+2\) in the following two patterns.
For far roots \(t,-a,-b\), \(tab=2\),
\(0<a,b<.617\), the trace is five,
\(6<t<6.234\), and the integer second
coefficient is \(-6\) or \(-7\).
The polynomials
\(X^3-5X^2-6X-2\),
\(X^3-5X^2-7X-2\)
have discriminants \(-424,229\).
The first is not totally real.
At \(-.617\), the second has value
\(.180669887>0\), inconsistent with both
negative roots above that lower endpoint.
For near roots \(v,a,b\) with \(vab=2\),
the trace is seven, \(5.766<v<6\),
and the second coefficient is seven.
Its polynomial \(X^3-7X^2+7X-2\)
has discriminant \(-59\), impossible.
Thus a remaining cubic end-helper is a unit,
so \(U/2\) integral.

Here is the full EVEN-cubic-head conclusion,
also reused below.
For integral \(Y=U/2\), physical \(2<Y<11/2\)
and other roots in \((0,1)\) give
\[
                       0<|\mathrm N(Y(1-Y))|<99/64<2.
\]
Thus \(\mathrm N(Y)=1,\mathrm N(1-Y)=-1\).
Its polynomial is
\(X^3-TX^2+(T-1)X-1\), \(T=3,\ldots,7\).
The COMPLETE discriminants are
\(-23,-31,-23,49,257\).
The first three are not totally real;
at \(T=7\), evaluation at six is \(-1\),
forcing physical root \(>6\).
Only \(T=6\) remains.
Its reciprocal generates the cyclic field
\(K_7\), with full basis \(1,u,u^2\) for
\(u^3-5u^2+6u-1=0\).
The square discriminant gives cyclicity, so
choose its PHYSICAL largest root as generator.
The full-order proof and root/derivative bounds
are those just proved above and are algebraic
facts of this field, not a finite-head hypothesis.

For a cubic end-helper \(v=A+Bu+Cu^2\),
the full far/near strips and
\eqref{eq:arithD17interpolation} imply
\(-1<C<2\), hence \(C=0,1\), and
\(|B+C(\beta+\chi)|<.738\).
At \(C=0\), \(B=0\), and the small strip
forces \(A=0\), contrary to physical \(v\ge1\).
At \(C=1\), only \(B=-2\) is possible.
Since \(2\chi-\chi^2\in(.356,.359)\),
the small strip gives \(A=0\) in either role.
But \(v_\beta=\beta(\beta-2)<-.689\),
outside BOTH strips.
This exhausts every coefficient triple.
Cubic end helpers are impossible, so BOTH
actual cubic forks are isolated dotted.
Finite own in degree three would violate
\(\mathrm N(V)<4(8/17)^2<1\).
Each integral \(W\) satisfies
\eqref{eq:arithD17forkstrip},
with norm one or two; no equality of the two
fork parameters has been assumed.

\subsection{Whole \(D_{17}\): the entire isolated cubic field domain}
\label{app:arithD17cubic}

Choose ONE actual isolated cubic fork; the other
may have a different parameter.
Write its monic polynomial
\(P(X)=X^3-AX^2+BX-n\),
\(n=\mathrm N(W)=1,2\).
Its two other positive roots lie in
\((0,q)\), \(q=4/17\), physical
\(W<676/17\), hence integer trace \(A\le40\).
Since \(P(q)<0\),
\[
                   64-272A+1156B-4913n<0 .
\]
For \(n=2\), physical \(W>289/8\), and the
small-root arithmetic-geometric inequality gives
\(B\ge2\sqrt{2W}+2/W>17\).
The polynomial inequality gives
\(B<20642/1156<18\), impossible.
Thus \(W\) is a UNIT.

For the actual head/fork products in the
normalization \(r_H^2=2U,r_f^2=4W\),
the LONG and TWO cycles give integral \(g,L\),
with
\[
                   L-g=(8U-17)W-2U+4=:d_0,\qquad
                   g^2=2UW(g+d_0).
\]
At every dyadic prime, \(d_0\equiv-W\) is a unit.
The equation forces \(g\) into the prime and
then
\(2\operatorname{ord}(g)
 =\operatorname{ord}(2)+\operatorname{ord}(U)\).
If two is unramified this makes \(U/2\) integral.
The even-cubic-head conclusion above gives
actual field \(K_7\).
Apply the SAME full-order interpolation as for
finite head seven, now merely as algebra of the
field and the CURRENT unit-\(W\) strip.
Its complete \(C=4,\ldots,7\) tests leave only
\(C=4,B=-7\), then no integer \(A\).
Thus the unramified case is impossible, without
assuming the head itself is finite.

If two ramifies, the unit polynomial has a
repeated factor modulo two.
Its derivative \(X^2+B\) shows this requires
BOTH \(A,B\) odd.
The physical root and two small-root bounds give
\[
 \begin{gathered}
 A=19,21,\ldots,39,\qquad
 B^2>4(A-8/17),\qquad
 B<{272A+4849\over1156}.
 \end{gathered}
\]
The COMPLETE lower/upper odd-integer tests are
\[
\begin{array}{c|rrrrrrrrrrr}
 A&19&21&23&25&27&29&31&33&35&37&39\\ \hline
 B_{\min}&9&11&11&11&11&11&13&13&13&13&13\\
 B_{\max}&7&9&9&9&9&11&11&11&11&11&13 .
\end{array}
\]
Only
\[
                   P_{29}=X^3-29X^2+11X-1,\qquad
                   P_{39}=X^3-39X^2+13X-1
\]
remain.
Their discriminants are \(4(3)(383)\),
\(4(5021)\), with squarefree odd parts.
For an elementary full-order certificate,
trial primes below \(\sqrt{383}\) are
\(2,3,5,7,11,13,17,19\), none dividing383.
For5021, all primes below its square root are
\[
2,3,5,7,11,13,17,19,23,29,31,37,41,43,47,53,59,61,67;
\]
the respective remainders are
\[
1,2,1,2,5,3,6,5,7,4,30,26,19,33,39,39,6,19,63 .
\]
The shifted polynomials are two-Eisenstein:
\[
 P_{29}(X+1)=X^3-26X^2-44X-18,\qquad
 P_{39}(X+1)=X^3-36X^2-62X-26 .
\]
Only two can divide the power-order index.
At its uniformizer \(W-1\), valuations of
\(a_0+a_1\pi+a_2\pi^2\) are distinct modulo
three; a numerator not all divisible by two
has unique least valuation at most two.
Dividing by two cannot be integral.
Thus \(\mathcal O_K=\mathbb Z[W]\) in BOTH cases.

The exact small-root brackets for \(P_{29}\)
are \(.14<u<.15\), \(.23<v<.235\);
its physical root satisfies \(28.615<w<28.63\).
They follow from values
\(-.025656,.000875,.008067,-.003547125\)
at those four endpoints.
The physical derivative is \(>800\), both
small derivative magnitudes \(>2.27\).
For \(P_{39}\), values
\(-.040569,.000128,.019361,-.016952\)
give \(.11<u<.12\), \(.21<v<.22\),
\(38.66<w<38.68\), physical derivative \(>1480\),
small derivative magnitudes \(>3.45\).

Write the ACTUAL head \(U=a+bW+cW^2\) in
these proved FULL rings.
Its current strips are physical \((4,11)\),
other values \((0,2)\).
Lagrange interpolation bounds \(c\) respectively by
\[
             -2/2.27<c<11/800+2/2.27,\qquad
             -2/3.45<c<11/1480+2/3.45 .
\]
Both intervals are strictly inside \((-1,1)\),
so \(c=0\).
Then \(b=(U(w)-U(u))/(w-u)\) satisfies
\(0<b<11/(w-u)<1\), impossible for an integer.
Every ramified cubic case has been excluded.
The ACTUAL field is QUADRATIC.

\subsection{Whole \(D_{17}\): complete quadratic end-helper exclusion}
\label{app:arithD17quadratichelpers}

The dotted-head fork table already excludes
all four non-end white roles.
For the remaining far/near end-three roles
write \(v\)'s polynomial \(X^2-TX+n\).
The COMPLETE far ranges are
\(T=1,\ldots,6\), \(n=-3,-2,-1,1,2\);
near ranges \(T=2,\ldots,7\),
\(n=-2,-1,1,2,3,4\).
The quadratic formula and EXACT time signs leave
the following fourteen and fifteen rows:
\[
\begin{array}{c|c}
 (T,n,d)\text{ far}&(T,n,d)\text{ near}\\ \hline
(2,-1,2)&(3,-1,13)\\
(3,-1,13)&(3,1,5)\\
(3,-2,17)&(4,-1,5)\\
(4,-1,5)&(4,1,3)\\
(4,-2,6)&(4,2,2)\\
(4,1,3)&(5,-1,29)\\
(5,-1,29)&(5,-2,33)\\
(5,-2,33)&(5,1,21)\\
(5,-3,37)&(5,2,17)\\
(5,1,21)&(6,-1,10)\\
(6,-1,10)&(6,-2,11)\\
(6,-2,11)&(6,1,2)\\
(6,1,2)&(6,2,7)\\
(6,2,7)&(6,3,6)\\
 &(7,3,37).
\end{array}
\]
Here actual \(K=\mathbb Q(\sqrt d)\).
At each bounded pair \(T,n\), substitute the larger
root and its conjugate into the FULL strips.
Every missing pair is rational, nonreal or
violates an endpoint.
The safe far decimal overstrip admits one extra
\((3,1)\), \(v=\varphi^2\).
Its exact conjugate time numerator is
\((123-55\sqrt5)/2>0\), excluded by
\(123^2-5(55)^2=4\).
Thus the exact, rather than decimal, domain
is fully enforced.

For odd \(n\), \(v\) is a unit at EVERY dyadic
prime, so parity~\eqref{eq:arithD17helperparity}
makes \(U/2\) integral even with ramification.
The FULL quadratic head strip for \(Y=U/2\)
is \(2<Y<11/2\), \(0<Y'<1\).
Its integer trace is three through six and
norm \(1\le k\le T-2\); at trace six the
physical bound removes \(k=1,2\).
The COMPLETE eight actual heads and
\(h=-\mathrm N(U(U-2))\) are
\[
\begin{array}{c|c|c}
 d&U&h\\ \hline
5&2\varphi^2,\ 4\varphi^2&16,64\\
2&4+2\sqrt2&32\\
3&4+2\sqrt3&32\\
13&5+\sqrt{13}&48\\
17&5+\sqrt{17}&64\\
21&5+\sqrt{21}&48\\
6&6+2\sqrt6&96 .
\end{array}
\]
Odd-norm helper rows in fields29/10/37 are
immediately incompatible.
For all others, put
\(p=-16\mathrm N(H_c^2)>0\).
The complete time-norm tests are
\[
\begin{array}{c|c|r|l}
 \text{role}&(T,n,d)&p&h\\ \hline
 \text{far}&(2,-1,2)&287&32\\
 &(3,-1,13)&711&48\\
 &(4,-1,5)&1271&16,64\\
 &(4,1,3)&423&32\\
 &(5,1,21)&983&48\\
 &(6,1,2)&1679&32\\
 \text{near}&(3,-1,13)&207&48\\
 &(3,1,5)&311&16,64\\
 &(4,-1,5)&599&16,64\\
 &(4,1,3)&839&32\\
 &(5,1,21)&1503&48\\
 &(6,1,2)&2303&32\\
 &(6,3,6)&503&96 .
\end{array}
\]
All \(p/(16h)\) are nonsquares.
At \(h=32,96\) the odd numerator gives odd
two-adic denominator valuation.
The four \(h=48\) reduced cases are
\(237/256,983/768,69/256,501/256\);
the middle denominator has odd three valuation,
the other numerators are nonsquares.
For \(h=16,64\), the numerators
\(1271,311,599\) are nonsquares.
These are NECESSARY actual common-time-line
square tests, not sufficiency of a norm-square.

For even-norm rows DO NOT assume \(U/2\)
integral.
Their complete actual fields, full head strips
and \(h\)'s are
\[
\begin{array}{c|l|l}
 d&U&h\\ \hline
2&3+\sqrt2,\ 3+2\sqrt2,\ 4+2\sqrt2,\
     5+3\sqrt2,\ 6+3\sqrt2&7,7,32,63,36\\
6&3+\sqrt6,\ 4+\sqrt6,\ 5+2\sqrt6,\ 6+2\sqrt6&
     15,20,15,96\\
7&3+\sqrt7,\ 4+\sqrt7&12,27\\
11&4+\sqrt{11},\ 5+\sqrt{11}&35,28\\
17&2+\eta,\ 3+\eta,\ 4+2\eta,\ 5+2\eta&
      8,16,64,19\\
33&3+\xi,\ 4+\xi&24,24 .
\end{array}
\]
Here \(\eta=(1+\sqrt{17})/2\),
\(\xi=(1+\sqrt{33})/2\).
The FULL integral bases and \(0<U'<2,4<U<11\)
give coefficient bounds respectively
\(3,2,2,1,2,1\).
At fixed positive coefficient, the width-two
conjugate interval gives two consecutive constants;
the physical strip removes the excess.
For example at \(\sqrt7\), coefficient two
already has least physical \(6+2\sqrt7>11\).
Thus no full-order head was omitted.

The COMPLETE remaining even-norm time rows are
\[
\begin{array}{c|c|r|l}
 \text{role}&(T,n,d)&p&\text{all }h\\ \hline
 \text{far}&(3,-2,17)&200&8,16,64,19\\
 &(4,-2,6)&828&15,20,15,96\\
 &(5,-2,33)&1592&24\\
 &(6,-2,11)&2492&35,28\\
 &(6,2,7)&252&12,27\\
 \text{near}&(4,2,2)&92&7,7,32,63,36\\
 &(5,-2,33)&72&24\\
 &(5,2,17)&824&8,16,64,19\\
 &(6,-2,11)&668&35,28\\
 &(6,2,7)&1692&12,27 .
\end{array}
\]
Every entry follows by expanding
\[
16\mathrm N(H_c^2)
 =(17T\pm4)^2n+(17T\pm4)(-17n-4)T+(-17n-4)^2 .
\]
Every ratio \(p/(16h)\) is nonsquare EXCEPT
far \((T,n,d)=(3,-2,17)\), \(h=8\).
The far nonexceptional rows retain odd valuations
at \(23,199,89,7\).
In the exceptional field the other ratios are
\(25/32,25/128,25/38\), nonsquares.
The near rows retain odd valuations at
\(23,3,103,167,47\).
Thus a genuine square-norm exception is kept.

For that exception
\[
 t=(3+\sqrt{17})/2,\quad U=(5+\sqrt{17})/2,
\]
the ACTUAL time squares are
\[
 H_c^2=(225+55\sqrt{17})/8,\qquad
 H_0^2=(11+3\sqrt{17})/2 .
\]
Their ratio
\(R=(165+35\sqrt{17})/64\) has norm \(25/16\).
If \(R=z^2\) in \(K\), then
\(\mathrm N(z)=\pm5/4\), and
\[
                    \operatorname{Tr}(z)^2=245/32
                                      \quad\hbox{or}\quad85/32 .
\]
Neither is a rational square, by its odd
two-adic denominator valuation.
But the trace is rational.
The last actual common-time square fails.
All quadratic end-helper rows are now excluded,
in both signed orientations.

Both ACTUAL forks are therefore isolated under
the dotted quadratic head.
Their original LONG+TWO cycles, applied BEFORE
assuming own dotted, give integral \(W=V/2\)
with \(0<W<676/17\) and \(0<W'<4/17\).
The lower bound \(2<W\) is deferred until
finite owns have been excluded.
Finite own would have \(0<W<2\), other
\(0<W'<4/17\), hence nonzero integer norm
\(0<\mathrm N(W)<8/17\), impossible.
Parallel has its actual null pair.
Both forks are isolated DOTTED, without
equating their two \(W\) parameters.

\subsection{Whole \(D_{17}\)'s distinct actual inclusion and final budget}
\label{app:arithD17domain}

Choose ONE of the two actual isolated dotted forks.
Its raw square is \(4W\), sole old product
\(-2W\), spatial \(W(P_{16}\pm e_{17})\),
time square \(W(17W-4)\).
The actual head has square \(2U\), spatial
\(Ue_1\), time square \(U(U-2)\).
The original \(K\)-time line gives
\[
 R_0=U(U-2)W(17W-4)\in K^2,\qquad
 g=UW-\sqrt{R_0},
\]
with PHYSICAL positive square root.
The original LONG path and TWO cycles give
\[
                   g,\ L={g^2\over2UW}\in\mathcal O_K .
\]
All physical/conjugate strips are now
\[
 K=\mathbb Q(U)=\mathbb Q(W),\quad
 4<U<11,\quad0<U'<2,\quad
 2<W<676/17,\quad0<W'<4/17 .
\]
The actual independent conjugate root lines give
\(0<L'<4\) when \(g\ne0\).
Here the current physical strips force
\(d_0=(8U-17)W-2U+4>0\), hence \(g<0\),
so no orthogonal exception is silently lost.
The true physical mutual \(L\) is finite
native or dotted, while parallel four has
its forbidden actual null pair.
Thus EVERY predicate of
Proposition~\ref{prop:pair-output}(iii) holds
for this ORIGINAL D17 geometry.

This inclusion is separate from
Section~\ref{app:arithB17domain}:
D17 raw fork norm is \(4W\);
B17 raw short norm was \(V\).
No equation of different root scales or
arithmeticity of a face supplied either inclusion.
The existing full-order algorithm covers every
\(\Delta<121\), every strict-strip model and
ordered pair with both field-square norm signs.
Its 36-field/120-pair output, explicitly ASSUMED
as the recorded computational proposition, has
the same TWO necessary survivors.
The field/counter/survivor combine does not
independently recompute the118 rejected tests.

For either actual survivor, use its actual
head point and ONE chosen signed fork.
The head gain at cost \(.072\) from \(P_5\)
is \(>.043\), by the full rational seals already
printed in Section~\ref{app:arithB17domain}.
The actual fork has
\(C_f^2=17W/4>442/5\), giving
\[
                 x\cdot(P_{16}\pm e_{17})>3.7 .
\]
Pay \(.13\) from the DISTINCT \(P_{14}\) price:
\[
             y=w_{14}+.36(Z-P_5/5)
                          +.13(P_{16}\pm e_{17}-P_{14}).
\]
The signed residual has norm three and
vanishes on all reference coordinates;
it is orthogonal to the head residual.
Both individual payments are legal BEFORE
using their original floors.
ONE square on the SAME original \(x\) gives
\[
              \|x\|^2>.9084+.043+.26(3.7-3.311)-.0169(3)
                                     =1.00184>1 .
\]
Every actual field and actual fork type has
now been covered.
This closes whole \(D_{17}\) at the stated
external and explicitly computed premises.
The argument needed one chosen fork for the
last contradiction, while preserving both
signed walls and all original constraints.

\section{Smaller chains: the same pressure method with fresh local domains}
\label{app:arithsmallerchains}

We now finish the smaller whole-\(A\) consumers.
The shared mechanism is unchanged: classify an
ACTUAL central opposite, constrain a pure helper
arithmetically, and absorb EVERY retained isolated
foot in ONE reference.
The index ranges, finite side exceptions, helper
packing and coefficient capacities are proved at
their own \(N\); no numerical theorem is reranked.

\subsection{A spectator-safe \(A_{14}\) bound greater than \(4/5\)}
\label{app:arithA14fourfifths}

Use the zero-sum space in \(\mathbb R^{15}\),
\[
 n_i=(e_i-e_{i+1})/\sqrt2,\quad
 \Omega_i=\sqrt2(P_i-iP_{15}/15),\quad
 D_i=2i(15-i)/15,\quad h_i=x_A\cdot\Omega_i .
\]
All spectators are kept; the SAME original
farthest vector has \(\|x\|<1\).
For \(i=5,\ldots,10\) take loss
\(\ell=.597\).
Here \(D_5=D_{10}=20/3\),
\(D_6=D_9=36/5\), \(D_7=D_8=112/15\).
The exact ordinary sufficient ratio
\[
 R_\ell(q)=1/\ell-(q-\ell)/(q(2q-\ell))
\]
increases on this interval.
Since \(q<11/4\),
\[
 R_\ell(q)<1000/597-8612/53933<761/500 ;
\]
the last positive cleared numerator is \(106860761\).
The older coarse \(q<3\) bound is not used.

Deleting a central node retains two paths
of rank at least four.
The FULL side-order-\(m\ge5\) finite-white
list is no contact, permitted endpoint/D/E
simple three, endpoint four, or two end-threes
joining the paths.
The connected replacement rank is fourteen,
so no high-rank exceptional type is missing.
The complete own-zero repairs retain internal
affine arms, stars and two-contact cycles;
far/next-three use the ENTIRE foreign cost
and reciprocal original-wall law;
endpoint-four has no positive-time foreign
completion.
All isolated own-three/four/six/parallel nulls
apply for EVERY foreign cost BEFORE splitting it,
as proved in the common chain section.

The ordinary nonisolated ratios exceed \(761/500\),
as do the complete foreign far/near bounds
\(227/(25\sqrt{35})\) and \(\sqrt{121/35}\).
The only extra near-four orientations are
\[
                   v_A=\Omega_i/\sqrt2+\Omega_{i-1}/2,
                         \qquad i=5,6,7
\]
and reflections, with zero foreign cost.
Their conjugate time numerators at \(N=15\)
are
\(114-80\sqrt2,128-90\sqrt2,136-96\sqrt2>0\).
The actual simple-three link puts \(\sqrt2\)
in \(K\); its sign-reversing embedding makes
the original Schur complement negative.
Near-four of side order at least eight is
already strong, so both helper orientations
are covered.

For an isolated-within-\(A\) cap write
\(c=Ap-\gamma\Omega_i-v_S\),
\(Q=\|v_S\|^2<1\), \(u=(1-Q)/\gamma^2\).
For \(Q>0,u>1\), the actual positive clearing
has loss \(<.202+.382=.584<\ell\).
At \(u=1\) it is a feasible null.
For pure native five the actual \(H_4\)
loss is \(<25/76+5/24<\ell\).
For EVERY pure finite native at least seven
the following original negative core gives
a fresh stronger bound, including native eight:
\[
 \xi=c+\gamma^{-1}n_i+
                    {2(1-\gamma^2)\over\gamma}n_j,
 \qquad j\text{ either old neighbour}.
\]
Put \(d=1-\gamma^2<.19\).
The cap and \(n_i\) products vanish;
\(\xi\cdot n_j=(2d-1/2)/\gamma<0\);
its square is \(-d(1-4d)/\gamma^2<0\).
Positive cap time and actual acuteness check
EVERY omitted original wall.
Its normalized comparison gives loss
\[
 <{10000\over39771}+{25\over81}
       ={22275\over39771}<.561<\ell .
\]
Here \(\sqrt{d(1-4d)}\le1/4\),
\(\gamma>.9\), \(q>2.58\), and
\(\sqrt{D_i-100/81}>2.33\).
Every remaining isolated foot is retained in
the FULL interval \(0<u<1\), pure or foreign;
none is assigned an unsupported strong floor.

Every remaining weak nonisolated cap is a PURE
endpoint helper with
\[
 202/125<t=2\gamma<3139/1458<54/25.
\]
The fresh cheapest negative/positive conjugate
Schur costs are
\[
 Q_-(t)=2it^2+4t+2-2(it+1)^2/15,\quad
 Q_+(t)=2it^2+2-2(it-1)^2/15 .
\]
At \(i=5\), \(Q_-(-4/5)=4\) EXACTLY,
\(Q_+(12/25)=4+16/375\).
They increase at the relevant signs over the
half-window, so strict positivity gives every
other conjugate in \((-4/5,12/25)\).
The COMPLETE all-degree norm/trace argument
of Section~\ref{app:arithgoldenfield} uses exactly
these physical/conjugate bounds and proves
\(t=\varphi\), actual \(K=\mathbb Q(\varphi)\).
No restriction is imposed in a branch without
a helper.

The COMPLETE six-type \(30A^2=M\) table is
\[
\begin{array}{c|c|r}
 \text{type}&M&\mathrm N(M)\\ \hline
 \text{near }5&34+70\varphi&-1364\\
 \text{near }6&38+72\varphi&-1004\\
 \text{near }7&40+72\varphi&-704\\
 \text{far }5&34+60\varphi&-404\\
 \text{far }6&38+66\varphi&-404\\
 \text{far }7&40+70\varphi&-500 .
\end{array}
\]
Different norm-squareclasses cannot coexist on
the SAME original connected-block \(K\)-time line.
For the repeated far-five/six class, the ratio is
\[
                         {97+9\varphi\over101}
                              ={203+9\sqrt5\over202}.
\]
Its norm is one.
A full \(K\)-square would have rational
coefficient square \(405/404\) or \(1/404\),
both impossible at prime101.
For reflected copies, the actual cyclic parameter
has constant numerators \(8,1,-10\) modulo fifteen,
so is nonintegral in \(\mathbb Z[\varphi]\).
Explicitly the near/far formulas are
\[
 { (15-2i)(i+(i\pm2)\varphi)-17\over15}.
\]
Every different and repeated type is covered;
at a fixed index the ACTUAL opposite is unique.
Thus at most ONE central index has a weak helper.
The number and adjacency of isolated feet are
UNRESTRICTED.

For each retained \(0<u<1\) foot the positive
original combination
\(c+v_S+\gamma u n_i\) gives
\[
 x_A\cdot(\Omega_i-un_i)\ge
                    \sqrt{D_i-u}-\sqrt{u(1-u)} .
\]
The raw loss satisfies
\(L_i\le a_i u+\sqrt{u(1-u)}\), where
\[
 a_{5,10}<.202,\quad a_{6,9}<.194,\quad
 a_{7,8}<.190 .
\]
Each follows from the ACTUAL
\(a_i=(\sqrt{D_i}+\sqrt{D_i-1})^{-1}\).
The full adjusted-loss inequality
\[
                 L_i-{\beta_i\over2}(2u-u^2)<.597
\]
holds for capacities
\[
 \beta_{5,10}\ge.06,\qquad
 \beta_{6,9}\ge.026,\qquad
 \beta_{7,8}\ge.0246 .
\]
For outer sites, \(u\le.4\) has loss \(<.5708\);
above \(.4\) the circle bound gives
\(.6111-.03(.64)=.5919\).
For sites six/nine, \(u\le.5\) gives \(<.597\);
above it use
\((.194+1.01865)/2-3(.026)/8=.596575\).
For seven/eight, the corresponding bounds are
\(.595\) and
\((.190+1.018)/2-3(.0246)/8=.594775\).
Thus EVERY pure/foreign foot parameter is covered.

Use symmetric rational heights
\[
 H_5=H_{10}=A=1.98498,\quad
 H_6=H_9=B=2.08628,\quad
 H_7=H_8=C=2.13551,
\]
linear on the outer tails.
Their exact floor squares are below the true
\(\ell\)-strong heights.
The central dual coefficients are
\[
                  .147848,.026035,.024615,
                  .024615,.026035,.147848 .
\]
A single internal helper is omitted by its
neighbour chord, making its coefficient zero
and increasing neighbouring ones.
The largest energy loss is
\((2B-A-C)^2/4=.000677821225\).

For \(w=\sum\beta_i\Omega_i\), and EVERY retained
foot with positive coefficient, put
\(Z=w-\sum\beta_i u_i n_i\).
Primitive actual pressures give
\(x_A\cdot Z\ge E_H-\sum\beta_i(L_i-\ell)\).
ALL adjacent-foot crosses have the correct sign
\(n_i\cdot n_j\le0\), so
\[
 Z^2\le E_H-\sum\beta_i^2(2u_i-u_i^2).
\]
ONE square and the adjusted losses give \(e_A>E_H\).
Strong constraints supply strictness when no
feet remain.
The full energy is
\(A^2/5+(B-A)^2+(C-B)^2=.80071440298\);
after the largest internal omission it is
\[
                              >.800036581755>4/5 .
\]

It remains to cover an OUTER helper at five/ten.
Reverse labels if necessary, without assuming
a chamber symmetry.
The actual \(H_4\) positive correction is
\[
 b_c=\gamma n_5+(2/3)n_e+(1/3)n_{\rm adj},\quad
 \delta=(c+b_c)^2=(5-3\varphi)/12<.111^2 .
\]
The positive FINITE projected point is used;
the spacelike \(c+b_c\) alone is not a point.
For helper end \(e=1\) or fourteen, the signed
endpoint \(y=x_1\) or \(-x_{15}\) obeys
\[
                    y>\sqrt2(A_c-.111-\gamma h_5).
\]
The smaller actual time has
\(A_c^2=(34+60\varphi)/30>2.09^2\);
weakness gives \(h_5<1.985\), \(\gamma<.81\).
Thus BOTH orientations have
\(y>1.4(2.09-.111-.81(1.985))=.51961>.49\).

Omit knot five by the chord zero-to-six.
Its reference energy is
\[
 E_{\rm end}={1\over2}
       [B^2/6+A^2/5+(B-A)^2+2(C-B)^2]>.764 .
\]
At the actual helper end there is a flat signed
block of length six or five, mean
\(b_0=B/(6\sqrt2)\in(.24,.25)\) or
\(A/(5\sqrt2)\in(.25,.3)\).
Graft the ACTUAL endpoint to \(.49\),
keeping the block sum fixed.
The other entries
\((Lb_0-.49)/(L-1)\) lie in \((.19,.2525)\),
above the next old coordinate \(<.073\).
All duals remain positive.
The changed central boundary coefficient is
\(>.117/1.415>.06\), meeting its own capacity;
all others retain theirs.
The missing-five coefficient remains zero.
The gain in this SAME reference is
\[
                  {L\over L-1}(.49-b_0)^2
                                  >(6/5)(.19)^2=.04332 .
\]
Absorb ALL other isolated feet with their
individual parameters as above.
Then \(e_A>.764+.04332=.80732>4/5\).
No second helper exists by packing.
The proof covers every actual field and
arbitrary finite spectators:
\[
                                  e_{A_{14}}>4/5 .
\tag{A14U}\label{eq:arithA14universal}
\]
In particular its ordinary-two-\(A_1\) and
actual-\(I\)-plus-\(A_1\) rank-sixteen consumers
close by adding respectively \(>.2\) or the
rank-free JOINT cost \(>.20204\), with \(I\)
counted once.
In rank seventeen, every rank-three remainder
has cost \(>.21\) by
\eqref{eq:arithD13spectator}, so the ENTIRE
older \(A_{14}\) row closes at once.

\subsection{A fresh spectator-safe \(A_{13}\) bound}
\label{app:arithA13universal}

Now \(N=14\), central sites five through nine,
with
\[
 D_5=D_9=45/7,\quad D_6=D_8=48/7,\quad D_7=7,\qquad
 \ell=1199/2000=.5995 .
\]
Deleting a central site again retains two
paths of rank at least four.
The complete side-order-\(m\ge5\) white list
is isolation, permitted simple-three/D/E,
endpoint-four, or one endpoint-three on EACH
retained path.
Two contacts within one path form a cycle;
the connected rank-thirteen completion is
classical, with no missing exceptional type.
The fresh sufficient ratio satisfies
\[
 R_\ell(q)<2000/1199-4801/32403<761/500,\qquad 5/2<q<3 .
\]
This follows from the same positive derivative
and a direct positive rational cross product.

Own zero is dispatched BEFORE positive-own
ratios.
Isolation has \(A^2=Q-1<0\).
Internal-three has the actual affine arms
\(1,2,5\) or \(1,3,3\); internal-four has its
weighted star, two threes their cycle.
The only residual endpoint/next-three/four
times are
\[
 \begin{array}{c|c}
 A^2&\text{white bound}\\ \hline
 Q-15/28&Q<3/5\\
 Q-1/7&Q<2/5\\
 Q-1/14&Q<1/5 .
 \end{array}
\]
The first positive time is \(<9/140<1/4\),
below any foreign contribution.
The second has at most one foreign contribution,
larger than \(A^2\).
Both violate the actual orthogonal reciprocal
opposite law; the last admits no positive
foreign cost.
At \(Q=0\) all three times are negative.
Every affine/null vector has positive cap time
and correct signs against EVERY original wall.

The complete ordinary ratios for positive own
exceed \(761/500\), including
\[
 { .808+.4\over\sqrt{.6}},\quad
 { .808+.2\over\sqrt{.4}},\quad
 { .808\over\sqrt{.25}},\quad q/2+1/q .
\]
The complete foreign far/near bounds are still
\(227/(25\sqrt{35})\), \(\sqrt{121/35}\).
Thus each weak nonisolated cap is PURE:
an endpoint helper or own-four/simple-three
near helper of side order five/six/seven.
Near-four of side order at least eight is
already strong.
We retain the smaller near-four roles rather
than use an invalid N15 conjugate exclusion.

For a pure endpoint helper, the exact square
ratio identity gives \(t=2\gamma<54/25<9/4\).
If its own is not five, \(t>9/5\).
Actual simple-three normalization makes \(t\)
integral.
At every lifted nonphysical embedding, the
cheapest signed old-\(A\) Schur costs are
\[
 2it^2+4t+2-2(it+1)^2/14,\qquad
 2it^2+2-2(it-1)^2/14 .
\]
At the relevant negative/positive endpoints
they increase with \(i\) and \(N\).
Already at \(N=13,i=5\), costs at \(-.83,.5\)
exceed four by \(69/1625,2/13\).
Hence every other conjugate belongs to
\((-.83,.5)\).
Degree at least six contradicts
\((9/4)(.83)^5<1\).
In degrees three/four/five the norm magnitude
is \(<2\); a positive small conjugate would
make it \(<1\), so all are negative, of product
magnitude \(1/t\).
AM--GM bounds the integer trace respectively in
\[
                         (.14,1),\quad(-.69,0),\quad(-1.52,-1),
\]
none containing an integer.
The exact underlying comparisons are
\(t(t-1)^2<(9/4)(5/4)^2<4\),
\(t^4<(9/4)^4<27\),
\(t(t+1)^4<(9/4)(13/4)^4<256\).
In degree two trace two gives a nonzero norm
strictly between \(-1,1\);
trace one leaves only norm \(-2\), whose
polynomial \(X^2-X-2\) is reducible.
Degree one leaves parallel two, an actual null.
Therefore every weak endpoint helper is golden,
with ACTUAL \(K=\mathbb Q(\varphi)\) by its
original Lorentzian principal.

The COMPLETE \(28A^2=M\) golden table is
\[
\begin{array}{c|c|r}
 \text{type}&M&\mathrm N(M)\\ \hline
 \text{near }5&30+63\varphi&-1179\\
 \text{near }6&33+64\varphi&-895\\
 \text{middle }7&34+63\varphi&-671\\
 \text{far }5&30+55\varphi&-475\\
 \text{far }6&33+60\varphi&-531 .
\end{array}
\]
All five norm-squareclasses differ, excluding
unlike types on the common actual \(K\)-time line.
Reflected same-type cyclic parameters are
\[
                  { (14-2i)(i+(i\pm2)\varphi)-16\over14}.
\]
At \(i=5,6\) their constant coefficients are
\(2/7,-2/7\), nonintegral.
At the middle there is only one actual index.
Thus at most ONE weak golden helper exists.

If such a helper occurs, EVERY other central
site is \(.65\)-strong.
Indeed \(R_{.65}(3)<1.4\), below all ordinary
and near-four roles; the isolated raw loss
bound below is \(<.6135\).
Any \(.65\)-weak endpoint helper is also
\(.5995\)-weak, so packing excludes a second one.
For the actual finite \(H_4\) point use
\[
 U=v_A-\gamma n_i-(2/3)n_e-(1/3)n_{\rm adj},
 \qquad\delta=(5-3\varphi)/12<.111^2,
 \qquad x_A\cdot U>A-\sqrt\delta .
\]
Positive finite projection, not the spacelike
corrected normal alone, checks all original
walls.

For helper index five, set \(a=1.885,b=1.968\).
Use heights \(H_6=H_7=H_8=b,H_9=a\), linear
on its two outer tails.
The energy is
\(E_5=(b^2/6+(b-a)^2+a^2/5)/2=.681519\).
Project \(U\) onto its flat blocks.
For far/near helper the projections are
\[
 (5\gamma/6)\Omega_6+\Omega_9/10,\qquad
 (5\gamma/6+1/12)\Omega_6 .
\]
Their reference products are \(<1.5169\).
The smaller time is \(A>2.061\); the deficit
is \(>.4331\).
The FULL orthogonal residual norms are
\((2/3)\gamma^2+1/15\) and
\((2/3)\gamma^2+1/12+\gamma/3\), both \(<.8\).
Pay \(\mu=.2\) from the INDIVIDUAL strong
coefficients: the largest debit is
\(.2(5(.81)/6+1/12)<.151667<b/12=.164\),
and the far other debit \(.02\) is below
\((a/5-(b-a))/2=.147\).
ONE square gives
\(e_A>E_5+.4(.4331)-.04(.8)=.822759\).

For helper at six/seven, use only
\(H_5=H_9=a\), with flat middle block of four
coordinates.
Energy \(E_0=a^2/5=.710645\), both individual
coefficients \(.1885\).
Writing \(l=i-5\), the flat-block projection is
\[
 \gamma[(4-l)\Omega_5+l\Omega_9]/4
                         +\Omega_{\rm helper\ end}/10 .
\]
Its reference product is
\(a(\gamma+.1)<1.71535\).
The smaller time is \(A>2.155\), so the
deficit is \(>.328\).
The FULL residual norm is
\(\gamma^2[2l(4-l)/4-1]+1/15<.724\).
At \(\mu=.1\), EACH actual debit is
\(<.091<.1885\).
ONE square gives
\[
                           e_A>.710645+.2(.328)-.01(.724)
                                      =.769005 .
\]
Every golden case therefore already exceeds
the target energy below.

A weak near-four also costs too much, without
a false arithmetic conjugate ban.
Its actual simple-three mate \(j=4,5,6\) has
the reciprocal finite endpoint data
\[
 \delta_j={1\over14-j},\quad
 \sigma_j={1\over2}+{13-j\over(14-j)\sqrt2},\quad
 C_j=\sigma_j\sqrt{(14-j)D_j}>31/4 .
\]
Positive finite inverse projection proves an
ACTUAL original endpoint, with every omitted
wall strict/nonpositive as required.
Thus its normalized dual height is
\(>\sqrt{27/35}>.878\).
At site nine every nongolden nonweak role has
raw loss \(<.618\), hence normalized height
\(>.75\).
If nine itself is weak near-four, use its
ACTUAL mate ten, whose height is \(>.878\).
The other orientation, with mate eight, has
side order nine and was already strong.
The two chosen unit duals have \(j\le6,k\ge9\),
correlation
\[
                         \rho\le\sqrt{5/12}<2/3 .
\]
Cauchy on their sum gives
\(e_A>(.878+.75)^2/[2(1+2/3)]=.7951152\).
Reflection only reverses index names.

It remains to absorb ALL isolated feet.
Own-three/four/six/parallel have their ALL-foreign
actual nulls before the \(Q\) split.
Pure five's actual \(H_4\) loss is
\(<595/1056<.57\).
Every pure finite native at least seven uses
the positive adjacent core already proved in
Section~\ref{app:arithA14fourfifths};
at this rank its loss is
\[
        <{4/3\over q+\sqrt{D-4/3}}+1/\sqrt{12}
                       <.279+.289=.568<\ell .
\]
For \(Q>0,u>1\), only own five remains,
and positive finite clearing gives loss
\(<.206+.383=.589<\ell\).
The \(u=1\) vector is a forbidden feasible null.
Retain EVERY \(0<u<1\), pure or foreign,
with its actual pressure
\[
 x_A\cdot(\Omega_i-un_i)\ge
                      \sqrt{D_i-u}-\sqrt{u(1-u)}.
\]
Its raw loss is
\(<a_i u+\sqrt{u(1-u)}\),
with \(a_{5,9}<.206\), \(a_{6,7,8}<.199\).
The raw circle maximum is \(<.6135\),
justifying its earlier use.

Choose symmetric heights, linear on the ends,
\[
                         H_5=1.9359,\quad H_6=2.0191,\quad H_7=2.0462.
\]
The exact floor squares after adding \(\ell\) are
\(2.5354^2<45/7\), \(2.6186^2<48/7\),
\(2.6457^2<7\).
Its central coefficients are
\[
                    \beta_5=\beta_9=.15199,\quad
                    \beta_6=\beta_8=.02805,\quad\beta_7=.0271 .
\]
The energy is
\[
 E=1.9359^2/5+.0832^2+.0271^2=.757198412,\qquad
 E-53/70={97221\over1750000000}>0 .
\]
Assume \(e_A\le E\).
The golden and near-four contradictions leave
only genuine strong constraints or isolated
ACTUAL feet at every paid index.
Their FULL adjusted losses satisfy
\[
               a_i u+\sqrt{u(1-u)}
                          -(\beta_i/2)(2u-u^2)<.5995 .
\]
For outer sites the proof works already at
\(\beta=.08\):
below \(.4\) raw loss \(<.5824\), above it
\(.6135-.04(.64)=.5879\).
For inner sites below \(.5\), raw loss \(<.5995\);
above it use
\[
              (.199+1.01961)/2-.0271(3/8)
                                   =.5991425<.5995 .
\]
All root comparisons are positive rational
squares, covering the ENTIRE intervals.

Form ONE vector \(Z=w-\sum\beta_i u_i n_i\)
over ALL retained weak feet.
The actual pressures and every adjacent-foot
cross give exactly the simultaneous-foot
inequalities proved above.
The single square implies \(e_A>E\),
contradicting the assumed \(e_A\le E\).
Strictness comes from strong floors if no
foot is present, otherwise from each
own-coefficient adjusted loss.
Thus
\[
                             e_{A_{13}}>.757198412>53/70
\tag{A13U}\label{eq:arithA13universal}
\]
with every actual field and arbitrary finite
spectators.
Only now add the disjoint original ordinary
costs: \(A_3>1/4\), or \(T_2>1/7,A_1>.1\).
They close BOTH ordinary rank-three alternatives.
The actual-\(I\) and actual-\(O_7\) alternatives
retain their separate hypotheses and proof.

\subsection{The actual seven-field \(A_{13}\) consumer}
\label{app:arithA13seven}

For the actual \(O_7+A_1\) alternative,
the ACTUAL field is \(K_7\).
We use a direct full-order integer strip,
rather than any computed field census.
Put \(\vartheta=2\cos(\pi/7)\),
\(f(T)=T^3-T^2-2T+1\).
Its full ring is \(\mathbb Z[\vartheta]\).
This follows already from the proved full ring
\(\mathbb Z[u]\), \(u^3-5u^2+6u-1=0\):
\[
                    \vartheta=u^2-3u+1,\qquad u=\vartheta^2 .
\]
Thus no new external integral-basis assumption
is required here.

There is NO algebraic integer \(w\in K_7\)
with physical \(0<w<256/25\) and both other
values in \((0,28/45)\).
To prove this, exact polynomial signs give roots
\[
 1.8019<\vartheta<1.802,\quad .445<\alpha<.446,\quad
                         -1.248<\beta<-1.246 .
\]
Write \(w=a+b\vartheta+c\vartheta^2\) in
the FULL ring.
For \(d=b-(\vartheta-1)c\),
\[
             w_\alpha-w_\beta=(\alpha-\beta)d,\qquad |d|<.369 ,
\]
using \(\alpha-\beta>1.691\).
Put \(\Delta=w-(w_\alpha+w_\beta)/2\).
The exact trace identities give
\[
 \Delta={(3\vartheta-1)\over2}d+
                     (3\vartheta^2-2\vartheta-2)c.
\]
The first coefficient is \(<2.203\), its product
with \(d\) has magnitude \(<.812907\);
the second coefficient is \(>4.136\).
Since \(-28/45<\Delta<10.24\), this forces
\(c=0,1,2\).
For \(c=0\), \(b=0\), leaving a positive integer
\(a<28/45\).
For \(c=2\), the \(d\) interval gives
\(1.2348<b<1.973\), no integer.
For \(c=1\), it gives \(b=1\).
Now
\(.643025<\alpha+\alpha^2<.644916\);
positivity of \(w_\alpha=a+\alpha+\alpha^2\)
forces \(a\ge0\), but then
\(w_\alpha>.643025>28/45\).
This exhausts every full-order coefficient.

Reclassify the SAME N14 central sites at the
NEW loss \(\ell=.574\), not the previous
\(.5995\) list.
The exact ratio is
\[
 R_\ell(q)<500/287-1213/8139<8/5 .
\]
The finite-white list and own-zero reciprocal
repairs are unchanged, with their full foreign
costs.
Native five is absent in \(K_7\).
Every own-four/simple-three near role, at
any position and with any foreign contacts,
has the actual \(\sqrt2\) cycle and is absent
in the cubic field.
The remaining COMPLETE nonisolated comparisons
against the larger threshold \(8/5\) are:
near-three \(>(1.3)/\sqrt{.6}\);
D-position-three \(>(1.1)/\sqrt{.4}\);
E-position \(>1.8\);
far-four \(>.9\sqrt5\);
near-four white contact has
\(\sigma>53/50,\delta\le1/5\);
two-white-three has \(q/2+1/q>8/5\).
The full FOREIGN far-three bound is
\[
                {1.8m+1\over\sqrt{m(m+2)}}\ge10/\sqrt{35}>8/5
                    \quad(m\ge5).
\]
It increases with \(m\), and positive foreign
cost only strengthens every other ratio.
All own-three/four/six/parallel affine cases
retain the original nulls in their proper roles.
Thus any weak nonisolated cap is pure endpoint
helper with \(\gamma>.9\),
\(9/5<t=2\gamma<57/25\).
The signed Schur costs proved above still give
every other conjugate in \((-.83,.5)\).
In actual \(K_7\), \(t\) has degree one or three.
Degree three gives unit norm, both small
conjugates negative, and integer trace
\[
                      .14<\operatorname{Tr}(t)
                              \le t-2/\sqrt t<1,
\]
using \(t(t-1)^2<58368/15625<4\).
Degree one leaves only parallel two.
No weak nonisolated cap remains.

For an isolated cap keep its ENTIRE spectator
projection.
Its actual own cycle \(w=4\gamma^2\) is integral
in \(K_7\).
Original \(K\)-rescaling gives new diagonal
\(2/w\), marked old product \(-1\).
At both other embeddings, the old-\(A\)
Schur complement requires
\[
                       0<w_\sigma<4/D_i\le28/45 .
\]
Weakness would give
\(\gamma/\sqrt{1-Q}\le R_\ell(q)<8/5\),
hence physical \(0<w<256/25\).
This is precisely the impossible FULL-order
strip just proved.
It covers EVERY finite/dotted/pure/foreign
isolated cap, not only generating \(w\)'s.
No weak feet remain.

All sites five through nine are now
\(.574\)-strong.
The symmetric positive reference with
\[
 H_5=H_9=1.9614,\quad H_6=H_8=2.0446,\quad H_7=2.0717
\]
has genuine floor squares
\(2.5354^2<45/7\), \(2.6186^2<48/7\),
\(2.6457^2<7\).
Its duals are
\(.15454,.02805,.0271,.02805,.15454\),
and ONE square gives
\[
 e_A>1.9614^2/5+.0832^2+.0271^2
                         =.777074642>233/300 .
\]
The ACTUAL \(O_7\) spectator costs \(>.09\)
and propagates the last ordinary \(A_1>2/15\).
Adding these disjoint original projections
ONCE gives \(233/300+.09+2/15=1\), strictly
exceeded.
A bare old seven bond is never substituted
for the actual \(O_7\) profile.

\subsection{The \(A_{13}+I+T_2\) consumer with every remaining foot}
\label{app:arithA13I}

Here \(I\) is the ACTUAL isolated short wall,
and \(T_2\) an ORDINARY irreducible dihedral.
The strengthened ordinary theorem proved in
Section~\ref{app:arithdihedral}, together with
the separate ordinary \(A_2\) bound, gives
\(e_T>3/19\).
Use the actual rank-free \(I\) parameters
from~\eqref{eq:Ibarrier}:
\[
 e_I\ge u^2,\quad0<u<1/3,\quad t+3u/2>.767 .
\]
Both signed \(A_{13}\) endpoint coordinates
satisfy
\[
 x_1,-x_{14}\ge\alpha t,\qquad
                     \alpha=\sqrt{13/14}>.9636.
\]
All these pressures concern the SAME original
farthest vector.
The \(I\) energy is counted once.

For endpoint means \(b<.2738\), or \(b<.2666\),
the positive linear lower bounds are
\[
 \alpha t-b>.4652-1.446u,\qquad
 \alpha t-b>.4724-1.446u,
\]
respectively.
They follow from
\(.9636(.767-1.5u)=.7390812-1.4454u\).
The following exact completion applies on
the FULL continuous domain:
\[
 (1+cB^2)[u^2+c(A-Bu)^2]-cA^2
                        =[(1+cB^2)u-cAB]^2 .
\]
For \(B=1.446\), the positive rational margins
\[
 (5/2)(.4652)^2-.0868[1+(5/2)B^2]=.000498828,
\]
\[
 (5/4)(.4724)^2-.077[1+(5/4)B^2]=.000701535
\]
therefore give the uniform bounds
\(u^2+(5/2)(.4652-Bu)^2>.0868\) and
\(u^2+(5/4)(.4724-Bu)^2>.077\).

The COMPLETE N14 helper/near-four classifications
of Section~\ref{app:arithA13universal} do not
depend on its contrary energy assumption.
A golden helper at five gives \(e_A>.822759\),
which closes with the joint \(I+T_2>.21\).
For a helper at six/seven, retain its SAME
one-reference \(H_4\) square of energy
\(L>.769005\), with price \(<.091\) at each
outer coefficient \(.1885\).
Choose the five-coordinate block OPPOSITE its
helper endpoint.
The actual \(H_4\) displacement and its flat
block projection are constant there; their
residual vanishes POINTWISE on that block.
Thus an endpoint residual \(G\) has
\[
 G^2=4/5,\quad w\cdot G=Y\cdot G=0,\quad
             b=1.885/(5\sqrt2)<.2666 .
\]
Its individual price is
\(\kappa/(5\sqrt2)\), with remaining capacity
\(>.0975\).
One JOINT pressure square gives
\[
                e_A>L+2\kappa(\alpha t-b)-(4/5)\kappa^2 .
\]
For \(u\ge.3\), \(L+3/19+.09>1\).
For \(u\le.2\), choose \(\kappa=.2\);
its price is \(<.0283\).
The total is
\(>L+3/19+u^2+.4(.4724-Bu)-.032\),
decreasing to an endpoint additive value
\(.08128\), already contradictory.
For \(.2\le u\le.3\), put
\(d=.4724-Bu>0\), \(\kappa=5d/4\);
its price is \(<.033<.0975\).
The total exceeds \(L+3/19+.077>1\)
by the completed-square certificate.
Both helper ends and reflected indices are
covered with ACTUAL orthogonal residuals.

Every weak near-four gives \(e_A>.7951152\);
adding joint \(I+T_2>.21\) also exceeds one.
We are left with ONLY true strong sites or
the actual isolated feet of
Section~\ref{app:arithA13universal}.

Use that section's full rational reference,
energy \(E=.757198412\), outer coefficients
\(.15199\), inner \(.02805,.0271,.02805\).
Its two endpoint residuals are
\[
 G_L=e_1-P_5/5,\qquad
 G_R=-e_{14}+(e_{10}+\cdots+e_{14})/5 .
\]
They are disjoint, have norm \(4/5\), and
are orthogonal to \(w\).
For \(W=w+\kappa(G_L+G_R)\), EACH outer
central coefficient becomes
\[
                  \beta'_5=\beta'_9=.15199-\kappa/(5\sqrt2).
\]
The others are unchanged.
For EVERY remaining foot parameter \(v_i\in(0,1)\),
form ONE vector
\(T=W-\sum\beta'_i v_i n_i\).
The primitive endpoint/strong/foot pressures,
paid BEFORE taking its square, give
\[
 \begin{gathered}
 x_A\cdot T\ge E+2\kappa(\alpha t-b)
                         -\sum\beta'_i(L_i-\ell),\\
 T^2\le E+(8/5)\kappa^2
                         -\sum(\beta'_i)^2(2v_i-v_i^2).
 \end{gathered}
\]
Here \(b=1.9359/(5\sqrt2)<.2738\).
Every endpoint/foot cross is incorporated
through the REMAINING coefficient \(\beta'_i\),
and every adjacent-foot cross has the actual
nonpositive old-root sign.
The proved continuous adjusted-loss inequalities
hold at outer coefficients at least \(.08\)
and the unchanged inner coefficients.
Consequently
\[
                    e_A>E+4\kappa(\alpha t-b)-(8/5)\kappa^2 .
\]
For \(u\ge.3\), \(E+3/19+.09>1\).
For \(u\le.2\), choose \(\kappa=.2\);
both outer coefficients remain \(>.1236>.08\).
The total is
\(>E+3/19+u^2+.8(.4652-Bu)-.064\),
decreasing to endpoint additive value \(.1168\).
For \(.2\le u\le.3\), use
\(\kappa=5(.4652-Bu)/4\le.22\);
both outer coefficients remain \(>.1199>.08\).
The exact completion gives total
\(>E+3/19+.0868>1\).
All actual \(I\) parameters, ALL foreign/adjacent
feet and both endpoints have been covered.
This closes the last \(A_{13}\) rank-sixteen
alternative with one original vector and one
\(I\) energy.

\subsection{A universal \(A_{12}\) bound from one paid endpoint}
\label{app:arithA12universal}

Set \(N=13\), central sites five through eight,
\[
 D_5=D_8=80/13,\quad D_6=D_7=84/13,\qquad
                         \ell=.6005 .
\]
Keep arbitrary finite spectators and EVERY
actual field.
Assume \(e_A\le5/7\).
The fresh exact sufficient ratio is
\[
 R_\ell(q)<2000/1201-77980/458949<3/2,\qquad q<51/20,
\]
with positive margin \(5105207/(2402\cdot458949)\).
Deleting a central site leaves two paths of rank
at least four.
The COMPLETE white list is isolation,
permitted endpoint/D/E simple three, endpoint
four, or one endpoint-three on EACH path.
Connected rank twelve precludes exceptional
completions; its two long arms preclude a
fork completion.
Every spectator contribution remains in the
actual retained finite Gram.

The complete own-zero residual repairs are
\[
 \begin{array}{c|c}
 A^2&Q\text{ bound}\\ \hline
 Q-7/13&Q<3/5\\
 Q-2/13&Q<2/5\\
 Q-1/13&Q<1/5 .
 \end{array}
\]
The first positive time is \(<4/65<1/4\),
below any foreign contribution.
The second permits at most one foreign component,
with contribution \(Q>A^2\).
Both violate the actual reciprocal orthogonal
opposite law.
The last has no possible foreign contact.
At \(Q=0\) their times are negative.
Internal-three/four/two-contact cases have the
actual affine arms/star/cycle.
ALL isolated owns three/four/six and parallel
retain their positive future nulls before
splitting the foreign cost.

The ordinary nonisolated ratios exceed
\(38/25>3/2\), including the same first three
side-order-five comparisons printed above,
and two-path \(q/2+1/q>97/60\).
Foreign far/near bounds are still
\(227/(25\sqrt{35})\), \(\sqrt{121/35}\).
Thus every weak nonisolated cap is PURE:
an endpoint helper, or OWN-four/SIMPLE-three
near helper, with BOTH orientations retained.
The roles are not interchanged.

For an endpoint helper the exact ratio-square
identity gives \(t=2\gamma<17/8\).
Non-five own has \(t>9/5\).
At \(N=13,i=5\), the signed conjugate Schur
costs at \(-.83,.5\) exceed four by
\(69/1625,2/13\), and they increase across
the half-window.
Hence every other conjugate lies in
\((-.83,.5)\).
The COMPLETE all-degree trace/norm proof
of Section~\ref{app:arithA13universal} uses
only these bounds, \(9/5<t<9/4\), actual
integrality and all lifted embeddings.
It excludes every non-five helper.
A weak endpoint helper is therefore golden,
actual \(K=\mathbb Q(\varphi)\).
A near-four simple-three normalization instead
forces actual \(K=\mathbb Q(\sqrt2)\), so
it cannot coexist with a golden helper.

The COMPLETE \(26A^2\) golden time rows are
\[
\begin{array}{c|c|r}
 \text{type}&26A^2&\mathrm N(26A^2)\\ \hline
 \text{near }5&26+56\varphi&-1004\\
 \text{near }6&28+56\varphi&-784\\
 \text{far }5&26+50\varphi&-524\\
 \text{far }6&28+54\varphi&-620 .
\end{array}
\]
The four norm-squareclasses differ.
Reflected cyclic parameters are
\(21\varphi/13,9\varphi/13\) at five,
and have constant \(-9/13\), coefficient
\(8/13\) or \(4/13\) at six.
All are nonintegral in the full golden ring.
Thus at most ONE golden helper occurs.
Every other central site is \(.65\)-strong,
by the complete geometric inventory at that
loss; another weak helper would also be
\(.6005\)-weak.

Use its ACTUAL \(H_4\) correction and positive
finite projected point as before, with
\(\delta=(5-3\varphi)/12<.111^2\).
For helper at five, take
\(a=1.830,b=1.891\), heights
\(H_6=H_7=b,H_8=a\), linear tails.
Its energy
\[
                 E=(b^2/6+a^2/5+.061^2)/2=.6347405833\ldots
\]
has individual coefficients
\(b/12>.15758,.0305,.1525\).
The far/near flat-block projections are
\((5\gamma/6)\Omega_6+\Omega_8/10\) and
\((5\gamma/6+1/12)\Omega_6\).
The reference product is \(<1.459425\),
smaller actual time \(>2.027\), deficit \(>.45\).
FULL residual norm is \(<.8\).
Pay \(\mu=.2\): largest own price
\(<.151667\), far other price \(.02\),
both individually affordable.
ONE square gives \(e_A>E+.148>5/7\).

For helper at six, take only \(H_5=H_8=a\).
Energy \(a^2/5=.66978\), both outer prices \(.183\).
The flat middle block has three coordinates.
The projections are
\((2\gamma/3)\Omega_5+(\gamma/3+.1)\Omega_8\)
or the reverse placement of the \(.1\).
The reference product is \(<1.6653\);
smaller actual time \(>2.106\), deficit \(>.329\).
FULL residual norm is \(\gamma^2/3+1/15<.286\).
At \(\mu=.1\), each price is \(<.064<.183\).
ONE square gives
\(.66978+.2(.329)-.01(.286)=.73272>5/7\).
Reflections merely reverse labels.
No golden helper occurs under the contrary
energy assumption.

For near-four side order \(m\), its exact ratio
\[
 R_m={1/\sqrt2+(m-1)/(2m)\over\sqrt{(m+1)/(2m)}}
\]
increases in \(m\).
At seven,
\(R_7^2=67/56+3\sqrt2/4>9/4\), by \(42\sqrt2>59\).
Weak near-four therefore has only own/helper
positions \(5/4,6/5\) or their reflections.
Deleting the ACTUAL helper \(j=4,5\) gives a
finite \(B\) completion and reciprocal endpoint:
\[
 \delta_j={1\over13-j},\quad
 \sigma_j={1\over2}+{12-j\over(13-j)\sqrt2},\quad
 C_j=\sigma_j\sqrt{(13-j)D_j}>31/4.
\]
Positive finite projection checks ALL original
walls and the reciprocal law identifies the
SAME actual opposite.
Its normalized height is \(>.878\).
On the opposite side, index eight has height
\(>.75\) from loss \(<.618\), or use its ACTUAL
near-four mate nine with height \(>.878\).
The two normalized duals have correlation
\(\le5/8\), so
\[
                       e_A>{(.878+.75)^2\over2(1+5/8)}>4/5.
\]
No weak near-four remains.

For isolated caps keep the entire \(Q\).
Pure five has the actual LEFT \(H_4\)
correction with
\(u=\gamma^{-2}<25/16\),
\(\sqrt{\delta_5}/\gamma<5/24\).
The divided displacement is the old concave
combination at barycentre \(i\), so it is
bounded above by \(h_i\).
At this ACTUAL \(D_i\ge80/13\),
\[
 q+\sqrt{D_i-25/16}>12/5+2=22/5,
\]
and its loss is
\[
                    <125/352+5/24=595/1056<.57 .
\]
This proves the exact local bound used in
the preceding N14 theorem; no larger-rank
denominator is assumed.

EVERY pure finite native at least seven uses
the positive adjacent core
\(c+\gamma^{-1}n_i+2(\gamma^{-1}-\gamma)n_j\).
It is zero on cap and \(n_i\), negative on
\(n_j\), future timelike and nonpositive at
ALL omitted original walls.
Its pressure loses
\[
       <{4/3\over q+\sqrt{D_i-4/3}}+1/\sqrt{12}
                        <.286+.289=.575<\ell .
\]
Here \(q>2.48\),
\(\sqrt{D_i-4/3}>2.195\).
For \(Q>0,u>1\), only own five is possible,
since every other remaining finite \(\gamma>.9\)
has \(u\le.75/.81<1\).
Positive finite clearing gives loss
\(<.212+.383=.595\).
At \(u=1\) its actual vector is a forbidden
future feasible null.
Retain ALL \(0<u<1\), pure/foreign, with pressure
\[
 x_A\cdot(\Omega_i-un_i)\ge
                   \sqrt{D_i-u}-\sqrt{u(1-u)} .
\]
Its loss is \(<a_i u+\sqrt{u(1-u)}\), where
\(a_{5,8}<.212\), \(a_{6,7}<.205\);
the raw circle maximum \(<.6175\) supplies
the mate argument above.

Fix the symmetric reference
\[
 H_5=H_8=1.8795,\quad H_6=H_7=1.9405,\qquad
 \beta_{5,8}=.15745,\quad\beta_{6,7}=.0305,
 \quad E=.71022505 .
\]
All heights are strictly below the true
strong floors.
The FULL continuous adjusted-loss estimate
\[
 a_i u+\sqrt{u(1-u)}-(\beta_i/2)(2u-u^2)<.6005
\]
holds even for ANY outer \(\beta_i\ge.08\).
For outer sites below \(.4\), raw loss
\(\le.5848\); above it,
\(.6175-.04(.64)=.5919\).
For inner sites the expression is at most
\(.1745u+.01525u^2+\sqrt{u(1-u)}\).
On \((0,.5]\) it is \(\le.5910625\).
On \([.5,.7]\), \(u^2\le1.2u-.35\)
and the circle bound with coefficient \(.1928\)
give \(<.6003125\).
On \([.7,1)\),
\(\sqrt{u(1-u)}<(3-u)/5\),
by positive square difference
\((2u-1)(13u-9)/25\), giving \(<.6\).
Thus EVERY endpoint of the continuous split
and every foot parameter are covered.

We record the full JOINT endpoint/foot rule.
Let \(\bar Z\) sum to one, supported on the
first at most FOUR coordinates, and let its
ACTUAL pure pressure be \(x_A\cdot Z\ge h\),
where \(Z=\bar Z-P_{13}/13\).
Its flat-end projection is
\(\Pi Z=\Omega_5/(5\sqrt2)\).
For \(Y=Z-\Pi Z\),
\[
 w\cdot Y=0,\quad S=Y^2=\|\bar Z\|^2-1/5,\quad
                         b=w\cdot\Pi Z<.266 .
\]
For \(0\le\mu\le.4\), spend
\(\mu/(5\sqrt2)\) from INDIVIDUAL \(\beta_5\).
Its remainder is \(>.1008>.08\); the opposite
outer and both inner coefficients are unchanged.
For EVERY weak foot use this REMAINING
\(\beta_i'=(w+\mu Y)\cdot n_i\), and form
\[
                    T=w+\mu Y-\sum\beta_i' u_i n_i .
\]
Primitive pressures and ALL old-root crosses give
\[
 \begin{gathered}
 x_A\cdot T\ge E+\mu(h-b)-\sum\beta_i'(L_i-\ell),\\
 T^2\le E+\mu^2S-\sum(\beta_i')^2(2u_i-u_i^2).
 \end{gathered}
\]
ONE square and the just-proved adjusted losses
yield
\[
                           e_A>E+2\mu(h-b)-\mu^2S .
\tag{A12P}\label{eq:arithA12endpointpayment}
\]
Endpoint/foot crosses have not been discarded:
they are paid through the REMAINING individual
coefficients.
Reflection gives the same rule at the other end.

If actual \(x_1\ge.34\), take \(\bar Z=e_1\),
\(S=4/5,\mu=.1\).
Then \(e_A>.71702505>5/7\), impossible.
Likewise \(-x_{13}<.34\).
The two ACTUAL endpoint edges therefore have
half-tangent square
\[
                 <(13/12)(.34)^2=3757/30000<3/23 ,
\]
hence cosh \(<13/10\).
The strict rational margin is \(3589/690000\).

Deleting endpoint one retains \(A_{11}\) and
ALL spectators.
Its COMPLETE finite list is isolation, simple
three at global \(j=2,3,11,12\), or four at
\(j=2,12\).
The six pure coshes for positive own \(\gamma\)
are respectively
\[
 {24\gamma+11\over13},\
 {24\gamma+10\over2\sqrt{13}},\
 {24\gamma+2\over2\sqrt{13}},\
 {24\gamma+1\over13},\
 {24\gamma+11\sqrt2\over\sqrt{26}},\
 {24\gamma+\sqrt2\over\sqrt{26}} .
\]
Foreign \(Q\) only increases them.
At \(\gamma\ge1/2\), all but far-three exceed
\(1.3\); far-three/own-three has its actual
affine cycle, and every remaining own gives
\((12\sqrt2+1)/13>1.3\).
For own zero with white three, the spatial
directional square is \((13-j)^2/312\le121/312\).
The actual short length bound gives time square
\(<69(121)/(169\cdot312)<1/4\).
Any positive foreign contribution violates
the reciprocal rule.
At \(Q=0\), the remaining times are negative
or the actual affine \(1,2,5\) principal.
White-four at two has its affine star of
squared weights \(1/2,1/4,1/4\);
far twelve has \(Q<1/12\), hence \(Q=0\)
and negative time.
Every own-zero and nonisolated role is excluded.

Each ACTUAL endpoint is therefore isolated
within \(A_{12}\), with foreign contacts retained:
\[
 c=Ap-\gamma\Omega_1-v_S,\quad
 u=(1-Q)/\gamma^2,\quad D_1=24/13,\quad
 C^2=D_1/u<169/100 .
\]
Thus \(u>2400/2197>1\), excluding the
dotted/parallel/\(u\le1\) domain by an ACTUAL
edge bound.
Own three has \(Q<1-600/2197<3/4\),
at most two contacted spectator components.
Their largest contribution exceeds
\(Q/2>A^2=Q-7/13\), forbidden by reciprocity;
at \(Q=0\) time is negative.
Own four has \(Q<1-1200/2197<1/2\),
so at most one foreign contribution
\(Q>A^2=Q-1/13\).
Own six has the actual \(6,3\) null.
Only finite owns five or at least seven remain.

For pure five use the ACTUAL positive finite
\(H_4\) projection on \(c,n_1,n_2,n_3\),
with \(\delta=(5-3\varphi)/8>0\).
Its sum-one displacement is \(\bar Z=P_4/4\),
\(S=1/20\), and pressure
\[
 h={A^2+\delta-\sqrt{\delta(A^2+\delta)}
                    \over A\gamma\sqrt2},\quad
 A^2=(24/13)\gamma^2-1,\quad
 A^2+\delta=9\gamma^2/26 .
\]
The exact brackets
\(.6545<\gamma^2<.65475\),
\(.2083<A^2<.2088\), \(A<.457\),
\(.017875<\delta<.01825\),
\(A^2+\delta>.2265\)
give \(\sqrt{\delta(A^2+\delta)}<.0644\)
and \(A\gamma\sqrt2<.525\).
Hence \(h>.308\).
Payment~\eqref{eq:arithA12endpointpayment}
at \(\mu=.4\) gains \(>.0256\).

For EVERY pure finite own at least seven use
the ACTUAL adjacent negative core at one/two.
Its sum-one displacement and pressure are
\[
 \bar Z=(1-u/2)(e_1+e_2)+(u-1)e_3,\quad
 h={\sqrt{D_1-u}-\sqrt{7u-4-3u^2}\over\sqrt2},
 \quad S\le.3 .
\]
Since \(u<100/81\),
\(D_1-u>644/1053>.782^2\);
the other radical is \(<.289\).
Thus \(h>.348\), and \(\mu=.4\) gains
\(>.8(.348-.266)-.16(.3)=.0176\).
This covers native eight and the whole finite
tail without a label cutoff.

For EVERY positive-foreign endpoint, \(u>1\)
forces native five.
Clear ALL contacted spectator normals positively
with \(n_1\).
The actual finite projection gives the simpler
pure pressure
\[
 x_A\cdot(\Omega_1-n_1)>
                  \sqrt{D_1-1}-\sqrt{u-1}.
\]
Its stronger exact pressure has additional
positive factor
\(\sqrt{D_1-1}/\sqrt{D_1-u}>1\), so the
displayed lower bound is valid.
Here \(\bar Z=(e_1+e_2)/2\), \(S=.3\),
\(\sqrt{11/13}>.919\), \(\sqrt{u-1}<.383\),
so \(h>.378\).
At \(\mu=.3\), gain \(>.0402\).

All these points have positive ACTUAL wall
coefficients, correct selected products, positive
future time and nonpositive products at EVERY
omitted original wall.
Every \(\bar Z\) has the support required
by the JOINT payment rule.
Even the weakest endpoint gain gives the
CONTRARY-branch bound
\(.71022505+.0176=.72782505>5/7\).
Thus the assumed \(e_A\le5/7\) is impossible:
\[
                                  e_{A_{12}}>5/7
\tag{A12U}\label{eq:arithA12universal}
\]
with every actual field and arbitrary spectators.
The number \(.72782505\) is not asserted
as an unconditional component rate.

This one stronger theorem closes the rank-sixteen
ordinary-rank-four and actual-\(I+T_3\) rows,
whose disjoint or JOINT remainder costs \(>2/7\).
It also closes BOTH typed-seven rank-two rows:
the actual defect costs \(>.09\) and propagates
the remaining ordinary \(T_2>.2\), giving
\(5/7+.09+.2>1\).
The weaker earlier typed-only component estimate
therefore needs no separate chapter.
In rank seventeen it closes every older
\(A_{12}\) row by the actual rank-five spectator
partitions and rank-free joint costs.

\subsection{A rank-eleven chain: the complete local alternatives}
\label{app:arithA11local}

We next prove a uniform estimate at the smaller
path size, and then a joint estimate when an
ACTUAL short isolated \(I\) is present.
Both arguments use the original vector \(x\);
the second will rebuild its reference rather
than add a new gain to an optimized first one.
Write \(N=12\), \(A=A_{11}\), and use
\[
 n_i=(e_i-e_{i+1})/\sqrt2,\qquad
 D_i={2i(12-i)\over12},\qquad
 C_{ij}={2i(12-j)\over12}\quad(i\le j).
\]
Only sites \(5,6,7\) will be paid.
Put \(\ell=.61\), \(q_i=\sqrt{D_i}\).
Then
\[
 D_5=D_7=35/6,\quad D_6=6,\qquad
 a_i={1\over q_i+\sqrt{D_i-1}}<.217.
\tag{A11a}\label{eq:arithA11a}
\]
Indeed \(q_i>2.415\),
\(\sqrt{D_i-1}>2.198\), and
\((2.415+2.198)(.217)=1.001021>1\).
A site is strong if \(h_i>q_i-\ell\).

Deleting one of these sites retains two paths
of ranks at least four.
The complete one-node finite extensions are:
no white contact; one admissible endpoint,
\(D\)- or \(E\)-position three; one terminal
four; or two endpoint threes joining the paths.
Two contacts to one path make a forbidden
cycle.  A connected completion has rank eleven,
so exceptional finite diagrams cannot add a
missing joined-path role.
We retain the ENTIRE contacted spectator forest,
writing its squared projection as \(Q\);
every positive whole-component contribution
is at least \(1/4\).

The own-zero cases are excluded before any
positive-own ratio is used.
An internal three contains the actual affine
arms \(1,2,5\) or \(1,3,3\); an internal four
contains the affine \(3,3,4\) star; two
white threes contain the affine cycle.
The remaining far-three role, with side
order \(m\ge5\), has time square
\(<1/(2m)\), below any necessary positive
foreign contribution.
The next-to-far three has \(Q<2/m\), hence
at most one contacted whole spectator, and
time square \(Q-1/6<Q\).
Either would make this same actual wall
close an edge in that orthogonal spectator,
contrary to the actual reciprocal-opposite
law.  With \(Q=0\) its time is negative.
Far-four has \(Q<1/m\le1/5\), so \(Q=0\)
and negative time.
These are all residual own-zero roles.
The affine vectors have positive time through
the actual cap and nonpositive product with
every omitted original wall.
They are forbidden feasible future nulls,
even in the presence of foreign contacts.

For \(q\ge12/5\), the sufficient strong-edge
ratio is
\[
 R_\ell(q)={100\over61}
              -{q-.61\over q(2q-.61)}.
\]
It increases with \(q\), and
\[
 R_\ell(q)\le R_\ell(3)<3/2,\qquad
 {3\over2}-R_\ell(3)={1669\over197274}>0.
\]
The complete near-three, \(D/E\), terminal-four
and two-white bridge formulas have, respectively,
the lower comparisons
\[
 { .808+.4\over\sqrt{.6}},\quad
 { .808+.2\over\sqrt{.4}},\quad
 { .808\over\sqrt{.25}},\quad
 q/2+1/q>97/60 .
\]
All exceed \(38/25>3/2\).
The smaller positive owns have their actual
affine cycle, star or \(6\)--\(3\) obstruction.
For the two exceptional profiles WITH \(Q>0\),
clearing the entire spectator gives ratio
\[
 >{227\over25\sqrt{35}}>38/25
       \quad\hbox{or ratio square }>121/35.
\]
Consequently a weak NONISOLATED cap is PURE,
and is either a simple-three END-helper
\[
 c=A_cp-\gamma\Omega_i-\Omega_e/2,\qquad
 e=1\ \hbox{or }11,
\tag{A11helper}\label{eq:arithA11helper}
\]
or own-four/simple-three near helper
\[
 c=A_cp-\Omega_i/\sqrt2-\Omega_j/2,\qquad |i-j|=1.
\tag{A11near4}\label{eq:arithA11near4}
\]
Both orientations are retained.
The near-four ratio increases with side order;
from order eight it already exceeds \(3/2\).
Thus its complete list, up to reflection, is
\((i,j)=(5,4),(6,5),(5,6)\).
It would be incorrect to exclude these three
roles by the larger-path conjugate test.

For an isolated-within-\(A\) cap retain
\[
 c=A_cp-\gamma\Omega_i-v_S,\quad
 Q=\|v_S\|^2,\quad u=(1-Q)/\gamma^2 .
\]
Before splitting \(Q\), own-three, four, six
and parallel have their actual affine
\(E_7\) arms \(1,3,3\), weighted star,
\(6\)--\(3\) path or positive null pair.
These nulls are feasible against ALL foreign
walls by acuteness.
Hence every remaining positive own has
\(\gamma\ge\varphi/2\), \(\varphi=(1+\sqrt5)/2\).
For \(Q>0,u>1\), positive clearing gives loss
\[
 <a_i+\sqrt{u-1}<.217+(2-\varphi)<.599<\ell.
\]
The actual \(u=1\) vector is a forbidden
future feasible null.

The PURE finite cases with \(u>1\) are
handled by actual points at this rank.
For own five the same positive \(H_4\)
projection used above has loss
\[
 <{25/16\over q_i+\sqrt{D_i-25/16}}+5/24
 <{625\over1792}+{5\over24}
 ={2995\over5376}<.558.
\]
Here the denominator is \(>2.415+2.066>4.48\);
\(2.066^2<205/48=D_{\min}-25/16\).
For EVERY finite native at least seven put
\(d=1-\gamma^2\) and choose an old adjacent
root \(n_j\).  The positive actual combination
\[
 \xi=c+\gamma^{-1}n_i+(2d/\gamma)n_j
\]
is zero on \(c,n_i\), negative on \(n_j\),
and has square \(-d(1-4d)/\gamma^2<0\).
It is future timelike and all omitted
ORIGINAL wall products are nonpositive.
Its comparison gives loss
\[
 <{100/81\over q_i+\sqrt{D_i-100/81}}
                   +{25\over81}
 <{213975\over369279}<.58<\ell.
\]
We used \(\gamma>.9\),
\(\sqrt{D_i-100/81}>2.144\) and
\(\sqrt{d(1-4d)}\le1/4\).
This includes native eight and the entire
unbounded finite-native tail.

Retain EVERY remaining \(0<u<1\), pure or
foreign.  Positive inverse-Gram clearing
expresses \(v_S\) as a positive combination
of all contacted old spectator roots.
The actual vector
\[
 \eta=c+v_S+\gamma un_i
   =\gamma\{\sqrt{D_i-u}\,p-(\Omega_i-un_i)\}
\]
has positive time, square
\(-\gamma^2u(1-u)\), and zero products on
the cap and cleared spectators.
Its product on \(n_i\) is
\(-\gamma(1-u)<0\); all other original
walls have nonpositive product termwise.
It therefore supplies the SAME-\(x\) pressure
\[
 x_A\cdot(\Omega_i-un_i)
       \ge\sqrt{D_i-u}-\sqrt{u(1-u)}
       =q_i-L_i,\qquad
 L_i<.217u+\sqrt{u(1-u)} .
\tag{A11foot}\label{eq:arithA11foot}
\]
No restriction on the number or adjacency
of these isolated feet is made.

\subsection{The all-degree end-helper reduction at \(N=12\)}
\label{app:arithA11helperfield}

For~\eqref{eq:arithA11helper}, the exact
side-order identity is
\[
 {(\gamma+1/(2m))^2\over(m+1)/(2m)}
                 -(4\gamma-2)
 ={(\gamma-1-1/(2m))^2\over(m+1)/(2m)}\ge0 .
\]
Weakness gives \(\varphi\le t=2\gamma<17/8\);
if the own is not five then \(t>9/5\).
The actual simple-three normalization makes
\(t\) integral in \(K\).
For a lifted nonphysical embedding of
\(\mathbb Q(t)\), strict old-\(A\) Schur
positivity gives the cheapest signed costs
\[
 Q_-(s)=2is^2+4s+2-2(is+1)^2/12,\qquad
 Q_+(s)=2is^2+2-2(is-1)^2/12 .
\]
They increase with \(i\) on the relevant
half-window.  At \(i=5\),
\[
 Q_-(-17/20)-4=31/480>0,\qquad
 Q_+(1/2)-4=1/8>0 .
\]
Their values at zero are below four, and
the outer quadratic branches increase away
from their allowed interval.
Thus EVERY other conjugate obeys
\[
                    -17/20<t_\sigma<1/2 .
\tag{A11strip}\label{eq:arithA11strip}
\]

We give the full degree argument, including
degree five.  In degree at least six a
nonzero integral norm has magnitude
\(<(17/8)(17/20)^5<1\), impossible.
In degrees three through five its magnitude
is less than two.
One positive small conjugate would instead
make it \(<(17/8)(1/2)(17/20)^{a-2}<1\).
All small conjugates are therefore negative
and their absolute product is \(1/t\).
AM--GM bounds their trace above by
\(t-(a-1)t^{-1/(a-1)}\).
The resulting strict trace intervals are
\[
 \begin{array}{c|c|c}
 a&\text{trace interval}&\text{upper comparison}\\ \hline
 3&(0,1)&t(t-1)^2<(17/8)(9/8)^2<4\\
 4&(-1,0)&t^4<(17/8)^4<27\\
 5&(-2,-1)&t(t+1)^4<(17/8)(25/8)^4<256 .
 \end{array}
\]
The lower bounds follow from
\(t>9/5\) and each small conjugate \(>-.85\).
None contains an integer.
In degree two the norm is \(1\) or \(-1\).
A negative small conjugate gives trace
\(t-1/t\in(1,2)\).
A positive one forces \(t>2\) and trace
\(t+1/t\in(2,3)\).
Degree one leaves only \(t=2\), the
forbidden actual parallel pair.
Thus every weak end-helper is native five.

Its actual connected old-path-plus-cap Gram,
normalized by its simple three, is Lorentzian
over \(\mathbb Q(\varphi)\).
A proper totally real extension would have
a nonphysical embedding fixing this negative
direction.  Admissibility forbids it:
the ACTUAL field is \(\mathbb Q(\varphi)\).
A near-four simple-three cap similarly forces
ACTUAL \(K=\mathbb Q(\sqrt2)\).
These restrictions apply only in the stated
exceptional branches; no field list is assumed
for the other caps.

\subsection{Uniform \(A_{11}\) energy, paying a helper and all feet}
\label{app:arithA11uniform}

We prove \(e_{A_{11}}>13/20\) with arbitrary
finite spectators and every actual field.
Assume contrariwise \(e_A\le.65\).
First remove the retained near-four cases
by their OWN actual mate, not a field ban.
For~\eqref{eq:arithA11near4} the three complete
time squares, up to reflection, are
\[
 {39+28\sqrt2\over12},\qquad
 {83+60\sqrt2\over24},\qquad
 {41+30\sqrt2\over12}.
\tag{A11silver}\label{eq:arithA11silvertimes}
\]
Each exceeds six.
With \(v=\Omega_i/\sqrt2+\Omega_j/2\) and
\(T_j=(v\cdot\Omega_j)/(A_c\sqrt{D_j})\),
direct inverse-Gram multiplication gives
\[
 T_j^2-1=
 \begin{cases}
  1/((12-j)A_c^2),&j=i-1,\\
  1/(jA_c^2),&j=i+1 .
 \end{cases}
\]
In all three roles and their reflections the
integer factor is at least six.
Consequently \(1<T_j^2<37/36<(51/50)^2\).
The full old-root replacement at \(j\) is
positive definite, its final Schur complement
being \((v\cdot\Omega_j/\sqrt{D_j})^2-A_c^2>0\).
Its actual off-diagonal entries are nonpositive.
Positive finite inverse projection of the
strict interior point gives a common point
strict against every other ORIGINAL wall.
Hence this cap is the ACTUAL opposite at
\(j\), rather than a formal cutter behind
another facet.
Its half-tangent pressure gives
\[
 x_A\cdot{\Omega_j\over\sqrt{D_j}}
 \ge T_j-\sqrt{T_j^2-1}
 >{51-\sqrt{101}\over50}>.81 .
\]
Thus \(e_A>.6561>.65\), a contradiction.
Every weak nonisolated exception is now golden.

The complete golden time table is
\[
 \begin{array}{c|c|r}
 \text{type}&24A_c^2&\mathrm N(24A_c^2)\\ \hline
 5,\text{ near}&22+49\varphi&-839\\
 6,\text{ either end}&23+48\varphi&-671=-11\cdot61\\
 5,\text{ far}&22+45\varphi&-551=-19\cdot29 .
 \end{array}
\tag{A11gold}\label{eq:arithA11goldtimes}
\]
Both actual norm-two helper caps lie on the
same \(K\)-time line; their time-square ratio
must be a \(K\)-square.
The three norm squareclasses are distinct,
so unlike rows cannot coexist.
Two center-six caps cannot be the distinct
actual opposites of the same edge.
For the reflected near-five and far-five
pairs the actual closed helper paths give
raw cyclic products
\((2-7\varphi)/6\) and \((2-3\varphi)/6\).
Neither is integral in the full ring
\(\mathbb Z[\varphi]\).
Thus at most ONE central index has a weak
golden helper.  This does not limit the
isolated feet of~\eqref{eq:arithA11foot}.

Set \(a=1.805,b=1.839\).
With heights linear on the unmentioned
outer tails, there are three reference cases:
\[
 \begin{array}{c|c|c}
 \text{helper}&\text{positive }\beta_i=w\cdot n_i
                      &E=w^2\\ \hline
 \text{none}&\beta_5=\beta_7=.1635,\ \beta_6=.034&.652761\\
 6&\beta_5=\beta_7=.1805&.651605\\
 5&\beta_6=.17025,\ \beta_7=.1635&.60820725 .
 \end{array}
\tag{A11refs}\label{eq:arithA11references}
\]
In the first case \(H_5=H_7=a,H_6=b\).
For the second replace \(H_6\) by \(a\),
so its helper coefficient is zero.
For the third flatten heights on \(0,\ldots,6\),
with \(H_6=b,H_7=a\); reflect a helper at seven.
Its energy is exactly
\((b^2/6+(b-a)^2+a^2/5)/2\).
All paid heights are strictly below the strong
floors, since \(2.415^2<35/6\) and \(2.449^2<6\).

Only the end-helper case needs a helper payment.
For the actual helper at five put \(g=\varphi/2\),
let \(e=1\) or \(11\) and let \(j=2\) or \(10\)
be its adjacent old root.
The ACTUAL finite \(H_4\) projection uses
\[
 w_{\rm old}=gn_5+(2/3)n_e+(1/3)n_j,\qquad
 \delta=2/3-g^2>0,\qquad U=v-w_{\rm old}.
\]
The positive vector
\[
       \xi=p+(A_c/\delta)(c+w_{\rm old})
\]
has selected products zero, square
\(-1-A_c^2/\delta\), and positive time.
Every unselected original wall has nonpositive
product, by positivity and acuteness.
Comparison with the ORIGINAL \(p,x\) gives
\[
 x_A\cdot U>A_c-\sqrt\delta,\qquad
             U\cdot n_k\ge0\quad\hbox{for ALL old }k .
\tag{A11U}\label{eq:arithA11helperpressure}
\]
The second conclusion follows directly from
the selected zero and unselected nonpositive
products of the actual projected point.
The correction \(U\) alone is not claimed
to be a hyperbolic point.

Let \(\Pi\) average the first six and last five
coordinates, leaving coordinate seven fixed.
Then \(w\) of the last row of
\eqref{eq:arithA11references} lies in its range.
The correction roots have zero block averages.
For \(Y=U-\Pi U\), direct block variance gives
\[
 \begin{array}{c|c|c}
 &\Pi U&Y^2\\ \hline
 \text{near}&(5g/6+1/12)\Omega_6&
           (2/3)g^2+g/3+1/12\\
 \text{far}&(5g/6)\Omega_6+\Omega_7/10&
           (2/3)g^2+1/15 .
 \end{array}
\]
In particular \(w\cdot Y=0\), \(Y^2<.8\).
For example the far variances of
\(\Omega_5,\Omega_{11}\) are \(5/3,8/5\),
their cross variance is zero, and the old
correction contributes \(-g^2-1/3\);
the near cross variance contributes \(g/3\).
These are FULL residual norms.
The smaller actual time square is
\((35g^2+5g)/6-13/24\).
Using \(.809<g<.81\) gives
\[
 A_c>1.987,\quad\sqrt\delta<.111,\qquad
 w\cdot\Pi U\le.675b+.1a=1.421825
\]
in the far case; in the near case the product
is \(<(91/120)b=1.394575\).
Both actual helper deficits therefore satisfy
\[
 \Delta=A_c-\sqrt\delta-w\cdot\Pi U
                         >.454175>.45 .
\]

Pay \(\mu=.1\).
Writing \(\Pi U=\sum a_i\Omega_i\), define
\[
 W=w+\mu Y,\qquad b_i=\beta_i-\mu a_i .
\]
The near debit at six is \(<.075834\).
The far debits are \(<.0675\) at six and
\(.01\) at seven.
Each remaining positive \(b_i\) exceeds
\(.034\), and no helper index is paid.
Crucially~\eqref{eq:arithA11helperpressure} gives
\[
                  W\cdot n_i=b_i+\mu U\cdot n_i\ge b_i.
\tag{A11price}\label{eq:arithA11individualprice}
\]
This sign retains every helper/foot cross.
It is not replaced by an unsupported
orthogonality assertion.
In the no-helper and center-helper cases set
\(W=w,b_i=\beta_i,\mu=0\).

Let \(S\) contain ALL retained isolated feet
at indices with positive \(b_i\), and set
\[
                  Z=W-\sum_{i\in S}b_i u_i n_i .
\]
Before assigning any floor the helper debit
has been removed from each individual
coefficient.  Combining its pressure with
all strong and actual foot pressures gives
\[
 x_A\cdot Z\ge E+\mu\Delta
                    -\sum_{i\in S}b_i(L_i-\ell).
\]
Since distinct old normals have nonpositive
products, \eqref{eq:arithA11individualprice}
gives, INCLUDING adjacent-foot crosses,
\[
 Z^2\le E+\mu^2Y^2
                    -\sum_{i\in S}b_i^2(2u_i-u_i^2).
\]
For \(0<u\le1/2\), raw loss is
\(<.217/2+1/2=.6085\).
For \(1/2\le u<1\), the circle bound and
\(\sqrt{1+.217^2}<1.024\) give
\[
 L_i-\frac{b_i}{2}(2u-u^2)
 <{.217+1.024\over2}-{3(.034)\over8}
 =.60775<\ell .
\]
One completed square therefore yields
\[
 e_A\ge2x_A\cdot Z-Z^2
                      >E+2\mu\Delta-\mu^2Y^2 .
\]
If no foot occurs, strict paid strong/helper
pressures give the same strictness.
The three cases in~\eqref{eq:arithA11references}
give respectively
\(.652761,.651605\), and
\(.60820725+.2(.45)-.01(.8)=.69020725\).
Each exceeds \(.65\), contradicting the
contrary assumption.  We have proved
\[
                            e_{A_{11}}>13/20 .
\tag{A11Urate}\label{eq:arithA11uniform}
\]
The ordinary-rank-five remainder costs \(>5/14\),
so \(13/20+5/14=141/140>1\).
The typed-seven plus ordinary \(A_3\) alternative
also closes directly: its distinct whole
blocks cost \(>.09\) and \(>.26\).
The actual-\(I+T_4\) case needs the joint
argument that follows.

\subsection{Actual \(I\): six time-denominator obstructions}
\label{app:arithA11Itimes}

Suppose an ACTUAL isolated short \(I\) spectator
has opposite
\[
 c_I=\sqrt h\,p-\sqrt{1+h}\,e_I,\qquad 0<h<1/2 .
\]
The common \(N=12\) local inventory above
was proved before assuming this wall exists.
Its only nonisolated weak exceptions are
exactly the three golden times in
\eqref{eq:arithA11goldtimes} and the three
silver times in~\eqref{eq:arithA11silvertimes}.
We now exclude ALL six by actual cycles.

For completeness write \(A_c^2=M/L\).
Their fields, denominators and numerator
norms are
\[
 \begin{array}{c|c|c|r}
 \text{type}&K&M/L&\mathrm N(M)\\ \hline
 5,\text{ near}&\mathbb Q(\varphi)&(22+49\varphi)/24&-839\\
 6&\mathbb Q(\varphi)&(23+48\varphi)/24&-671\\
 5,\text{ far}&\mathbb Q(\varphi)&(22+45\varphi)/24&-551\\
 (5,4)&\mathbb Q(\sqrt2)&(39+28\sqrt2)/12&-47\\
 (6,5)&\mathbb Q(\sqrt2)&(83+60\sqrt2)/24&-311\\
 (5,6)&\mathbb Q(\sqrt2)&(41+30\sqrt2)/12&-119 .
 \end{array}
\]
The golden numerators follow from
\[
 M_N=i(12-i)-13+(i+2)(12-i)\varphi,\quad
 M_F=i(12-i)-13+i(14-i)\varphi;
\]
the silver formula is
\(A_c^2=D_i/2+C_{ij}/\sqrt2+D_j/4-1\).
The actual fields equal the displayed quadratic
fields by the connected Lorentzian-principal
argument already proved.
The rings of integers are respectively
\(\mathbb Z[\varphi]\), \(\mathbb Z[\sqrt2]\).
In particular this is not a field assumption
in the remaining no-exception branch.

The own two-cycle of \(c_I\) makes
\(v=4h\) integral.
Strict positivity of its two-root Gram at the
other embedding gives \(-4<v'<0\);
physically \(0<v<2\).
Thus
\[
                        0<|\mathrm N(v)|<8 .
\]
Every listed cap is PURE \(A_{11}\), whereas
\(c_I\) has spatial part only in \(I\).
Their product is \(-A_c\sqrt h\).
The scale-invariant actual two-cycle therefore
makes \(A_c^2v=Mv/L\) integral.
Each displayed \(\mathrm N(M)\) is coprime
to \(6\), while \(L\) is \(12\) or \(24\).
No principal-ideal hypothesis is needed:
multiply by the conjugate of \(M\) and use
an integer Bezout relation for
\(\mathrm N(M)\) and \(L\).
It follows that \(v\in L\mathcal O_K\).
Its nonzero norm is then divisible by
\(L^2\ge144\), contradicting the bound eight.
All six types and their reflected orientations
are excluded.
Consequently every paid central site is
strong or has its actual directional foot
\eqref{eq:arithA11foot}, with no restriction
on the actual field in this remaining branch.

\subsection{One actual endpoint and a joint \(A_{11}+I\) square}
\label{app:arithA11Ijoint}

The rank-free actual-\(I\) edge parameters are
\[
 u^2={\sqrt{1+h}-1\over\sqrt{1+h}+1},\qquad
 t=\sqrt{1+4h}-2\sqrt h,\qquad 0<u<1/3 .
\]
The original \(I\) edge gives \(E_I\ge u^2\).
Let \(d=A_dp-v_A-\beta e_I-v_{\rm other}\)
be the ACTUAL opposite on the first \(A_{11}\)
edge.  It differs from \(c_I\), which stays
strictly negative on the entire outgoing
\(A_{11}\) ray.
Acuteness gives
\[
       -d\cdot c_I=A_d\sqrt h-\beta\sqrt{1+h}\ge0 .
\]
If \(\beta=0\), this positive actual-wall
mutual magnitude is at least \(1/2\)
by the native gap.
If \(\beta>0\), its finite contact with the
retained old \(I\) root has \(\beta\ge1/2\),
so \(A_d\sqrt h>1/2\).
Thus \(A_d\sqrt h\ge1/2\) in every white,
own and foreign role.
For the actual inward unit first-edge direction
\(U_1=\Omega_1/\sqrt{11/6}\), the ENTIRE
spatial vector of \(d\) has norm
\(\sqrt{1+A_d^2}\).
Therefore
\[
 \tanh L={A_d\over v_{\rm total}\cdot U_1}
       \ge {A_d\over\sqrt{1+A_d^2}}
       \ge {1\over\sqrt{1+4h}} .
\]
Taking the increasing half-tangent yields
\[
 x_A\cdot U_1\ge t,\qquad
 x_1\ge\alpha t,\quad\alpha=\sqrt{11/12}.
\tag{A11Iend}\label{eq:arithA11Iendpoint}
\]
Only physical unit-normal geometry is used;
no disconnected raw-root scales are identified.

The exact relation between \(u,t\) is
\[
                    (1-u^2)(1-t^2)=8ut .
\]
It implies \(t+(3/2)u>3/4\).
Indeed substituting \(t_0=3/4-3u/2>0\)
into left-minus-right and multiplying by
sixteen gives
\[
 7-60u+149u^2-36u^3+36u^4
       >7-60u+137u^2>0 .
\]
The last quadratic has discriminant \(-236\),
and the relation decreases strictly in
positive \(t\).
Also \(t<1\) and
\(t>\sqrt3-\sqrt2>.3\).
These estimates hold throughout the full
interval \(0<h<1/2\).

Keep the EXACT central heights
\(H_5=H_7=1.805,H_6=1.839\) and
\(E=.652761\).
Put
\[
 b={361\over1000\sqrt2},\quad
 c={.034\over\sqrt2},\quad
 \eta=\alpha t,\quad r=(5b-\eta)/4,
\]
and choose the NEW zero-sum reference
\[
                  w=(\eta,r,r,r,r,c,-c,-b,-b,-b,-b,-b).
\]
Its first-five sum is still \(5b\), so its
central dual heights are the three fixed
\(H_i\), while its first height
\(\sqrt2\eta\) is paid by the ACTUAL
endpoint~\eqref{eq:arithA11Iendpoint}.
The safe rational bounds are
\[
 .95<\alpha<.958,\quad .255<b<.256,\quad
 \eta>b,\quad\eta<.958,\quad
 r>.07925,\quad c<.025 .
\]
All follow by positive squares and \(t>.3\).
Hence \(\eta>r>c>0\), so all dual coefficients
are nonnegative.  The ONLY positive ones are
\[
 \beta_1=(\eta-r)/\sqrt2,\quad
 \beta_5=(r-c)/\sqrt2>217/6000>.034,\quad
 \beta_6=.034,\quad\beta_7=.1635 .
\]
In particular \(\beta_5\) is the coefficient
of this NEW reference, including its full
endpoint/foot cross.

Let \(S\) be ANY set of retained feet among
\(5,6,7\) and form
\(Z=w-\sum_{i\in S}\beta_i s_i n_i\).
Every new central coefficient is at least
\(.034\).
The continuous adjusted-loss bound proved
above therefore applies afresh:
\[
 L_i-\frac{\beta_i}{2}(2s_i-s_i^2)<.61 .
\]
Combining all foot pressures, strict strong
central heights and the actual endpoint gives
\[
 \begin{split}
 x_A\cdot Z&\ge w^2-\sum_{i\in S}\beta_i(L_i-.61),\\
 Z^2&\le w^2-\sum_{i\in S}\beta_i^2(2s_i-s_i^2).
 \end{split}
\]
Every adjacent-foot cross is nonpositive.
The single completed square proves \(e_A>w^2\),
including the empty-foot case by strict
central floors.
At the REFERENCE level, the change from
the constant first-five block is
\((\eta-b,-(\eta-b)/4,\ldots,-(\eta-b)/4)\);
it has zero sum and zero cross with that
old constant block.
Thus exactly
\[
                      w^2=E+(5/4)(\eta-b)^2 .
\]
This identity is not added to an independently
optimized old-foot norm.

Write \(D=\eta-b>0\).
The exact \(I\) barrier and \(\alpha<1\) give
\[
 D+(3/2)u>(3/4)\alpha-b>913/2000 .
\]
Weighted Cauchy gives
\[
 (D+3u/2)^2
       \le\{(5/4)D^2+u^2\}(4/5+9/4)
       ={61\over20}\{(5/4)D^2+u^2\}.
\]
Since the ACTUAL \(I\) energy is counted
ONCE, on the original vector,
\[
 e_{A_{11}}+E_I
 >{652761\over1000000}+{833569\over12200000}
 ={43986266\over61000000}>{18\over25}.
\tag{A11Irate}\label{eq:arithA11Ijoint}
\]
The positive final margin is
\(66266/61000000\).
Every ordinary \(T_4=A_4,B_4,D_4,F_4,H_4\)
costs \(>2/7\), so
\(18/25+2/7=176/175>1\).
This closes the entire actual-\(I+T_4\)
alternative with arbitrary other spectators
and every actual field.

\subsection{The fresh \(N=11\) window and its small-side repairs}
\label{app:arithA10window}

The remaining rank-ten chain cases require a
different finite-white list.
At \(A_{10}\), use zero-sum \(\mathbb R^{11}\),
\[
 D_i={2i(11-i)\over11},\qquad
 C_{ij}={2i(11-j)\over11}\quad(i\le j),
\]
and ONLY sites \(4,5,6,7\).
Here \(D_i=56/11\) or \(60/11\), so
\(9/4<q_i=\sqrt{D_i}<12/5\).
Write an actual cap as
\(c=A_cp-v_A-v_S\), with full spectator
cost \(Q=\|v_S\|^2<1\).
Each positive whole-component contribution
is at least \(1/4\).
The actual reciprocal-opposite law gives the
stronger individual restriction
\[
                         Q_C<A_c^2
\tag{A10foreign}\label{eq:arithA10foreign}
\]
for each contacted whole spectator \(C\).
If \(Q_C>A_c^2\), the original finite
replacement on its contacted node has
positive Schur residual \(Q_C-A_c^2\).
Positive finite projection checks EVERY
original wall and would make \(c\) close
two orthogonal components at the original
farthest vertex.
Equality instead gives the positive-cleared
future null.
Neither is possible.
The cleared projection is never treated as
an additional native facet normal.

Take loss \(\ell=13/20=.65\).
The sufficient ratio
\[
 T(q)={20\over13}
       -{q-13/20\over q(2q-13/20)}
\]
increases on the stated \(q\)-interval.
Indeed its derivative has positive numerator
\(2q^2-(13/5)q+169/400\).
The exact upper comparison is
\[
 T(q)<{17645\over12948}<{1363\over1000},
\]
with cleared positive gap \(3124\) in the
last cross product.
Thus ratio at least \(1.363\) is strong.

Deleting a selected site leaves path ORDERS
in \(\{4,5,6,7\}\).
The COMPLETE finite one-node extensions are:
isolation; one admissible endpoint/\(D/E\)
white-three; terminal white-four; terminal
white-five on an \(A_3\), giving \(H_4\);
one endpoint-three in each path, giving
\(A_{10}\); or, at site four or its
reflection seven, a middle-\(A_3\)/end-three
join, giving \(D_{10}\).
This last role is a genuine new small-side
case.  Two contacts to one retained path
make a cycle.  Rank ten precludes exceptional
joined completions; a weighted \(B_{10}\)
join would need a rank-one end side, which
is absent.  An \(A_3\) has no missing
one-node \(F_4\) extension with these
unweighted old bonds.

At own zero, isolation has \(A_c^2=Q-1<0\).
Internal three has the affine \(E_8\) arms
\(1,2,5\), or \(E_7\) arms \(1,3,3\);
internal four has the weighted \(4,3,3\)
star; either two-three join has its actual
all-three cycle.
The positive \(E_7\) kernel has central
coefficient four, short arm two and the
two other arms \(3,2,1\).
The \(E_8\) kernel has center six and
arms \(3\), \(4,2\), \(5,4,3,2,1\).
The weighted star has center \(\sqrt2\),
four-leaf one and three-leaves \(1/\sqrt2\);
cycle coefficients are all one.
Every kernel is future through the positive
cap time and nonpositive against all
unselected ORIGINAL walls.

The residual single-contact rows are
\[
 \begin{array}{c|c}
 A_c^2&\text{strict finite-white bound}\\ \hline
 Q-6/11&Q<5/8\\
 Q-2/11&Q<1/2\\
 Q-1/11&Q<1/4 .
 \end{array}
\]
In the first, positive time is \(<7/88<1/4\),
contradicting~\eqref{eq:arithA10foreign}.
The second permits at most one whole
foreign contribution, larger than its time.
The third has \(Q=0\) and negative time.
The white-five/own-zero case will instead
be directly strong below.

Isolated owns three, four and six have
the actual \(E_7\), weighted-star or
\(6\)--\(3\) null, BEFORE splitting \(Q\).
A parallel actual pair has its positive
future null in every profile.
For SINGLE white-three, own-three closes
the actual all-three cycle; own-six uses
an old neighbor different from the helper.
Own-four has its star unless that helper
is adjacent to the own site.
The adjacent near-four role is retained
in the next calculation.
All other positive owns have
\(\gamma\ge\varphi/2>4/5\).

For a white-three at position \(t\) in a
side of order \(l\), discarding foreign
cost only decreases the ratio, giving
\[
                  {\gamma+t/(2l)\over
                            \sqrt{1-t(l-t)/(2l)}} .
\]
Near-three exceeds
\((4/5+3/8)/\sqrt{5/8}>1.363\).
The \(D\)-position exceeds
\((4/5+1/4)/\sqrt{1/2}>1.363\);
its formula increases with \(l\ge4\).
Allowed \(E\)-positions have denominator
at most \(1/2\), hence ratio \(>1.6\).
Near own-four/simple-three has minimum
at \(l=4\), with
\[
                     R_4^2=41/40+3\sqrt2/5
                             >1.8734>1.363^2 .
\]
Its numerator increases and denominator
decreases with \(l\), so ALL these roles
are strong without using a quadratic-field
exclusion.

White-four has ratio \(\ge\gamma\sqrt l\).
For own at least four this is \(\ge\sqrt2\).
Near white-four/own-three already has
ratio \(>2\) at \(l=4\).
Far white-four/own-three has
\(\sqrt l/2+1/\sqrt{2l}\), increasing for
\(l\ge4\).
At \(l=5\) its square is
\(5/4+1/\sqrt2+1/10>1.363^2\).
Only the far role at \(i=4\), or its
reflection seven, remains.
It has \(Q=0\) and
\[
 v_A=\Omega_4/2+\Omega_1/\sqrt2,\qquad
                   A_c^2=(13+7\sqrt2)/11 .
\]
The actual own-three normalization and
closed cycle put \(\sqrt2\) in the field.
At any embedding extending its sign reversal,
the connected principal still has positive
time square \((13-7\sqrt2)/11>0\).
It therefore retains a negative direction,
contradicting strict nonphysical positivity.
This argument needs only the contained
quadratic field, not equality with it.
It is the required small-side repair;
the unavailable longer \(E_7\) arm is not
used to discard this role.

For white-five on \(A_3\), put
\(\beta=\varphi/2>4/5\).
Its white residual is
\(\delta_H=(7-3\sqrt5)/16<1/50\).
Even own zero has ratio
\(\beta/(4\sqrt{\delta_H})\), whose square
exceeds two.  All such profiles are strong.
The \(A_{10}\) two-end-three join has ratio
\(\ge\gamma q_i+1/q_i>113/72\) for
positive own.
The \(D_{10}\) join has white cost
\(1/2+3/7=13/14\), so \(Q=0\).
Its two offsets are \(19/28,9/28\);
the positive-own ratio is at least
\(23\sqrt{14}/28>3\).
All their own-zero cycles were handled above.

Finally a far-three helper WITH foreign cost
has residual at most \(3/8\) and numerator
\(>4/5+1/14=61/70\).
Its ratio square exceeds
\(29768/14700>1.363^2\).
Every possible weak NONISOLATED cap is
therefore a PURE end-helper
\[
              v_A=\gamma\Omega_i+\Omega_e/2,
               \qquad e=1\ \hbox{or }10 .
\tag{A10helper}\label{eq:arithA10helper}
\]
This conclusion is independent of both
the actual-\(I\) and typed-seven hypotheses.

\subsection{A small conjugate polynomial replaces the degree census}
\label{app:arithA10helperfield}

For~\eqref{eq:arithA10helper}, the exact
side-order square identity gives
\(4\gamma-2\le R^2\).
Weakness and \(T(q)<1.363<1.37\) imply
\[
             \varphi\le t=2\gamma
                     <38769/20000<39/20 .
\]
Actual simple-three normalization makes
\(t\) integral.
The old roots and cap have a common
totally positive scale times the norm-two
matrix with helper product \(-1\) and
own product \(-t\).
This full-rank connected principal is
physically Lorentzian.
Hence a nonphysical embedding fixing
\(\mathbb Q(t)\) would preserve a negative
direction: ACTUAL \(K=\mathbb Q(t)\).

At every other embedding the exact Schur
inequality is
\[
                  D_i t_\sigma^2+2C_{ie}t_\sigma+20/11<4 .
\]
At \(i=4\), its cheapest signed costs are
\((56s^2+28s+20)/11\) for negative \(s\),
and \((56s^2+16s+20)/11\) for positive \(s\).
At five the increases are
\(4(s^2-s)/11\), \(4(s^2+s)/11\), respectively.
They are positive in the stated half-windows.
Boundary excesses \(456/6875\) at \(-24/25\)
and \(263/13750\) at \(53/100\) imply
\[
                   -24/25<t_\sigma<53/100 .
\]
Let \(f(T)=T(T^2-1)\).
The physical \(f(t)\) is nonzero integral,
with \(0<f(t)<43719/8000\).
On the ENTIRE small-conjugate interval,
\[
 |f(s)|\le {2\over3\sqrt3}<{77\over200};
\]
the last comparison squares to
\(160000<160083\).
Degree at least three would give a nonzero
integral norm of magnitude
\[
 <{43719\over8000}\left({77\over200}\right)^2
 ={259209951\over320000000}<1 ,
\]
and every further factor is smaller than one.
The degree is at most two.
Degree one has no integer in
\([\varphi,39/20)\).
In degree two the nonzero norm has magnitude
\(<1.95(.96)<2\), hence is \(1\) or \(-1\).
A positive small conjugate gives integral
trace \(t+1/t\in(2,3)\), impossible.
A negative one gives trace
\(t-1/t\in[1,2)\), hence exactly one.
Thus
\[
                 t=\varphi,\qquad K=\mathbb Q(\varphi).
\tag{A10gold}\label{eq:arithA10gold}
\]
This is a complete all-degree argument.
No field census or assumed cubic bound is used.

\subsection{Every isolated \(A_{10}\) cap is strong}
\label{app:arithA10isolated}

This part also needs neither spectator
hypothesis.  Keep
\(u=(1-Q)/\gamma^2>0\).
The small-own/null cases were excluded for
all \(Q\), so \(\gamma\ge\varphi/2\).
For \(0<u<1\), the positive ACTUAL vector
\[
        \xi=c+v_S+{1-Q\over\gamma}n_i
\]
has square \(-(1-Q)(1-u)<0\), positive
time, zero selected products and negative
product on \(n_i\).
The entire spectator correction is positive,
so every original omitted wall is nonpositive.
The original farthest comparison gives
\[
 h_i>\sqrt{D_i-u}-\sqrt{u(1-u)} .
\]
Here
\[
 a_i={1\over q_i+\sqrt{D_i-1}}<{100\over427}<{47\over200},
\]
since \(q_i>9/4\) and
\(\sqrt{D_i-1}>101/50\).
The full continuous loss is
\[
 <{.235+\sqrt{1+.235^2}\over2}
 <{.235+1.028\over2}=.6315<.65 .
\]
Equality \(u=1\) is the feasible future null.
For \(Q>0,u>1\), positive finite clearing
has loss \(<a_i+\sqrt{u-1}\).
Since \(Q\ge1/4\) and \(\gamma\ge\varphi/2\),
\(u\le3/\varphi^2\), hence
\(\sqrt{u-1}\le2-\varphi<.4\).
The loss is \(<.635<.65\).

For PURE native five the actual finite
\(H_4\) point has loss
\[
 <{25/16\over q_i+\sqrt{D_i-25/16}}+5/24
              <625/1648+5/24<.588 .
\]
The denominator is \(>2.25+1.87=4.12\).
It is the positive finite projected point,
not its spacelike correction, that checks
all original walls.
For EVERY finite native at least eight
the finite two-root projected point gives,
with \(B=\sqrt{D_i-1}\), \(v=\sqrt{u-1}\),
\[
 h_i>{\sqrt{B^2-v^2}\over1+v/B}>B-v .
\]
Its bracket is positive, and
\(v\le\tan(\pi/8)=\sqrt2-1<.415\).
Thus the full loss is \(<.235+.415=.65\),
covering the entire infinite native tail.

For native seven, the elementary cubic
\(8\gamma^3-4\gamma^2-4\gamma+1\) and
\(\pi<22/7\) give \(.9<\gamma<10/11\).
Consequently \(121/100<u<100/81<4/3\).
For an adjacent old root \(n_j\), the
positive actual core
\[
             \xi=c+\gamma un_i+2\gamma(u-1)n_j
\]
is zero on \(c,n_i\), negative on \(n_j\),
and has square
\(-\gamma^2(7u-4-3u^2)<0\).
It is future and nonpositive against all
other ORIGINAL walls.
Its pressure gives
\[
 h_i>\sqrt{D_i-u}-\sqrt{7u-4-3u^2}.
\]
The first loss is \(<.294\), using
\((100/81)/4.21<.294\),
\(q_i>2.25\), and
\(\sqrt{D_i-100/81}>1.96\).
The second radicand decreases for \(u>7/6\);
at \(121/100\) it is \(.0777<.28^2\).
Thus the total loss is \(<.574<.65\).
All isolated finite, foreign and dotted
profiles are strong; no foot is left unpaid.

\subsection{The actual-\(I\) and typed-seven consumers of the same window}
\label{app:arithA10consumers}

In the actual-\(I\) case, only the golden
helpers~\eqref{eq:arithA10gold} still need
to be excluded.
The fresh COMPLETE time table, with reflected
indices covering six and seven, is
\[
 \begin{array}{c|c|c|r}
 i&e&11A_c^2&\mathrm N(11A_c^2)\\ \hline
 4&1&8+21\varphi&-209\\
 4&10&8+18\varphi&-116\\
 5&1&9+21\varphi&-171\\
 5&10&9+20\varphi&-139 .
 \end{array}
\]
These are direct substitutions in
\[
 A_c^2=(\varphi^2/4)D_i+(\varphi/2)C_{ie}+D_e/4-1,
                        \qquad D_e=20/11 .
\]
The first numerator norm has exactly one
factor eleven; none of the others has one.
For the ACTUAL isolated \(I\), its own
two-cycle gives \(V=4h\in\mathcal O_K\),
\(0<V<2\), and \(-4<V'<0\).
We are now in the actual quadratic field
\(K=\mathbb Q(\varphi)\), so
\(0<|\mathrm N(V)|<8\).
The helper and \(I\) cap have disjoint
spatial supports.
Their actual two-cycle makes \(A_c^2V\)
integral, hence
\(\mathrm N(A_c^2)\mathrm N(V)\) integral.
Every displayed row forces eleven to divide
\(\mathrm N(V)\): its denominators are
eleven or eleven squared, with numerators
coprime to eleven after cancelling the one
factor in the first row.
This contradicts the nonzero norm bound eight.
No relative primitive scale, PID assertion
or extra \(I\)-energy charge is involved.

For EITHER actual typed-seven defect, the
same generic local geometry and all-degree
helper proof remain valid.
The previously proved ACTUAL \([7,3]\)
principal instead gives \(K=K_7\), a cubic
field with no quadratic subfield.
Thus~\eqref{eq:arithA10gold} is impossible
in this branch as well.
This reuses the complete generic window;
it does not transport the \(N=12\) inventory
or infer arithmeticity of an arbitrary face.
It also avoids a duplicated cubic-only
trace argument for the same already excluded
helper domain.

In both cases every site \(4,\ldots,7\)
is genuinely \(.65\)-strong.
Take symmetric heights, linear on the tails,
\[
                  H_4=H_7=8/5,\qquad H_5=H_6=42/25 .
\]
They lie strictly below the proved floors,
because \((9/4)^2<56/11\) and
\((233/100)^2<60/11\), with positive
margins \(5/176,2821/110000\).
The only positive dual coefficients are
\(\beta_4=\beta_7=4/25\),
\(\beta_5=\beta_6=1/25\).
Every one is paid by an actual strong height.
One square on the original vector yields
\[
                       e_{A_{10}}>
                   (8/5)^2/4+(42/25-8/5)^2
                   ={404\over625}.
\tag{A10rate}\label{eq:arithA10conditionalrate}
\]
This rate has the stated actual-\(I\) OR
actual typed-seven hypothesis; it is not
asserted for a bare arbitrary-field \(A_{10}\).

For the actual-\(I+(A_5\hbox{ or }D_5)\)
alternative, add only the ordinary whole
rank-five component:
\(404/625+5/14=8781/8750>1\).
The positive \(I\) energy is not needed.
For actual typed-seven plus ordinary \(T_4\),
count the defect once and use
\[
                  404/625+9/100+2/7
                         ={17887\over17500}>1 .
\]
These are sums of disjoint whole-component
energies of the SAME original \(x\).
They exclude both entire rank-sixteen
\(A_{10}\) alternatives, without a field
census or a newly executed finite calculation.

\subsection{The two actual typed-seven \(A_{14}\) forests}
\label{app:arithA14typed}

Here the WHOLE incident forest is \(A_{14}\)
plus one actual rank-two defect, with no
further spectators.
Put \(c=\cos(\pi/7)\), \(\theta=2c\).
The actual \([7,3]\) principal proves
\(K=K_7=\mathbb Q(\theta)\).
Its full integer ring and root brackets
are those proved in
Section~\ref{app:arithA13seven}.
The defect wall is isolated from the ENTIRE
old \(A_{14}\), and has the form
\[
 \begin{array}{c|c|c}
 &d=\sqrt{h_D}\,p-v_D&h_D\\ \hline
 N_7&v_D=c\,\Omega_e^{A_2}&(4/3)c^2-1\\
 O_7&v_D=\Omega_e^{I_2(7)}/2&
                       1/(4(1-c^2))-1 .
 \end{array}
\]
Both defects have disjoint energy \(>.09\).
We first strengthen the central floors,
then use this ACTUAL wall to pay both ends.

At \(N=15\), use sites \(5,\ldots,10\)
and fresh loss \(\ell=.55\).
Their dual squares are \(20/3,36/5,112/15\),
in reflected order, with \(5/2<q<11/4\).
The increasing sufficient ratio obeys
\[
                R_\ell(q)<R_\ell(11/4)=164/99<5/3 .
\]
The full finite-white list remains isolation,
admissible endpoint/\(D/E\) three, endpoint
four, or two endpoint threes joining paths.
The retained paths have ranks at least four.
The target-independent own-zero residuals
are \(Q-8/15,Q-2/15,Q-1/15\), with bounds
\(Q<3/5,2/5,1/5\).
The first positive time is below any whole
foreign contribution; the second permits
at most one contribution greater than its
time; the last permits none.
The actual reciprocal law excludes those
positive cases; at \(Q=0\) times are negative.
Internal-three/four and two-three roles
retain their all-wall positive affine kernels.

The actual field \(K_7\) has no quadratic
subfield.  Native five is absent by its
two-cycle; own-four/simple-three is absent
by its UNSQUARED connected \(\sqrt2\) cycle.
All appropriate smaller owns and parallel
retain their all-foreign actual nulls.
For the other ordinary roles, the FRESH
lower ratios are
\[
 \begin{array}{c|c}
 \text{role}&\text{lower ratio}\\ \hline
 \text{near-three}&13/\sqrt{60}>5/3\\
 D\text{-position}&11/\sqrt{40}>5/3\\
 E\text{-position}&1.8\\
 \text{far-four}&.9\sqrt5\\
 \text{near white-four}&53\sqrt5/50\\
 \text{two end-threes}&q/2+1/q\ge13/\sqrt{60}>5/3 .
 \end{array}
\]
These use \(\gamma>.9\), side order
\(m\ge5\), white-three residual
\(1-j(m-j)/(2m)\) and offset \(j/(2m)\);
white-four has residual \(1/m\), offset
\(j/(\sqrt2m)\).
The last comparison uses the ACTUAL minimum
\(q^2=20/3\), not merely \(q>5/2\).
Positive foreign cost strengthens all
these ratios.
For a foreign far-three helper complete
clearing gives
\[
        {1.8m+1\over\sqrt{m(m+2)}}\ge10/\sqrt{35}>5/3 ,
\]
increasing for \(m\ge5\).
Thus a weak nonisolated cap is PURE
end-helper with \(\gamma>.9\).

Put \(t=2\gamma\).
Its exact square-ratio identity gives
\[
                      9/5<t<23249/9801<19/8 .
\]
It is integral in the ACTUAL cubic field.
The signed Schur costs are
\[
 2it^2+4t+2-2(it+1)^2/15,\qquad
 2it^2+2-2(it-1)^2/15 .
\]
At \(i=5\), their values at \(-4/5\),
\(12/25\) are four and \(4+16/375\).
They increase across the half-window,
so both other conjugates lie in
\((-.8,.48)\).
In degree three the nonzero norm has
magnitude \(<(19/8)(.8)^2<2\).
A positive small conjugate would make it
\(<(19/8)(.48)(.8)=.912<1\).
Both are negative with absolute product
\(1/t\); the integral trace is one,
being \(>t-1.6>.2\) and
\(\le t-2/\sqrt t<2\).
Its second coefficient is
\[
 B=t-t^2+1/t,\qquad
                    -3459/1216<B<-199/225 .
\]
The function decreases, so \(B=-2\) or
\(-1\).
The candidate cubics
\(X^3-X^2-2X-1\), \(X^3-X^2-X-1\)
have discriminants \(-31,-44\), contradicting
total reality.
Degree one leaves only \(t=2\), its
actual parallel null.
No weak helper remains.

We now prove directly that no
\(w\in\mathcal O_{K_7}\) satisfies
\[
                    0<w<11,\qquad0<w_\alpha,w_\beta<3/5 .
\tag{A14strip}\label{eq:arithA14typedstrip}
\]
Use the proved FULL basis
\(w=a+b\theta+c_0\theta^2\), with roots
\[
 1.8019<\theta<1.802,\quad .445<\alpha<.446,\quad
                            -1.248<\beta<-1.246 .
\]
Set \(d_0=b-(\theta-1)c_0\) and
\(\Delta=w-(w_\alpha+w_\beta)/2\).
The exact identity is
\[
 \Delta=A_0d_0+B_0c_0,\quad
 A_0=(3\theta-1)/2<2.203,\quad
 B_0=3\theta^2-2\theta-2>4.136 .
\]
The small-value difference gives
\(|d_0|<.6/(\alpha-\beta)<.355\),
hence \(|A_0d_0|<.782065\), while
\(-.6<\Delta<11\).
Values \(c_0\le-1\) force
\(\Delta<-3.353935\);
\(c_0\ge3\) forces \(\Delta>11.625935\).
Only \(c_0=0,1,2\) remain.
At zero, \(b=0\), then \(a\) would be an
integer in \((0,.6)\).
At two, \(1.2488<b<1.959\), impossible.
At one, \(.4469<b<1.157\), so \(b=1\).
But \(.643025<\alpha+\alpha^2<.644916\);
positivity forces \(a\ge0\), contradicting
\(w_\alpha<.6\).
This proves the entire strip.

For ANY isolated-within-\(A\) cap retain
its entire foreign cost and put \(w=4\gamma^2\).
The actual own two-cycle makes \(w\)
nonzero integral.
Normalize its original root to own raw
product \(-1\), square \(2/w\).
The old inverse entry \(D_i/2\) gives
\[
                     0<w_\sigma<4/D_i\le3/5
\]
at BOTH other embeddings, even with
foreign contacts.
Weakness would imply
\[
           0<w<4(164/99)^2=107584/9801<11,
\]
contradicting~\eqref{eq:arithA14typedstrip}.
Every site \(5,\ldots,10\) is therefore
genuinely \(.55\)-strong, with no helper
or weak foot left to omit.

\subsection{The actual typed wall forces both \(A_{14}\) ends}
\label{app:arithA14typedends}

The increasing cubic \(8c^3-4c^2-4c+1\)
has values \(-.000569738168,.000257608\)
at \(.9009,.901\).
Thus
\[
 .9009<c<.901,\quad
 .08<h_N<.083,\quad .32<h_O<.3284=821/2500 .
\]
Substitution into the displayed height
formulas gives these rational seals.
The shared exact gap is
\[
 \psi={15(3c^2-2)\over28(4c^2-3)},\qquad
                         c^2<\psi<(1+c)/2<1 .
\tag{A14gap}\label{eq:arithA14typedgap}
\]
For the lower bound,
\(\psi=(1+1/(4h_O))/(28/15)>21690/22988>.9>c^2\).
For the upper, clearing positive denominators
and using the cubic reduces to
\(39c^2-14c-19>0\), increasing for \(c>.9\)
and equal to \(.04061159>0\) at \(.9009\).
In an \(N_7\) branch, \(h_O\) in this
identity is only its scalar function of \(c\).

The COMPLETE finite actual own-square
alphabet in \(K_7\) is
\[
                 \gamma^2\in\{0,1/4,1/2,3/4,c^2,(1+c)/2\}.
\]
Cyclotomic degree first permits
\(2,3,4,6,7,9,14,18\).
The \(9/18\) cubic field differs:
\(2\cos(\pi/9)\) has polynomial
\(X^3-3X-1\), whose shift by one is
Eisenstein at three, whereas \(K_7\)
has discriminant \(49\).
If \(\gamma\) itself is \(K_7\)-valued,
its finite list further reduces to
\(0,1/2,c\); \(\cos(\pi/14)\) has degree six.

For EITHER actual endpoint cap
\(c_A=A_cp-v_A-v_D'\), let \(\delta>0\)
be its Schur residual against ALL retained
old roots.  The actual edge gives
\[
 C_{\rm end}^2=1+A_c^2/\delta,\quad0<\delta\le1,\qquad
 y\ge\sqrt{14/15}\sqrt{{C_{\rm end}-1\over C_{\rm end}+1}} .
\tag{A14end}\label{eq:arithA14endpointlaw}
\]
The two signed coordinates are \(x_{A,1}\),
\(-x_{A,15}\).
This cap differs from the actual defect
wall \(d\), which cannot close its edge.
The two endpoint caps need not differ
from each other.

For \(O_7\), positive retained
\(I_2(7)\) contact would cost at least
\(\csc^2(\pi/7)/4>1\), impossible.
Thus \(v_D'=0\), and the positive ACTUAL
mutual magnitude
\(\lambda=-c_A\cdot d=A_c\sqrt{h_O}\)
is at least \(1/2\).
We exclude equality in every endpoint role.
An isolated cap has
\(A_c^2=(28/15)\gamma^2-1\);
equality forces \(\gamma^2=\psi\).
The finite alphabet, parallel one and
dotted \(\gamma>1\) all avoid
\eqref{eq:arithA14typedgap}.

For a nonisolated endpoint, retained
\(A_{13}\) has the COMPLETE finite list:
helper-three at \(j=2,3,13,14\), or
helper-four at \(j=2,14\).
The exact fresh time polynomials are
\[
 \begin{array}{c|c}
 (j,\text{helper label})&15A_c^2\\ \hline
 (2,3)&28\gamma^2+26\gamma-2\\
 (3,3)&28\gamma^2+24\gamma+3\\
 (13,3)&28\gamma^2+4\gamma-2\\
 (14,3)&28\gamma^2+2\gamma-8\\
 (2,4)&28\gamma^2+26\sqrt2\gamma+11\\
 (14,4)&28\gamma^2+2\sqrt2\gamma-1 .
 \end{array}
\tag{A14endtimes}\label{eq:arithA14endpointtimes}
\]
Simple-three normalization makes
\(\gamma\) itself \(K_7\)-valued.
At own zero only \(j=3\) is future,
with time \(1/5\), giving \(\lambda^2<1/4\).
At own \(1/2\) the four times are
\(6/5,22/15,7/15,0\).
The first two have mutual square \(>1/4\),
the third \(<1/4\), and the last is not future.
For every \(\gamma\ge c>.9\) all four
polynomials increase and their minimum
exceeds \(16.48/15>1\).
Thus their mutual squares exceed \(1/4\).
This covers every remaining finite and
every dotted own; parallel has its actual null.

For helper-four, connected normalization
gives \(2\sqrt2\gamma\in K_7\).
At finite own labels \(3,6,7,14\), the
norms of \(2w\), \(w=4\gamma^2\), are
\(8,216,8,56\), all nonsquares.
Only own four remains, besides zero,
parallel or dotted.
At zero, near time \(11/15\) has mutual
square \(<1/4\), and far time is negative.
At own four, times \(17/5,1\) have mutual
squares \(>1/4\).
At \(\gamma\ge1\), both times exceed one.
Hence equality \(\lambda=1/2\) is impossible
in EVERY role.
The next positive native magnitude is
at least \(1/\sqrt2\), so
\[
 C_{\rm end}^2\ge1+1/(2h_O)>2071/821>(397/250)^2 .
\]
Equation~\eqref{eq:arithA14endpointlaw} gives
\[
                       y^2>{2058\over9705}>(23/50)^2,\qquad y>.46 .
\]
The two cleared rational margins are
\(40511,11055\).
Only connected products within \(A_{14}\)
were used to normalize helpers; no equal
scale of its disconnected defect is assumed.

For \(N_7\), retain the ENTIRE foreign
old \(A_2\) block and write
\[
 \lambda=A_c\sqrt{h_N}-\rho\ge0,\qquad
                              \rho=v_D'\cdot v_D .
\]
At \(Q=0\), the native gap gives
\(A_c^2>250/83\), hence \(C_{\rm end}>2\).
Positive foreign contact has the COMPLETE
finite list
\[
 \begin{array}{c|c|c}
 \text{label}&Q&\rho\\ \hline
 3&1/3&c/3\ \hbox{or }2c/3\\
 4&2/3&\sqrt2c/3\ \hbox{or }2\sqrt2c/3 .
 \end{array}
\]
Two contacts make a cycle; six is affine;
larger labels are nonfinite, and five is
forbidden in \(K_7\).
If an \(A_{13}\) contact is also present,
its full inverse cost is at least \(13/28\).
Thus \(\delta\le17/84\).
Acuteness gives \(A_c\sqrt{h_N}\ge\rho>.3\),
so \(A_c^2>90/83\) and
\(A_c^2/\delta>7560/1411>3\).
Without an \(A_{13}\) contact, \(\delta=1-Q\).
At \(Q=2/3\), \(\rho>.3\sqrt2\) gives
\(A_c^2/\delta>540/83>3\).
At \(Q=1/3,\rho=2c/3\), the same bound
follows from \(\rho>.6\).
These arguments retain EVERY own parameter.

Only \(Q=1/3,\rho=c/3\), with no
\(A_{13}\) contact, remains.
Positive \(\lambda\) gives
\(A_c\sqrt{h_N}\ge c/3+1/2>.8\)
and \(A_c^2/\delta>960/83>3\).
At \(\lambda=0\), the exact equation
\(A_c^2=(28/15)\gamma^2-2/3\) instead gives
\(\gamma^2=\psi\).
Its ACTUAL own two-cycle puts
\(4\gamma^2\) in \(K_7\);
the finite alphabet, parallel and dotted
ranges again avoid the gap.
Thus every \(N_7\) endpoint has \(C_{\rm end}>2\),
giving \(y>\sqrt{14/45}>.55>.46\).
The two strict \(.46\) signed demands
therefore hold for BOTH actual defect types.

\subsection{Both typed endpoints in one legal reference}
\label{app:arithA14typedcomposition}

All central sites are \(.55\)-strong.
Choose symmetric heights, linear on the tails,
\[
 H_5=H_{10}=a=2.03198,\quad
 H_6=H_9=b=2.13328,\quad
 H_7=H_8=c_0=2.18251 .
\]
The floor tests are
\(2.58198^2<20/3\),
\(2.68328^2<36/5\),
\(2.73251^2<112/15\).
The six central coefficients are
\(.152548,.026035,.024615,.024615,.026035,.152548\).
The base norm is
\[
                   E=a^2/5+(b-a)^2+(c_0-b)^2=.83847382698 .
\]
The first and last five coordinates are
flat with signed mean \(b_0=a/(5\sqrt2)\),
\(.28<b_0<.2875\).
In EACH block replace the signed endpoint
by \(.46\), and redistribute its remaining
sum constantly over the other four coordinates.
Their value \(r=(5b_0-.46)/4\) lies in
\((.235,.245)\), below the endpoint and
above the next old coordinate \(<.073\).
The NEW reference stays concave, with
positive endpoint and changed boundary
coefficients.  Every central dual height
is unchanged and genuinely paid.
There are no omitted weak feet or helpers.

The two residuals have disjoint block
supports, zero block sums and zero cross
with the original constant blocks.
Each reference gain is exactly
\((5/4)(.46-b_0)^2>.0371953125\).
All actual paid pressures are strict.
One square therefore gives
\[
                    e_{A_{14}}>
                  .83847382698+2(.0371953125)=.91286445198 .
\]
Only NOW add the disjoint defect energy
\(>.09\) of the SAME original vector,
counted ONCE.  The full norm exceeds
\(1.00286445198\), a contradiction.
Both actual typed-seven \(A_{14}\) forests
are excluded.
No \(K_7\) condition is imposed on the
separate ordinary rank-two branch.

\subsection{The ordinary \(A_{14}+T_2\) branch has one actual three-wall core}
\label{app:arithA14ordinarycore}

Now the WHOLE forest is \(A_{14}+S\), where
\(S=I_2(m)\), \(m\ge3\), is ORDINARY:
neither actual typed defect is present.
The universal bound
\eqref{eq:arithA14universal} gives
\[
                         0<e_S<1/5 .
\]
Let \(k=\cos(\pi/m)\ge1/2\),
\(d_0=(1-k^2)^{-1}\).
The old unit Gram inverse is
\(d_0\left(\begin{smallmatrix}1&k\\k&1\end{smallmatrix}\right)\).
The positive unit incident duals \(u_1,u_2\)
have correlation \(k\).
For their ACTUAL edges put
\(\tau_i=\tanh(L_i/2)\), \(r_i=\tanh L_i\).
Original farthest comparison gives
\(x\cdot u_i\ge\tau_i>0\);
hence \(\tau_i^2<1/5\) and each
actual edge cosine is \(<3/2\).

The two edges have ONE actual opposite.
Indeed Cauchy on \(u_1+u_2\) gives
\(\tau_1+\tau_2\le\sqrt{2e_S(1+k)}\),
so \(\tau_1\tau_2<(1+k)/10\).
For \(0<a,b<1\),
\[
 {2a\over1+a^2}{2b\over1+b^2}
          \le {4ab\over(1+ab)^2}.
\]
The right side increases in \(ab<1\);
therefore
\[
                         r_1r_2<{40(1+k)\over(11+k)^2}<k .
\]
The last cleared polynomial is
\(k^3+22k^2+81k-40\), increasing for
\(k\ge1/2\) and equal to \(49/8>0\)
at \(1/2\).
The actual two-wall law for distinct
opposites requires \(r_1r_2\ge k\).
Thus one wall \(d\) closes BOTH directions,
including the old \(A_2\) case.

Write
\[
 d=A_dp-v_S-v_A,\quad h=A_d^2>0,\quad
 b=-d\cdot n_1,\quad g=-d\cdot n_2,\quad Q=\|v_A\|^2,
\]
so
\[
             h=d_0(b^2+g^2+2kbg)+Q-1 .
\]
Every positive \(b,g\) is finite native,
being in a retained positive two-root Gram
at the other actual endpoint.
The whole \(A_{14}\) is retained there,
so \(Q<1\); a positive \(Q\) is at least
\(7/15\).
Both \(b=g=0\) would give negative time.
If both are positive and \(Q>0\), the
actual edge formula gives
\[
 C_1^2={d_0(b+kg)^2\over1-g^2-Q}
       \ge{(4/3)(3/4)^2\over3/4-7/15}
       ={45\over17}>{9\over4},
\]
contrary to \(C_1<3/2\).

If precisely one contact is positive,
write \(b=\beta,g=0\).
At the other endpoint the connected finite
rank-sixteen diagram is classical.
Its only possible foreign attachment is
a SINGLE old \(A_{14}\) END-three, with
\(\beta=1/2\) or \(1/\sqrt2\).
An internal attachment would give arms
of lengths two and two others totaling
thirteen, which are neither \(D_{16}\)
arms nor exceptional finite arms.
An end-four would become an internal
weighted bond; a five bond cannot occur.
Thus \(Q=7/15\) and
\(h=d_0\beta^2-8/15\).
The smaller normalized derivative square is
\(k^2d_0\beta^2\).
Since \(r_i^2<5/9\), we get
\[
                          \beta^2(4d_0+5)\le24/5 .
\]
At \(\beta=1/\sqrt2\), this forces
\(d_0\le23/20<4/3\), impossible.
At \(\beta=1/2\), \(d_0\le71/20<4\),
so \(m=3,4,5\).
The first two times are \(-1/5,-1/30\).
At five, \(d_0=2+2/\sqrt5<3\), giving
\(0<h<13/60<1/4\).
The actual small-height foreign-isolation
lemma then forbids this remaining contact:
the same wall would close an \(A_{14}\)
direction in addition to its \(S\) directions.
Consequently
\[
                         Q=0 .
\tag{A14coreiso}\label{eq:arithA14spectatorisolated}
\]
This is isolation of an ACTUAL wall from
the entire old component, not face heredity.

We classify every native \(S\)-contact.
With one contact the same derivative law
now gives \(\beta^2(4d_0+5)\le9\).
Thus \(d_0\le31/4\).
For \(m\ge9\), \(\sin(\pi/m)<\pi/m\) and
\(\pi^2<10\) give \(d_0>81/10>31/4\).
Hence \(m=3,\ldots,8\) is complete.
If \(\beta\ge1/\sqrt2\), then
\(d_0\le13/4\), so \(m\le5\).
At four, positive time and the inequality
leave only new five.
At five, new three has negative time,
new four remains, and new five or larger
violates \(\beta^2(4d_0+5)\le9\).
At six only new three remains, with
zero time.
At seven/eight only new three remains;
old seven/new three is ACTUAL \(O_7\)
and excluded by ordinary meaning.
At old three, positivity and the inequality
give \(3/4<\beta^2\le27/31\).
Only new seven or eight remains:
for all native nine or larger,
\(\cos(\pi/9)>76/81\) and
\((76/81)^2>27/31\).
Old three/new seven is ACTUAL \(N_7\)
and is excluded.
Thus precisely four one-contact rows survive.

With two positive contacts, every old
\(m\ge5\) has
\[
                   C_i\ge{\cot(\pi/(2m))\over\sqrt3}
                       \ge{\cot(\pi/10)\over\sqrt3}
                       >\sqrt3>3/2 .
\]
At old four, either contact of at least
four gives \(C_i\ge\sqrt3\);
both must be three.
At old three, both contacts at least four
give the same contradiction.
One is therefore three; if the other
\(\beta\ge\cos(\pi/5)>4/5\), the opposite
edge has \(C^2=(1+\beta)/(3(1-\beta))>3\).
The remaining possibilities are three/three,
with zero time, or three/four in either
direction orientation.
This completes the classification.

Put \(\varphi=(1+\sqrt5)/2\), \(s=\sqrt2\).
The SIX necessary old-role models are
\[
 \begin{array}{c|c|c|c}
 S&d\text{ contacts in }S&h&K\\ \hline
 I_2(4)&\text{one }5&(\varphi-1)/2&\mathbb Q(\varphi)\\
 I_2(5)&\text{one }4&1/\sqrt5&\mathbb Q(\varphi)\\
 A_2&\text{one }8&(s-1)/3&\mathbb Q(s)\\
 I_2(8)&\text{one }3&s/2&\mathbb Q(s)\\
 A_2&3,4\text{ in either orientation}&s/3&\mathbb Q(s)\\
 I_2(4)&3,3&s/2&\mathbb Q(s).
 \end{array}
\tag{A14six}\label{eq:arithA14sixmodels}
\]
They are necessary actual Gram types,
never realized chambers.

We justify the ACTUAL field column.
The native two-cycles put the quadratic
field inside \(K\) in the tree types
\((2,4,5)\), \((2,3,8)\).
For \((3,3,4)\), the essential actual
THREE-cycle is \(8kbg=\sqrt2\).
Choose the degree-two root as base in a
tree, and any root with both products
nonzero in the triangle.
Let its original \(K\)-root square be
the totally positive \(q_0\).
Rescale the two other original roots by
ratios of \(K\)-valued products so their
products with the base are \(-q_0\).
Their squares become \(q_0/k^2,q_0/b^2\);
the remaining product is zero or
\(-q_0g/(kb)\), whose coefficient is in
the displayed field \(F\) because
\(g/(kb)=kbg/(k^2b^2)\).
Thus the three ACTUAL root lines have
Gram \(q_0M\), \(M\) over \(F\).
It is Lorentzian by \(Q=0,h>0\).
A proper totally real extension \(K/F\)
would have a nonphysical embedding fixing
this negative direction while preserving
the positive scalar.
Strict nonphysical positivity forbids it.
Hence \(K=F\), without identifying raw
scales across the old components.

\subsection{Every ordinary-core \(A_{14}\) endpoint exceeds \(.44\)}
\label{app:arithA14ordinaryends}

Name the five distinct heights in
\eqref{eq:arithA14sixmodels}
\[
 h_a=(\varphi-1)/2,\quad h_b=1/\sqrt5,\quad
 h_c=(s-1)/3,\quad h_d=s/2,\quad h_e=s/3 .
\]
Every one is positive and at most \(s/2\).
For an actual endpoint cap
\(c=A_cp-v_A-v_F\), let
\(\gamma=-c\cdot n_1\ge0\),
\(Q=\|v_F\|^2\), and \(\delta>0\)
be the full retained-white Schur residual.
Its actual edge obeys
\eqref{eq:arithA14endpointlaw}.
It suffices to prove \(C_{\rm end}>31/20\),
because then its signed coordinate satisfies
\[
                           y^2>{154\over765}>.44^2 .
\]
This actual wall differs from the common
spectator wall \(d\), which is orthogonal
to all \(A_{14}\) directions.
Their mutual magnitude is
\[
                \lambda=-c\cdot d=A_c\sqrt h-\rho\ge0,
                         \qquad \rho=v_F\cdot v_S\ge0.
\tag{A14mut}\label{eq:arithA14ordinarymutual}
\]
A positive \(\lambda\) is at least \(1/2\).
No individual unit product is declared
\(K\)-valued: its invariant square is.
The finite positive mutual labels are
\[
 \begin{array}{c|c}
 K&\text{finite positive labels}\\ \hline
 \mathbb Q(\varphi)&3,4,5,6,10\\
 \mathbb Q(s)&3,4,6,8 .
 \end{array}
\]
Indeed cyclotomic degree at most two first
permits \(3,4,5,6,8,10,12\);
the \(5/10,8,12\) fields are respectively
\(\mathbb Q(\sqrt5),\mathbb Q(\sqrt2),\mathbb Q(\sqrt3)\).
The full orders are \(\mathbb Z[\varphi]\),
\(\mathbb Z[s]\).

First suppose \(c\) is PURE isolated:
no retained \(A_{13}\) or \(S\) contact.
Then \(\delta=1\),
\(A_c^2=(28/15)\gamma^2-1\),
and \(\lambda=A_c\sqrt h>0\).
Its own two-cycle requires
\[
                    4\gamma^2={15\over7}(1+\lambda^2/h)
                                     \in\mathcal O_K .
\]
Here are ALL values \(z=1+\lambda^2/h\)
for the finite mutual alphabet:
\[
 \begin{array}{c|ccccc}
 &3&4&5&6&10\\ \hline
 h_a&(2+\varphi)/2&1+\varphi&(3+2\varphi)/2&
                    (2+3\varphi)/2&(3+3\varphi)/2\\
 h_b&(3+2\varphi)/4&(1+2\varphi)/2&(5+3\varphi)/4&
                    (1+6\varphi)/4&(4+5\varphi)/4
 \end{array}
\]
and
\[
 \begin{array}{c|cccc}
 &3&4&6&8\\ \hline
 h_c&(7+3s)/4&(5+3s)/2&(13+9s)/4&(16+9s)/4\\
 h_d&(4+s)/4&(2+s)/2&(4+3s)/4&(3+s)/2\\
 h_e&(8+3s)/8&(4+3s)/4&(8+9s)/8&(7+3s)/4 .
 \end{array}
\]
Each dyadic numerator has a coefficient
not divisible by seven.
Multiplication by \(15/7\) cannot make
both full-order coefficients integral.
Thus EVERY finite mutual is impossible.
Retain parallel/dotted mutuals by the
weaker conclusion \(\lambda\ge1\).
Then \(A_c^2\ge1/h\ge s\), so
\[
                  C_{\rm end}^2\ge1+s>(31/20)^2 .
\]

For PURE nonisolated caps, the complete
endpoint roles and times are
\eqref{eq:arithA14endpointtimes}.
Their residuals, in row order, are
\[
                         15/28,\quad1/7,\quad1/7,\quad15/28,
                                      \quad1/14,\quad1/14 .
\]
At own zero only row \(j=3\), helper-three,
and near helper-four have positive time.
The first has \(\lambda^2=h/5<1/4\),
impossible.
For the second \(\lambda^2=11h/15\):
it lies in \((0,1/4)\) at \(h_a,h_c\),
in \((1/4,1/2)\) at \(h_b,h_e\),
and in \((1/2,\cos^2(\pi/5))\) at \(h_d\).
These are all native gaps.
The coarse exact brackets
\(1.6<\varphi<5/3\),
\(2<\sqrt5<5/2\), \(1.4<s<1.5\)
seal every strict interval.
Every other own-zero time is nonpositive.

For ANY positive own \(\gamma\ge1/2\),
the first three helper-three rows already
give \(C_{\rm end}\ge9/5\),
\(C_{\rm end}^2\ge169/15\),
\(C_{\rm end}^2\ge64/15\).
Both helper-four rows give
\(C_{\rm end}^2>113/15\).
Only far helper-three at \(j=14\) remains.
Its actual simple-three normalization gives
\(2\gamma\in K\).
Thus finite own is
\(0,1/2,\varphi/2\) in \(\mathbb Q(\varphi)\),
or \(0,1/2,s/2\) in \(\mathbb Q(s)\).
Zero is negative time and \(1/2\) has zero
time.  At \(\varphi/2\), time
\((8\varphi-1)/15\) gives mutual squares
\[
        (9-\varphi)/30\in(0,1/4),\qquad
       (8\varphi-1)/(15\sqrt5)\in(1/4,1/2) .
\]
At \(s/2\), time \((6+s)/15\) gives,
for \(h_c,h_d,h_e\),
\[
 {5s-4\over45}\in(0,1/4),\quad
 {3s+1\over15}\in(1/4,1/2),\quad
 {6s+2\over45}\in(0,1/4) .
\]
All are forbidden native gaps.
Parallel own has its positive actual null.
For EVERY dotted \(\gamma>1\) all time
polynomials increase and the minimum far
time exceeds \(22/15\).
With residual at most \(15/28\) this gives
\(C_{\rm end}>29/15\).
The entire unbounded own range is covered.

We now keep the ENTIRE foreign \(S\) block.
At most one \(S\) contact is possible in
the retained finite Gram; two would make
a nonpositive cycle through its old bond.
For old \(I_2(4)\), \(I_2(5)\), only native
three is finite, of costs \(1/2,d_5/4\),
where \(d_5=2+2/\sqrt5\).
For old \(I_2(8)\), any contact costs \(>1\).
For old \(A_2\), the finite labels are
three, four, five; its actual \(\mathbb Q(s)\)
excludes five, leaving costs \(1/3,2/3\).
White parallel/dotted products cannot occur.

If there is also an \(A_{13}\) contact,
the full connected finite rank-sixteen
diagram is classical.
An old five/eight bond is impossible there.
For old \(S=A_2,I_2(4)\), the ONLY join is
an \(A_{13}\) END-three and a native-three
\(S\) contact; an internal or weighted
attachment gives excessive arms, an internal
double bond, or a cycle.
The exact residuals are
\[
                    S=I_2(4):\ \delta=1/28,\qquad
                    S=A_2:\ \delta=17/84 .
\]
At old four, \(\rho>1/2,h<3/4\), so
acuteness alone gives \(A_c^2>1/3\) and
\(C_{\rm end}^2>31/3\).
At old \(A_2\), the complete inverse products
give
\[
                   {\rho_{\min}^2\over h}
                       ={1\over3}+{s\over4}>{41\over60}.
\]
Thus \(C_{\rm end}^2>1+(41/60)(84/17)>4\).
These estimates retain all own parameters
and require no finite mutual assumption.

It remains to consider foreign caps isolated
WITHIN \(A_{14}\), with
\[
                 A_c^2=(28/15)\gamma^2+Q-1,\qquad \delta=1-Q .
\]
The full inverse products yield
\[
 \begin{array}{c|c|c|c}
 S,h&\text{foreign label}&\rho_{\min}^2/h&1-Q\\ \hline
 I_2(4),h_a&3&\varphi^3/4>1&1/2\\
 I_2(4),h_d&3&(4+3s)/8>1&1/2\\
 I_2(5),h_b&3&d_5(d_5-1)/(4(d_5-2))>3/2&1-d_5/4<3/10\\
 A_2,h_c\text{ or }h_e&4&2(1/3+s/4)>4/3&1/3 .
 \end{array}
\]
The \(d_5\) inequality reduces to
\((d_5-3)(d_5-4)>0\), with \(2<d_5<3\).
Acuteness gives respectively
\(C_{\rm end}^2>3,3,6,5\), sufficient
even when the mutual is zero.

Only old \(A_2\), foreign-three, remains.
At \(h_c\), the two inverse products are
\(a_8/3,2a_8/3\), \(a_8=\cos(\pi/8)\);
at \(h_e\), they are
\((s+2)/6,(2s+1)/6\).
These include BOTH old-node orientations.
If \(\lambda=0\), the exact time equation
forces the own two-cycle \(4\gamma^2\)
to be one of
\[
                  {15(4+s)\over28},\qquad
                  {15(2+s)\over7},\qquad
                  {15(8+3s)\over56}.
\]
The first occurs for the smaller product
in both models; the others are their
respective larger products.
None is integral in \(\mathbb Z[s]\),
by its nonzero coefficient residue at seven.
Thus every zero-mutual exception is impossible
for ANY finite or dotted own.
If \(\lambda>0\), then \(\lambda\ge1/2\),
while \(h<1/2,\rho>.3\).
Hence \(A_c^2>(.5+.3)^2/.5=32/25\) and
\[
                    C_{\rm end}^2>
                      1+{32/25\over2/3}=73/25>(31/20)^2 .
\]
We have covered every white, own, foreign
and zero-mutual role.
Both actual endpoints therefore satisfy
\[
                     x_{A,1}>.44,\qquad -x_{A,15}>.44 .
\tag{A14ordends}\label{eq:arithA14ordinaryendpoint}
\]
Reversing index names proves the second
demand without requiring a chamber symmetry
or distinct opposite walls at the two ends.

\subsection{Ordinary \(A_{14}\): pay both ends and all feet afresh}
\label{app:arithA14ordinarycomposition}

We return to the COMPLETE \(.597\) local
inventory and directional pressures in
Section~\ref{app:arithA14fourfifths}.
There is at most ONE weak golden helper.
All other weak central sites are the
retained pure/foreign \(0<u_i<1\) feet;
no adjacency or packing restriction is
imposed on them.
Their full continuous adjusted-loss bounds
hold whenever the ACTUAL remaining prices
satisfy
\[
             \beta_{5,10}\ge.06,\quad
             \beta_{6,9}\ge.026,\quad
             \beta_{7,8}\ge.0246 .
\]
These are individual coefficient thresholds,
not a total debit budget.

The reference must be rebuilt BEFORE any
new energy is used.
Let \(v=\sum\beta_j\Omega_j\) be a concave
reference whose endpoint prices are paid
by actual endpoint demands, and whose
central heights retain their strong floors.
For ALL paid isolated feet put
\[
                         Z=v-\sum\beta_i u_i n_i .
\]
Its actual pressure and full norm obey
\[
 \begin{split}
 x_A\cdot Z&\ge v^2-\sum\beta_i(L_i-.597),\\
 Z^2&\le v^2-\sum\beta_i^2(2u_i-u_i^2).
 \end{split}
\]
Here \(v\cdot n_i=\beta_i\) is the NEW
coefficient, including every endpoint/foot
cross.  Every adjacent-foot cross is
nonpositive.
The previously proved full-interval loss
bounds give \(e_A>v^2\) by ONE square.
Strict endpoint pressures give strictness
also when no foot remains.
We use orthogonality of the block changes
only to compute the NEW reference norm,
never to ignore their crosses with feet.

Use the base heights
\[
 H_5=H_{10}=a=1.98498,\quad
 H_6=H_9=b=2.08628,\quad
 H_7=H_8=c_0=2.13551 .
\]
The full base norm is \(.80071440298\).
If one internal helper lies at six through
nine, omit ONLY its knot by the neighbor
chord, losing at most \(.000677821225\).
The neighboring coefficients increase.
Thus the base norm is at least
\(.800036581755\), with its first and
last five coordinate blocks still flat
at signed mean
\[
                           b_0={a\over5\sqrt2},\qquad .28<b_0<.3 .
\]
In each block replace the signed endpoint
by its actual paid value \(.44\), keeping
the block sum fixed.
The other four become
\(r=(5b_0-.44)/4\), with
\(.24<r<.265<.44\).
They exceed the next old signed coordinate,
which is \(<.073\).
The modified reference is concave, its
new endpoint coefficients are positive,
and EACH changed outer central price is
\[
                    >.167/1.415>.118>.06 .
\]
Every other paid central coefficient stays
unchanged or increases; the helper price
remains zero.
All central heights are preserved.
The two disjoint block gains give
\[
              v^2>.800036581755+
                            2(5/4)(.44-.3)^2=.849036581755 .
\tag{A14ord0}\label{eq:arithA14ordinarynoouter}
\]
Apply the simultaneous-foot square to
THIS new reference.

For ONE outer helper, relabel it as site
five.  Its ACTUAL \(H_4\) point already
proved in Section~\ref{app:arithA14fourfifths}
gives the signed helper-end demand \(>.51961\),
in BOTH helper orientations.
Pay that endpoint at \(.51\), and the
OTHER endpoint at \(.44\) by
\eqref{eq:arithA14ordinaryendpoint}.
Omit only knot five by the chord zero-to-six.
Its base norm is \(>.764\).
The left six-block mean lies in \((.24,.25)\);
the right five-block mean lies in \((.25,.3)\).
Distribute each remaining block sum
constantly after replacing its endpoint.

If the helper end is near, the five other
left coordinates exceed
\((1.44-.51)/5=.186\), while the other
block's four exceed
\((1.25-.44)/4=.2025\).
If it is far, the right four exceed
\((1.25-.51)/4=.185\), and the left
five exceed \((1.44-.44)/5=.2\).
Each is positive, below its endpoint
target and above the next old coordinate
\(<.073\).
Thus both new boundary prices exceed
\(.112/1.415>.079\).
Site six needs only \(.026\);
site ten needs \(.06\).
All other prices retain their own thresholds;
the missing-five price remains zero.
No strong-five floor has been used.

For near helper the exact block-gain
formula gives gains \(>.08112,.0245\).
For far helper they are \(>.055125,.04332\).
Consequently
\[
        v_{\rm near}^2>.86962,\qquad
        v_{\rm far}^2>.862445 .
\]
The SAME simultaneous-foot square applies,
retaining every foot and all boundary
crosses with these end changes.

The minimum of the three complete branches
is~\eqref{eq:arithA14ordinarynoouter}, which
exceeds \(16/19\):
it exceeds \(.849\), and \(19(.849)>16\).
Only NOW add the disjoint ORDINARY \(T_2\)
energy \(>3/19\) proved in
Proposition~\ref{prop:ordinarydihedralthreenineteenths}.
The original norm exceeds one.
Thus the entire ordinary \(A_{14}+T_2\)
branch is excluded.
Together with the actual typed-seven proof
and the already stated two-\(A_1\) and
joint-\(I+A_1\) consumers, this covers every
rank-sixteen \(A_{14}\) alternative.
The proof never increases the ordinary
dihedral floor or transfers a \(K_7\)
assumption to this quadratic-field branch.

\subsection{The rank-seventeen \(A_{15}\) typed alternatives}
\label{app:arithA15typed}

Here the WHOLE forest is \(A_{15}\) plus
one actual rank-two typed defect.
Thus \(K=K_7\), its disjoint defect energy
is \(>.09\), and the ENTIRE foreign cost is
\[
                         Q\in\{0,1/3,2/3\}.
\]
Positive values occur only for old \(A_2\):
finite foreign labels are three, four, five,
and \(K_7\) excludes five.
For old \(I_2(7)\), even a three costs
\(\csc^2(\pi/7)/4>1\).
There are no further spectators.

Use \(\vartheta=2\cos(2\pi/7)\), satisfying
\(p(X)=X^3+X^2-2X-1\).
The earlier generator \(\theta=2\cos(\pi/7)\)
satisfies
\[
           \vartheta=\theta^2-2,\qquad
           \theta=\vartheta^2+\vartheta-1 .
\]
Hence the already proved FULL integer ring
is also \(\mathbb Z[\vartheta]\).
Its roots and derivatives satisfy
\[
\begin{gathered}
 1.246<\vartheta_1<1.247,\quad
 -.446<\vartheta_2<-.445,\quad
 -1.802<\vartheta_3<-1.801,\\
 5<p'(\vartheta_1)<6,\quad-3<p'(\vartheta_2)<-2,\quad
 4<p'(\vartheta_3)<5,\quad\vartheta_2-\vartheta_3>1.355 .
\end{gathered}
\]
The three paired endpoint signs are
\((-.005049064,.000105223)\),
\((.002199464,-.000096125)\),
\((-.000257608,.003875599)\);
derivative bounds then follow on the brackets.
For ANY integer \(y=a+b\vartheta+c_0\vartheta^2\),
Lagrange interpolation gives
\[
 c_0=\sum_j{y_j\over p'(\vartheta_j)},\qquad
 {y_2-y_3\over\vartheta_2-\vartheta_3}
                 =b-(1+\vartheta_1)c_0 .
\]
We prove three COMPLETE integer-strip exclusions.

There is no integer with
\[
              4<y_1<16,\qquad0<y_2,y_3<32/55 .
\tag{A15stripA}\label{eq:arithA15typedstripA}
\]
Interpolation gives \(0<c_0<4\) and
\(|b-(1+\vartheta_1)c_0|<.43\).
At \(c_0=1,2,3\), respectively,
\(b=2\), no integer, or \(b=7\).
For the first case
\[
          -.358399<\vartheta_3^2+2\vartheta_3<-.356796
\]
forces \(.356796<a<32/55+.358399<1\).
For the last case
\[
           -2.876197<3\vartheta_3^2+7\vartheta_3<-2.872388
\]
forces \(a=3\), but the physical value then
exceeds \(3+7(1.246)+3(1.246)^2=16.379548\).

There is no integer with
\[
                9/5<y_1<68/25,\qquad-.8<y_2,y_3<.48 .
\tag{A15stripB}\label{eq:arithA15typedstripB}
\]
Interpolation gives \(-.14<c_0<1.064\);
hence \(c_0=0,1\).
The difference bound is \(<.945\).
At zero \(b=0\), giving incompatible
physical and small rational-integer strips.
At one \(b=2,3\); their values at
\(\vartheta_2\), before adding \(a\), lie in
\((-.693084,-.691975)\),
\((-1.139084,-1.136975)\).
The small-value strip permits \(a=0,1\)
in the first case, \(a=1\) in the second.
Their physical values exceed \(4.044516\),
\(6.290516\), respectively.

Finally no integer has
\[
                   2<y_1<21/10,\qquad-7/5<y_2,y_3<-1/5 .
\tag{A15stripC}\label{eq:arithA15typedstripC}
\]
Now \(-1<c_0<2\):
the interpolation lower and upper bounds
are \(-1/60\), \(1.12\).
The difference bound is \(<.886\).
At zero the rational case fails;
at one \(b=2,3\).
The same \(\vartheta_2\) intervals force
\(a=0\), with physical values
\(>4.044516,5.290516\).
All three full-ring strips are proved.

At \(N=16\) use sites \(5,\ldots,11\)
and loss \(1/2\).
Here \(D_i=i(16-i)/8\ge55/8\), \(2.6<q_i<3\).
The full white list is the two paths
of ranks at least four, with the usual
endpoint/\(D/E\) three, endpoint four,
or two-end-three join.
Rank fifteen excludes exceptional completions.
The increasing sufficient ratio is
\[
             T(q)=2-{q-1/2\over q(2q-1/2)}<61/33 .
\]
Own-zero global end-three has time
\(Q-17/32\), \(Q<3/5\), negative for
the actual foreign-cost set.
Next-end-three has time \(Q-1/8\),
\(Q=0,1/3\); its positive time is
less than its foreign cost.
End-four has \(Q=0\) and negative time.
The actual reciprocal law, affine internal
three/four and two-three kernels, and
parallel nulls dispose of all those roles.

For a simple-three helper \(t=2\gamma\)
is integral in \(K_7\); the only positive
finite own cosines themselves in this
field are labels three and seven.
Three closes the actual affine cycle.
For native seven at a NONEND helper node
\(j\), \(D_j\ge7/2\), and its positive
conjugate \(t_\sigma=2\cos(3\pi/7)>2/5\)
violates Schur positivity already through
\[
                 D_it_\sigma^2+D_j>
                       (55/8)(4/25)+7/2=23/5>4 .
\]
The omitted cross and foreign costs are
nonnegative.
For dotted owns the remaining lower
ratios are \(6/\sqrt{10}>61/33\) at
\(D\)-positions, \(2\gamma+3/8>19/8\)
at \(E\)-positions, and
\[
 {3m-1\over\sqrt{2m(m+1)}}\ge17/\sqrt{84}>61/33
\]
for near-three of order \(m\ge6\).
The last function increases; its cleared
square margin at six is \(2157\).
Endpoint-four/own-three would put
\(\sqrt2\) in \(K_7\); other owns have
ratio at least \(6/\sqrt{10}\).
The two-path bridge has
\(\gamma q+1/q>.9q+1/q>61/33\).
Foreign far-three has \(Q\ge1/3\)
and ratio square
\[
          \ge{3(2m\gamma+1)^2\over2m(m+3)}
                         >15/4>(61/33)^2 .
\]
It increases with \(m\ge5\) at \(\gamma>.9\).
Only pure far helper or near dotted
side-order-five helper remains.

For the pure far role the square identity
gives \(9/5<t<5899/2178<68/25\).
The fresh \(N=16,i=5\) signed Schur
costs exceed four at \(-.8,.48\)
by \(3/40,59/1000\).
Their monotonicity puts both other values
in \((-.8,.48)\).
Strip~\eqref{eq:arithA15typedstripB}
excludes every such cap.
For near dotted order five,
\(2<t<1363/660<21/10\).
Its necessary conjugate inequality is
\(55t_\sigma^2+88t_\sigma+16<0\).
That quadratic equals \(3/5>0\) at
\(-7/5,-1/5\), so both values lie
between these endpoints.
Strip~\eqref{eq:arithA15typedstripC}
excludes the last nonisolated weak cap.

For isolation within \(A\), the actual
own two-cycle \(B=4\gamma^2\) is integral,
and full conjugate Schur positivity gives
\[
                         0<B_\sigma<4/D_i\le32/55 .
\]
Use the COMPLETE finite-square alphabet
already proved in
Section~\ref{app:arithA14typedends}.
Small owns three/four/six have all-foreign
actual nulls.
The other finite labels seven/fourteen
each have a conjugate \(>1\), using
the stated brackets for \(\theta^2\),
\(2+\theta\).
Thus no finite isolated own survives.
For a dotted own, \(0<u=(1-Q)/\gamma^2<1\).
Its actual all-wall foot has loss
\[
                  <u/5+\sqrt{u(1-u)},
\]
since \((\sqrt D+\sqrt{D-1})^{-1}<1/5\).
At \(u\le1/4\) this is
\(<1/20+\sqrt3/4<.484<1/2\).
Weakness therefore implies \(u>1/4\)
and \(4<B=4(1-Q)/u<16\).
Strip~\eqref{eq:arithA15typedstripA}
excludes every such dotted foot, including
both positive foreign costs.

Every site \(5,\ldots,11\) is now actually
half-loss strong.
Choose symmetric heights, linear on the tails,
\[
           H_5=2.122,\quad H_6=2.238,\quad
           H_7=2.306,\quad H_8=2.328 .
\]
The strict floor tests are
\(2.622^2<55/8\), \(2.738^2<15/2\),
\(2.806^2<63/8\), \(2.828^2<8\).
The successive slopes
\(.4244,.116,.068,.022\) and their
reflection decrease at all paid knots.
One positive reference square gives
\[
        e_{A_{15}}>2.122^2/5+.116^2+.068^2+.022^2=.9191408 .
\]
Only NOW add the disjoint actual defect
energy \(>.09\), obtaining \(>1.0091408\).
Both actual typed-seven \(A_{15}\) alternatives
are excluded, with no remaining weak foot,
field census or independently added optimum.

































